% Vectors I Homework 3 (Vector Product): HOUSE FORMAT (house_format/SPEC.md), converted 2026-11-09.
% Questions, marks, diagrams and solution working copied word for word from sources/*.tex
% (backups: house_format/backups/pre_phase2_20261109, house_format/phase2/vec_backup/hw3).
% The same file feeds the worked solutions (HW3_Solutions.tex).
\documentclass[a4paper,11pt,oneside]{book}
\providecommand\housemode{hw}
% =====================================================================
% house.tex: ONE preamble for every Maximus Pung Jun H2 Math document,
% modelled on Ch10 Definite Integrals + Ch11 Differential Equations (see SPEC.md).
%
% Before % =====================================================================
% house.tex: ONE preamble for every Maximus Pung Jun H2 Math document,
% modelled on Ch10 Definite Integrals + Ch11 Differential Equations (see SPEC.md).
%
% Before % =====================================================================
% house.tex: ONE preamble for every Maximus Pung Jun H2 Math document,
% modelled on Ch10 Definite Integrals + Ch11 Differential Equations (see SPEC.md).
%
% Before \input{house}, set:
%   \def\housemode{notes}    lecture notes: book class, fncychap Lenny chapter page, boxes
%   \def\housemode{hw}       homework / additional / workshop QUESTIONS (+ Answers section)
%   \def\housemode{sol}      worked SOLUTIONS of a hw (RI "Solution | Comments" table)
%   \def\housemode{test}     test question paper (prelim cover, prelim layout)
%   \def\housemode{testsol}  test worked solutions (prelim layout + RI table)
%   \def\housechapterno{10}  (notes: chapter number; optional)
%   \def\housenotesLM{1}     (notes only, optional) keep Latin Modern as in the current
%                            DI/DE notes instead of Times New Roman + TeX Gyre Termes Math
% then \housechapter{Definite Integrals} and \ditle{Homework 1}.
% Document class (put BEFORE \input{house}):
%   notes: \documentclass[a4paper]{book}
%   others: \documentclass[a4paper,11pt,oneside]{book}
% A HW file uses \providecommand\housemode{hw}; its solutions file has NO \documentclass:
%   \def\housemode{sol}\def\diquestionsonly{1}\input{HW1}  (as DI/DE HWn_solutions.tex)
% Compile: TEXINPUTS=.:/workspace/house_format: xelatex (2 runs; solutions until .aux stable)
% =====================================================================
\providecommand\housemode{hw}
\providecommand\housedir{/workspace/house_format}
\newcommand\housechap{}
\newcommand\housechapter[1]{\renewcommand\housechap{#1}}
\def\housexnotes{notes}\def\housexhw{hw}\def\housexsol{sol}\def\housextest{test}\def\housextestsol{testsol}
\edef\housexm{\housemode}
\input{house_base}
\input{house_local}
\ifx\housexm\housexnotes
  \ifdefined\housenotesLM\else\input{house_fonts}\fi
  \input{house_tf}
  \providecommand{\units}[1]{\ \text{units}^{#1}}
  \raggedbottom
\else
  \input{house_fonts}
  \input{house_hw}
  \input{house_solmacros}
  \input{house_tf}
  \ifx\housexm\housexsol \def\prelimkeepcolour{1}\input{house_prelim}\input{house_sol}\tfanswerstrue\fi
  \ifx\housexm\housextest \input{house_prelim}\fi
  \ifx\housexm\housextestsol \def\prelimkeepcolour{1}\input{house_prelim}\input{house_sol}\tfanswerstrue\fi
\fi
, set:
%   \def\housemode{notes}    lecture notes: book class, fncychap Lenny chapter page, boxes
%   \def\housemode{hw}       homework / additional / workshop QUESTIONS (+ Answers section)
%   \def\housemode{sol}      worked SOLUTIONS of a hw (RI "Solution | Comments" table)
%   \def\housemode{test}     test question paper (prelim cover, prelim layout)
%   \def\housemode{testsol}  test worked solutions (prelim layout + RI table)
%   \def\housechapterno{10}  (notes: chapter number; optional)
%   \def\housenotesLM{1}     (notes only, optional) keep Latin Modern as in the current
%                            DI/DE notes instead of Times New Roman + TeX Gyre Termes Math
% then \housechapter{Definite Integrals} and \ditle{Homework 1}.
% Document class (put BEFORE % =====================================================================
% house.tex: ONE preamble for every Maximus Pung Jun H2 Math document,
% modelled on Ch10 Definite Integrals + Ch11 Differential Equations (see SPEC.md).
%
% Before \input{house}, set:
%   \def\housemode{notes}    lecture notes: book class, fncychap Lenny chapter page, boxes
%   \def\housemode{hw}       homework / additional / workshop QUESTIONS (+ Answers section)
%   \def\housemode{sol}      worked SOLUTIONS of a hw (RI "Solution | Comments" table)
%   \def\housemode{test}     test question paper (prelim cover, prelim layout)
%   \def\housemode{testsol}  test worked solutions (prelim layout + RI table)
%   \def\housechapterno{10}  (notes: chapter number; optional)
%   \def\housenotesLM{1}     (notes only, optional) keep Latin Modern as in the current
%                            DI/DE notes instead of Times New Roman + TeX Gyre Termes Math
% then \housechapter{Definite Integrals} and \ditle{Homework 1}.
% Document class (put BEFORE \input{house}):
%   notes: \documentclass[a4paper]{book}
%   others: \documentclass[a4paper,11pt,oneside]{book}
% A HW file uses \providecommand\housemode{hw}; its solutions file has NO \documentclass:
%   \def\housemode{sol}\def\diquestionsonly{1}\input{HW1}  (as DI/DE HWn_solutions.tex)
% Compile: TEXINPUTS=.:/workspace/house_format: xelatex (2 runs; solutions until .aux stable)
% =====================================================================
\providecommand\housemode{hw}
\providecommand\housedir{/workspace/house_format}
\newcommand\housechap{}
\newcommand\housechapter[1]{\renewcommand\housechap{#1}}
\def\housexnotes{notes}\def\housexhw{hw}\def\housexsol{sol}\def\housextest{test}\def\housextestsol{testsol}
\edef\housexm{\housemode}
\input{house_base}
\input{house_local}
\ifx\housexm\housexnotes
  \ifdefined\housenotesLM\else\input{house_fonts}\fi
  \input{house_tf}
  \providecommand{\units}[1]{\ \text{units}^{#1}}
  \raggedbottom
\else
  \input{house_fonts}
  \input{house_hw}
  \input{house_solmacros}
  \input{house_tf}
  \ifx\housexm\housexsol \def\prelimkeepcolour{1}\input{house_prelim}\input{house_sol}\tfanswerstrue\fi
  \ifx\housexm\housextest \input{house_prelim}\fi
  \ifx\housexm\housextestsol \def\prelimkeepcolour{1}\input{house_prelim}\input{house_sol}\tfanswerstrue\fi
\fi
):
%   notes: \documentclass[a4paper]{book}
%   others: \documentclass[a4paper,11pt,oneside]{book}
% A HW file uses \providecommand\housemode{hw}; its solutions file has NO \documentclass:
%   \def\housemode{sol}\def\diquestionsonly{1}\input{HW1}  (as DI/DE HWn_solutions.tex)
% Compile: TEXINPUTS=.:/workspace/house_format: xelatex (2 runs; solutions until .aux stable)
% =====================================================================
\providecommand\housemode{hw}
\providecommand\housedir{/workspace/house_format}
\newcommand\housechap{}
\newcommand\housechapter[1]{\renewcommand\housechap{#1}}
\def\housexnotes{notes}\def\housexhw{hw}\def\housexsol{sol}\def\housextest{test}\def\housextestsol{testsol}
\edef\housexm{\housemode}
% house_base.tex: HOUSE COPY of definite_integrals/common/ws_preamble_build.tex
% (= Maximus's own preamble.tex + the DE v3 / AoI additions). Only change: GC font path. [HOUSE]
% WORKSHEET COPY of ../notes/preamble_build.tex (2026-10-09 version), unchanged.
% preamble_build.tex: copy of Maximus's preamble.tex for the DE v3 build.
% The blocks marked [BUILD COPY CHANGE] are the earlier build fixes.
% The blocks marked [V3 ADDITION] are new in this version (example box,
% spacing). Everything else is unchanged from preamble.tex.

% For margins
\usepackage[a4paper, margin = 20mm]{geometry}

\usepackage[Lenny]{fncychap}
% Rejne, Glenn, Lenny

\usepackage{etoc}

% [V3 ADDITION] cancel: used to show the dx's "cancelling" in Section 11.3
\usepackage{cancel}
\renewcommand{\CancelColor}{\color{red}}

% For enumeration & math symbols
\usepackage{enumerate}
\usepackage{amsmath,amsfonts,amssymb,amsthm,graphicx,lipsum}

% For diagrams, plots and boxes
\usepackage{pgfplots}
\pgfplotsset{compat=1.18}
\usepackage{tikz}
\usepackage[most]{tcolorbox}
\tcbuselibrary{skins}
\usepackage{xcolor}
\pgfplotsset{
    standard/.style={clip = false,
    axis line style = thick,
    trig format=rad,
    axis lines = middle,
    every axis x label/.style={at={(current axis.right of origin)},anchor= west},
    every axis y label/.style={at={(current axis.above origin)},anchor=south}
    }
}

% For GC font
\usepackage{fontspec}
% [BUILD COPY CHANGE 1] The TI-84 key font is loaded only if the file exists
% (checked in ./Fonts/... and ../Fonts/...). Otherwise \ti{} falls back to a
% small boxed sans-serif key. Known letters of the TI font are mapped to their
% key names in the fallback: o = [y=], z = [mode], f = [L3].
\newcommand{\tikey}[1]{\tcbox[on line, enhanced, boxsep=0pt, left=2pt, right=2pt,
  top=1pt, bottom=1pt, boxrule=0.5pt, arc=1.5pt, colback=white,
  colframe=black!70]{\sffamily\footnotesize\strut #1}}
\ExplSyntaxOn
\NewDocumentCommand{\tifallback}{m}{
  \str_case:nnF {#1}
    { {o}{\tikey{y=}} {z}{\tikey{mode}} {f}{\tikey{L3}} }
    { \tikey{#1} }
}
\ExplSyntaxOff
\def\tifontcase{0}
\IfFileExists{Fonts/TI-84 Plus CE Key Fonts/TI84PlusCEKeys.ttf}
  {\def\tifontcase{1}}{}
\IfFileExists{../Fonts/TI-84 Plus CE Key Fonts/TI84PlusCEKeys.ttf}
  {\def\tifontcase{2}}{}
\ifcase\tifontcase\relax
  % 0: font not found, use boxed fallback keys
  \newcommand{\ti}[1]{\tifallback{#1}}
\or
  % 1: original location (as in preamble.tex)
  \newfontfamily\tifont[
    Path = {Fonts/TI-84 Plus CE Key Fonts/},
    UprightFont = {TI84PlusCEKeys.ttf}
  ]{}
  \newcommand{\ti}[1]{{\tifont #1}}
\or
  % 2: one level up (when compiling inside draft-ch1/)
  \newfontfamily\tifont[
    Path = {../Fonts/TI-84 Plus CE Key Fonts/},
    UprightFont = {TI84PlusCEKeys.ttf}
  ]{}
  \newcommand{\ti}[1]{{\tifont #1}}
\fi

% [AoI items 17/18/37] GC DISPLAY FONT. Every GC key, command, menu item and
% screen line goes through \gcfont, so the font can be swapped HERE in one place.
%   \gc{fnInt(}      command / menu item / stored variable in the display font
%   \gckey{2nd}      a calculator key (display font in a small key outline)
%   \gcscreen{...}   a calculator screen (display font on a grey-green LCD panel)
% The font file lives in ../Fonts (shared by notes, slides and worksheet).
\def\gcfontfile{VT323-Regular.ttf}
% [HOUSE] GC font: always the absolute copy in \housedir/Fonts (relative Path= breaks under TEXINPUTS)
\providecommand\housedir{/workspace/house_format}
\IfFileExists{\housedir/Fonts/\gcfontfile}{\newfontfamily\gcfont[Path=\housedir/Fonts/, Scale=1.32]{\gcfontfile}}{%
  \newcommand\gcfont{\ttfamily}}
\newcommand{\gc}[1]{\mbox{\gcfont #1}}
\newcommand{\gckey}[1]{\tcbox[on line, enhanced, boxsep=0pt, left=2pt, right=2pt,
  top=0.6pt, bottom=0.2pt, boxrule=0.5pt, arc=1.5pt, colback=black!4,
  colframe=black!65]{\gcfont\strut #1}}
\definecolor{lcd}{RGB}{214,222,204}
\newcommand{\gcscreen}[1]{\tcbox[on line, enhanced, boxsep=1pt, left=4pt, right=4pt,
  top=2pt, bottom=2pt, boxrule=1.2pt, arc=1pt, colback=lcd, colframe=black!70]{%
  \gcfont\begin{tabular}{@{}l@{}}#1\end{tabular}}}
% integral template as shown on the TI-84 Plus CE (MathPrint): \gcint{0}{2}{e^(-X²)}{X}
\newcommand{\gcint}[4]{{\gcfont\raisebox{-0.35ex}{\Large ∫}\hspace{-0.15em}%
  \raisebox{-0.2ex}{\begin{tabular}{@{}l@{}}\footnotesize #2\\[-1.3ex]\footnotesize #1\end{tabular}}%
  \hspace{0.1em}(#3)\mathrm{d}#4}}
% old key macro kept as an alias so nothing breaks
\renewcommand{\ti}[1]{\gckey{#1}}

% SI units
\usepackage{siunitx}

% setspace
%\usepackage{setspace}
%\onehalfspacing

% For hyperlinks within and out of document
\usepackage{hyperref}
\hypersetup{
    colorlinks = true,
    linkcolor = black,
    urlcolor = blue,
}   

% For captions for pics
\usepackage{capt-of}


% tcolorbox custom environments
% Definition ========================================================
\newtcolorbox[auto counter,number within=section]{definition}[2][drop shadow]{%
enhanced, before skip=5mm, after skip=6mm, % [V3 ADDITION] spacing
colback=blue!5!white,colframe=blue!50!black,fonttitle=\bfseries,
title= Definition~\thetcbcounter. #2,#1}

% Theorem ==========================================================
\newtcolorbox[auto counter,number within=section]{theorem}[2][drop shadow]{%
enhanced, before skip=5mm, after skip=6mm, % [V3 ADDITION] spacing
colback=black!5!white,colframe=black,fonttitle=\bfseries,
title= Theorem~\thetcbcounter. #2,#1}

% Caution box ======================================================
\newtcolorbox{caution}[1][]{enhanced,
before skip=5mm,after skip=6mm, % [V3 ADDITION] was 2mm/3mm
boxrule=1pt,left=11mm,right=2mm,top=3mm,bottom=3mm,
colback=yellow!50,
colframe=yellow!20!black,
sharp corners,rounded corners=southeast,arc is angular,arc=3mm,
underlay={%
\path[fill=tcbcolback!80!black] ([yshift=3mm]interior.south east)--++(-0.4,-0.1)--++(0.1,-0.2);
\path[draw=tcbcolframe,shorten <=-0.05mm,shorten >=-0.05mm] ([yshift=3mm]interior.south east)--++(-0.4,-0.1)--++(0.1,-0.2);
\path[fill=yellow!50!black,draw=none] (interior.south west) rectangle node[white]{\Huge\bfseries !} ([xshift=9mm]interior.north west);
},
drop fuzzy shadow,#1}

% Question box ======================================================
\newtcolorbox{question}[1][]{enhanced,
before skip=5mm,after skip=6mm, % [V3 ADDITION] was 2mm/3mm
boxrule=1pt,left=11mm,right=2mm,top=3mm,bottom=3mm,
colback=blue!20!white,
colframe=blue!20!black,
sharp corners,rounded corners=southeast,arc is angular,arc=3mm,
underlay={%
\path[fill=tcbcolback!80!black] ([yshift=3mm]interior.south east)--++(-0.4,-0.1)--++(0.1,-0.2);
\path[draw=tcbcolframe,shorten <=-0.05mm,shorten >=-0.05mm] ([yshift=3mm]interior.south east)--++(-0.4,-0.1)--++(0.1,-0.2);
\path[fill=blue!50!black,draw=none] (interior.south west) rectangle node[white]{\Huge\bfseries ?} ([xshift=9mm]interior.north west);
},
drop fuzzy shadow,#1}

% Information box ======================================================
\newtcolorbox{info}[1][]{enhanced,
before skip=5mm,after skip=6mm, % [V3 ADDITION] was 2mm/3mm
boxrule=1pt,left=11mm,right=2mm,top=3mm,bottom=3mm,
colback=cyan!20!white,
colframe=cyan!20!black,
sharp corners,rounded corners=southeast,arc is angular,arc=3mm,
underlay={%
\path[fill=tcbcolback!80!black] ([yshift=3mm]interior.south east)--++(-0.4,-0.1)--++(0.1,-0.2);
\path[draw=tcbcolframe,shorten <=-0.05mm,shorten >=-0.05mm] ([yshift=3mm]interior.south east)--++(-0.4,-0.1)--++(0.1,-0.2);
\path[fill=cyan!50!black,draw=none] (interior.south west) rectangle node[white]{\Huge\bfseries i} ([xshift=9mm]interior.north west);
},
drop fuzzy shadow,#1}

% Mistake box ======================================================
\newtcolorbox{mistake}[1][]{enhanced,
before skip=5mm,after skip=6mm, % [V3 ADDITION] was 2mm/3mm
boxrule=1pt,left=11mm,right=2mm,top=3mm,bottom=3mm,
colback=red!20!white,
colframe=red!20!black,
sharp corners,rounded corners=southeast,arc is angular,arc=3mm,
underlay={%
\path[fill=tcbcolback!80!black] ([yshift=3mm]interior.south east)--++(-0.4,-0.1)--++(0.1,-0.2);
\path[draw=tcbcolframe,shorten <=-0.05mm,shorten >=-0.05mm] ([yshift=3mm]interior.south east)--++(-0.4,-0.1)--++(0.1,-0.2);
\path[fill=red!50!black,draw=none] (interior.south west) rectangle node[white]{\Huge\bfseries $\times$} ([xshift=9mm]interior.north west);
},
drop fuzzy shadow,#1}



% [V3 ADDITION] Worked example box ==================================
% Own counter, numbered within the section (Example 11.2.1, 11.2.2, ...).
% For per-subsection numbering (11.2.1.1) change `section' to `subsection'. Blue `question' boxes are kept for thought-provoking questions
% only; worked examples always use this box. Usage:
%   \begin{example} ...question text... \end{example}
%   \solution{...}
\newtcolorbox[auto counter,number within=section]{example}[1][]{%
enhanced, breakable=false, before skip=5mm, after skip=4mm,
colback=teal!5!white, colframe=teal!55!black, boxrule=0.8pt,
left=3mm, right=3mm, top=2mm, bottom=2mm,
sharp corners, rounded corners=northwest, arc=3mm,
fonttitle=\bfseries, coltitle=white,
title=Example~\thetcbcounter, attach boxed title to top left={yshift=-2mm, xshift=3mm},
boxed title style={colback=teal!55!black, sharp corners, rounded corners=northwest, arc=2mm},
drop fuzzy shadow, #1}

% [V3 ADDITION] Spacing between sections, subsections, parts, paragraphs =
% (fncychap's chapter heading is untouched; only \section and below.)
\makeatletter
\renewcommand\section{\@startsection{section}{1}{\z@}%
  {-2ex \@plus -.5ex \@minus -.2ex}{3.2ex \@plus .3ex}%
  {\normalfont\LARGE\bfseries\color{blue!50!black}}}
\renewcommand\subsection{\@startsection{subsection}{2}{\z@}%
  {-1.5ex \@plus -.5ex \@minus -.2ex}{2.6ex \@plus .3ex}%
  {\normalfont\Large\bfseries}}
\renewcommand\subsubsection{\@startsection{subsubsection}{3}{\z@}%
  {-4ex \@plus -1ex \@minus -.5ex}{1.6ex \@plus .2ex}%
  {\normalfont\large\bfseries\color{teal!50!black}}}
\makeatother
\setlength{\parindent}{0pt}
\setlength{\parskip}{0.9ex plus 0.3ex minus 0.2ex}

% Headers and Footers
\usepackage{fancyhdr,lastpage}
% [V3 R2] grey left-margin line on the out-of-syllabus pages (item 13)
\usepackage{eso-pic}
% [AoI item 63] thicker line (2.6pt, was 1.6pt) that starts just BELOW the header rule
% (top 14mm of the page left clear). Flag: copy this line into every chapter's preamble.
\newcommand{\oosmargin}{\AddToShipoutPictureBG{\put(12mm,0){\color{gray!75}\rule{2.6pt}{\dimexpr\paperheight-14mm\relax}}}}
\pagestyle{fancy}
\lhead{Maximus Pung Jun}
% [BUILD COPY CHANGE 2] Header bug fix: \chaptername is the word "Chapter",
% so the original header printed "Chapter 1: Chapter". \chaptermark now stores
% the bare chapter title in \leftmark, and the header uses \leftmark.
\renewcommand{\chaptermark}[1]{\markboth{#1}{}}
\rhead{Chapter \thechapter: \leftmark}
% \cfoot{Page \thepage\  of \pageref{LastPage}} % buggy, does not appear on chapter pages.

% [V3 R2] tick and cross markers for mistake boxes (G3)
\newcommand{\wrong}{\textcolor{red!70!black}{\ensuremath{\boldsymbol{\times}}}}
\newcommand{\cmark}{\textcolor{green!50!black}{\ensuremath{\checkmark}}}

\newcommand{\xmark}{\wrong}   % alias: red cross
% [V3 R2] out-of-syllabus helpers (DE notes item 13)
%   \oosmargin    : grey line in the left margin on every following page
%   \oosmarginoff : stop the margin line
%   \oosnote      : grey banner "Out of syllabus. For your interest only."
%   outofsyllabus : environment = new page + margin line + banner; ends with
%                   \clearpage and switches the margin line off again
\newcommand{\oosmarginoff}{\ClearShipoutPictureBG}
\newcommand{\oosnote}[1][Nothing in this section is examined in H2 Mathematics.]{\par\noindent\colorbox{gray!15}{\parbox{\dimexpr\textwidth-2\fboxsep}{\textbf{Out of syllabus. For your interest only.} #1}}\par}
\newenvironment{outofsyllabus}{\clearpage\oosmargin}{\clearpage\oosmarginoff}

% Notation as in Maximus's notes: italic d, display style everywhere (G1).
% Fractions inside exponents stay small, so never use these macros inside ^{}.
\newcommand{\dydx}{\dfrac{\mathrm{d}y}{\mathrm{d}x}}
\newcommand{\dvdx}{\dfrac{\mathrm{d}v}{\mathrm{d}x}}
\newcommand{\dudx}{\dfrac{\mathrm{d}u}{\mathrm{d}x}}
\newcommand{\ddydx}{\dfrac{\mathrm{d}^{2}y}{\mathrm{d}x^{2}}}
\newcommand{\ddt}[1]{\dfrac{\mathrm{d}#1}{\mathrm{d}t}}
% Small syllabus flag (no longer used in the DE notes, kept for other chapters)
\newcommand{\syllabusflag}[1]{\par\noindent{\footnotesize\color{red!70!black}\textbf{Syllabus note.} #1}\par}


% Teacher's copy toggle ================================================
% [BUILD COPY CHANGE 3] \ifx compares the whole word (the original
% \if\conform test was always false). DE_notes_student.tex predefines
% \conform as HideSolutions.
% A student copy defines \conform as HideSolutions BEFORE loading this preamble.
\providecommand\conform{ShowSolutions}
\newcommand\ShowSolutionsKey{ShowSolutions}
% [V3 ADDITION] In the student copy the hidden solution leaves stretchable
% blank space (\vfill), so the free space on each page is shared out as
% working space below the examples instead of collecting at the bottom.
% [AoI item 10] "Solution." on its own line, so that (a), (b), (c) line up below it.
\newcommand\solution[1]{\ifx\conform\ShowSolutionsKey
  \par\noindent\textbf{Solution.}\par\nobreak\vspace{0.5ex}\noindent #1\par\medskip\else\par\vspace{\workspace}\fi}
% [V3 R2] minimum working space left under each example in the student copy
\newlength\workspace\setlength\workspace{3cm plus 1fill}


\newcommand{\NN}{\mathbb{N}}
\newcommand{\ZZ}{\mathbb{Z}}
\newcommand{\QQ}{\mathbb{Q}}
\newcommand{\RR}{\mathbb{R}}
\newcommand{\CC}{\mathbb{C}}


% house_local.tex: HOUSE COPY of definite_integrals/common/ws_preamble_local.tex ([HOUSE] figstyles renamed)
% WORKSHEET COPY of ../notes/preamble_local.tex (version the worksheet was last built with,
% from ../_backup_pre_limitsum, 2026-10-09 17:31). Local so the worksheet no longer depends
% on ../notes; changes: \input{house_figstyles}; [2026-10-09 follow-up] \mk marks now bold, as in the notes.
% =====================================================================
% preamble_local.tex: extra packages, styles and macros for the
% Definite Integrals chapter ONLY. Loaded after preamble_build.tex.
% Plot styles and the mini 3D macros are in figstyles.tex (shared with the
% slides and worksheets).
% =====================================================================
\usepgfplotslibrary{fillbetween}
\usepackage{adjustbox} % mini 3D pictures (figstyles.tex)
\usetikzlibrary{arrows.meta, decorations.pathreplacing, calc}
\input{house_figstyles}

% Small fractions are allowed only in exponents and in the limits of an
% integral (which are superscripts/subscripts); everywhere else use \dfrac.
\newcommand{\dx}{\,\mathrm{d}x}
\newcommand{\dy}{\,\mathrm{d}y}
\newcommand{\evalat}[3]{\left[\,#1\,\right]_{#2}^{#3}}

% [item 78] right-aligned A-level style marks at the end of a question part
\newcommand{\mk}[1]{{\unskip\nobreak\hfil\penalty50\hskip1.5em\hbox{}\nobreak\hfill\textbf{[#1]}%
  \parfillskip=0pt\finalhyphendemerits=0\par}}
% [item 11] visually emphasised step labels in worked solutions
\newcommand{\step}[1]{\textcolor{teal!55!black}{\textbf{Step #1}}\quad}
% (a), (b), (c) part lists inside questions and solutions
\newenvironment{parts}{\begin{enumerate}[(a)]\setlength{\itemsep}{0.4ex}}{\end{enumerate}}
% SEAB-style solution layout (approved 10.5.8 template): working on the
% left, region sketch on the right (mini 3D pictures go under the sketch)
\newcommand{\solpair}[3][0.56]{\noindent\begin{minipage}[t]{#1\textwidth}\vspace{0pt}#2\end{minipage}\hfill
  \begin{minipage}[t]{\dimexpr0.98\textwidth-#1\textwidth\relax}\vspace{0pt}\centering #3\end{minipage}\par}
% minus-diagram style (drafts/minus_style_template.tex, items 29/31/62)
\pgfplotsset{minus/.style={axis lines=middle, width=5.0cm, height=4.2cm, ticks=none,
  xlabel={$x$}, ylabel={$y$}, every axis x label/.style={at={(ticklabel* cs:1)},anchor=west},
  every axis y label/.style={at={(ticklabel* cs:1)},anchor=south}, clip=false, axis on top,
  axis line style={gray!70}, every axis label/.append style={gray},
  every node/.append style={text=gray!80!black, font=\footnotesize}}}
\tikzset{hlreg/.style={fill=hl!35, draw=hl, thick}}

\ifx\housexm\housexnotes
  \ifdefined\housenotesLM\else% house_fonts.tex: HOUSE COPY of definite_integrals/common/prelim_fonts.tex (unchanged)
% =====================================================================
% prelim_fonts.tex: fonts of the real 9758 prelim papers (Word + MathType; pdffonts on
% ACJC 2026 P1/P2, ACJC 2025 P1, JPJC 2026, RVHS 2026, CJC/DHS/TJC/NYJC/HCI/EJC 2025):
%   text Times New Roman; cover/footer Arial; RI page footer Calibri (here Carlito);
%   maths = MathType: Times New Roman letters/digits (italic variables; upright digits,
%   function names, d, e, f), SymbolMT for the integral, sigma, pi and = + < >,
%   scripts 7 pt on 12 pt (about 58 %), tight operator spacing.
% Candidate C (teacher-approved 13 Oct 2026): the real Times New Roman and Arial
% (ttf-mscorefonts-installer) + URW Standard Symbols PS (fonts-urw-base35, a SymbolMT
% clone with the same Symbol encoding) for the integral, sigma, pi, = + < >;
% TeX Gyre Termes Math only for what MathType draws itself (radicals, stretchy brackets,
% fraction bars, arrows). House style: upright d (\mathrm{d}x), upright e and f.
% Loaded once by di_style / de_core (HW sheets) and prelim_style (Tests, solutions).
% =====================================================================
\ifdefined\prelimfontsloaded\endinput\fi
\def\prelimfontsloaded{1}
\edef\pfrestoreat{\catcode`\noexpand\@=\the\catcode`\@\relax}\makeatletter
% MathType script sizes: 7 pt superscripts on 12 pt text (58 %), 5 pt second level
\DeclareMathSizes{12}{12}{7}{5}
\DeclareMathSizes{11}{11}{6.5}{5}
\DeclareMathSizes{10.95}{10.95}{6.5}{5}
\DeclareMathSizes{10}{10}{6}{5}
\DeclareMathSizes{9}{9}{5.5}{5}
\DeclareMathSizes{14.4}{14.4}{8.5}{6}
\setmainfont{Times New Roman}
\setsansfont{Arial}
\newfontfamily\coverfont{Arial}
\newfontfamily\calibrifont{Carlito}
\usepackage{unicode-math}
\setmathfont{TeX Gyre Termes Math}
\setmathfont[range={up/{latin,Latin,num}}]{Times New Roman}
\setmathfont[range={it/{latin,Latin}}]{Times New Roman Italic}
\setmathfont[range={bfup/{latin,Latin,num}}]{Times New Roman Bold}
\setmathfont[range={"002B,"003D,"003C,"003E}]{Standard Symbols PS}
\newfontfamily\pf@sym{Standard Symbols PS}
\newfontfamily\pf@symi[FakeSlant=0.21]{Standard Symbols PS}
% MathType-like spacing (tighter than TeX around = + -)
\AtBeginDocument{\thickmuskip=2.5mu plus 1mu\medmuskip=1.5mu plus 1mu\thinmuskip=3mu
% Symbol integral and sigma (18 pt at 12 pt text, as MathType), slanted Symbol pi
 \renewcommand{\int}{\mathop{\vcenter{\hbox{\pf@sym\fontsize{\fpeval{1.5*\f@size}}{0pt}\selectfont\char"F2}}}\nolimits}%
 \renewcommand{\sum}{\mathop{\vcenter{\hbox{\pf@sym\fontsize{\fpeval{1.5*\f@size}}{0pt}\selectfont\char"E5}}}\limits}%
 \renewcommand{\pi}{\mathord{\text{\pf@symi\char"70}}}}
\pfrestoreat
\fi
  % =====================================================================
% house_tf.tex: True/False table (standing rules 2026-11-05/06/07):
%   * a T/F set is a NUMBERED QUESTION (no "Check Your Understanding" heading);
%   * table: row label (a),(b),... | Statement (left) | T/F column (right);
%   * row spacing slightly larger than default (extrarowheight 6pt, arraystretch 1.18,
%     as the Vectors Basics Worksheet after the 2026-11-07 change), font size unchanged;
%   * question doc: T/F column blank; solutions doc (\tfanswerstrue): bold T or F,
%     optional short reason under the statement in small grey.
% Usage:
%   \begin{tftable}  \tfrow{statement}{T}{optional reason}  ...  \end{tftable}
% Plain tabular: a question never splits across pages (house rule), so the table is kept whole.
% =====================================================================
\RequirePackage{array}
\makeatletter
\newif\iftfanswers \tfanswersfalse
\newcounter{tfrow}
\newlength{\tf@lab}\setlength{\tf@lab}{0.9cm}
\newlength{\tf@ans}\setlength{\tf@ans}{1.6cm}
\newenvironment{tftable}{\setcounter{tfrow}{0}\par\smallskip
  \setlength{\extrarowheight}{6pt}\renewcommand{\arraystretch}{1.18}%
  \noindent\begin{tabular}{|>{\centering\arraybackslash}p{\tf@lab}|p{\dimexpr\linewidth-\tf@lab-\tf@ans-6\tabcolsep-4\arrayrulewidth\relax}|>{\centering\arraybackslash}p{\tf@ans}|}
  \hline &\textbf{Statement}&\textbf{T/F}\\\hline}
 {\end{tabular}\par\smallskip}
\newcommand{\tfrow}[3]{\stepcounter{tfrow}\textbf{(\alph{tfrow})} & #1%
  \iftfanswers\if\relax\detokenize{#3}\relax\else\newline{\small\color{black!60}#3}\fi\fi
  & \iftfanswers\textbf{#2}\fi\\\hline}
% explicit row label (a set that keeps the Basics Worksheet letters, e.g. (e) omitted)
\newcommand{\tfrowl}[4]{\textbf{#1} & #2%
  \iftfanswers\if\relax\detokenize{#4}\relax\else\newline{\small\color{black!60}#4}\fi\fi
  & \iftfanswers\textbf{#3}\fi\\\hline}
% \tfcompact: allowed exception (Vectors I HW2/HW3) - compressed rows so the whole set
% stays on page 1 of the question sheet (extrarowheight 1.5pt, tabcolsep 4pt, 0.85/1.45 cm)
\newif\iftf@compact
\newcommand{\tfcompact}{\tf@compacttrue\setlength{\tf@lab}{0.85cm}\setlength{\tf@ans}{1.45cm}}
\renewenvironment{tftable}{\setcounter{tfrow}{0}\par\smallskip
  \iftf@compact\setlength{\extrarowheight}{1.5pt}\setlength{\tabcolsep}{4pt}\else\setlength{\extrarowheight}{6pt}\fi
  \renewcommand{\arraystretch}{1.18}%
  \noindent\begin{tabular}{|>{\centering\arraybackslash}p{\tf@lab}|p{\dimexpr\linewidth-\tf@lab-\tf@ans-6\tabcolsep-4\arrayrulewidth\relax}|>{\centering\arraybackslash}p{\tf@ans}|}
  \hline &\textbf{Statement}&\textbf{T/F}\\\hline}
 {\end{tabular}\par\smallskip}
% tflongtable: same table, may break across pages with the header repeated (solutions
% copies whose reasons make the set longer than a page). Needs longtable.
\RequirePackage{longtable}
\newenvironment{tflongtable}{\setcounter{tfrow}{0}\par\smallskip
  \setlength{\extrarowheight}{6pt}\renewcommand{\arraystretch}{1.18}%
  \setlength{\LTleft}{0pt}\setlength{\LTright}{0pt}%
  \begin{longtable}{|>{\centering\arraybackslash}p{\tf@lab}|p{\dimexpr\linewidth-\tf@lab-\tf@ans-6\tabcolsep-4\arrayrulewidth\relax}|>{\centering\arraybackslash}p{\tf@ans}|}
  \hline &\textbf{Statement}&\textbf{T/F}\\\hline\endhead}
 {\end{longtable}\par}
% \putTFlong{n}: True/False question n (body of \defQ{n}: intro, then tftable) printed with
% tflongtable, so a long set may run over a page (house_sol redefines it for solutions).
\newcommand{\putTFlong}[1]{\begingroup\let\tftable\tflongtable\let\endtftable\endtflongtable
  \stepcounter{diq}\expandafter\xdef\csname diLab@#1\endcsname{\thediq}%
  \Qhead{\thediq}{\csname diQtag@#1\endcsname}\csname diQbody@#1\endcsname\par\endgroup}
\makeatother

  \providecommand{\units}[1]{\ \text{units}^{#1}}
  \raggedbottom
\else
  % house_fonts.tex: HOUSE COPY of definite_integrals/common/prelim_fonts.tex (unchanged)
% =====================================================================
% prelim_fonts.tex: fonts of the real 9758 prelim papers (Word + MathType; pdffonts on
% ACJC 2026 P1/P2, ACJC 2025 P1, JPJC 2026, RVHS 2026, CJC/DHS/TJC/NYJC/HCI/EJC 2025):
%   text Times New Roman; cover/footer Arial; RI page footer Calibri (here Carlito);
%   maths = MathType: Times New Roman letters/digits (italic variables; upright digits,
%   function names, d, e, f), SymbolMT for the integral, sigma, pi and = + < >,
%   scripts 7 pt on 12 pt (about 58 %), tight operator spacing.
% Candidate C (teacher-approved 13 Oct 2026): the real Times New Roman and Arial
% (ttf-mscorefonts-installer) + URW Standard Symbols PS (fonts-urw-base35, a SymbolMT
% clone with the same Symbol encoding) for the integral, sigma, pi, = + < >;
% TeX Gyre Termes Math only for what MathType draws itself (radicals, stretchy brackets,
% fraction bars, arrows). House style: upright d (\mathrm{d}x), upright e and f.
% Loaded once by di_style / de_core (HW sheets) and prelim_style (Tests, solutions).
% =====================================================================
\ifdefined\prelimfontsloaded\endinput\fi
\def\prelimfontsloaded{1}
\edef\pfrestoreat{\catcode`\noexpand\@=\the\catcode`\@\relax}\makeatletter
% MathType script sizes: 7 pt superscripts on 12 pt text (58 %), 5 pt second level
\DeclareMathSizes{12}{12}{7}{5}
\DeclareMathSizes{11}{11}{6.5}{5}
\DeclareMathSizes{10.95}{10.95}{6.5}{5}
\DeclareMathSizes{10}{10}{6}{5}
\DeclareMathSizes{9}{9}{5.5}{5}
\DeclareMathSizes{14.4}{14.4}{8.5}{6}
\setmainfont{Times New Roman}
\setsansfont{Arial}
\newfontfamily\coverfont{Arial}
\newfontfamily\calibrifont{Carlito}
\usepackage{unicode-math}
\setmathfont{TeX Gyre Termes Math}
\setmathfont[range={up/{latin,Latin,num}}]{Times New Roman}
\setmathfont[range={it/{latin,Latin}}]{Times New Roman Italic}
\setmathfont[range={bfup/{latin,Latin,num}}]{Times New Roman Bold}
\setmathfont[range={"002B,"003D,"003C,"003E}]{Standard Symbols PS}
\newfontfamily\pf@sym{Standard Symbols PS}
\newfontfamily\pf@symi[FakeSlant=0.21]{Standard Symbols PS}
% MathType-like spacing (tighter than TeX around = + -)
\AtBeginDocument{\thickmuskip=2.5mu plus 1mu\medmuskip=1.5mu plus 1mu\thinmuskip=3mu
% Symbol integral and sigma (18 pt at 12 pt text, as MathType), slanted Symbol pi
 \renewcommand{\int}{\mathop{\vcenter{\hbox{\pf@sym\fontsize{\fpeval{1.5*\f@size}}{0pt}\selectfont\char"F2}}}\nolimits}%
 \renewcommand{\sum}{\mathop{\vcenter{\hbox{\pf@sym\fontsize{\fpeval{1.5*\f@size}}{0pt}\selectfont\char"E5}}}\limits}%
 \renewcommand{\pi}{\mathord{\text{\pf@symi\char"70}}}}
\pfrestoreat

  % house_hw.tex: HOUSE COPY of definite_integrals/common/di_style.tex; [HOUSE] chapter name from \housechapter{...}
% =====================================================================
% di_style.tex: Definite Integrals homework / test layout.
% Loads the house preamble (local copies of the AoI worksheet files:
% ws_preamble_build.tex, ws_preamble_local.tex -> ws_figstyles.tex),
% then adds the homework layout used by HW1-HW3 and the Test.
% Usage (before \begin{document}):  \ditle{Homework 1}
% =====================================================================
% [HOUSE] base, local and fonts are loaded by house.tex
\usepackage{needspace}
\usepackage{zref-savepos,zref-abspage}\makeatletter\zref@addprop{savepos}{abspage}\makeatother
\usepackage{array}

% ---- shortlist toggle ------------------------------------------------
% \ifshortlist true  -> only the questions in SHORTLIST.md (+ challenges)
% A wrapper file sets it with \def\setshortlist{1} before \input-ing the
% main file; the main file can also be edited to say \shortlisttrue.
\newif\ifshortlist
% Task 5: the FINAL homework (= former shortlist) is the default build. HWn.tex -> HWn.pdf.
% A wrapper that sets \def\setfull{1} builds the archive with all original questions.
\ifdefined\setfull\else\shortlisttrue\fi

\makeatletter
% ---- REVIEW build: highlight amended content ---------------------------
% A wrapper sets \def\setreview{1}. In review mode amended text is shaded
% yellow and carries a hover note (PDF tooltip / sticky note); otherwise the
% macros print the content plainly.
%   \chg{note}{inline content}        short inline change (no line breaks), text mode
%   \chgm{note}{math}                 the same inside math mode
%   \begin{chgblock}{note} ... \end{chgblock}   paragraph/display-level change
\newif\ifreview
\ifdefined\setreview \reviewtrue \fi
% hover notes are raw PDF annotations (xdvipdfmx \special; no extra packages).
% Keep notes plain ASCII; balanced parentheses are fine.
\newsavebox{\chg@box}
\newcommand{\chg@ann}[4]{% width, height, depth, note -> invisible markup annotation with a hover note
  \special{pdf:ann width #1 height #2 depth #3 << /Type /Annot /Subtype /Square /Border [0 0 0] /BS << /W 0 >> /C [1 0.85 0] /CA 0.01
    /T (Review) /Subj (Amendment) /Contents (#4) >>}}
\newcommand{\chg@sticky}[1]{\special{pdf:ann width 12pt height 12pt depth 0pt << /Type /Annot /Subtype /Text /Name /Comment
    /C [1 0.85 0] /T (Review) /Subj (Amendment) /Contents (#1) >>}}
\newcommand{\chg}[2]{\ifreview\sbox\chg@box{\colorbox{yellow!60}{#2}}%
  \leavevmode\rlap{\chg@ann{\the\wd\chg@box}{\the\ht\chg@box}{\the\dp\chg@box}{#1}}\usebox\chg@box\else#2\fi}
% \chgm{note}{math}: the same, for use INSIDE math mode (display style kept)
\newcommand{\chgm}[2]{\ifreview\mathchoice{\chg{#1}{$\displaystyle#2$}}{\chg{#1}{$\textstyle#2$}}{\chg{#1}{$\scriptstyle#2$}}{\chg{#1}{$\scriptscriptstyle#2$}}\else#2\fi}
\newtcolorbox{chgbox}[1]{enhanced, breakable, colback=yellow!35, colframe=yellow!35, boxrule=0pt, sharp corners,
  left=2pt, right=14pt, top=1pt, bottom=1pt, before skip=2pt, after skip=2pt,
  overlay first and middle={}, overlay last={}, overlay unbroken and first={\node[anchor=north east, inner sep=1pt] at (frame.north east) {\chg@sticky{#1}\hbox to 12pt{\vrule width 0pt height 12pt}};}}
% \chgnote{note}: just the sticky-note icon (e.g. inside a figure); nothing outside review mode
\newcommand{\chgnote}[1]{\ifreview\chg@sticky{#1}\hbox to 12pt{\vrule width 0pt height 12pt}\fi}
\newenvironment{chgblock}[1]{\ifreview\begin{chgbox}{#1}\fi}{\ifreview\end{chgbox}\fi}

% ---- header ---------------------------------------------------------
\newcommand\di@title{}
\makeatletter
\newcommand{\ditle}[1]{\gdef\di@title{#1}}
\fancypagestyle{di}{\fancyhf{}%
  \lhead{\small \housechap: \di@title}%
  \rhead{\small Compiled by Maximus Pung Jun}%
  \cfoot{\small\thepage}%
  \renewcommand{\headrulewidth}{0.4pt}}
\pagestyle{di}
% title block on page 1, as in his Typst sheets
\newcommand{\dititleblock}[1][]{\thispagestyle{di}%
  {\LARGE\scshape \housechap: \di@title}\par\smallskip
  {Compiled by Maximus Pung Jun}\par\ifx&#1&\else\smallskip{\small\textcolor{black!60}{#1}}\par\fi\medskip}
\makeatother
\newcommand{\disection}[1]{\par\bigskip\noindent{\Large\scshape #1}\par\medskip}

% ---- questions --------------------------------------------------------
% \Qhead{label}{source tag}: question heading; grey [tag] straight after the number, on the left (rule 3B, 9 Nov 2026)
\newcommand{\Qhead}[2]{\par\needspace{8\baselineskip}\bigskip\noindent
  {\large\bfseries #1}\setbox0=\hbox{#2}\ifdim\wd0=0pt\else\enspace{\small\textcolor{black!60}{[#2]}}\fi\par\smallskip}
% part lists with bold labels as he writes them: \pt{(a)}, \pt{(i)}, \pt{(b) (i)}
\newenvironment{qparts}[1][2.4em]{\begin{list}{}{\setlength\leftmargin{#1}%
  \setlength\labelwidth{\dimexpr#1-0.5em\relax}\setlength\labelsep{0.5em}%
  \setlength\itemsep{0.8ex}\setlength\topsep{0.6ex}\setlength\parsep{0.6ex}%
  \renewcommand\makelabel[1]{\textbf{##1}\hfil}}}{\end{list}}
\newcommand{\pt}[1]{\item[#1]}
% thin horizontal label for a challenge question
\newcommand{\challengebanner}{\par\bigskip\noindent\textcolor{black!50}{\rule[0.5ex]{\linewidth}{0.4pt}}\par
  \noindent{\large\scshape Challenge (optional)}\par}
% upright "units"
\newcommand{\units}[1]{\ \text{units}^{#1}}
% END line
\newcommand{\diend}[1][END OF WORKSHEET]{\par\vspace{2\baselineskip}\begin{center}\textbf{#1}\end{center}}

% ---- question registry (full / shortlist from one source) -------------
% \defQ{n}{tag}{body}   question n of the original sheet
% \defA{n}{answer body} its final answer, as printed in his answer section
% \putQ{n} / \putC{n}   print question n as the next (challenge) question
% \putA{n}              print its answer under the same new label
\newcounter{diq}
\makeatletter
\newcommand{\defQ}[3]{\expandafter\gdef\csname diQtag@#1\endcsname{#2}%
  \expandafter\long\expandafter\gdef\csname diQbody@#1\endcsname{#3}}
\newcommand{\defA}[2]{\expandafter\long\expandafter\gdef\csname diA@#1\endcsname{#2}}
\newcommand{\di@tag}[1]{\csname diQtag@#1\endcsname}% source tag only (no original Q number)
% A question never breaks across pages: it is measured and starts on a new page unless it
% fits whole on the rest of this one (only a question taller than a page may break, and then
% only between parts). Page positions of its start/end (zref labels q-s<n>/q-e<n> in the .aux)
% let the checker verify this.
\newsavebox{\di@qmeas}
\newcommand{\di@placeQ}[1]{\par\setbox\di@qmeas=\vbox{#1}%
  \ifdim\dimexpr\ht\di@qmeas+\dp\di@qmeas\relax<0.95\textheight
    \needspace{\dimexpr\ht\di@qmeas+\dp\di@qmeas+\bigskipamount\relax}\else\needspace{14\baselineskip}\fi
  \zsavepos{q-s\thediq}#1\par\zsavepos{q-e\thediq}}
\newcommand{\putQ}[1]{\stepcounter{diq}%
  \expandafter\xdef\csname diLab@#1\endcsname{\thediq}%
  \di@placeQ{\Qhead{\thediq}{\di@tag{#1}}\csname diQbody@#1\endcsname\par}}
\newcommand{\putC}[1]{\stepcounter{diq}%
  \expandafter\xdef\csname diLab@#1\endcsname{\thediq}%
  \expandafter\gdef\csname diChal@#1\endcsname{1}%
  \di@placeQ{\Qhead{\thediq\quad\normalsize\mdseries\scshape Challenge}{\di@tag{#1}}\csname diQbody@#1\endcsname\par}}
\newcommand{\putA}[1]{\par\needspace{4\baselineskip}%
  \begin{qparts}[2.4em]\item[{\csname diLab@#1\endcsname}]\csname diA@#1\endcsname\end{qparts}\medskip}
% ---- worked solutions: \defS{n}{body}, \putQS{n} prints question n + its solution ----
\newcommand{\defS}[2]{\expandafter\long\expandafter\gdef\csname diS@#1\endcsname{#2}}
\newcommand{\disolhead}{\par\medskip\noindent{\bfseries Solution}\par\nobreak\vspace{-0.6ex}\noindent\textcolor{black!35}{\rule{\linewidth}{0.4pt}}\par\nobreak}
\newcommand{\putQS}[1]{\putQ{#1}\disolhead\csname diS@#1\endcsname\par}
% source of a worked solution; visible only in REVIEW builds. #1: his / pub / pip
\newcommand{\solsrc}[2]{\ifreview\par\noindent\def\di@a{#1}\def\di@b{pip}%
  \ifx\di@a\di@b\colorbox{yellow!60}{\parbox{\dimexpr\linewidth-2\fboxsep-14pt}{\small\textbf{Pip-worked:} #2}}\chgnote{Pip-worked solution. No worked solution by the teacher and no published solution was found. Please check it.}%
  \else{\small\textcolor{black!55}{Source of this solution: #2}}\fi\par\smallskip\fi}
\makeatother
% compact sketch for worked solutions: \solfig[axis options]{plot commands}
\newcommand{\solfig}[2][]{\par\begin{center}\begin{tikzpicture}\begin{axis}[qfig, width=6.6cm, height=5cm, #1]#2\end{axis}\end{tikzpicture}\end{center}}
% a sketch added to an existing worked solution (teacher rule), shaded in REVIEW
\newenvironment{addedsketch}{\begin{chgblock}{Sketch added: teacher rule - every area or volume solution has a sketch.}}{\end{chgblock}}
% answer part list (compact)
\newenvironment{aparts}[1][2.4em]{\begin{qparts}[#1]\setlength\itemsep{0.3ex}}{\end{qparts}}

% ---- figure helpers ---------------------------------------------------------
\tikzset{rect/.style={fill=blue!10, draw=blue!55!black, thin}}
\pgfplotsset{qfig/.style={sketch, every node/.append style={font=\scriptsize}}}
% aligned working that starts on the same line as a part label: \pt{(iii)} \begin{topalign} ... \end{topalign}
\newenvironment{topalign}{\hfill$\displaystyle\begin{aligned}[t]}{\end{aligned}$\hfill\null\par\smallskip}
% one display that starts on the same line as a part label
\newcommand{\topdisplay}[1]{\hfill$\displaystyle #1$\hfill\null\par\smallskip}

  % house_solmacros.tex: HOUSE COPY of de-docs/common/de_style.tex minus \input{de_core} and the DE test titles.
% Adds the Ch11 solution helpers: soln/spart tables, sbs/sbsplain side-by-side, \note, \award, \tagm, \dstart.
% =====================================================================
% de_style.tex: Differential Equations (Ch11) homework / test layout.
% = Ch10 di_style (copied as de_core.tex; header text changed) + the DE
% notation and the framed solution boxes of the v1 DE documents
% (soln / spart, used by the Test worked solutions).
% =====================================================================
\providecommand{\e}{\mathrm{e}}
\providecommand{\cosec}{\operatorname{cosec}}
\providecommand{\dt}{\,\mathrm{d}t}
\providecommand{\abs}[1]{\left|#1\right|}
\providecommand{\dd}[2]{\dfrac{\mathrm{d}#1}{\mathrm{d}#2}}
\providecommand{\dvdx}{\dfrac{\mathrm{d}v}{\mathrm{d}x}}
\providecommand{\dudx}{\dfrac{\mathrm{d}u}{\mathrm{d}x}}
\colorlet{notegrey}{black!60}
\makeatletter
\protected\def\note#1{\ifmmode\text{\small\color{notegrey}\itshape #1}\else{\small\color{notegrey}\itshape #1}\fi}
\newcommand{\award}[1]{\unskip\nobreak\hfill\textbf{\small #1}}
\newcommand{\tagm}[1]{\tag*{\textbf{\small #1}}}
\newcommand{\mfAminus}{\[\begin{aligned}&\int\dfrac{1}{a^{2}-x^{2}}\dx\\&\quad=\dfrac{1}{2a}\ln\abs{\dfrac{a+x}{a-x}}+c\end{aligned}\]}
\newcommand{\mfXminus}{\[\begin{aligned}&\int\dfrac{1}{x^{2}-a^{2}}\dx\\&\quad=\dfrac{1}{2a}\ln\abs{\dfrac{x-a}{x+a}}+c\end{aligned}\]}
\newcommand{\sbs@start}{\if@inlabel\def\sbs@shift{\vspace{-\ht\strutbox}}\else
  \ifvmode\else\par\fi\addvspace{0.4ex}\def\sbs@shift{}\fi\noindent}
\newsavebox{\sbs@box}
\newenvironment{sbs}[1]{%
  \sbs@start\begin{minipage}[t]{#1\linewidth}\vspace{0pt}\sbs@shift\raggedright%
  \def\sbsnext{\end{minipage}\hfill\begin{minipage}[t]{\dimexpr0.97\linewidth-#1\linewidth\relax}\vspace{0pt}\sbs@shift\setlength{\fboxsep}{1.5mm}\begin{lrbox}{\sbs@box}\begin{minipage}[t]{\dimexpr\linewidth-2\fboxsep-1pt\relax}\small\raggedright}%
}{\end{minipage}\end{lrbox}\setlength{\fboxrule}{0.5pt}\fbox{\usebox{\sbs@box}}\end{minipage}\par}
\newenvironment{sbsplain}[1]{%
  \sbs@start\begin{minipage}[t]{#1\linewidth}\vspace{0pt}\sbs@shift\raggedright%
  \def\sbsnext{\end{minipage}\hfill\begin{minipage}[t]{\dimexpr0.97\linewidth-#1\linewidth\relax}\vspace{0pt}\sbs@shift\centering}%
}{\end{minipage}\par}
% ---- framed worked-solution box (v1 DE mark-scheme layout, kept for the Test) ----
\newlength{\solcol}\setlength{\solcol}{15mm}
\newcounter{spartno}
\newcommand{\solskips}{\abovedisplayskip=4pt plus 2pt\belowdisplayskip=4pt plus 2pt\abovedisplayshortskip=2pt\belowdisplayshortskip=3pt}
\newtcolorbox{soln}[1][]{enhanced, breakable, sharp corners,
  colback=white, colframe=black, boxrule=0.6pt,
  boxsep=0pt, left=0pt, right=3mm, top=1.5mm, bottom=2mm,
  before skip=1.2ex, after skip=1.5ex,
  before upper={\setcounter{spartno}{0}\parskip=0.5ex\solskips\@fleqntrue\@mathmargin=0pt\raggedright},
  overlay={\draw[line width=0.4pt] ([xshift=\solcol]frame.north west) -- ([xshift=\solcol]frame.south west);},
  #1}
\newenvironment{spart}[1]{%
  \stepcounter{spartno}%
  \ifnum\value{spartno}>1 \par\vspace{0.2ex}\noindent
    \makebox[\linewidth][l]{\rule{\dimexpr\linewidth+3mm\relax}{0.4pt}}\par\nobreak\vspace{0.3ex}\fi
  \begin{list}{}{\leftmargin=\dimexpr\solcol+2.5mm\relax
    \labelwidth=\solcol \labelsep=2.5mm \itemindent=0pt
    \topsep=0.3ex \parsep=0.5ex \rightmargin=0pt \@beginparpenalty=\@M}%
  \item[{\makebox[\solcol][l]{\strut\hspace{1.5mm}\textbf{#1}}}]%
  \ignorespaces}{\end{list}}
\makeatother
% ---- DE sketches: thin light axes (teacher rule), as the Ch10 'sketch' style ----
\pgfplotsset{dsketch/.style={sketch, axis lines=middle, axis line style={thin, black!45, -{Stealth[length=1.6mm]}},
    ticks=none, clip=false, width=6.5cm, height=4.6cm,
    every axis x label/.style={at={(current axis.right of origin)},anchor=west},
    every axis y label/.style={at={(current axis.above origin)},anchor=south},
    every node/.append style={font=\small}},
  dcurve/.style={black, semithick, smooth, samples=120}}
\tikzset{asym/.style={black, dashed, thin}}
\makeatletter
\newcommand{\dstart}{\leavevmode\vadjust{\nobreak\vskip\dimexpr-\baselineskip-\abovedisplayskip+1pt\relax}}
\newcommand{\sdstart}{\leavevmode\vadjust{\nobreak\vskip\dimexpr-\baselineskip-\abovedisplayskip+4pt\relax}}
\makeatother

  % =====================================================================
% house_tf.tex: True/False table (standing rules 2026-11-05/06/07):
%   * a T/F set is a NUMBERED QUESTION (no "Check Your Understanding" heading);
%   * table: row label (a),(b),... | Statement (left) | T/F column (right);
%   * row spacing slightly larger than default (extrarowheight 6pt, arraystretch 1.18,
%     as the Vectors Basics Worksheet after the 2026-11-07 change), font size unchanged;
%   * question doc: T/F column blank; solutions doc (\tfanswerstrue): bold T or F,
%     optional short reason under the statement in small grey.
% Usage:
%   \begin{tftable}  \tfrow{statement}{T}{optional reason}  ...  \end{tftable}
% Plain tabular: a question never splits across pages (house rule), so the table is kept whole.
% =====================================================================
\RequirePackage{array}
\makeatletter
\newif\iftfanswers \tfanswersfalse
\newcounter{tfrow}
\newlength{\tf@lab}\setlength{\tf@lab}{0.9cm}
\newlength{\tf@ans}\setlength{\tf@ans}{1.6cm}
\newenvironment{tftable}{\setcounter{tfrow}{0}\par\smallskip
  \setlength{\extrarowheight}{6pt}\renewcommand{\arraystretch}{1.18}%
  \noindent\begin{tabular}{|>{\centering\arraybackslash}p{\tf@lab}|p{\dimexpr\linewidth-\tf@lab-\tf@ans-6\tabcolsep-4\arrayrulewidth\relax}|>{\centering\arraybackslash}p{\tf@ans}|}
  \hline &\textbf{Statement}&\textbf{T/F}\\\hline}
 {\end{tabular}\par\smallskip}
\newcommand{\tfrow}[3]{\stepcounter{tfrow}\textbf{(\alph{tfrow})} & #1%
  \iftfanswers\if\relax\detokenize{#3}\relax\else\newline{\small\color{black!60}#3}\fi\fi
  & \iftfanswers\textbf{#2}\fi\\\hline}
% explicit row label (a set that keeps the Basics Worksheet letters, e.g. (e) omitted)
\newcommand{\tfrowl}[4]{\textbf{#1} & #2%
  \iftfanswers\if\relax\detokenize{#4}\relax\else\newline{\small\color{black!60}#4}\fi\fi
  & \iftfanswers\textbf{#3}\fi\\\hline}
% \tfcompact: allowed exception (Vectors I HW2/HW3) - compressed rows so the whole set
% stays on page 1 of the question sheet (extrarowheight 1.5pt, tabcolsep 4pt, 0.85/1.45 cm)
\newif\iftf@compact
\newcommand{\tfcompact}{\tf@compacttrue\setlength{\tf@lab}{0.85cm}\setlength{\tf@ans}{1.45cm}}
\renewenvironment{tftable}{\setcounter{tfrow}{0}\par\smallskip
  \iftf@compact\setlength{\extrarowheight}{1.5pt}\setlength{\tabcolsep}{4pt}\else\setlength{\extrarowheight}{6pt}\fi
  \renewcommand{\arraystretch}{1.18}%
  \noindent\begin{tabular}{|>{\centering\arraybackslash}p{\tf@lab}|p{\dimexpr\linewidth-\tf@lab-\tf@ans-6\tabcolsep-4\arrayrulewidth\relax}|>{\centering\arraybackslash}p{\tf@ans}|}
  \hline &\textbf{Statement}&\textbf{T/F}\\\hline}
 {\end{tabular}\par\smallskip}
% tflongtable: same table, may break across pages with the header repeated (solutions
% copies whose reasons make the set longer than a page). Needs longtable.
\RequirePackage{longtable}
\newenvironment{tflongtable}{\setcounter{tfrow}{0}\par\smallskip
  \setlength{\extrarowheight}{6pt}\renewcommand{\arraystretch}{1.18}%
  \setlength{\LTleft}{0pt}\setlength{\LTright}{0pt}%
  \begin{longtable}{|>{\centering\arraybackslash}p{\tf@lab}|p{\dimexpr\linewidth-\tf@lab-\tf@ans-6\tabcolsep-4\arrayrulewidth\relax}|>{\centering\arraybackslash}p{\tf@ans}|}
  \hline &\textbf{Statement}&\textbf{T/F}\\\hline\endhead}
 {\end{longtable}\par}
% \putTFlong{n}: True/False question n (body of \defQ{n}: intro, then tftable) printed with
% tflongtable, so a long set may run over a page (house_sol redefines it for solutions).
\newcommand{\putTFlong}[1]{\begingroup\let\tftable\tflongtable\let\endtftable\endtflongtable
  \stepcounter{diq}\expandafter\xdef\csname diLab@#1\endcsname{\thediq}%
  \Qhead{\thediq}{\csname diQtag@#1\endcsname}\csname diQbody@#1\endcsname\par\endgroup}
\makeatother

  \ifx\housexm\housexsol \def\prelimkeepcolour{1}% house_prelim.tex: HOUSE COPY of definite_integrals/common/prelim_style.tex; [HOUSE] titles from \housechapter
% =====================================================================
% prelim_style.tex: makes the Test look like a Singapore JC H2 Mathematics
% 9758 prelim paper (studied from ACJC/RI 2026 P1/P2 on the box:
% body Times New Roman + Times-style maths, cover in Arial).
% Fonts: prelim_fonts.tex (real Times New Roman/Arial + Times-style maths).
% Shared by Ch10 (di_style) and Ch11 (de_style; sets \prelimtitle/\prelimsubject). Loaded AFTER the chapter style (overrides its layout for the test only).
% =====================================================================
% [HOUSE] fonts loaded by house.tex
\providecommand{\prelimtitle}{\MakeUppercase{\housechap} TEST}
\providecommand{\prelimsubject}{\housechap}
\geometry{a4paper, top=22mm, bottom=24mm, left=25mm, right=25mm, headsep=6mm, footskip=12mm}
\setlength{\parindent}{0pt}
\setlength{\parskip}{0.9ex plus 0.3ex minus 0.2ex}
\linespread{1.3}
\setlength{\emergencystretch}{3em}
\clubpenalty=10000 \widowpenalty=10000 \interlinepenalty=10000
\AtBeginDocument{\fontsize{12}{14.4}\selectfont}
\renewcommand{\normalsize}{\fontsize{12}{14.4}\selectfont}

% ---- figures: plain black/grey like a printed paper ----------------------
\ifdefined\prelimkeepcolour\else % question paper: black/grey figures; solutions keep colour
\pgfplotsset{curve/.style={black, semithick, smooth, samples=150},
  line2/.style={black, semithick, samples=2},
  sketch/.append style={axis line style={thin, black}, every axis label/.append style={black}}}
\tikzset{region/.style={fill=black!15, draw=none}, regionB/.style={fill=black!30, draw=none},
  rect/.style={fill=black!8, draw=black, thin}, ax3/.style={-{Stealth[length=1.5mm]}, thin, black}}
\colorlet{blue}{black}\colorlet{red}{black}\colorlet{orange}{black}
\fi

% ---- header / footer: page number centred top (bold), [Turn over at bottom right
\makeatletter
\providecommand{\prelimfootlabel}{\prelimtitle\ 2026}
\newcommand{\pp@footer}{{\coverfont\scriptsize \prelimfootlabel\hspace{3em}H2 MATHEMATICS 9758}}
\fancypagestyle{prelim}{\fancyhf{}%
  \chead{\bfseries\thepage}%
  \lfoot{\pp@footer}%
  \rfoot{\ifnum\value{page}<\getpagerefnumber{LastPage}\relax\textbf{[Turn over}\fi}%
  \renewcommand{\headrulewidth}{0pt}\renewcommand{\footrulewidth}{0pt}}
\fancypagestyle{prelimcover}{\fancyhf{}%
  \rfoot{\coverfont\textbf{[Turn over}}%
  \renewcommand{\headrulewidth}{0pt}\renewcommand{\footrulewidth}{0pt}}
\pagestyle{prelim}

% ---- marks: [n] regular weight, right-aligned at the end of the part ----------
\renewcommand{\mk}[1]{{\unskip\nobreak\hfil\penalty50\hskip1.5em\hbox{}\nobreak\hfill[#1]%
  \parfillskip=0pt\finalhyphendemerits=0\par}}

% ---- questions: bold number at the left margin, text indented; grey [tag] straight after the number (3B)
\newsavebox{\pp@qmeas}
\newcommand{\pp@qlist}[1]{%
  \begin{list}{}{\setlength\leftmargin{2.6em}\setlength\labelwidth{2.1em}\setlength\labelsep{0.5em}%
     \setlength\topsep{0pt}\setlength\itemsep{0pt}\setlength\parsep{0.9ex}%
     \renewcommand\makelabel[1]{\textbf{##1}\hfil}}%
  \item[\thediq]\@endparpenalty\@M\setbox0=\hbox{\di@tag{#1}}\ifdim\wd0=0pt\else{\small\textcolor{black!60}{[\di@tag{#1}]}}\par\nobreak\fi\csname diQbody@#1\endcsname\par\penalty\@M
  \end{list}}
% a question starts on a new page unless it fits whole on the current one
% (questions taller than a page may break, but only between parts)
\renewcommand{\putQ}[1]{\stepcounter{diq}%
  \expandafter\xdef\csname diLab@#1\endcsname{\thediq}%
  \par\setbox\pp@qmeas=\vbox{\pp@qlist{#1}}%
  \ifdim\dimexpr\ht\pp@qmeas+\dp\pp@qmeas\relax<0.95\textheight
    \needspace{\dimexpr\ht\pp@qmeas+\dp\pp@qmeas\relax}\else\needspace{14\baselineskip}\fi
  \zsavepos{q-s\thediq}\pp@qlist{#1}\par\zsavepos{q-e\thediq}\vspace{5ex plus 2ex minus 1ex}}
% part labels (a)/(i) bold, indented under the question text
\renewenvironment{qparts}[1][2.6em]{\begin{list}{}{\setlength\leftmargin{#1}%
  \setlength\labelwidth{\dimexpr#1-0.5em\relax}\setlength\labelsep{0.5em}%
  \setlength\itemsep{1.4ex}\setlength\topsep{1.2ex}\setlength\parsep{0.8ex}%
  \renewcommand\makelabel[1]{\textbf{##1}\hfil}}}{\end{list}}
\makeatother

% ---- cover page --------------------------------------------------------------
% A-level H2 style cover (9 Nov 2026; reference house_format/ref/alevel_cover_ref.png), no logos/barcodes.
% \prelimcover{total marks}{duration}{date}; heading block from \housechap; override with
% \prelimcoverhead (4 lines) and \prelimcoverline (the line under MATHEMATICS).
\makeatletter
\providecommand{\prelimcoverhead}{MAXIMUS PUNG JUN\\ H2 Mathematics (9758)\\ \housechap\\ Chapter Test}
\providecommand{\prelimcoverline}{\prelimsubject\ Test}
\newcommand{\pp@htwo}{\begin{tikzpicture}[baseline=(b.base)]
  \node[draw, line width=0.9pt, minimum width=19mm, minimum height=15mm, inner sep=0pt] (b)
    {\fontsize{38}{38}\selectfont\special{pdf:literal direct 1 Tr 0.7 w}H2\special{pdf:literal direct 0 Tr}};
  \end{tikzpicture}}
\newcommand{\prelimcover}[3]{% #1 total marks, #2 duration, #3 date line
\thispagestyle{prelimcover}%
{\coverfont\fontsize{11}{13.5}\selectfont\linespread{1}\selectfont\setlength{\parskip}{0pt}%
\vspace*{4mm}
\noindent\begin{tabular}{@{}l@{\hspace{5mm}}l@{}}\raisebox{-0.5\height}{\pp@htwo}&
  \raisebox{-0.5\height}{\begin{minipage}{120mm}\fontsize{12}{14.5}\selectfont\prelimcoverhead\end{minipage}}\end{tabular}\par
\vspace{9mm}
\noindent\rule{\linewidth}{0.8pt}\par\vspace{1.2mm}
\noindent{\fontsize{12}{14.5}\selectfont\bfseries MATHEMATICS\hfill 9758}\par\vspace{1.2mm}
\noindent\prelimcoverline\hfill\textbf{#3}\par\vspace{1.2mm}
\setbox0=\hbox{#2}\ifdim\wd0>0pt \noindent\hfill\textbf{#2}\par\fi
\vspace{3mm}
\noindent You must answer on writing paper.\par\vspace{3mm}
\noindent\begin{tabular}{@{}l@{\hspace{2mm}}l@{}}You will need:&Writing paper\\ &List of formulae and results (MF27)\end{tabular}\par
\vspace{2mm}
\noindent\rule{\linewidth}{0.8pt}\par\vspace{1.5mm}
\noindent{\bfseries INSTRUCTIONS}\par\vspace{0.5mm}
\begin{list}{\textbullet}{\setlength\leftmargin{7mm}\setlength\labelwidth{4mm}\setlength\labelsep{3mm}\setlength\itemsep{0.6mm}\setlength\topsep{0pt}\setlength\parsep{0pt}}
\item Answer \textbf{all} questions.
\item Give non-exact numerical answers correct to 3 significant figures, or 1 decimal place for angles in degrees, unless a different level of accuracy is specified in the question.
\item You should use an approved graphing calculator.
\item Unsupported answers from a graphing calculator are allowed unless a question specifically states otherwise.
\item Where unsupported answers from a graphing calculator are \textbf{not} allowed in a question, you must present the mathematical steps using mathematical notations and not calculator commands.
\item You must show all working clearly.
\end{list}
\vspace{6mm}
\noindent{\bfseries INFORMATION}\par\vspace{0.5mm}
\begin{list}{\textbullet}{\setlength\leftmargin{7mm}\setlength\labelwidth{4mm}\setlength\labelsep{3mm}\setlength\itemsep{0.6mm}\setlength\topsep{0pt}\setlength\parsep{0pt}}
\item The total mark for this paper is #1.
\item The number of marks for each question or part question is shown in brackets [\ ].
\end{list}
\vfill
\noindent\rule{\linewidth}{0.8pt}\par\vspace{1mm}
\centerline{This document has \textbf{\pageref{LastPage}} pages.}\par\vspace{2mm}
}\clearpage}
\makeatother

% ---- solutions header block (page 1 of the worked solutions; same cover style) ----
\newcommand{\prelimsolhead}[1]{% #1 total marks
{\coverfont\fontsize{10}{12}\selectfont\linespread{1}\selectfont
\noindent{\bfseries \prelimtitle}\\[0.2ex]{\bfseries Compiled by Maximus Pung Jun}\\[1.6ex]Higher 2\par\vspace{2ex}
\noindent\rule{\linewidth}{1.2pt}\par\vspace{0.4ex}
\noindent{\Large\bfseries MATHEMATICS}\hfill{\Large\bfseries 9758}\par\vspace{1.2ex}
\noindent \prelimsubject\hfill\textbf{2026}\\
WORKED SOLUTIONS\hfill\textbf{Total marks: #1}\par\vspace{1ex}
\noindent\rule{\linewidth}{0.6pt}\par}\vspace{4ex}}
% house_sol.tex: HOUSE COPY of definite_integrals/common/ri_solstyle.tex (unchanged)
% =====================================================================
% ri_solstyle.tex: worked solutions in the layout of the RI 2026 Year 6 Prelim
% "Solutions with comments" (grail.moe 26612/26613), loaded AFTER prelim_style.
%   * page 1: bold school-style block + bold underlined title; later pages: running
%     title top-left with a rule under it; footer "Page N of M" bottom right (Calibri).
%   * each question: prelim-format question text (bold number, marks right, source tag
%     bottom right), then a bordered table:  Solution | Comments  (bold header row);
%     body rows  label + [marks] (bold, narrow left column) | working | comments,
%     a full-width rule between parts, working left-aligned (fleqn).
% Rows come from the existing sources: \pt in qparts (Ch10 HW), \slabel (Ch10 Test),
% spart (Ch11). Marks under each label are taken from the question's \mk.
% =====================================================================
\makeatletter
\usepackage{lastpage}
\usepackage{adjustbox}
\usepackage{zref-savepos,zref-abspage}
\zref@addprop{savepos}{abspage}
\geometry{a4paper, top=22mm, bottom=24mm, left=22mm, right=22mm, headsep=8mm, footskip=12mm, headheight=24pt}
\providecommand{\risoltitle}{\prelimtitle: Solutions with comments}
\providecommand{\risolschool}{Compiled by Maximus Pung Jun}
\fancypagestyle{risol}{\fancyhf{}%
  \lhead{\fontsize{11}{12}\selectfont 2026 \risoltitle\\[-1.2ex]\rule{\textwidth}{0.6pt}}%
  \rfoot{\calibrifont\fontsize{11}{13}\selectfont Page \textbf{\thepage} of \textbf{\pageref*{LastPage}}}%
  \renewcommand{\headrulewidth}{0pt}\renewcommand{\footrulewidth}{0pt}}
\fancypagestyle{risolfirst}{\fancyhf{}%
  \rfoot{\calibrifont\fontsize{11}{13}\selectfont Page \textbf{\thepage} of \textbf{\pageref*{LastPage}}}%
  \renewcommand{\headrulewidth}{0pt}\renewcommand{\footrulewidth}{0pt}}
\pagestyle{risol}
\newif\ifri@fresh
\newcommand{\risolfront}{\thispagestyle{risolfirst}\global\ri@freshtrue%
  {\noindent\bfseries\linespread{1}\selectfont \risolschool\\ H2 Mathematics (9758)\\ 2026\par}%
  \vspace{4ex}{\noindent\bfseries\underline{2026 \risoltitle}\par}\vspace{3ex}}
% ---- table geometry (RI: label 42pt | working 284pt | comments 142pt of 468pt) ----
\newlength{\ri@L}\setlength{\ri@L}{15mm}
\newlength{\ri@C}\setlength{\ri@C}{40mm}
\newlength{\ri@pad}\setlength{\ri@pad}{1.8mm}
\newif\ifri@sol \newif\ifri@first
\newcount\ri@depth
\def\ri@outer{}
\newtcolorbox{risol}{enhanced, breakable, sharp corners, frame hidden, colback=white,
  boxrule=0pt, boxsep=0pt, left=\dimexpr\ri@L+\ri@pad\relax, right=\dimexpr\ri@C+\ri@pad\relax,
  top=1.2mm, bottom=2.5mm, toptitle=0.6mm, bottomtitle=0.6mm, lefttitle=\ri@pad, righttitle=0pt,
  colbacktitle=white, coltitle=black, fonttitle=\bfseries,
  title={Solution\hfill\makebox[\ri@C][l]{\hspace{\ri@pad}Comments}},
  before skip=2ex, after skip=2ex,
  overlay unbroken and first={\draw[line width=0.5pt] (frame.south west) rectangle (frame.north east);
     \draw[line width=0.5pt] (title.south west) -- (title.south east);
     \draw[line width=0.5pt] ([xshift=\ri@L]interior.north west) -- ([xshift=\ri@L]frame.south west);
     \draw[line width=0.5pt] ([xshift=-\ri@C]frame.north east) -- ([xshift=-\ri@C]frame.south east);},
  overlay middle and last={\draw[line width=0.5pt] (frame.south west) rectangle (frame.north east);
     \draw[line width=0.5pt] ([xshift=\ri@L]frame.north west) -- ([xshift=\ri@L]frame.south west);
     \draw[line width=0.5pt] ([xshift=-\ri@C]frame.north east) -- ([xshift=-\ri@C]frame.south east);},
  after upper={\ri@endrow\protected@write\@auxout{}{\string\ri@settotal{\ri@qid}{\the\ri@sumH}}},
  before upper={\global\ri@sumH=14mm\ri@soltrue\ri@firsttrue\ri@depth=0 \parindent=0pt \parskip=0.6ex
     \predisplaypenalty=100 
     \linespread{1.12}\selectfont\raggedright\@fleqntrue\@mathmargin=0pt
     \abovedisplayskip=5pt plus 2pt\belowdisplayskip=5pt plus 2pt
     \abovedisplayshortskip=3pt plus 1pt\belowdisplayshortskip=4pt plus 2pt
     \ifcsname ri@hp@\ri@qid\endcsname\else\ri@row{}\fi}}
\def\ri@qid{}
% one row: full-width rule (glue leaders, so it vanishes at a page break where the frame
% edge takes its place; not before the first row), then bold label + bold [marks]
% in the left column, on the baseline of the first line of the working
\ExplSyntaxOn
\tl_new:N \g__ri_tmp_tl
\cs_new_protected:Npn \ri_setkey:Nn #1#2
  { \group_begin: \cs_set_eq:NN \chg \use_ii:nn \cs_set_eq:NN \chgm \use_ii:nn
    \cs_set_eq:NN \textbf \use:n \cs_set_eq:NN \mathrm \use:n
    \tl_gset:Nx \g__ri_tmp_tl {#2} \group_end:
    \tl_set_eq:NN #1 \g__ri_tmp_tl \tl_remove_all:Nn #1 {~} }
\cs_set_eq:NN \ri@setkey \ri_setkey:Nn
\ExplSyntaxOff
% ---- page-break rules (positions from the previous run, zref labels in the .aux) ----
% * a row (one part) never breaks across pages: if it does not fit on the rest of the
%   page it starts the next page (forced break between rows); only a row taller than a
%   whole page may break inside;
% * a whole solution that fits on one page is kept together (needspace before the table);
%   a longer one starts where at least its first row fits;
% * a row is at least as tall as its label + [marks].
\newcount\ri@n \newbox\ri@labbox \def\ri@open{0}
\newdimen\ri@ybot \newdimen\ri@ytop \newdimen\ri@h \newdimen\ri@sumH
\AtBeginDocument{\ri@ybot=\dimexpr\paperheight-1in-\voffset-\topmargin-\headheight-\headsep-\textheight\relax
  \ri@ytop=\dimexpr\ri@ybot+\textheight\relax}
\def\ri@pg#1{\zref@extractdefault{#1}{abspage}{-1}}
% \ri@rowh{n}: height of row n in the previous run -> \ri@h (0pt if unknown)
\newcommand{\ri@rowh}[1]{\ri@h=0pt
  \ifnum\ri@pg{ri-s#1}>0 \ifnum\ri@pg{ri-e#1}>0
    \ifnum\ri@pg{ri-s#1}=\ri@pg{ri-e#1}\relax \ri@h=\dimexpr\zposy{ri-s#1}sp-\zposy{ri-e#1}sp\relax
    \else \ri@h=\dimexpr\zposy{ri-s#1}sp-\ri@ybot+\ri@ytop-\zposy{ri-e#1}sp-6mm\relax\fi\fi\fi}
\newcommand{\ri@setfirst}[2]{\expandafter\gdef\csname ri@first@#1\endcsname{#2}}
\newcommand{\ri@setkept}[1]{\expandafter\gdef\csname ri@kept@#1\endcsname{}}
\newcommand{\ri@settotal}[2]{\expandafter\gdef\csname ri@total@#1\endcsname{#2}}
\newcommand{\ri@endrow}{\ifnum\ri@open=1 \par\zsavepos{ri-e\the\ri@n}%
  \ifnum\ri@pg{ri-s\the\ri@n}=\ri@pg{ri-e\the\ri@n}\relax
    \dimen@=\dimexpr\ri@curH sp+1.6mm-\zposy{ri-s\the\ri@n}sp+\zposy{ri-e\the\ri@n}sp\relax
    \ifdim\dimen@>0pt \nobreak\kern\dimen@\penalty\z@\fi\fi
  \ri@rowh{\the\ri@n}\global\advance\ri@sumH\dimexpr\ri@h+3.5mm\relax
  \gdef\ri@open{0}\fi}
\newcommand{\ri@row}[1]{\par\ri@endrow\global\advance\ri@n\@ne\gdef\ri@open{1}\@beginparpenalty\@M
  \ifri@first\global\ri@firstfalse\ri@firstfalse
    \protected@write\@auxout{}{\string\ri@setfirst{\ri@qid}{\the\ri@n}}%
  \else
    \zsavepos{ri-p\the\ri@n}%
    % push the row to the next page if it would not fit on the rest of this one
    \ri@rowh{\the\ri@n}%
    \ifnum\ri@pg{ri-p\the\ri@n}>0 \ifdim\ri@h>0pt
      \ifdim\dimexpr\ri@h+4mm\relax<\dimexpr\textheight-10mm\relax
        \ifdim\dimexpr\ri@h+7mm\relax>\dimexpr\zposy{ri-p\the\ri@n}sp-\ri@ybot\relax
          \penalty-\@M\fi\fi\fi\fi
    \vskip0.6ex\nointerlineskip
    \leaders\hbox{\hskip-\dimexpr\ri@L+\ri@pad\relax\vrule width\dimexpr\linewidth+\ri@L+\ri@C+2\ri@pad\relax height0.5pt depth0pt}\vskip0.5pt
    \nobreak\vskip1.4mm\fi
  \ri@setkey\ri@key{\ri@qid#1}\def\ri@lab{#1}%
  \setbox\ri@labbox=\vtop{\hsize=\dimexpr\ri@L-1.6mm\relax\bfseries\raggedright\parskip=0pt
      \strut\ri@lab\ifcsname ri@mk@\ri@key\endcsname\ifx\ri@lab\@empty\else\\\fi
      [\csname ri@mk@\ri@key\endcsname]\fi\strut\par}%
  \xdef\ri@curH{\number\dimexpr\ht\ri@labbox+\dp\ri@labbox\relax}%
  \nointerlineskip\zsavepos{ri-s\the\ri@n}\nointerlineskip\nobreak
  \vbox to 0pt{\hbox{\hskip-\dimexpr\ri@L+\ri@pad-1.2mm\relax\box\ri@labbox}\vss}\nobreak\prevdepth=\dp\strutbox
  \vskip-\parskip\ignorespaces}
% ---- qparts / \pt: question mode records labels (for the marks); solution mode makes rows
\NewCommandCopy\ri@qparts@begin\qparts \let\ri@qparts@end\endqparts
\renewenvironment{qparts}[1][2.6em]{\advance\ri@depth\@ne
  \ifri@sol\else\ri@qparts@begin[#1]\fi}{\ifri@sol\par\else\ri@qparts@end\fi}
\def\ri@qpartsname{qparts}
\renewcommand{\pt}[1]{\ifnum\ri@depth<2 \gdef\ri@outer{#1}\gdef\ri@cur{#1}\else\expandafter\gdef\expandafter\ri@cur\expandafter{\ri@outer#1}\fi
  \ifri@sol\expandafter\ri@solpt\else\expandafter\ri@qpt\fi{#1}}
\def\ri@qpt#1{\expandafter\gdef\csname ri@hp@\ri@qid\endcsname{}\item[#1]}
\def\ri@solpt#1{\@ifnextchar\begin\ri@ptpeek{\ri@row{\ri@cur}}}
\def\ri@ptpeek\begin#1{\def\ri@tmp{#1}\ifx\ri@tmp\ri@qpartsname
   \expandafter\@firstoftwo\else\expandafter\@secondoftwo\fi
  {\begin{qparts}}{\ri@row{\ri@cur}\begin{#1}}}
% marks: record [n] against the current question/part
\let\ri@mk\mk
\renewcommand{\mk}[1]{\ri@setkey\ri@k{\ri@qid\ri@cur}\ri@setkey\ri@mv{#1}\expandafter\xdef\csname ri@mk@\ri@k\endcsname{\ri@mv}\ri@mk{#1}}
% \slabel (Ch10 Test answers): letter labels set the outer part, roman ones are appended;
% consecutive \slabel calls merge into one row label
\def\ri@romans{(i)(ii)(iii)(iv)(v)(vi)(vii)}
\def\ri@firstpart(#1)#2\@nil{(#1)}
\def\slabel#1{\ifri@sol\expandafter\ri@solslabel\else\expandafter\ri@qslabel\fi{#1}}
\def\ri@qslabel#1{\par\smallskip\noindent\textbf{#1}\par\nobreak}
\def\ri@solslabel#1{\@expandtwoargs\in@{#1}{\ri@romans}%
  \ifin@\expandafter\gdef\expandafter\ri@cur\expandafter{\ri@outer#1}\else\gdef\ri@cur{#1}\xdef\ri@outer{\ri@firstpart#1\@nil}\fi
  \@ifnextchar\slabel{}{\ri@row{\ri@cur}}}
% Ch11 spart{label}: a row; soln: the table
\def\spart#1{\gdef\ri@cur{#1}\ri@row{#1}}\def\endspart{\par}
\def\soln{\@ifnextchar[{\ri@solnopt}{\risol}}\def\ri@solnopt[#1]{\risol}\def\endsoln{\endrisol}
% ---- working that started on the label line is left-aligned in the working column
\AtBeginDocument{\@ifundefined{topalign}{}{%
  \renewenvironment{topalign}{\par\noindent$\displaystyle\begin{aligned}[t]}{\end{aligned}$\par\smallskip}%
  \renewcommand{\topdisplay}[1]{\par\noindent$\displaystyle #1$\par\smallskip}}}
% ---- Ch11 side-by-side (text | sketch): the sketch shrinks to fit its half of the working column
\AtBeginDocument{\@ifundefined{sbsplain}{}{%
  \renewenvironment{sbsplain}[1]{%
    \sbs@start\begin{minipage}[t]{#1\linewidth}\vspace{0pt}\sbs@shift\raggedright%
    \def\sbsnext{\end{minipage}\hfill\begin{minipage}[t]{\dimexpr0.97\linewidth-#1\linewidth\relax}\vspace{0pt}\sbs@shift\centering
      \begin{adjustbox}{max width=\linewidth}}%
  }{\end{adjustbox}\end{minipage}\par}}}
% ---- figures in the working column shrink to fit it (only if wider), colours kept
\newcommand{\rifitfig}[1]{\expandafter\let\csname ri@orig@\string#1\endcsname#1%
  \def#1{\@ifnextchar\bgroup{\ri@fitarg#1}{\ri@fitnoarg#1}}}
\def\ri@fitarg#1#2{\ifri@sol\adjustbox{max width=\linewidth}{\csname ri@orig@\string#1\endcsname{#2}}\else\csname ri@orig@\string#1\endcsname{#2}\fi}
\def\ri@fitnoarg#1{\ifri@sol\adjustbox{max width=\linewidth}{\csname ri@orig@\string#1\endcsname}\else\csname ri@orig@\string#1\endcsname\fi}
\@ifundefined{solfig}{}{\renewcommand{\solfig}[2][]{\par\begin{center}\adjustbox{max width=\linewidth}{\begin{tikzpicture}\begin{axis}[qfig, width=6.6cm, height=5cm, #1]#2\end{axis}\end{tikzpicture}}\end{center}}}
% \ricomment{text}: text in the Comments column of the current row (call right after \pt,
% \slabel or \spart). The row is made at least as tall as the comment.
\newbox\ri@cmtbox
\newcommand{\ricomment}[1]{\par\setbox\ri@cmtbox=\vtop{\hsize=\dimexpr\ri@C-2\ri@pad\relax
    \parskip=0pt\raggedright\small{\color{black!60}\strut#1\strut\par}}%
  \ifdim\dimexpr\ht\ri@cmtbox+\dp\ri@cmtbox\relax>\ri@curH sp\xdef\ri@curH{\number\dimexpr\ht\ri@cmtbox+\dp\ri@cmtbox\relax}\fi
  \nointerlineskip\vbox to 0pt{\hbox to\linewidth{\hskip\dimexpr\linewidth+2\ri@pad\relax\box\ri@cmtbox\hss}\vss}\nobreak\prevdepth=\dp\strutbox
  \vskip-\parskip\ignorespaces}
% \putTFlong{n}: the True/False question n with its answer table allowed to break across
% pages (the body of \defQ{n} must be: intro text, then \begin{tftable}...\end{tftable});
% the solutions copy prints it with tflongtable.
\renewcommand{\putTFlong}[1]{\ri@newq\begingroup\let\tftable\tflongtable\let\endtftable\endtflongtable
  \stepcounter{diq}\expandafter\xdef\csname diLab@#1\endcsname{\thediq}%
  \par\noindent\hbox to 2.6em{\textbf{\thediq}\hfil}\setbox0=\hbox{\di@tag{#1}}\ifdim\wd0=0pt\else{\small\textcolor{black!60}{[\di@tag{#1}]}}\par\noindent\hskip2.6em\fi\ignorespaces\csname diQbody@#1\endcsname\par
  \endgroup\vspace{3ex}}
% ---- question + solution; each question starts on a new page (as RI)
\def\ri@qid{}
\let\ri@putQ\putQ
\renewcommand{\putQ}[1]{\gdef\ri@qid{#1}\global\ri@depth=0 \gdef\ri@cur{}\ri@putQ{#1}}
% \risolQ{n}{body macro name}: question n + its solution body in the table;
% \risolQraw: body already contains its own soln environment (Ch11 Test)
\newcommand{\ri@newq}{\ifri@fresh\global\ri@freshfalse\else\clearpage\fi}
% keep the whole table together if it fits on a page; otherwise make sure at least its
% header and first row fit before starting it
\newcommand{\ri@keep}[1]{\par
  \ifcsname ri@total@#1\endcsname
    % hysteresis: a table kept whole in the previous run stays whole up to textheight-4.5mm,
    % so a total near the limit cannot flip between runs
    \ifdim\csname ri@total@#1\endcsname<\dimexpr\textheight-\ifcsname ri@kept@#1\endcsname4.5\else8\fi mm\relax
      \protected@write\@auxout{}{\string\ri@setkept{#1}}%
      \needspace{\dimexpr\csname ri@total@#1\endcsname+4mm\relax}%
    \else\ifcsname ri@first@#1\endcsname
      \ri@rowh{\csname ri@first@#1\endcsname}%
      \ifdim\ri@h<\dimexpr\textheight-20mm\relax\needspace{\dimexpr\ri@h+18mm\relax}\fi\fi\fi\fi}
\newcommand{\risolQ}[2]{\ri@newq\putQ{#1}\vspace{-3ex}\ri@keep{#1}\begin{risol}\csname #2@#1\endcsname\par\end{risol}}
\newcommand{\risolQraw}[2]{\ri@newq\putQ{#1}\vspace{-3ex}\ri@keep{#1}\csname #2@#1\endcsname\par}
\renewcommand{\putQS}[1]{\risolQ{#1}{diS}}
% end marker only if it fits on the last page (no page holding just the marker)
\renewcommand{\diend}[1][END OF SOLUTIONS]{\par\ifdim\dimexpr\pagegoal-\pagetotal\relax<4\baselineskip\else\vspace{1.5ex}\begin{center}\textbf{#1}\end{center}\fi}
\makeatother
\tfanswerstrue\fi
  \ifx\housexm\housextest % house_prelim.tex: HOUSE COPY of definite_integrals/common/prelim_style.tex; [HOUSE] titles from \housechapter
% =====================================================================
% prelim_style.tex: makes the Test look like a Singapore JC H2 Mathematics
% 9758 prelim paper (studied from ACJC/RI 2026 P1/P2 on the box:
% body Times New Roman + Times-style maths, cover in Arial).
% Fonts: prelim_fonts.tex (real Times New Roman/Arial + Times-style maths).
% Shared by Ch10 (di_style) and Ch11 (de_style; sets \prelimtitle/\prelimsubject). Loaded AFTER the chapter style (overrides its layout for the test only).
% =====================================================================
% [HOUSE] fonts loaded by house.tex
\providecommand{\prelimtitle}{\MakeUppercase{\housechap} TEST}
\providecommand{\prelimsubject}{\housechap}
\geometry{a4paper, top=22mm, bottom=24mm, left=25mm, right=25mm, headsep=6mm, footskip=12mm}
\setlength{\parindent}{0pt}
\setlength{\parskip}{0.9ex plus 0.3ex minus 0.2ex}
\linespread{1.3}
\setlength{\emergencystretch}{3em}
\clubpenalty=10000 \widowpenalty=10000 \interlinepenalty=10000
\AtBeginDocument{\fontsize{12}{14.4}\selectfont}
\renewcommand{\normalsize}{\fontsize{12}{14.4}\selectfont}

% ---- figures: plain black/grey like a printed paper ----------------------
\ifdefined\prelimkeepcolour\else % question paper: black/grey figures; solutions keep colour
\pgfplotsset{curve/.style={black, semithick, smooth, samples=150},
  line2/.style={black, semithick, samples=2},
  sketch/.append style={axis line style={thin, black}, every axis label/.append style={black}}}
\tikzset{region/.style={fill=black!15, draw=none}, regionB/.style={fill=black!30, draw=none},
  rect/.style={fill=black!8, draw=black, thin}, ax3/.style={-{Stealth[length=1.5mm]}, thin, black}}
\colorlet{blue}{black}\colorlet{red}{black}\colorlet{orange}{black}
\fi

% ---- header / footer: page number centred top (bold), [Turn over at bottom right
\makeatletter
\providecommand{\prelimfootlabel}{\prelimtitle\ 2026}
\newcommand{\pp@footer}{{\coverfont\scriptsize \prelimfootlabel\hspace{3em}H2 MATHEMATICS 9758}}
\fancypagestyle{prelim}{\fancyhf{}%
  \chead{\bfseries\thepage}%
  \lfoot{\pp@footer}%
  \rfoot{\ifnum\value{page}<\getpagerefnumber{LastPage}\relax\textbf{[Turn over}\fi}%
  \renewcommand{\headrulewidth}{0pt}\renewcommand{\footrulewidth}{0pt}}
\fancypagestyle{prelimcover}{\fancyhf{}%
  \rfoot{\coverfont\textbf{[Turn over}}%
  \renewcommand{\headrulewidth}{0pt}\renewcommand{\footrulewidth}{0pt}}
\pagestyle{prelim}

% ---- marks: [n] regular weight, right-aligned at the end of the part ----------
\renewcommand{\mk}[1]{{\unskip\nobreak\hfil\penalty50\hskip1.5em\hbox{}\nobreak\hfill[#1]%
  \parfillskip=0pt\finalhyphendemerits=0\par}}

% ---- questions: bold number at the left margin, text indented; grey [tag] straight after the number (3B)
\newsavebox{\pp@qmeas}
\newcommand{\pp@qlist}[1]{%
  \begin{list}{}{\setlength\leftmargin{2.6em}\setlength\labelwidth{2.1em}\setlength\labelsep{0.5em}%
     \setlength\topsep{0pt}\setlength\itemsep{0pt}\setlength\parsep{0.9ex}%
     \renewcommand\makelabel[1]{\textbf{##1}\hfil}}%
  \item[\thediq]\@endparpenalty\@M\setbox0=\hbox{\di@tag{#1}}\ifdim\wd0=0pt\else{\small\textcolor{black!60}{[\di@tag{#1}]}}\par\nobreak\fi\csname diQbody@#1\endcsname\par\penalty\@M
  \end{list}}
% a question starts on a new page unless it fits whole on the current one
% (questions taller than a page may break, but only between parts)
\renewcommand{\putQ}[1]{\stepcounter{diq}%
  \expandafter\xdef\csname diLab@#1\endcsname{\thediq}%
  \par\setbox\pp@qmeas=\vbox{\pp@qlist{#1}}%
  \ifdim\dimexpr\ht\pp@qmeas+\dp\pp@qmeas\relax<0.95\textheight
    \needspace{\dimexpr\ht\pp@qmeas+\dp\pp@qmeas\relax}\else\needspace{14\baselineskip}\fi
  \zsavepos{q-s\thediq}\pp@qlist{#1}\par\zsavepos{q-e\thediq}\vspace{5ex plus 2ex minus 1ex}}
% part labels (a)/(i) bold, indented under the question text
\renewenvironment{qparts}[1][2.6em]{\begin{list}{}{\setlength\leftmargin{#1}%
  \setlength\labelwidth{\dimexpr#1-0.5em\relax}\setlength\labelsep{0.5em}%
  \setlength\itemsep{1.4ex}\setlength\topsep{1.2ex}\setlength\parsep{0.8ex}%
  \renewcommand\makelabel[1]{\textbf{##1}\hfil}}}{\end{list}}
\makeatother

% ---- cover page --------------------------------------------------------------
% A-level H2 style cover (9 Nov 2026; reference house_format/ref/alevel_cover_ref.png), no logos/barcodes.
% \prelimcover{total marks}{duration}{date}; heading block from \housechap; override with
% \prelimcoverhead (4 lines) and \prelimcoverline (the line under MATHEMATICS).
\makeatletter
\providecommand{\prelimcoverhead}{MAXIMUS PUNG JUN\\ H2 Mathematics (9758)\\ \housechap\\ Chapter Test}
\providecommand{\prelimcoverline}{\prelimsubject\ Test}
\newcommand{\pp@htwo}{\begin{tikzpicture}[baseline=(b.base)]
  \node[draw, line width=0.9pt, minimum width=19mm, minimum height=15mm, inner sep=0pt] (b)
    {\fontsize{38}{38}\selectfont\special{pdf:literal direct 1 Tr 0.7 w}H2\special{pdf:literal direct 0 Tr}};
  \end{tikzpicture}}
\newcommand{\prelimcover}[3]{% #1 total marks, #2 duration, #3 date line
\thispagestyle{prelimcover}%
{\coverfont\fontsize{11}{13.5}\selectfont\linespread{1}\selectfont\setlength{\parskip}{0pt}%
\vspace*{4mm}
\noindent\begin{tabular}{@{}l@{\hspace{5mm}}l@{}}\raisebox{-0.5\height}{\pp@htwo}&
  \raisebox{-0.5\height}{\begin{minipage}{120mm}\fontsize{12}{14.5}\selectfont\prelimcoverhead\end{minipage}}\end{tabular}\par
\vspace{9mm}
\noindent\rule{\linewidth}{0.8pt}\par\vspace{1.2mm}
\noindent{\fontsize{12}{14.5}\selectfont\bfseries MATHEMATICS\hfill 9758}\par\vspace{1.2mm}
\noindent\prelimcoverline\hfill\textbf{#3}\par\vspace{1.2mm}
\setbox0=\hbox{#2}\ifdim\wd0>0pt \noindent\hfill\textbf{#2}\par\fi
\vspace{3mm}
\noindent You must answer on writing paper.\par\vspace{3mm}
\noindent\begin{tabular}{@{}l@{\hspace{2mm}}l@{}}You will need:&Writing paper\\ &List of formulae and results (MF27)\end{tabular}\par
\vspace{2mm}
\noindent\rule{\linewidth}{0.8pt}\par\vspace{1.5mm}
\noindent{\bfseries INSTRUCTIONS}\par\vspace{0.5mm}
\begin{list}{\textbullet}{\setlength\leftmargin{7mm}\setlength\labelwidth{4mm}\setlength\labelsep{3mm}\setlength\itemsep{0.6mm}\setlength\topsep{0pt}\setlength\parsep{0pt}}
\item Answer \textbf{all} questions.
\item Give non-exact numerical answers correct to 3 significant figures, or 1 decimal place for angles in degrees, unless a different level of accuracy is specified in the question.
\item You should use an approved graphing calculator.
\item Unsupported answers from a graphing calculator are allowed unless a question specifically states otherwise.
\item Where unsupported answers from a graphing calculator are \textbf{not} allowed in a question, you must present the mathematical steps using mathematical notations and not calculator commands.
\item You must show all working clearly.
\end{list}
\vspace{6mm}
\noindent{\bfseries INFORMATION}\par\vspace{0.5mm}
\begin{list}{\textbullet}{\setlength\leftmargin{7mm}\setlength\labelwidth{4mm}\setlength\labelsep{3mm}\setlength\itemsep{0.6mm}\setlength\topsep{0pt}\setlength\parsep{0pt}}
\item The total mark for this paper is #1.
\item The number of marks for each question or part question is shown in brackets [\ ].
\end{list}
\vfill
\noindent\rule{\linewidth}{0.8pt}\par\vspace{1mm}
\centerline{This document has \textbf{\pageref{LastPage}} pages.}\par\vspace{2mm}
}\clearpage}
\makeatother

% ---- solutions header block (page 1 of the worked solutions; same cover style) ----
\newcommand{\prelimsolhead}[1]{% #1 total marks
{\coverfont\fontsize{10}{12}\selectfont\linespread{1}\selectfont
\noindent{\bfseries \prelimtitle}\\[0.2ex]{\bfseries Compiled by Maximus Pung Jun}\\[1.6ex]Higher 2\par\vspace{2ex}
\noindent\rule{\linewidth}{1.2pt}\par\vspace{0.4ex}
\noindent{\Large\bfseries MATHEMATICS}\hfill{\Large\bfseries 9758}\par\vspace{1.2ex}
\noindent \prelimsubject\hfill\textbf{2026}\\
WORKED SOLUTIONS\hfill\textbf{Total marks: #1}\par\vspace{1ex}
\noindent\rule{\linewidth}{0.6pt}\par}\vspace{4ex}}
\fi
  \ifx\housexm\housextestsol \def\prelimkeepcolour{1}% house_prelim.tex: HOUSE COPY of definite_integrals/common/prelim_style.tex; [HOUSE] titles from \housechapter
% =====================================================================
% prelim_style.tex: makes the Test look like a Singapore JC H2 Mathematics
% 9758 prelim paper (studied from ACJC/RI 2026 P1/P2 on the box:
% body Times New Roman + Times-style maths, cover in Arial).
% Fonts: prelim_fonts.tex (real Times New Roman/Arial + Times-style maths).
% Shared by Ch10 (di_style) and Ch11 (de_style; sets \prelimtitle/\prelimsubject). Loaded AFTER the chapter style (overrides its layout for the test only).
% =====================================================================
% [HOUSE] fonts loaded by house.tex
\providecommand{\prelimtitle}{\MakeUppercase{\housechap} TEST}
\providecommand{\prelimsubject}{\housechap}
\geometry{a4paper, top=22mm, bottom=24mm, left=25mm, right=25mm, headsep=6mm, footskip=12mm}
\setlength{\parindent}{0pt}
\setlength{\parskip}{0.9ex plus 0.3ex minus 0.2ex}
\linespread{1.3}
\setlength{\emergencystretch}{3em}
\clubpenalty=10000 \widowpenalty=10000 \interlinepenalty=10000
\AtBeginDocument{\fontsize{12}{14.4}\selectfont}
\renewcommand{\normalsize}{\fontsize{12}{14.4}\selectfont}

% ---- figures: plain black/grey like a printed paper ----------------------
\ifdefined\prelimkeepcolour\else % question paper: black/grey figures; solutions keep colour
\pgfplotsset{curve/.style={black, semithick, smooth, samples=150},
  line2/.style={black, semithick, samples=2},
  sketch/.append style={axis line style={thin, black}, every axis label/.append style={black}}}
\tikzset{region/.style={fill=black!15, draw=none}, regionB/.style={fill=black!30, draw=none},
  rect/.style={fill=black!8, draw=black, thin}, ax3/.style={-{Stealth[length=1.5mm]}, thin, black}}
\colorlet{blue}{black}\colorlet{red}{black}\colorlet{orange}{black}
\fi

% ---- header / footer: page number centred top (bold), [Turn over at bottom right
\makeatletter
\providecommand{\prelimfootlabel}{\prelimtitle\ 2026}
\newcommand{\pp@footer}{{\coverfont\scriptsize \prelimfootlabel\hspace{3em}H2 MATHEMATICS 9758}}
\fancypagestyle{prelim}{\fancyhf{}%
  \chead{\bfseries\thepage}%
  \lfoot{\pp@footer}%
  \rfoot{\ifnum\value{page}<\getpagerefnumber{LastPage}\relax\textbf{[Turn over}\fi}%
  \renewcommand{\headrulewidth}{0pt}\renewcommand{\footrulewidth}{0pt}}
\fancypagestyle{prelimcover}{\fancyhf{}%
  \rfoot{\coverfont\textbf{[Turn over}}%
  \renewcommand{\headrulewidth}{0pt}\renewcommand{\footrulewidth}{0pt}}
\pagestyle{prelim}

% ---- marks: [n] regular weight, right-aligned at the end of the part ----------
\renewcommand{\mk}[1]{{\unskip\nobreak\hfil\penalty50\hskip1.5em\hbox{}\nobreak\hfill[#1]%
  \parfillskip=0pt\finalhyphendemerits=0\par}}

% ---- questions: bold number at the left margin, text indented; grey [tag] straight after the number (3B)
\newsavebox{\pp@qmeas}
\newcommand{\pp@qlist}[1]{%
  \begin{list}{}{\setlength\leftmargin{2.6em}\setlength\labelwidth{2.1em}\setlength\labelsep{0.5em}%
     \setlength\topsep{0pt}\setlength\itemsep{0pt}\setlength\parsep{0.9ex}%
     \renewcommand\makelabel[1]{\textbf{##1}\hfil}}%
  \item[\thediq]\@endparpenalty\@M\setbox0=\hbox{\di@tag{#1}}\ifdim\wd0=0pt\else{\small\textcolor{black!60}{[\di@tag{#1}]}}\par\nobreak\fi\csname diQbody@#1\endcsname\par\penalty\@M
  \end{list}}
% a question starts on a new page unless it fits whole on the current one
% (questions taller than a page may break, but only between parts)
\renewcommand{\putQ}[1]{\stepcounter{diq}%
  \expandafter\xdef\csname diLab@#1\endcsname{\thediq}%
  \par\setbox\pp@qmeas=\vbox{\pp@qlist{#1}}%
  \ifdim\dimexpr\ht\pp@qmeas+\dp\pp@qmeas\relax<0.95\textheight
    \needspace{\dimexpr\ht\pp@qmeas+\dp\pp@qmeas\relax}\else\needspace{14\baselineskip}\fi
  \zsavepos{q-s\thediq}\pp@qlist{#1}\par\zsavepos{q-e\thediq}\vspace{5ex plus 2ex minus 1ex}}
% part labels (a)/(i) bold, indented under the question text
\renewenvironment{qparts}[1][2.6em]{\begin{list}{}{\setlength\leftmargin{#1}%
  \setlength\labelwidth{\dimexpr#1-0.5em\relax}\setlength\labelsep{0.5em}%
  \setlength\itemsep{1.4ex}\setlength\topsep{1.2ex}\setlength\parsep{0.8ex}%
  \renewcommand\makelabel[1]{\textbf{##1}\hfil}}}{\end{list}}
\makeatother

% ---- cover page --------------------------------------------------------------
% A-level H2 style cover (9 Nov 2026; reference house_format/ref/alevel_cover_ref.png), no logos/barcodes.
% \prelimcover{total marks}{duration}{date}; heading block from \housechap; override with
% \prelimcoverhead (4 lines) and \prelimcoverline (the line under MATHEMATICS).
\makeatletter
\providecommand{\prelimcoverhead}{MAXIMUS PUNG JUN\\ H2 Mathematics (9758)\\ \housechap\\ Chapter Test}
\providecommand{\prelimcoverline}{\prelimsubject\ Test}
\newcommand{\pp@htwo}{\begin{tikzpicture}[baseline=(b.base)]
  \node[draw, line width=0.9pt, minimum width=19mm, minimum height=15mm, inner sep=0pt] (b)
    {\fontsize{38}{38}\selectfont\special{pdf:literal direct 1 Tr 0.7 w}H2\special{pdf:literal direct 0 Tr}};
  \end{tikzpicture}}
\newcommand{\prelimcover}[3]{% #1 total marks, #2 duration, #3 date line
\thispagestyle{prelimcover}%
{\coverfont\fontsize{11}{13.5}\selectfont\linespread{1}\selectfont\setlength{\parskip}{0pt}%
\vspace*{4mm}
\noindent\begin{tabular}{@{}l@{\hspace{5mm}}l@{}}\raisebox{-0.5\height}{\pp@htwo}&
  \raisebox{-0.5\height}{\begin{minipage}{120mm}\fontsize{12}{14.5}\selectfont\prelimcoverhead\end{minipage}}\end{tabular}\par
\vspace{9mm}
\noindent\rule{\linewidth}{0.8pt}\par\vspace{1.2mm}
\noindent{\fontsize{12}{14.5}\selectfont\bfseries MATHEMATICS\hfill 9758}\par\vspace{1.2mm}
\noindent\prelimcoverline\hfill\textbf{#3}\par\vspace{1.2mm}
\setbox0=\hbox{#2}\ifdim\wd0>0pt \noindent\hfill\textbf{#2}\par\fi
\vspace{3mm}
\noindent You must answer on writing paper.\par\vspace{3mm}
\noindent\begin{tabular}{@{}l@{\hspace{2mm}}l@{}}You will need:&Writing paper\\ &List of formulae and results (MF27)\end{tabular}\par
\vspace{2mm}
\noindent\rule{\linewidth}{0.8pt}\par\vspace{1.5mm}
\noindent{\bfseries INSTRUCTIONS}\par\vspace{0.5mm}
\begin{list}{\textbullet}{\setlength\leftmargin{7mm}\setlength\labelwidth{4mm}\setlength\labelsep{3mm}\setlength\itemsep{0.6mm}\setlength\topsep{0pt}\setlength\parsep{0pt}}
\item Answer \textbf{all} questions.
\item Give non-exact numerical answers correct to 3 significant figures, or 1 decimal place for angles in degrees, unless a different level of accuracy is specified in the question.
\item You should use an approved graphing calculator.
\item Unsupported answers from a graphing calculator are allowed unless a question specifically states otherwise.
\item Where unsupported answers from a graphing calculator are \textbf{not} allowed in a question, you must present the mathematical steps using mathematical notations and not calculator commands.
\item You must show all working clearly.
\end{list}
\vspace{6mm}
\noindent{\bfseries INFORMATION}\par\vspace{0.5mm}
\begin{list}{\textbullet}{\setlength\leftmargin{7mm}\setlength\labelwidth{4mm}\setlength\labelsep{3mm}\setlength\itemsep{0.6mm}\setlength\topsep{0pt}\setlength\parsep{0pt}}
\item The total mark for this paper is #1.
\item The number of marks for each question or part question is shown in brackets [\ ].
\end{list}
\vfill
\noindent\rule{\linewidth}{0.8pt}\par\vspace{1mm}
\centerline{This document has \textbf{\pageref{LastPage}} pages.}\par\vspace{2mm}
}\clearpage}
\makeatother

% ---- solutions header block (page 1 of the worked solutions; same cover style) ----
\newcommand{\prelimsolhead}[1]{% #1 total marks
{\coverfont\fontsize{10}{12}\selectfont\linespread{1}\selectfont
\noindent{\bfseries \prelimtitle}\\[0.2ex]{\bfseries Compiled by Maximus Pung Jun}\\[1.6ex]Higher 2\par\vspace{2ex}
\noindent\rule{\linewidth}{1.2pt}\par\vspace{0.4ex}
\noindent{\Large\bfseries MATHEMATICS}\hfill{\Large\bfseries 9758}\par\vspace{1.2ex}
\noindent \prelimsubject\hfill\textbf{2026}\\
WORKED SOLUTIONS\hfill\textbf{Total marks: #1}\par\vspace{1ex}
\noindent\rule{\linewidth}{0.6pt}\par}\vspace{4ex}}
% house_sol.tex: HOUSE COPY of definite_integrals/common/ri_solstyle.tex (unchanged)
% =====================================================================
% ri_solstyle.tex: worked solutions in the layout of the RI 2026 Year 6 Prelim
% "Solutions with comments" (grail.moe 26612/26613), loaded AFTER prelim_style.
%   * page 1: bold school-style block + bold underlined title; later pages: running
%     title top-left with a rule under it; footer "Page N of M" bottom right (Calibri).
%   * each question: prelim-format question text (bold number, marks right, source tag
%     bottom right), then a bordered table:  Solution | Comments  (bold header row);
%     body rows  label + [marks] (bold, narrow left column) | working | comments,
%     a full-width rule between parts, working left-aligned (fleqn).
% Rows come from the existing sources: \pt in qparts (Ch10 HW), \slabel (Ch10 Test),
% spart (Ch11). Marks under each label are taken from the question's \mk.
% =====================================================================
\makeatletter
\usepackage{lastpage}
\usepackage{adjustbox}
\usepackage{zref-savepos,zref-abspage}
\zref@addprop{savepos}{abspage}
\geometry{a4paper, top=22mm, bottom=24mm, left=22mm, right=22mm, headsep=8mm, footskip=12mm, headheight=24pt}
\providecommand{\risoltitle}{\prelimtitle: Solutions with comments}
\providecommand{\risolschool}{Compiled by Maximus Pung Jun}
\fancypagestyle{risol}{\fancyhf{}%
  \lhead{\fontsize{11}{12}\selectfont 2026 \risoltitle\\[-1.2ex]\rule{\textwidth}{0.6pt}}%
  \rfoot{\calibrifont\fontsize{11}{13}\selectfont Page \textbf{\thepage} of \textbf{\pageref*{LastPage}}}%
  \renewcommand{\headrulewidth}{0pt}\renewcommand{\footrulewidth}{0pt}}
\fancypagestyle{risolfirst}{\fancyhf{}%
  \rfoot{\calibrifont\fontsize{11}{13}\selectfont Page \textbf{\thepage} of \textbf{\pageref*{LastPage}}}%
  \renewcommand{\headrulewidth}{0pt}\renewcommand{\footrulewidth}{0pt}}
\pagestyle{risol}
\newif\ifri@fresh
\newcommand{\risolfront}{\thispagestyle{risolfirst}\global\ri@freshtrue%
  {\noindent\bfseries\linespread{1}\selectfont \risolschool\\ H2 Mathematics (9758)\\ 2026\par}%
  \vspace{4ex}{\noindent\bfseries\underline{2026 \risoltitle}\par}\vspace{3ex}}
% ---- table geometry (RI: label 42pt | working 284pt | comments 142pt of 468pt) ----
\newlength{\ri@L}\setlength{\ri@L}{15mm}
\newlength{\ri@C}\setlength{\ri@C}{40mm}
\newlength{\ri@pad}\setlength{\ri@pad}{1.8mm}
\newif\ifri@sol \newif\ifri@first
\newcount\ri@depth
\def\ri@outer{}
\newtcolorbox{risol}{enhanced, breakable, sharp corners, frame hidden, colback=white,
  boxrule=0pt, boxsep=0pt, left=\dimexpr\ri@L+\ri@pad\relax, right=\dimexpr\ri@C+\ri@pad\relax,
  top=1.2mm, bottom=2.5mm, toptitle=0.6mm, bottomtitle=0.6mm, lefttitle=\ri@pad, righttitle=0pt,
  colbacktitle=white, coltitle=black, fonttitle=\bfseries,
  title={Solution\hfill\makebox[\ri@C][l]{\hspace{\ri@pad}Comments}},
  before skip=2ex, after skip=2ex,
  overlay unbroken and first={\draw[line width=0.5pt] (frame.south west) rectangle (frame.north east);
     \draw[line width=0.5pt] (title.south west) -- (title.south east);
     \draw[line width=0.5pt] ([xshift=\ri@L]interior.north west) -- ([xshift=\ri@L]frame.south west);
     \draw[line width=0.5pt] ([xshift=-\ri@C]frame.north east) -- ([xshift=-\ri@C]frame.south east);},
  overlay middle and last={\draw[line width=0.5pt] (frame.south west) rectangle (frame.north east);
     \draw[line width=0.5pt] ([xshift=\ri@L]frame.north west) -- ([xshift=\ri@L]frame.south west);
     \draw[line width=0.5pt] ([xshift=-\ri@C]frame.north east) -- ([xshift=-\ri@C]frame.south east);},
  after upper={\ri@endrow\protected@write\@auxout{}{\string\ri@settotal{\ri@qid}{\the\ri@sumH}}},
  before upper={\global\ri@sumH=14mm\ri@soltrue\ri@firsttrue\ri@depth=0 \parindent=0pt \parskip=0.6ex
     \predisplaypenalty=100 
     \linespread{1.12}\selectfont\raggedright\@fleqntrue\@mathmargin=0pt
     \abovedisplayskip=5pt plus 2pt\belowdisplayskip=5pt plus 2pt
     \abovedisplayshortskip=3pt plus 1pt\belowdisplayshortskip=4pt plus 2pt
     \ifcsname ri@hp@\ri@qid\endcsname\else\ri@row{}\fi}}
\def\ri@qid{}
% one row: full-width rule (glue leaders, so it vanishes at a page break where the frame
% edge takes its place; not before the first row), then bold label + bold [marks]
% in the left column, on the baseline of the first line of the working
\ExplSyntaxOn
\tl_new:N \g__ri_tmp_tl
\cs_new_protected:Npn \ri_setkey:Nn #1#2
  { \group_begin: \cs_set_eq:NN \chg \use_ii:nn \cs_set_eq:NN \chgm \use_ii:nn
    \cs_set_eq:NN \textbf \use:n \cs_set_eq:NN \mathrm \use:n
    \tl_gset:Nx \g__ri_tmp_tl {#2} \group_end:
    \tl_set_eq:NN #1 \g__ri_tmp_tl \tl_remove_all:Nn #1 {~} }
\cs_set_eq:NN \ri@setkey \ri_setkey:Nn
\ExplSyntaxOff
% ---- page-break rules (positions from the previous run, zref labels in the .aux) ----
% * a row (one part) never breaks across pages: if it does not fit on the rest of the
%   page it starts the next page (forced break between rows); only a row taller than a
%   whole page may break inside;
% * a whole solution that fits on one page is kept together (needspace before the table);
%   a longer one starts where at least its first row fits;
% * a row is at least as tall as its label + [marks].
\newcount\ri@n \newbox\ri@labbox \def\ri@open{0}
\newdimen\ri@ybot \newdimen\ri@ytop \newdimen\ri@h \newdimen\ri@sumH
\AtBeginDocument{\ri@ybot=\dimexpr\paperheight-1in-\voffset-\topmargin-\headheight-\headsep-\textheight\relax
  \ri@ytop=\dimexpr\ri@ybot+\textheight\relax}
\def\ri@pg#1{\zref@extractdefault{#1}{abspage}{-1}}
% \ri@rowh{n}: height of row n in the previous run -> \ri@h (0pt if unknown)
\newcommand{\ri@rowh}[1]{\ri@h=0pt
  \ifnum\ri@pg{ri-s#1}>0 \ifnum\ri@pg{ri-e#1}>0
    \ifnum\ri@pg{ri-s#1}=\ri@pg{ri-e#1}\relax \ri@h=\dimexpr\zposy{ri-s#1}sp-\zposy{ri-e#1}sp\relax
    \else \ri@h=\dimexpr\zposy{ri-s#1}sp-\ri@ybot+\ri@ytop-\zposy{ri-e#1}sp-6mm\relax\fi\fi\fi}
\newcommand{\ri@setfirst}[2]{\expandafter\gdef\csname ri@first@#1\endcsname{#2}}
\newcommand{\ri@setkept}[1]{\expandafter\gdef\csname ri@kept@#1\endcsname{}}
\newcommand{\ri@settotal}[2]{\expandafter\gdef\csname ri@total@#1\endcsname{#2}}
\newcommand{\ri@endrow}{\ifnum\ri@open=1 \par\zsavepos{ri-e\the\ri@n}%
  \ifnum\ri@pg{ri-s\the\ri@n}=\ri@pg{ri-e\the\ri@n}\relax
    \dimen@=\dimexpr\ri@curH sp+1.6mm-\zposy{ri-s\the\ri@n}sp+\zposy{ri-e\the\ri@n}sp\relax
    \ifdim\dimen@>0pt \nobreak\kern\dimen@\penalty\z@\fi\fi
  \ri@rowh{\the\ri@n}\global\advance\ri@sumH\dimexpr\ri@h+3.5mm\relax
  \gdef\ri@open{0}\fi}
\newcommand{\ri@row}[1]{\par\ri@endrow\global\advance\ri@n\@ne\gdef\ri@open{1}\@beginparpenalty\@M
  \ifri@first\global\ri@firstfalse\ri@firstfalse
    \protected@write\@auxout{}{\string\ri@setfirst{\ri@qid}{\the\ri@n}}%
  \else
    \zsavepos{ri-p\the\ri@n}%
    % push the row to the next page if it would not fit on the rest of this one
    \ri@rowh{\the\ri@n}%
    \ifnum\ri@pg{ri-p\the\ri@n}>0 \ifdim\ri@h>0pt
      \ifdim\dimexpr\ri@h+4mm\relax<\dimexpr\textheight-10mm\relax
        \ifdim\dimexpr\ri@h+7mm\relax>\dimexpr\zposy{ri-p\the\ri@n}sp-\ri@ybot\relax
          \penalty-\@M\fi\fi\fi\fi
    \vskip0.6ex\nointerlineskip
    \leaders\hbox{\hskip-\dimexpr\ri@L+\ri@pad\relax\vrule width\dimexpr\linewidth+\ri@L+\ri@C+2\ri@pad\relax height0.5pt depth0pt}\vskip0.5pt
    \nobreak\vskip1.4mm\fi
  \ri@setkey\ri@key{\ri@qid#1}\def\ri@lab{#1}%
  \setbox\ri@labbox=\vtop{\hsize=\dimexpr\ri@L-1.6mm\relax\bfseries\raggedright\parskip=0pt
      \strut\ri@lab\ifcsname ri@mk@\ri@key\endcsname\ifx\ri@lab\@empty\else\\\fi
      [\csname ri@mk@\ri@key\endcsname]\fi\strut\par}%
  \xdef\ri@curH{\number\dimexpr\ht\ri@labbox+\dp\ri@labbox\relax}%
  \nointerlineskip\zsavepos{ri-s\the\ri@n}\nointerlineskip\nobreak
  \vbox to 0pt{\hbox{\hskip-\dimexpr\ri@L+\ri@pad-1.2mm\relax\box\ri@labbox}\vss}\nobreak\prevdepth=\dp\strutbox
  \vskip-\parskip\ignorespaces}
% ---- qparts / \pt: question mode records labels (for the marks); solution mode makes rows
\NewCommandCopy\ri@qparts@begin\qparts \let\ri@qparts@end\endqparts
\renewenvironment{qparts}[1][2.6em]{\advance\ri@depth\@ne
  \ifri@sol\else\ri@qparts@begin[#1]\fi}{\ifri@sol\par\else\ri@qparts@end\fi}
\def\ri@qpartsname{qparts}
\renewcommand{\pt}[1]{\ifnum\ri@depth<2 \gdef\ri@outer{#1}\gdef\ri@cur{#1}\else\expandafter\gdef\expandafter\ri@cur\expandafter{\ri@outer#1}\fi
  \ifri@sol\expandafter\ri@solpt\else\expandafter\ri@qpt\fi{#1}}
\def\ri@qpt#1{\expandafter\gdef\csname ri@hp@\ri@qid\endcsname{}\item[#1]}
\def\ri@solpt#1{\@ifnextchar\begin\ri@ptpeek{\ri@row{\ri@cur}}}
\def\ri@ptpeek\begin#1{\def\ri@tmp{#1}\ifx\ri@tmp\ri@qpartsname
   \expandafter\@firstoftwo\else\expandafter\@secondoftwo\fi
  {\begin{qparts}}{\ri@row{\ri@cur}\begin{#1}}}
% marks: record [n] against the current question/part
\let\ri@mk\mk
\renewcommand{\mk}[1]{\ri@setkey\ri@k{\ri@qid\ri@cur}\ri@setkey\ri@mv{#1}\expandafter\xdef\csname ri@mk@\ri@k\endcsname{\ri@mv}\ri@mk{#1}}
% \slabel (Ch10 Test answers): letter labels set the outer part, roman ones are appended;
% consecutive \slabel calls merge into one row label
\def\ri@romans{(i)(ii)(iii)(iv)(v)(vi)(vii)}
\def\ri@firstpart(#1)#2\@nil{(#1)}
\def\slabel#1{\ifri@sol\expandafter\ri@solslabel\else\expandafter\ri@qslabel\fi{#1}}
\def\ri@qslabel#1{\par\smallskip\noindent\textbf{#1}\par\nobreak}
\def\ri@solslabel#1{\@expandtwoargs\in@{#1}{\ri@romans}%
  \ifin@\expandafter\gdef\expandafter\ri@cur\expandafter{\ri@outer#1}\else\gdef\ri@cur{#1}\xdef\ri@outer{\ri@firstpart#1\@nil}\fi
  \@ifnextchar\slabel{}{\ri@row{\ri@cur}}}
% Ch11 spart{label}: a row; soln: the table
\def\spart#1{\gdef\ri@cur{#1}\ri@row{#1}}\def\endspart{\par}
\def\soln{\@ifnextchar[{\ri@solnopt}{\risol}}\def\ri@solnopt[#1]{\risol}\def\endsoln{\endrisol}
% ---- working that started on the label line is left-aligned in the working column
\AtBeginDocument{\@ifundefined{topalign}{}{%
  \renewenvironment{topalign}{\par\noindent$\displaystyle\begin{aligned}[t]}{\end{aligned}$\par\smallskip}%
  \renewcommand{\topdisplay}[1]{\par\noindent$\displaystyle #1$\par\smallskip}}}
% ---- Ch11 side-by-side (text | sketch): the sketch shrinks to fit its half of the working column
\AtBeginDocument{\@ifundefined{sbsplain}{}{%
  \renewenvironment{sbsplain}[1]{%
    \sbs@start\begin{minipage}[t]{#1\linewidth}\vspace{0pt}\sbs@shift\raggedright%
    \def\sbsnext{\end{minipage}\hfill\begin{minipage}[t]{\dimexpr0.97\linewidth-#1\linewidth\relax}\vspace{0pt}\sbs@shift\centering
      \begin{adjustbox}{max width=\linewidth}}%
  }{\end{adjustbox}\end{minipage}\par}}}
% ---- figures in the working column shrink to fit it (only if wider), colours kept
\newcommand{\rifitfig}[1]{\expandafter\let\csname ri@orig@\string#1\endcsname#1%
  \def#1{\@ifnextchar\bgroup{\ri@fitarg#1}{\ri@fitnoarg#1}}}
\def\ri@fitarg#1#2{\ifri@sol\adjustbox{max width=\linewidth}{\csname ri@orig@\string#1\endcsname{#2}}\else\csname ri@orig@\string#1\endcsname{#2}\fi}
\def\ri@fitnoarg#1{\ifri@sol\adjustbox{max width=\linewidth}{\csname ri@orig@\string#1\endcsname}\else\csname ri@orig@\string#1\endcsname\fi}
\@ifundefined{solfig}{}{\renewcommand{\solfig}[2][]{\par\begin{center}\adjustbox{max width=\linewidth}{\begin{tikzpicture}\begin{axis}[qfig, width=6.6cm, height=5cm, #1]#2\end{axis}\end{tikzpicture}}\end{center}}}
% \ricomment{text}: text in the Comments column of the current row (call right after \pt,
% \slabel or \spart). The row is made at least as tall as the comment.
\newbox\ri@cmtbox
\newcommand{\ricomment}[1]{\par\setbox\ri@cmtbox=\vtop{\hsize=\dimexpr\ri@C-2\ri@pad\relax
    \parskip=0pt\raggedright\small{\color{black!60}\strut#1\strut\par}}%
  \ifdim\dimexpr\ht\ri@cmtbox+\dp\ri@cmtbox\relax>\ri@curH sp\xdef\ri@curH{\number\dimexpr\ht\ri@cmtbox+\dp\ri@cmtbox\relax}\fi
  \nointerlineskip\vbox to 0pt{\hbox to\linewidth{\hskip\dimexpr\linewidth+2\ri@pad\relax\box\ri@cmtbox\hss}\vss}\nobreak\prevdepth=\dp\strutbox
  \vskip-\parskip\ignorespaces}
% \putTFlong{n}: the True/False question n with its answer table allowed to break across
% pages (the body of \defQ{n} must be: intro text, then \begin{tftable}...\end{tftable});
% the solutions copy prints it with tflongtable.
\renewcommand{\putTFlong}[1]{\ri@newq\begingroup\let\tftable\tflongtable\let\endtftable\endtflongtable
  \stepcounter{diq}\expandafter\xdef\csname diLab@#1\endcsname{\thediq}%
  \par\noindent\hbox to 2.6em{\textbf{\thediq}\hfil}\setbox0=\hbox{\di@tag{#1}}\ifdim\wd0=0pt\else{\small\textcolor{black!60}{[\di@tag{#1}]}}\par\noindent\hskip2.6em\fi\ignorespaces\csname diQbody@#1\endcsname\par
  \endgroup\vspace{3ex}}
% ---- question + solution; each question starts on a new page (as RI)
\def\ri@qid{}
\let\ri@putQ\putQ
\renewcommand{\putQ}[1]{\gdef\ri@qid{#1}\global\ri@depth=0 \gdef\ri@cur{}\ri@putQ{#1}}
% \risolQ{n}{body macro name}: question n + its solution body in the table;
% \risolQraw: body already contains its own soln environment (Ch11 Test)
\newcommand{\ri@newq}{\ifri@fresh\global\ri@freshfalse\else\clearpage\fi}
% keep the whole table together if it fits on a page; otherwise make sure at least its
% header and first row fit before starting it
\newcommand{\ri@keep}[1]{\par
  \ifcsname ri@total@#1\endcsname
    % hysteresis: a table kept whole in the previous run stays whole up to textheight-4.5mm,
    % so a total near the limit cannot flip between runs
    \ifdim\csname ri@total@#1\endcsname<\dimexpr\textheight-\ifcsname ri@kept@#1\endcsname4.5\else8\fi mm\relax
      \protected@write\@auxout{}{\string\ri@setkept{#1}}%
      \needspace{\dimexpr\csname ri@total@#1\endcsname+4mm\relax}%
    \else\ifcsname ri@first@#1\endcsname
      \ri@rowh{\csname ri@first@#1\endcsname}%
      \ifdim\ri@h<\dimexpr\textheight-20mm\relax\needspace{\dimexpr\ri@h+18mm\relax}\fi\fi\fi\fi}
\newcommand{\risolQ}[2]{\ri@newq\putQ{#1}\vspace{-3ex}\ri@keep{#1}\begin{risol}\csname #2@#1\endcsname\par\end{risol}}
\newcommand{\risolQraw}[2]{\ri@newq\putQ{#1}\vspace{-3ex}\ri@keep{#1}\csname #2@#1\endcsname\par}
\renewcommand{\putQS}[1]{\risolQ{#1}{diS}}
% end marker only if it fits on the last page (no page holding just the marker)
\renewcommand{\diend}[1][END OF SOLUTIONS]{\par\ifdim\dimexpr\pagegoal-\pagetotal\relax<4\baselineskip\else\vspace{1.5ex}\begin{center}\textbf{#1}\end{center}\fi}
\makeatother
\tfanswerstrue\fi
\fi
, set:
%   \def\housemode{notes}    lecture notes: book class, fncychap Lenny chapter page, boxes
%   \def\housemode{hw}       homework / additional / workshop QUESTIONS (+ Answers section)
%   \def\housemode{sol}      worked SOLUTIONS of a hw (RI "Solution | Comments" table)
%   \def\housemode{test}     test question paper (prelim cover, prelim layout)
%   \def\housemode{testsol}  test worked solutions (prelim layout + RI table)
%   \def\housechapterno{10}  (notes: chapter number; optional)
%   \def\housenotesLM{1}     (notes only, optional) keep Latin Modern as in the current
%                            DI/DE notes instead of Times New Roman + TeX Gyre Termes Math
% then \housechapter{Definite Integrals} and \ditle{Homework 1}.
% Document class (put BEFORE % =====================================================================
% house.tex: ONE preamble for every Maximus Pung Jun H2 Math document,
% modelled on Ch10 Definite Integrals + Ch11 Differential Equations (see SPEC.md).
%
% Before % =====================================================================
% house.tex: ONE preamble for every Maximus Pung Jun H2 Math document,
% modelled on Ch10 Definite Integrals + Ch11 Differential Equations (see SPEC.md).
%
% Before \input{house}, set:
%   \def\housemode{notes}    lecture notes: book class, fncychap Lenny chapter page, boxes
%   \def\housemode{hw}       homework / additional / workshop QUESTIONS (+ Answers section)
%   \def\housemode{sol}      worked SOLUTIONS of a hw (RI "Solution | Comments" table)
%   \def\housemode{test}     test question paper (prelim cover, prelim layout)
%   \def\housemode{testsol}  test worked solutions (prelim layout + RI table)
%   \def\housechapterno{10}  (notes: chapter number; optional)
%   \def\housenotesLM{1}     (notes only, optional) keep Latin Modern as in the current
%                            DI/DE notes instead of Times New Roman + TeX Gyre Termes Math
% then \housechapter{Definite Integrals} and \ditle{Homework 1}.
% Document class (put BEFORE \input{house}):
%   notes: \documentclass[a4paper]{book}
%   others: \documentclass[a4paper,11pt,oneside]{book}
% A HW file uses \providecommand\housemode{hw}; its solutions file has NO \documentclass:
%   \def\housemode{sol}\def\diquestionsonly{1}\input{HW1}  (as DI/DE HWn_solutions.tex)
% Compile: TEXINPUTS=.:/workspace/house_format: xelatex (2 runs; solutions until .aux stable)
% =====================================================================
\providecommand\housemode{hw}
\providecommand\housedir{/workspace/house_format}
\newcommand\housechap{}
\newcommand\housechapter[1]{\renewcommand\housechap{#1}}
\def\housexnotes{notes}\def\housexhw{hw}\def\housexsol{sol}\def\housextest{test}\def\housextestsol{testsol}
\edef\housexm{\housemode}
\input{house_base}
\input{house_local}
\ifx\housexm\housexnotes
  \ifdefined\housenotesLM\else\input{house_fonts}\fi
  \input{house_tf}
  \providecommand{\units}[1]{\ \text{units}^{#1}}
  \raggedbottom
\else
  \input{house_fonts}
  \input{house_hw}
  \input{house_solmacros}
  \input{house_tf}
  \ifx\housexm\housexsol \def\prelimkeepcolour{1}\input{house_prelim}\input{house_sol}\tfanswerstrue\fi
  \ifx\housexm\housextest \input{house_prelim}\fi
  \ifx\housexm\housextestsol \def\prelimkeepcolour{1}\input{house_prelim}\input{house_sol}\tfanswerstrue\fi
\fi
, set:
%   \def\housemode{notes}    lecture notes: book class, fncychap Lenny chapter page, boxes
%   \def\housemode{hw}       homework / additional / workshop QUESTIONS (+ Answers section)
%   \def\housemode{sol}      worked SOLUTIONS of a hw (RI "Solution | Comments" table)
%   \def\housemode{test}     test question paper (prelim cover, prelim layout)
%   \def\housemode{testsol}  test worked solutions (prelim layout + RI table)
%   \def\housechapterno{10}  (notes: chapter number; optional)
%   \def\housenotesLM{1}     (notes only, optional) keep Latin Modern as in the current
%                            DI/DE notes instead of Times New Roman + TeX Gyre Termes Math
% then \housechapter{Definite Integrals} and \ditle{Homework 1}.
% Document class (put BEFORE % =====================================================================
% house.tex: ONE preamble for every Maximus Pung Jun H2 Math document,
% modelled on Ch10 Definite Integrals + Ch11 Differential Equations (see SPEC.md).
%
% Before \input{house}, set:
%   \def\housemode{notes}    lecture notes: book class, fncychap Lenny chapter page, boxes
%   \def\housemode{hw}       homework / additional / workshop QUESTIONS (+ Answers section)
%   \def\housemode{sol}      worked SOLUTIONS of a hw (RI "Solution | Comments" table)
%   \def\housemode{test}     test question paper (prelim cover, prelim layout)
%   \def\housemode{testsol}  test worked solutions (prelim layout + RI table)
%   \def\housechapterno{10}  (notes: chapter number; optional)
%   \def\housenotesLM{1}     (notes only, optional) keep Latin Modern as in the current
%                            DI/DE notes instead of Times New Roman + TeX Gyre Termes Math
% then \housechapter{Definite Integrals} and \ditle{Homework 1}.
% Document class (put BEFORE \input{house}):
%   notes: \documentclass[a4paper]{book}
%   others: \documentclass[a4paper,11pt,oneside]{book}
% A HW file uses \providecommand\housemode{hw}; its solutions file has NO \documentclass:
%   \def\housemode{sol}\def\diquestionsonly{1}\input{HW1}  (as DI/DE HWn_solutions.tex)
% Compile: TEXINPUTS=.:/workspace/house_format: xelatex (2 runs; solutions until .aux stable)
% =====================================================================
\providecommand\housemode{hw}
\providecommand\housedir{/workspace/house_format}
\newcommand\housechap{}
\newcommand\housechapter[1]{\renewcommand\housechap{#1}}
\def\housexnotes{notes}\def\housexhw{hw}\def\housexsol{sol}\def\housextest{test}\def\housextestsol{testsol}
\edef\housexm{\housemode}
\input{house_base}
\input{house_local}
\ifx\housexm\housexnotes
  \ifdefined\housenotesLM\else\input{house_fonts}\fi
  \input{house_tf}
  \providecommand{\units}[1]{\ \text{units}^{#1}}
  \raggedbottom
\else
  \input{house_fonts}
  \input{house_hw}
  \input{house_solmacros}
  \input{house_tf}
  \ifx\housexm\housexsol \def\prelimkeepcolour{1}\input{house_prelim}\input{house_sol}\tfanswerstrue\fi
  \ifx\housexm\housextest \input{house_prelim}\fi
  \ifx\housexm\housextestsol \def\prelimkeepcolour{1}\input{house_prelim}\input{house_sol}\tfanswerstrue\fi
\fi
):
%   notes: \documentclass[a4paper]{book}
%   others: \documentclass[a4paper,11pt,oneside]{book}
% A HW file uses \providecommand\housemode{hw}; its solutions file has NO \documentclass:
%   \def\housemode{sol}\def\diquestionsonly{1}\input{HW1}  (as DI/DE HWn_solutions.tex)
% Compile: TEXINPUTS=.:/workspace/house_format: xelatex (2 runs; solutions until .aux stable)
% =====================================================================
\providecommand\housemode{hw}
\providecommand\housedir{/workspace/house_format}
\newcommand\housechap{}
\newcommand\housechapter[1]{\renewcommand\housechap{#1}}
\def\housexnotes{notes}\def\housexhw{hw}\def\housexsol{sol}\def\housextest{test}\def\housextestsol{testsol}
\edef\housexm{\housemode}
% house_base.tex: HOUSE COPY of definite_integrals/common/ws_preamble_build.tex
% (= Maximus's own preamble.tex + the DE v3 / AoI additions). Only change: GC font path. [HOUSE]
% WORKSHEET COPY of ../notes/preamble_build.tex (2026-10-09 version), unchanged.
% preamble_build.tex: copy of Maximus's preamble.tex for the DE v3 build.
% The blocks marked [BUILD COPY CHANGE] are the earlier build fixes.
% The blocks marked [V3 ADDITION] are new in this version (example box,
% spacing). Everything else is unchanged from preamble.tex.

% For margins
\usepackage[a4paper, margin = 20mm]{geometry}

\usepackage[Lenny]{fncychap}
% Rejne, Glenn, Lenny

\usepackage{etoc}

% [V3 ADDITION] cancel: used to show the dx's "cancelling" in Section 11.3
\usepackage{cancel}
\renewcommand{\CancelColor}{\color{red}}

% For enumeration & math symbols
\usepackage{enumerate}
\usepackage{amsmath,amsfonts,amssymb,amsthm,graphicx,lipsum}

% For diagrams, plots and boxes
\usepackage{pgfplots}
\pgfplotsset{compat=1.18}
\usepackage{tikz}
\usepackage[most]{tcolorbox}
\tcbuselibrary{skins}
\usepackage{xcolor}
\pgfplotsset{
    standard/.style={clip = false,
    axis line style = thick,
    trig format=rad,
    axis lines = middle,
    every axis x label/.style={at={(current axis.right of origin)},anchor= west},
    every axis y label/.style={at={(current axis.above origin)},anchor=south}
    }
}

% For GC font
\usepackage{fontspec}
% [BUILD COPY CHANGE 1] The TI-84 key font is loaded only if the file exists
% (checked in ./Fonts/... and ../Fonts/...). Otherwise \ti{} falls back to a
% small boxed sans-serif key. Known letters of the TI font are mapped to their
% key names in the fallback: o = [y=], z = [mode], f = [L3].
\newcommand{\tikey}[1]{\tcbox[on line, enhanced, boxsep=0pt, left=2pt, right=2pt,
  top=1pt, bottom=1pt, boxrule=0.5pt, arc=1.5pt, colback=white,
  colframe=black!70]{\sffamily\footnotesize\strut #1}}
\ExplSyntaxOn
\NewDocumentCommand{\tifallback}{m}{
  \str_case:nnF {#1}
    { {o}{\tikey{y=}} {z}{\tikey{mode}} {f}{\tikey{L3}} }
    { \tikey{#1} }
}
\ExplSyntaxOff
\def\tifontcase{0}
\IfFileExists{Fonts/TI-84 Plus CE Key Fonts/TI84PlusCEKeys.ttf}
  {\def\tifontcase{1}}{}
\IfFileExists{../Fonts/TI-84 Plus CE Key Fonts/TI84PlusCEKeys.ttf}
  {\def\tifontcase{2}}{}
\ifcase\tifontcase\relax
  % 0: font not found, use boxed fallback keys
  \newcommand{\ti}[1]{\tifallback{#1}}
\or
  % 1: original location (as in preamble.tex)
  \newfontfamily\tifont[
    Path = {Fonts/TI-84 Plus CE Key Fonts/},
    UprightFont = {TI84PlusCEKeys.ttf}
  ]{}
  \newcommand{\ti}[1]{{\tifont #1}}
\or
  % 2: one level up (when compiling inside draft-ch1/)
  \newfontfamily\tifont[
    Path = {../Fonts/TI-84 Plus CE Key Fonts/},
    UprightFont = {TI84PlusCEKeys.ttf}
  ]{}
  \newcommand{\ti}[1]{{\tifont #1}}
\fi

% [AoI items 17/18/37] GC DISPLAY FONT. Every GC key, command, menu item and
% screen line goes through \gcfont, so the font can be swapped HERE in one place.
%   \gc{fnInt(}      command / menu item / stored variable in the display font
%   \gckey{2nd}      a calculator key (display font in a small key outline)
%   \gcscreen{...}   a calculator screen (display font on a grey-green LCD panel)
% The font file lives in ../Fonts (shared by notes, slides and worksheet).
\def\gcfontfile{VT323-Regular.ttf}
% [HOUSE] GC font: always the absolute copy in \housedir/Fonts (relative Path= breaks under TEXINPUTS)
\providecommand\housedir{/workspace/house_format}
\IfFileExists{\housedir/Fonts/\gcfontfile}{\newfontfamily\gcfont[Path=\housedir/Fonts/, Scale=1.32]{\gcfontfile}}{%
  \newcommand\gcfont{\ttfamily}}
\newcommand{\gc}[1]{\mbox{\gcfont #1}}
\newcommand{\gckey}[1]{\tcbox[on line, enhanced, boxsep=0pt, left=2pt, right=2pt,
  top=0.6pt, bottom=0.2pt, boxrule=0.5pt, arc=1.5pt, colback=black!4,
  colframe=black!65]{\gcfont\strut #1}}
\definecolor{lcd}{RGB}{214,222,204}
\newcommand{\gcscreen}[1]{\tcbox[on line, enhanced, boxsep=1pt, left=4pt, right=4pt,
  top=2pt, bottom=2pt, boxrule=1.2pt, arc=1pt, colback=lcd, colframe=black!70]{%
  \gcfont\begin{tabular}{@{}l@{}}#1\end{tabular}}}
% integral template as shown on the TI-84 Plus CE (MathPrint): \gcint{0}{2}{e^(-X²)}{X}
\newcommand{\gcint}[4]{{\gcfont\raisebox{-0.35ex}{\Large ∫}\hspace{-0.15em}%
  \raisebox{-0.2ex}{\begin{tabular}{@{}l@{}}\footnotesize #2\\[-1.3ex]\footnotesize #1\end{tabular}}%
  \hspace{0.1em}(#3)\mathrm{d}#4}}
% old key macro kept as an alias so nothing breaks
\renewcommand{\ti}[1]{\gckey{#1}}

% SI units
\usepackage{siunitx}

% setspace
%\usepackage{setspace}
%\onehalfspacing

% For hyperlinks within and out of document
\usepackage{hyperref}
\hypersetup{
    colorlinks = true,
    linkcolor = black,
    urlcolor = blue,
}   

% For captions for pics
\usepackage{capt-of}


% tcolorbox custom environments
% Definition ========================================================
\newtcolorbox[auto counter,number within=section]{definition}[2][drop shadow]{%
enhanced, before skip=5mm, after skip=6mm, % [V3 ADDITION] spacing
colback=blue!5!white,colframe=blue!50!black,fonttitle=\bfseries,
title= Definition~\thetcbcounter. #2,#1}

% Theorem ==========================================================
\newtcolorbox[auto counter,number within=section]{theorem}[2][drop shadow]{%
enhanced, before skip=5mm, after skip=6mm, % [V3 ADDITION] spacing
colback=black!5!white,colframe=black,fonttitle=\bfseries,
title= Theorem~\thetcbcounter. #2,#1}

% Caution box ======================================================
\newtcolorbox{caution}[1][]{enhanced,
before skip=5mm,after skip=6mm, % [V3 ADDITION] was 2mm/3mm
boxrule=1pt,left=11mm,right=2mm,top=3mm,bottom=3mm,
colback=yellow!50,
colframe=yellow!20!black,
sharp corners,rounded corners=southeast,arc is angular,arc=3mm,
underlay={%
\path[fill=tcbcolback!80!black] ([yshift=3mm]interior.south east)--++(-0.4,-0.1)--++(0.1,-0.2);
\path[draw=tcbcolframe,shorten <=-0.05mm,shorten >=-0.05mm] ([yshift=3mm]interior.south east)--++(-0.4,-0.1)--++(0.1,-0.2);
\path[fill=yellow!50!black,draw=none] (interior.south west) rectangle node[white]{\Huge\bfseries !} ([xshift=9mm]interior.north west);
},
drop fuzzy shadow,#1}

% Question box ======================================================
\newtcolorbox{question}[1][]{enhanced,
before skip=5mm,after skip=6mm, % [V3 ADDITION] was 2mm/3mm
boxrule=1pt,left=11mm,right=2mm,top=3mm,bottom=3mm,
colback=blue!20!white,
colframe=blue!20!black,
sharp corners,rounded corners=southeast,arc is angular,arc=3mm,
underlay={%
\path[fill=tcbcolback!80!black] ([yshift=3mm]interior.south east)--++(-0.4,-0.1)--++(0.1,-0.2);
\path[draw=tcbcolframe,shorten <=-0.05mm,shorten >=-0.05mm] ([yshift=3mm]interior.south east)--++(-0.4,-0.1)--++(0.1,-0.2);
\path[fill=blue!50!black,draw=none] (interior.south west) rectangle node[white]{\Huge\bfseries ?} ([xshift=9mm]interior.north west);
},
drop fuzzy shadow,#1}

% Information box ======================================================
\newtcolorbox{info}[1][]{enhanced,
before skip=5mm,after skip=6mm, % [V3 ADDITION] was 2mm/3mm
boxrule=1pt,left=11mm,right=2mm,top=3mm,bottom=3mm,
colback=cyan!20!white,
colframe=cyan!20!black,
sharp corners,rounded corners=southeast,arc is angular,arc=3mm,
underlay={%
\path[fill=tcbcolback!80!black] ([yshift=3mm]interior.south east)--++(-0.4,-0.1)--++(0.1,-0.2);
\path[draw=tcbcolframe,shorten <=-0.05mm,shorten >=-0.05mm] ([yshift=3mm]interior.south east)--++(-0.4,-0.1)--++(0.1,-0.2);
\path[fill=cyan!50!black,draw=none] (interior.south west) rectangle node[white]{\Huge\bfseries i} ([xshift=9mm]interior.north west);
},
drop fuzzy shadow,#1}

% Mistake box ======================================================
\newtcolorbox{mistake}[1][]{enhanced,
before skip=5mm,after skip=6mm, % [V3 ADDITION] was 2mm/3mm
boxrule=1pt,left=11mm,right=2mm,top=3mm,bottom=3mm,
colback=red!20!white,
colframe=red!20!black,
sharp corners,rounded corners=southeast,arc is angular,arc=3mm,
underlay={%
\path[fill=tcbcolback!80!black] ([yshift=3mm]interior.south east)--++(-0.4,-0.1)--++(0.1,-0.2);
\path[draw=tcbcolframe,shorten <=-0.05mm,shorten >=-0.05mm] ([yshift=3mm]interior.south east)--++(-0.4,-0.1)--++(0.1,-0.2);
\path[fill=red!50!black,draw=none] (interior.south west) rectangle node[white]{\Huge\bfseries $\times$} ([xshift=9mm]interior.north west);
},
drop fuzzy shadow,#1}



% [V3 ADDITION] Worked example box ==================================
% Own counter, numbered within the section (Example 11.2.1, 11.2.2, ...).
% For per-subsection numbering (11.2.1.1) change `section' to `subsection'. Blue `question' boxes are kept for thought-provoking questions
% only; worked examples always use this box. Usage:
%   \begin{example} ...question text... \end{example}
%   \solution{...}
\newtcolorbox[auto counter,number within=section]{example}[1][]{%
enhanced, breakable=false, before skip=5mm, after skip=4mm,
colback=teal!5!white, colframe=teal!55!black, boxrule=0.8pt,
left=3mm, right=3mm, top=2mm, bottom=2mm,
sharp corners, rounded corners=northwest, arc=3mm,
fonttitle=\bfseries, coltitle=white,
title=Example~\thetcbcounter, attach boxed title to top left={yshift=-2mm, xshift=3mm},
boxed title style={colback=teal!55!black, sharp corners, rounded corners=northwest, arc=2mm},
drop fuzzy shadow, #1}

% [V3 ADDITION] Spacing between sections, subsections, parts, paragraphs =
% (fncychap's chapter heading is untouched; only \section and below.)
\makeatletter
\renewcommand\section{\@startsection{section}{1}{\z@}%
  {-2ex \@plus -.5ex \@minus -.2ex}{3.2ex \@plus .3ex}%
  {\normalfont\LARGE\bfseries\color{blue!50!black}}}
\renewcommand\subsection{\@startsection{subsection}{2}{\z@}%
  {-1.5ex \@plus -.5ex \@minus -.2ex}{2.6ex \@plus .3ex}%
  {\normalfont\Large\bfseries}}
\renewcommand\subsubsection{\@startsection{subsubsection}{3}{\z@}%
  {-4ex \@plus -1ex \@minus -.5ex}{1.6ex \@plus .2ex}%
  {\normalfont\large\bfseries\color{teal!50!black}}}
\makeatother
\setlength{\parindent}{0pt}
\setlength{\parskip}{0.9ex plus 0.3ex minus 0.2ex}

% Headers and Footers
\usepackage{fancyhdr,lastpage}
% [V3 R2] grey left-margin line on the out-of-syllabus pages (item 13)
\usepackage{eso-pic}
% [AoI item 63] thicker line (2.6pt, was 1.6pt) that starts just BELOW the header rule
% (top 14mm of the page left clear). Flag: copy this line into every chapter's preamble.
\newcommand{\oosmargin}{\AddToShipoutPictureBG{\put(12mm,0){\color{gray!75}\rule{2.6pt}{\dimexpr\paperheight-14mm\relax}}}}
\pagestyle{fancy}
\lhead{Maximus Pung Jun}
% [BUILD COPY CHANGE 2] Header bug fix: \chaptername is the word "Chapter",
% so the original header printed "Chapter 1: Chapter". \chaptermark now stores
% the bare chapter title in \leftmark, and the header uses \leftmark.
\renewcommand{\chaptermark}[1]{\markboth{#1}{}}
\rhead{Chapter \thechapter: \leftmark}
% \cfoot{Page \thepage\  of \pageref{LastPage}} % buggy, does not appear on chapter pages.

% [V3 R2] tick and cross markers for mistake boxes (G3)
\newcommand{\wrong}{\textcolor{red!70!black}{\ensuremath{\boldsymbol{\times}}}}
\newcommand{\cmark}{\textcolor{green!50!black}{\ensuremath{\checkmark}}}

\newcommand{\xmark}{\wrong}   % alias: red cross
% [V3 R2] out-of-syllabus helpers (DE notes item 13)
%   \oosmargin    : grey line in the left margin on every following page
%   \oosmarginoff : stop the margin line
%   \oosnote      : grey banner "Out of syllabus. For your interest only."
%   outofsyllabus : environment = new page + margin line + banner; ends with
%                   \clearpage and switches the margin line off again
\newcommand{\oosmarginoff}{\ClearShipoutPictureBG}
\newcommand{\oosnote}[1][Nothing in this section is examined in H2 Mathematics.]{\par\noindent\colorbox{gray!15}{\parbox{\dimexpr\textwidth-2\fboxsep}{\textbf{Out of syllabus. For your interest only.} #1}}\par}
\newenvironment{outofsyllabus}{\clearpage\oosmargin}{\clearpage\oosmarginoff}

% Notation as in Maximus's notes: italic d, display style everywhere (G1).
% Fractions inside exponents stay small, so never use these macros inside ^{}.
\newcommand{\dydx}{\dfrac{\mathrm{d}y}{\mathrm{d}x}}
\newcommand{\dvdx}{\dfrac{\mathrm{d}v}{\mathrm{d}x}}
\newcommand{\dudx}{\dfrac{\mathrm{d}u}{\mathrm{d}x}}
\newcommand{\ddydx}{\dfrac{\mathrm{d}^{2}y}{\mathrm{d}x^{2}}}
\newcommand{\ddt}[1]{\dfrac{\mathrm{d}#1}{\mathrm{d}t}}
% Small syllabus flag (no longer used in the DE notes, kept for other chapters)
\newcommand{\syllabusflag}[1]{\par\noindent{\footnotesize\color{red!70!black}\textbf{Syllabus note.} #1}\par}


% Teacher's copy toggle ================================================
% [BUILD COPY CHANGE 3] \ifx compares the whole word (the original
% \if\conform test was always false). DE_notes_student.tex predefines
% \conform as HideSolutions.
% A student copy defines \conform as HideSolutions BEFORE loading this preamble.
\providecommand\conform{ShowSolutions}
\newcommand\ShowSolutionsKey{ShowSolutions}
% [V3 ADDITION] In the student copy the hidden solution leaves stretchable
% blank space (\vfill), so the free space on each page is shared out as
% working space below the examples instead of collecting at the bottom.
% [AoI item 10] "Solution." on its own line, so that (a), (b), (c) line up below it.
\newcommand\solution[1]{\ifx\conform\ShowSolutionsKey
  \par\noindent\textbf{Solution.}\par\nobreak\vspace{0.5ex}\noindent #1\par\medskip\else\par\vspace{\workspace}\fi}
% [V3 R2] minimum working space left under each example in the student copy
\newlength\workspace\setlength\workspace{3cm plus 1fill}


\newcommand{\NN}{\mathbb{N}}
\newcommand{\ZZ}{\mathbb{Z}}
\newcommand{\QQ}{\mathbb{Q}}
\newcommand{\RR}{\mathbb{R}}
\newcommand{\CC}{\mathbb{C}}


% house_local.tex: HOUSE COPY of definite_integrals/common/ws_preamble_local.tex ([HOUSE] figstyles renamed)
% WORKSHEET COPY of ../notes/preamble_local.tex (version the worksheet was last built with,
% from ../_backup_pre_limitsum, 2026-10-09 17:31). Local so the worksheet no longer depends
% on ../notes; changes: \input{house_figstyles}; [2026-10-09 follow-up] \mk marks now bold, as in the notes.
% =====================================================================
% preamble_local.tex: extra packages, styles and macros for the
% Definite Integrals chapter ONLY. Loaded after preamble_build.tex.
% Plot styles and the mini 3D macros are in figstyles.tex (shared with the
% slides and worksheets).
% =====================================================================
\usepgfplotslibrary{fillbetween}
\usepackage{adjustbox} % mini 3D pictures (figstyles.tex)
\usetikzlibrary{arrows.meta, decorations.pathreplacing, calc}
\input{house_figstyles}

% Small fractions are allowed only in exponents and in the limits of an
% integral (which are superscripts/subscripts); everywhere else use \dfrac.
\newcommand{\dx}{\,\mathrm{d}x}
\newcommand{\dy}{\,\mathrm{d}y}
\newcommand{\evalat}[3]{\left[\,#1\,\right]_{#2}^{#3}}

% [item 78] right-aligned A-level style marks at the end of a question part
\newcommand{\mk}[1]{{\unskip\nobreak\hfil\penalty50\hskip1.5em\hbox{}\nobreak\hfill\textbf{[#1]}%
  \parfillskip=0pt\finalhyphendemerits=0\par}}
% [item 11] visually emphasised step labels in worked solutions
\newcommand{\step}[1]{\textcolor{teal!55!black}{\textbf{Step #1}}\quad}
% (a), (b), (c) part lists inside questions and solutions
\newenvironment{parts}{\begin{enumerate}[(a)]\setlength{\itemsep}{0.4ex}}{\end{enumerate}}
% SEAB-style solution layout (approved 10.5.8 template): working on the
% left, region sketch on the right (mini 3D pictures go under the sketch)
\newcommand{\solpair}[3][0.56]{\noindent\begin{minipage}[t]{#1\textwidth}\vspace{0pt}#2\end{minipage}\hfill
  \begin{minipage}[t]{\dimexpr0.98\textwidth-#1\textwidth\relax}\vspace{0pt}\centering #3\end{minipage}\par}
% minus-diagram style (drafts/minus_style_template.tex, items 29/31/62)
\pgfplotsset{minus/.style={axis lines=middle, width=5.0cm, height=4.2cm, ticks=none,
  xlabel={$x$}, ylabel={$y$}, every axis x label/.style={at={(ticklabel* cs:1)},anchor=west},
  every axis y label/.style={at={(ticklabel* cs:1)},anchor=south}, clip=false, axis on top,
  axis line style={gray!70}, every axis label/.append style={gray},
  every node/.append style={text=gray!80!black, font=\footnotesize}}}
\tikzset{hlreg/.style={fill=hl!35, draw=hl, thick}}

\ifx\housexm\housexnotes
  \ifdefined\housenotesLM\else% house_fonts.tex: HOUSE COPY of definite_integrals/common/prelim_fonts.tex (unchanged)
% =====================================================================
% prelim_fonts.tex: fonts of the real 9758 prelim papers (Word + MathType; pdffonts on
% ACJC 2026 P1/P2, ACJC 2025 P1, JPJC 2026, RVHS 2026, CJC/DHS/TJC/NYJC/HCI/EJC 2025):
%   text Times New Roman; cover/footer Arial; RI page footer Calibri (here Carlito);
%   maths = MathType: Times New Roman letters/digits (italic variables; upright digits,
%   function names, d, e, f), SymbolMT for the integral, sigma, pi and = + < >,
%   scripts 7 pt on 12 pt (about 58 %), tight operator spacing.
% Candidate C (teacher-approved 13 Oct 2026): the real Times New Roman and Arial
% (ttf-mscorefonts-installer) + URW Standard Symbols PS (fonts-urw-base35, a SymbolMT
% clone with the same Symbol encoding) for the integral, sigma, pi, = + < >;
% TeX Gyre Termes Math only for what MathType draws itself (radicals, stretchy brackets,
% fraction bars, arrows). House style: upright d (\mathrm{d}x), upright e and f.
% Loaded once by di_style / de_core (HW sheets) and prelim_style (Tests, solutions).
% =====================================================================
\ifdefined\prelimfontsloaded\endinput\fi
\def\prelimfontsloaded{1}
\edef\pfrestoreat{\catcode`\noexpand\@=\the\catcode`\@\relax}\makeatletter
% MathType script sizes: 7 pt superscripts on 12 pt text (58 %), 5 pt second level
\DeclareMathSizes{12}{12}{7}{5}
\DeclareMathSizes{11}{11}{6.5}{5}
\DeclareMathSizes{10.95}{10.95}{6.5}{5}
\DeclareMathSizes{10}{10}{6}{5}
\DeclareMathSizes{9}{9}{5.5}{5}
\DeclareMathSizes{14.4}{14.4}{8.5}{6}
\setmainfont{Times New Roman}
\setsansfont{Arial}
\newfontfamily\coverfont{Arial}
\newfontfamily\calibrifont{Carlito}
\usepackage{unicode-math}
\setmathfont{TeX Gyre Termes Math}
\setmathfont[range={up/{latin,Latin,num}}]{Times New Roman}
\setmathfont[range={it/{latin,Latin}}]{Times New Roman Italic}
\setmathfont[range={bfup/{latin,Latin,num}}]{Times New Roman Bold}
\setmathfont[range={"002B,"003D,"003C,"003E}]{Standard Symbols PS}
\newfontfamily\pf@sym{Standard Symbols PS}
\newfontfamily\pf@symi[FakeSlant=0.21]{Standard Symbols PS}
% MathType-like spacing (tighter than TeX around = + -)
\AtBeginDocument{\thickmuskip=2.5mu plus 1mu\medmuskip=1.5mu plus 1mu\thinmuskip=3mu
% Symbol integral and sigma (18 pt at 12 pt text, as MathType), slanted Symbol pi
 \renewcommand{\int}{\mathop{\vcenter{\hbox{\pf@sym\fontsize{\fpeval{1.5*\f@size}}{0pt}\selectfont\char"F2}}}\nolimits}%
 \renewcommand{\sum}{\mathop{\vcenter{\hbox{\pf@sym\fontsize{\fpeval{1.5*\f@size}}{0pt}\selectfont\char"E5}}}\limits}%
 \renewcommand{\pi}{\mathord{\text{\pf@symi\char"70}}}}
\pfrestoreat
\fi
  % =====================================================================
% house_tf.tex: True/False table (standing rules 2026-11-05/06/07):
%   * a T/F set is a NUMBERED QUESTION (no "Check Your Understanding" heading);
%   * table: row label (a),(b),... | Statement (left) | T/F column (right);
%   * row spacing slightly larger than default (extrarowheight 6pt, arraystretch 1.18,
%     as the Vectors Basics Worksheet after the 2026-11-07 change), font size unchanged;
%   * question doc: T/F column blank; solutions doc (\tfanswerstrue): bold T or F,
%     optional short reason under the statement in small grey.
% Usage:
%   \begin{tftable}  \tfrow{statement}{T}{optional reason}  ...  \end{tftable}
% Plain tabular: a question never splits across pages (house rule), so the table is kept whole.
% =====================================================================
\RequirePackage{array}
\makeatletter
\newif\iftfanswers \tfanswersfalse
\newcounter{tfrow}
\newlength{\tf@lab}\setlength{\tf@lab}{0.9cm}
\newlength{\tf@ans}\setlength{\tf@ans}{1.6cm}
\newenvironment{tftable}{\setcounter{tfrow}{0}\par\smallskip
  \setlength{\extrarowheight}{6pt}\renewcommand{\arraystretch}{1.18}%
  \noindent\begin{tabular}{|>{\centering\arraybackslash}p{\tf@lab}|p{\dimexpr\linewidth-\tf@lab-\tf@ans-6\tabcolsep-4\arrayrulewidth\relax}|>{\centering\arraybackslash}p{\tf@ans}|}
  \hline &\textbf{Statement}&\textbf{T/F}\\\hline}
 {\end{tabular}\par\smallskip}
\newcommand{\tfrow}[3]{\stepcounter{tfrow}\textbf{(\alph{tfrow})} & #1%
  \iftfanswers\if\relax\detokenize{#3}\relax\else\newline{\small\color{black!60}#3}\fi\fi
  & \iftfanswers\textbf{#2}\fi\\\hline}
% explicit row label (a set that keeps the Basics Worksheet letters, e.g. (e) omitted)
\newcommand{\tfrowl}[4]{\textbf{#1} & #2%
  \iftfanswers\if\relax\detokenize{#4}\relax\else\newline{\small\color{black!60}#4}\fi\fi
  & \iftfanswers\textbf{#3}\fi\\\hline}
% \tfcompact: allowed exception (Vectors I HW2/HW3) - compressed rows so the whole set
% stays on page 1 of the question sheet (extrarowheight 1.5pt, tabcolsep 4pt, 0.85/1.45 cm)
\newif\iftf@compact
\newcommand{\tfcompact}{\tf@compacttrue\setlength{\tf@lab}{0.85cm}\setlength{\tf@ans}{1.45cm}}
\renewenvironment{tftable}{\setcounter{tfrow}{0}\par\smallskip
  \iftf@compact\setlength{\extrarowheight}{1.5pt}\setlength{\tabcolsep}{4pt}\else\setlength{\extrarowheight}{6pt}\fi
  \renewcommand{\arraystretch}{1.18}%
  \noindent\begin{tabular}{|>{\centering\arraybackslash}p{\tf@lab}|p{\dimexpr\linewidth-\tf@lab-\tf@ans-6\tabcolsep-4\arrayrulewidth\relax}|>{\centering\arraybackslash}p{\tf@ans}|}
  \hline &\textbf{Statement}&\textbf{T/F}\\\hline}
 {\end{tabular}\par\smallskip}
% tflongtable: same table, may break across pages with the header repeated (solutions
% copies whose reasons make the set longer than a page). Needs longtable.
\RequirePackage{longtable}
\newenvironment{tflongtable}{\setcounter{tfrow}{0}\par\smallskip
  \setlength{\extrarowheight}{6pt}\renewcommand{\arraystretch}{1.18}%
  \setlength{\LTleft}{0pt}\setlength{\LTright}{0pt}%
  \begin{longtable}{|>{\centering\arraybackslash}p{\tf@lab}|p{\dimexpr\linewidth-\tf@lab-\tf@ans-6\tabcolsep-4\arrayrulewidth\relax}|>{\centering\arraybackslash}p{\tf@ans}|}
  \hline &\textbf{Statement}&\textbf{T/F}\\\hline\endhead}
 {\end{longtable}\par}
% \putTFlong{n}: True/False question n (body of \defQ{n}: intro, then tftable) printed with
% tflongtable, so a long set may run over a page (house_sol redefines it for solutions).
\newcommand{\putTFlong}[1]{\begingroup\let\tftable\tflongtable\let\endtftable\endtflongtable
  \stepcounter{diq}\expandafter\xdef\csname diLab@#1\endcsname{\thediq}%
  \Qhead{\thediq}{\csname diQtag@#1\endcsname}\csname diQbody@#1\endcsname\par\endgroup}
\makeatother

  \providecommand{\units}[1]{\ \text{units}^{#1}}
  \raggedbottom
\else
  % house_fonts.tex: HOUSE COPY of definite_integrals/common/prelim_fonts.tex (unchanged)
% =====================================================================
% prelim_fonts.tex: fonts of the real 9758 prelim papers (Word + MathType; pdffonts on
% ACJC 2026 P1/P2, ACJC 2025 P1, JPJC 2026, RVHS 2026, CJC/DHS/TJC/NYJC/HCI/EJC 2025):
%   text Times New Roman; cover/footer Arial; RI page footer Calibri (here Carlito);
%   maths = MathType: Times New Roman letters/digits (italic variables; upright digits,
%   function names, d, e, f), SymbolMT for the integral, sigma, pi and = + < >,
%   scripts 7 pt on 12 pt (about 58 %), tight operator spacing.
% Candidate C (teacher-approved 13 Oct 2026): the real Times New Roman and Arial
% (ttf-mscorefonts-installer) + URW Standard Symbols PS (fonts-urw-base35, a SymbolMT
% clone with the same Symbol encoding) for the integral, sigma, pi, = + < >;
% TeX Gyre Termes Math only for what MathType draws itself (radicals, stretchy brackets,
% fraction bars, arrows). House style: upright d (\mathrm{d}x), upright e and f.
% Loaded once by di_style / de_core (HW sheets) and prelim_style (Tests, solutions).
% =====================================================================
\ifdefined\prelimfontsloaded\endinput\fi
\def\prelimfontsloaded{1}
\edef\pfrestoreat{\catcode`\noexpand\@=\the\catcode`\@\relax}\makeatletter
% MathType script sizes: 7 pt superscripts on 12 pt text (58 %), 5 pt second level
\DeclareMathSizes{12}{12}{7}{5}
\DeclareMathSizes{11}{11}{6.5}{5}
\DeclareMathSizes{10.95}{10.95}{6.5}{5}
\DeclareMathSizes{10}{10}{6}{5}
\DeclareMathSizes{9}{9}{5.5}{5}
\DeclareMathSizes{14.4}{14.4}{8.5}{6}
\setmainfont{Times New Roman}
\setsansfont{Arial}
\newfontfamily\coverfont{Arial}
\newfontfamily\calibrifont{Carlito}
\usepackage{unicode-math}
\setmathfont{TeX Gyre Termes Math}
\setmathfont[range={up/{latin,Latin,num}}]{Times New Roman}
\setmathfont[range={it/{latin,Latin}}]{Times New Roman Italic}
\setmathfont[range={bfup/{latin,Latin,num}}]{Times New Roman Bold}
\setmathfont[range={"002B,"003D,"003C,"003E}]{Standard Symbols PS}
\newfontfamily\pf@sym{Standard Symbols PS}
\newfontfamily\pf@symi[FakeSlant=0.21]{Standard Symbols PS}
% MathType-like spacing (tighter than TeX around = + -)
\AtBeginDocument{\thickmuskip=2.5mu plus 1mu\medmuskip=1.5mu plus 1mu\thinmuskip=3mu
% Symbol integral and sigma (18 pt at 12 pt text, as MathType), slanted Symbol pi
 \renewcommand{\int}{\mathop{\vcenter{\hbox{\pf@sym\fontsize{\fpeval{1.5*\f@size}}{0pt}\selectfont\char"F2}}}\nolimits}%
 \renewcommand{\sum}{\mathop{\vcenter{\hbox{\pf@sym\fontsize{\fpeval{1.5*\f@size}}{0pt}\selectfont\char"E5}}}\limits}%
 \renewcommand{\pi}{\mathord{\text{\pf@symi\char"70}}}}
\pfrestoreat

  % house_hw.tex: HOUSE COPY of definite_integrals/common/di_style.tex; [HOUSE] chapter name from \housechapter{...}
% =====================================================================
% di_style.tex: Definite Integrals homework / test layout.
% Loads the house preamble (local copies of the AoI worksheet files:
% ws_preamble_build.tex, ws_preamble_local.tex -> ws_figstyles.tex),
% then adds the homework layout used by HW1-HW3 and the Test.
% Usage (before \begin{document}):  \ditle{Homework 1}
% =====================================================================
% [HOUSE] base, local and fonts are loaded by house.tex
\usepackage{needspace}
\usepackage{zref-savepos,zref-abspage}\makeatletter\zref@addprop{savepos}{abspage}\makeatother
\usepackage{array}

% ---- shortlist toggle ------------------------------------------------
% \ifshortlist true  -> only the questions in SHORTLIST.md (+ challenges)
% A wrapper file sets it with \def\setshortlist{1} before \input-ing the
% main file; the main file can also be edited to say \shortlisttrue.
\newif\ifshortlist
% Task 5: the FINAL homework (= former shortlist) is the default build. HWn.tex -> HWn.pdf.
% A wrapper that sets \def\setfull{1} builds the archive with all original questions.
\ifdefined\setfull\else\shortlisttrue\fi

\makeatletter
% ---- REVIEW build: highlight amended content ---------------------------
% A wrapper sets \def\setreview{1}. In review mode amended text is shaded
% yellow and carries a hover note (PDF tooltip / sticky note); otherwise the
% macros print the content plainly.
%   \chg{note}{inline content}        short inline change (no line breaks), text mode
%   \chgm{note}{math}                 the same inside math mode
%   \begin{chgblock}{note} ... \end{chgblock}   paragraph/display-level change
\newif\ifreview
\ifdefined\setreview \reviewtrue \fi
% hover notes are raw PDF annotations (xdvipdfmx \special; no extra packages).
% Keep notes plain ASCII; balanced parentheses are fine.
\newsavebox{\chg@box}
\newcommand{\chg@ann}[4]{% width, height, depth, note -> invisible markup annotation with a hover note
  \special{pdf:ann width #1 height #2 depth #3 << /Type /Annot /Subtype /Square /Border [0 0 0] /BS << /W 0 >> /C [1 0.85 0] /CA 0.01
    /T (Review) /Subj (Amendment) /Contents (#4) >>}}
\newcommand{\chg@sticky}[1]{\special{pdf:ann width 12pt height 12pt depth 0pt << /Type /Annot /Subtype /Text /Name /Comment
    /C [1 0.85 0] /T (Review) /Subj (Amendment) /Contents (#1) >>}}
\newcommand{\chg}[2]{\ifreview\sbox\chg@box{\colorbox{yellow!60}{#2}}%
  \leavevmode\rlap{\chg@ann{\the\wd\chg@box}{\the\ht\chg@box}{\the\dp\chg@box}{#1}}\usebox\chg@box\else#2\fi}
% \chgm{note}{math}: the same, for use INSIDE math mode (display style kept)
\newcommand{\chgm}[2]{\ifreview\mathchoice{\chg{#1}{$\displaystyle#2$}}{\chg{#1}{$\textstyle#2$}}{\chg{#1}{$\scriptstyle#2$}}{\chg{#1}{$\scriptscriptstyle#2$}}\else#2\fi}
\newtcolorbox{chgbox}[1]{enhanced, breakable, colback=yellow!35, colframe=yellow!35, boxrule=0pt, sharp corners,
  left=2pt, right=14pt, top=1pt, bottom=1pt, before skip=2pt, after skip=2pt,
  overlay first and middle={}, overlay last={}, overlay unbroken and first={\node[anchor=north east, inner sep=1pt] at (frame.north east) {\chg@sticky{#1}\hbox to 12pt{\vrule width 0pt height 12pt}};}}
% \chgnote{note}: just the sticky-note icon (e.g. inside a figure); nothing outside review mode
\newcommand{\chgnote}[1]{\ifreview\chg@sticky{#1}\hbox to 12pt{\vrule width 0pt height 12pt}\fi}
\newenvironment{chgblock}[1]{\ifreview\begin{chgbox}{#1}\fi}{\ifreview\end{chgbox}\fi}

% ---- header ---------------------------------------------------------
\newcommand\di@title{}
\makeatletter
\newcommand{\ditle}[1]{\gdef\di@title{#1}}
\fancypagestyle{di}{\fancyhf{}%
  \lhead{\small \housechap: \di@title}%
  \rhead{\small Compiled by Maximus Pung Jun}%
  \cfoot{\small\thepage}%
  \renewcommand{\headrulewidth}{0.4pt}}
\pagestyle{di}
% title block on page 1, as in his Typst sheets
\newcommand{\dititleblock}[1][]{\thispagestyle{di}%
  {\LARGE\scshape \housechap: \di@title}\par\smallskip
  {Compiled by Maximus Pung Jun}\par\ifx&#1&\else\smallskip{\small\textcolor{black!60}{#1}}\par\fi\medskip}
\makeatother
\newcommand{\disection}[1]{\par\bigskip\noindent{\Large\scshape #1}\par\medskip}

% ---- questions --------------------------------------------------------
% \Qhead{label}{source tag}: question heading; grey [tag] straight after the number, on the left (rule 3B, 9 Nov 2026)
\newcommand{\Qhead}[2]{\par\needspace{8\baselineskip}\bigskip\noindent
  {\large\bfseries #1}\setbox0=\hbox{#2}\ifdim\wd0=0pt\else\enspace{\small\textcolor{black!60}{[#2]}}\fi\par\smallskip}
% part lists with bold labels as he writes them: \pt{(a)}, \pt{(i)}, \pt{(b) (i)}
\newenvironment{qparts}[1][2.4em]{\begin{list}{}{\setlength\leftmargin{#1}%
  \setlength\labelwidth{\dimexpr#1-0.5em\relax}\setlength\labelsep{0.5em}%
  \setlength\itemsep{0.8ex}\setlength\topsep{0.6ex}\setlength\parsep{0.6ex}%
  \renewcommand\makelabel[1]{\textbf{##1}\hfil}}}{\end{list}}
\newcommand{\pt}[1]{\item[#1]}
% thin horizontal label for a challenge question
\newcommand{\challengebanner}{\par\bigskip\noindent\textcolor{black!50}{\rule[0.5ex]{\linewidth}{0.4pt}}\par
  \noindent{\large\scshape Challenge (optional)}\par}
% upright "units"
\newcommand{\units}[1]{\ \text{units}^{#1}}
% END line
\newcommand{\diend}[1][END OF WORKSHEET]{\par\vspace{2\baselineskip}\begin{center}\textbf{#1}\end{center}}

% ---- question registry (full / shortlist from one source) -------------
% \defQ{n}{tag}{body}   question n of the original sheet
% \defA{n}{answer body} its final answer, as printed in his answer section
% \putQ{n} / \putC{n}   print question n as the next (challenge) question
% \putA{n}              print its answer under the same new label
\newcounter{diq}
\makeatletter
\newcommand{\defQ}[3]{\expandafter\gdef\csname diQtag@#1\endcsname{#2}%
  \expandafter\long\expandafter\gdef\csname diQbody@#1\endcsname{#3}}
\newcommand{\defA}[2]{\expandafter\long\expandafter\gdef\csname diA@#1\endcsname{#2}}
\newcommand{\di@tag}[1]{\csname diQtag@#1\endcsname}% source tag only (no original Q number)
% A question never breaks across pages: it is measured and starts on a new page unless it
% fits whole on the rest of this one (only a question taller than a page may break, and then
% only between parts). Page positions of its start/end (zref labels q-s<n>/q-e<n> in the .aux)
% let the checker verify this.
\newsavebox{\di@qmeas}
\newcommand{\di@placeQ}[1]{\par\setbox\di@qmeas=\vbox{#1}%
  \ifdim\dimexpr\ht\di@qmeas+\dp\di@qmeas\relax<0.95\textheight
    \needspace{\dimexpr\ht\di@qmeas+\dp\di@qmeas+\bigskipamount\relax}\else\needspace{14\baselineskip}\fi
  \zsavepos{q-s\thediq}#1\par\zsavepos{q-e\thediq}}
\newcommand{\putQ}[1]{\stepcounter{diq}%
  \expandafter\xdef\csname diLab@#1\endcsname{\thediq}%
  \di@placeQ{\Qhead{\thediq}{\di@tag{#1}}\csname diQbody@#1\endcsname\par}}
\newcommand{\putC}[1]{\stepcounter{diq}%
  \expandafter\xdef\csname diLab@#1\endcsname{\thediq}%
  \expandafter\gdef\csname diChal@#1\endcsname{1}%
  \di@placeQ{\Qhead{\thediq\quad\normalsize\mdseries\scshape Challenge}{\di@tag{#1}}\csname diQbody@#1\endcsname\par}}
\newcommand{\putA}[1]{\par\needspace{4\baselineskip}%
  \begin{qparts}[2.4em]\item[{\csname diLab@#1\endcsname}]\csname diA@#1\endcsname\end{qparts}\medskip}
% ---- worked solutions: \defS{n}{body}, \putQS{n} prints question n + its solution ----
\newcommand{\defS}[2]{\expandafter\long\expandafter\gdef\csname diS@#1\endcsname{#2}}
\newcommand{\disolhead}{\par\medskip\noindent{\bfseries Solution}\par\nobreak\vspace{-0.6ex}\noindent\textcolor{black!35}{\rule{\linewidth}{0.4pt}}\par\nobreak}
\newcommand{\putQS}[1]{\putQ{#1}\disolhead\csname diS@#1\endcsname\par}
% source of a worked solution; visible only in REVIEW builds. #1: his / pub / pip
\newcommand{\solsrc}[2]{\ifreview\par\noindent\def\di@a{#1}\def\di@b{pip}%
  \ifx\di@a\di@b\colorbox{yellow!60}{\parbox{\dimexpr\linewidth-2\fboxsep-14pt}{\small\textbf{Pip-worked:} #2}}\chgnote{Pip-worked solution. No worked solution by the teacher and no published solution was found. Please check it.}%
  \else{\small\textcolor{black!55}{Source of this solution: #2}}\fi\par\smallskip\fi}
\makeatother
% compact sketch for worked solutions: \solfig[axis options]{plot commands}
\newcommand{\solfig}[2][]{\par\begin{center}\begin{tikzpicture}\begin{axis}[qfig, width=6.6cm, height=5cm, #1]#2\end{axis}\end{tikzpicture}\end{center}}
% a sketch added to an existing worked solution (teacher rule), shaded in REVIEW
\newenvironment{addedsketch}{\begin{chgblock}{Sketch added: teacher rule - every area or volume solution has a sketch.}}{\end{chgblock}}
% answer part list (compact)
\newenvironment{aparts}[1][2.4em]{\begin{qparts}[#1]\setlength\itemsep{0.3ex}}{\end{qparts}}

% ---- figure helpers ---------------------------------------------------------
\tikzset{rect/.style={fill=blue!10, draw=blue!55!black, thin}}
\pgfplotsset{qfig/.style={sketch, every node/.append style={font=\scriptsize}}}
% aligned working that starts on the same line as a part label: \pt{(iii)} \begin{topalign} ... \end{topalign}
\newenvironment{topalign}{\hfill$\displaystyle\begin{aligned}[t]}{\end{aligned}$\hfill\null\par\smallskip}
% one display that starts on the same line as a part label
\newcommand{\topdisplay}[1]{\hfill$\displaystyle #1$\hfill\null\par\smallskip}

  % house_solmacros.tex: HOUSE COPY of de-docs/common/de_style.tex minus \input{de_core} and the DE test titles.
% Adds the Ch11 solution helpers: soln/spart tables, sbs/sbsplain side-by-side, \note, \award, \tagm, \dstart.
% =====================================================================
% de_style.tex: Differential Equations (Ch11) homework / test layout.
% = Ch10 di_style (copied as de_core.tex; header text changed) + the DE
% notation and the framed solution boxes of the v1 DE documents
% (soln / spart, used by the Test worked solutions).
% =====================================================================
\providecommand{\e}{\mathrm{e}}
\providecommand{\cosec}{\operatorname{cosec}}
\providecommand{\dt}{\,\mathrm{d}t}
\providecommand{\abs}[1]{\left|#1\right|}
\providecommand{\dd}[2]{\dfrac{\mathrm{d}#1}{\mathrm{d}#2}}
\providecommand{\dvdx}{\dfrac{\mathrm{d}v}{\mathrm{d}x}}
\providecommand{\dudx}{\dfrac{\mathrm{d}u}{\mathrm{d}x}}
\colorlet{notegrey}{black!60}
\makeatletter
\protected\def\note#1{\ifmmode\text{\small\color{notegrey}\itshape #1}\else{\small\color{notegrey}\itshape #1}\fi}
\newcommand{\award}[1]{\unskip\nobreak\hfill\textbf{\small #1}}
\newcommand{\tagm}[1]{\tag*{\textbf{\small #1}}}
\newcommand{\mfAminus}{\[\begin{aligned}&\int\dfrac{1}{a^{2}-x^{2}}\dx\\&\quad=\dfrac{1}{2a}\ln\abs{\dfrac{a+x}{a-x}}+c\end{aligned}\]}
\newcommand{\mfXminus}{\[\begin{aligned}&\int\dfrac{1}{x^{2}-a^{2}}\dx\\&\quad=\dfrac{1}{2a}\ln\abs{\dfrac{x-a}{x+a}}+c\end{aligned}\]}
\newcommand{\sbs@start}{\if@inlabel\def\sbs@shift{\vspace{-\ht\strutbox}}\else
  \ifvmode\else\par\fi\addvspace{0.4ex}\def\sbs@shift{}\fi\noindent}
\newsavebox{\sbs@box}
\newenvironment{sbs}[1]{%
  \sbs@start\begin{minipage}[t]{#1\linewidth}\vspace{0pt}\sbs@shift\raggedright%
  \def\sbsnext{\end{minipage}\hfill\begin{minipage}[t]{\dimexpr0.97\linewidth-#1\linewidth\relax}\vspace{0pt}\sbs@shift\setlength{\fboxsep}{1.5mm}\begin{lrbox}{\sbs@box}\begin{minipage}[t]{\dimexpr\linewidth-2\fboxsep-1pt\relax}\small\raggedright}%
}{\end{minipage}\end{lrbox}\setlength{\fboxrule}{0.5pt}\fbox{\usebox{\sbs@box}}\end{minipage}\par}
\newenvironment{sbsplain}[1]{%
  \sbs@start\begin{minipage}[t]{#1\linewidth}\vspace{0pt}\sbs@shift\raggedright%
  \def\sbsnext{\end{minipage}\hfill\begin{minipage}[t]{\dimexpr0.97\linewidth-#1\linewidth\relax}\vspace{0pt}\sbs@shift\centering}%
}{\end{minipage}\par}
% ---- framed worked-solution box (v1 DE mark-scheme layout, kept for the Test) ----
\newlength{\solcol}\setlength{\solcol}{15mm}
\newcounter{spartno}
\newcommand{\solskips}{\abovedisplayskip=4pt plus 2pt\belowdisplayskip=4pt plus 2pt\abovedisplayshortskip=2pt\belowdisplayshortskip=3pt}
\newtcolorbox{soln}[1][]{enhanced, breakable, sharp corners,
  colback=white, colframe=black, boxrule=0.6pt,
  boxsep=0pt, left=0pt, right=3mm, top=1.5mm, bottom=2mm,
  before skip=1.2ex, after skip=1.5ex,
  before upper={\setcounter{spartno}{0}\parskip=0.5ex\solskips\@fleqntrue\@mathmargin=0pt\raggedright},
  overlay={\draw[line width=0.4pt] ([xshift=\solcol]frame.north west) -- ([xshift=\solcol]frame.south west);},
  #1}
\newenvironment{spart}[1]{%
  \stepcounter{spartno}%
  \ifnum\value{spartno}>1 \par\vspace{0.2ex}\noindent
    \makebox[\linewidth][l]{\rule{\dimexpr\linewidth+3mm\relax}{0.4pt}}\par\nobreak\vspace{0.3ex}\fi
  \begin{list}{}{\leftmargin=\dimexpr\solcol+2.5mm\relax
    \labelwidth=\solcol \labelsep=2.5mm \itemindent=0pt
    \topsep=0.3ex \parsep=0.5ex \rightmargin=0pt \@beginparpenalty=\@M}%
  \item[{\makebox[\solcol][l]{\strut\hspace{1.5mm}\textbf{#1}}}]%
  \ignorespaces}{\end{list}}
\makeatother
% ---- DE sketches: thin light axes (teacher rule), as the Ch10 'sketch' style ----
\pgfplotsset{dsketch/.style={sketch, axis lines=middle, axis line style={thin, black!45, -{Stealth[length=1.6mm]}},
    ticks=none, clip=false, width=6.5cm, height=4.6cm,
    every axis x label/.style={at={(current axis.right of origin)},anchor=west},
    every axis y label/.style={at={(current axis.above origin)},anchor=south},
    every node/.append style={font=\small}},
  dcurve/.style={black, semithick, smooth, samples=120}}
\tikzset{asym/.style={black, dashed, thin}}
\makeatletter
\newcommand{\dstart}{\leavevmode\vadjust{\nobreak\vskip\dimexpr-\baselineskip-\abovedisplayskip+1pt\relax}}
\newcommand{\sdstart}{\leavevmode\vadjust{\nobreak\vskip\dimexpr-\baselineskip-\abovedisplayskip+4pt\relax}}
\makeatother

  % =====================================================================
% house_tf.tex: True/False table (standing rules 2026-11-05/06/07):
%   * a T/F set is a NUMBERED QUESTION (no "Check Your Understanding" heading);
%   * table: row label (a),(b),... | Statement (left) | T/F column (right);
%   * row spacing slightly larger than default (extrarowheight 6pt, arraystretch 1.18,
%     as the Vectors Basics Worksheet after the 2026-11-07 change), font size unchanged;
%   * question doc: T/F column blank; solutions doc (\tfanswerstrue): bold T or F,
%     optional short reason under the statement in small grey.
% Usage:
%   \begin{tftable}  \tfrow{statement}{T}{optional reason}  ...  \end{tftable}
% Plain tabular: a question never splits across pages (house rule), so the table is kept whole.
% =====================================================================
\RequirePackage{array}
\makeatletter
\newif\iftfanswers \tfanswersfalse
\newcounter{tfrow}
\newlength{\tf@lab}\setlength{\tf@lab}{0.9cm}
\newlength{\tf@ans}\setlength{\tf@ans}{1.6cm}
\newenvironment{tftable}{\setcounter{tfrow}{0}\par\smallskip
  \setlength{\extrarowheight}{6pt}\renewcommand{\arraystretch}{1.18}%
  \noindent\begin{tabular}{|>{\centering\arraybackslash}p{\tf@lab}|p{\dimexpr\linewidth-\tf@lab-\tf@ans-6\tabcolsep-4\arrayrulewidth\relax}|>{\centering\arraybackslash}p{\tf@ans}|}
  \hline &\textbf{Statement}&\textbf{T/F}\\\hline}
 {\end{tabular}\par\smallskip}
\newcommand{\tfrow}[3]{\stepcounter{tfrow}\textbf{(\alph{tfrow})} & #1%
  \iftfanswers\if\relax\detokenize{#3}\relax\else\newline{\small\color{black!60}#3}\fi\fi
  & \iftfanswers\textbf{#2}\fi\\\hline}
% explicit row label (a set that keeps the Basics Worksheet letters, e.g. (e) omitted)
\newcommand{\tfrowl}[4]{\textbf{#1} & #2%
  \iftfanswers\if\relax\detokenize{#4}\relax\else\newline{\small\color{black!60}#4}\fi\fi
  & \iftfanswers\textbf{#3}\fi\\\hline}
% \tfcompact: allowed exception (Vectors I HW2/HW3) - compressed rows so the whole set
% stays on page 1 of the question sheet (extrarowheight 1.5pt, tabcolsep 4pt, 0.85/1.45 cm)
\newif\iftf@compact
\newcommand{\tfcompact}{\tf@compacttrue\setlength{\tf@lab}{0.85cm}\setlength{\tf@ans}{1.45cm}}
\renewenvironment{tftable}{\setcounter{tfrow}{0}\par\smallskip
  \iftf@compact\setlength{\extrarowheight}{1.5pt}\setlength{\tabcolsep}{4pt}\else\setlength{\extrarowheight}{6pt}\fi
  \renewcommand{\arraystretch}{1.18}%
  \noindent\begin{tabular}{|>{\centering\arraybackslash}p{\tf@lab}|p{\dimexpr\linewidth-\tf@lab-\tf@ans-6\tabcolsep-4\arrayrulewidth\relax}|>{\centering\arraybackslash}p{\tf@ans}|}
  \hline &\textbf{Statement}&\textbf{T/F}\\\hline}
 {\end{tabular}\par\smallskip}
% tflongtable: same table, may break across pages with the header repeated (solutions
% copies whose reasons make the set longer than a page). Needs longtable.
\RequirePackage{longtable}
\newenvironment{tflongtable}{\setcounter{tfrow}{0}\par\smallskip
  \setlength{\extrarowheight}{6pt}\renewcommand{\arraystretch}{1.18}%
  \setlength{\LTleft}{0pt}\setlength{\LTright}{0pt}%
  \begin{longtable}{|>{\centering\arraybackslash}p{\tf@lab}|p{\dimexpr\linewidth-\tf@lab-\tf@ans-6\tabcolsep-4\arrayrulewidth\relax}|>{\centering\arraybackslash}p{\tf@ans}|}
  \hline &\textbf{Statement}&\textbf{T/F}\\\hline\endhead}
 {\end{longtable}\par}
% \putTFlong{n}: True/False question n (body of \defQ{n}: intro, then tftable) printed with
% tflongtable, so a long set may run over a page (house_sol redefines it for solutions).
\newcommand{\putTFlong}[1]{\begingroup\let\tftable\tflongtable\let\endtftable\endtflongtable
  \stepcounter{diq}\expandafter\xdef\csname diLab@#1\endcsname{\thediq}%
  \Qhead{\thediq}{\csname diQtag@#1\endcsname}\csname diQbody@#1\endcsname\par\endgroup}
\makeatother

  \ifx\housexm\housexsol \def\prelimkeepcolour{1}% house_prelim.tex: HOUSE COPY of definite_integrals/common/prelim_style.tex; [HOUSE] titles from \housechapter
% =====================================================================
% prelim_style.tex: makes the Test look like a Singapore JC H2 Mathematics
% 9758 prelim paper (studied from ACJC/RI 2026 P1/P2 on the box:
% body Times New Roman + Times-style maths, cover in Arial).
% Fonts: prelim_fonts.tex (real Times New Roman/Arial + Times-style maths).
% Shared by Ch10 (di_style) and Ch11 (de_style; sets \prelimtitle/\prelimsubject). Loaded AFTER the chapter style (overrides its layout for the test only).
% =====================================================================
% [HOUSE] fonts loaded by house.tex
\providecommand{\prelimtitle}{\MakeUppercase{\housechap} TEST}
\providecommand{\prelimsubject}{\housechap}
\geometry{a4paper, top=22mm, bottom=24mm, left=25mm, right=25mm, headsep=6mm, footskip=12mm}
\setlength{\parindent}{0pt}
\setlength{\parskip}{0.9ex plus 0.3ex minus 0.2ex}
\linespread{1.3}
\setlength{\emergencystretch}{3em}
\clubpenalty=10000 \widowpenalty=10000 \interlinepenalty=10000
\AtBeginDocument{\fontsize{12}{14.4}\selectfont}
\renewcommand{\normalsize}{\fontsize{12}{14.4}\selectfont}

% ---- figures: plain black/grey like a printed paper ----------------------
\ifdefined\prelimkeepcolour\else % question paper: black/grey figures; solutions keep colour
\pgfplotsset{curve/.style={black, semithick, smooth, samples=150},
  line2/.style={black, semithick, samples=2},
  sketch/.append style={axis line style={thin, black}, every axis label/.append style={black}}}
\tikzset{region/.style={fill=black!15, draw=none}, regionB/.style={fill=black!30, draw=none},
  rect/.style={fill=black!8, draw=black, thin}, ax3/.style={-{Stealth[length=1.5mm]}, thin, black}}
\colorlet{blue}{black}\colorlet{red}{black}\colorlet{orange}{black}
\fi

% ---- header / footer: page number centred top (bold), [Turn over at bottom right
\makeatletter
\providecommand{\prelimfootlabel}{\prelimtitle\ 2026}
\newcommand{\pp@footer}{{\coverfont\scriptsize \prelimfootlabel\hspace{3em}H2 MATHEMATICS 9758}}
\fancypagestyle{prelim}{\fancyhf{}%
  \chead{\bfseries\thepage}%
  \lfoot{\pp@footer}%
  \rfoot{\ifnum\value{page}<\getpagerefnumber{LastPage}\relax\textbf{[Turn over}\fi}%
  \renewcommand{\headrulewidth}{0pt}\renewcommand{\footrulewidth}{0pt}}
\fancypagestyle{prelimcover}{\fancyhf{}%
  \rfoot{\coverfont\textbf{[Turn over}}%
  \renewcommand{\headrulewidth}{0pt}\renewcommand{\footrulewidth}{0pt}}
\pagestyle{prelim}

% ---- marks: [n] regular weight, right-aligned at the end of the part ----------
\renewcommand{\mk}[1]{{\unskip\nobreak\hfil\penalty50\hskip1.5em\hbox{}\nobreak\hfill[#1]%
  \parfillskip=0pt\finalhyphendemerits=0\par}}

% ---- questions: bold number at the left margin, text indented; grey [tag] straight after the number (3B)
\newsavebox{\pp@qmeas}
\newcommand{\pp@qlist}[1]{%
  \begin{list}{}{\setlength\leftmargin{2.6em}\setlength\labelwidth{2.1em}\setlength\labelsep{0.5em}%
     \setlength\topsep{0pt}\setlength\itemsep{0pt}\setlength\parsep{0.9ex}%
     \renewcommand\makelabel[1]{\textbf{##1}\hfil}}%
  \item[\thediq]\@endparpenalty\@M\setbox0=\hbox{\di@tag{#1}}\ifdim\wd0=0pt\else{\small\textcolor{black!60}{[\di@tag{#1}]}}\par\nobreak\fi\csname diQbody@#1\endcsname\par\penalty\@M
  \end{list}}
% a question starts on a new page unless it fits whole on the current one
% (questions taller than a page may break, but only between parts)
\renewcommand{\putQ}[1]{\stepcounter{diq}%
  \expandafter\xdef\csname diLab@#1\endcsname{\thediq}%
  \par\setbox\pp@qmeas=\vbox{\pp@qlist{#1}}%
  \ifdim\dimexpr\ht\pp@qmeas+\dp\pp@qmeas\relax<0.95\textheight
    \needspace{\dimexpr\ht\pp@qmeas+\dp\pp@qmeas\relax}\else\needspace{14\baselineskip}\fi
  \zsavepos{q-s\thediq}\pp@qlist{#1}\par\zsavepos{q-e\thediq}\vspace{5ex plus 2ex minus 1ex}}
% part labels (a)/(i) bold, indented under the question text
\renewenvironment{qparts}[1][2.6em]{\begin{list}{}{\setlength\leftmargin{#1}%
  \setlength\labelwidth{\dimexpr#1-0.5em\relax}\setlength\labelsep{0.5em}%
  \setlength\itemsep{1.4ex}\setlength\topsep{1.2ex}\setlength\parsep{0.8ex}%
  \renewcommand\makelabel[1]{\textbf{##1}\hfil}}}{\end{list}}
\makeatother

% ---- cover page --------------------------------------------------------------
% A-level H2 style cover (9 Nov 2026; reference house_format/ref/alevel_cover_ref.png), no logos/barcodes.
% \prelimcover{total marks}{duration}{date}; heading block from \housechap; override with
% \prelimcoverhead (4 lines) and \prelimcoverline (the line under MATHEMATICS).
\makeatletter
\providecommand{\prelimcoverhead}{MAXIMUS PUNG JUN\\ H2 Mathematics (9758)\\ \housechap\\ Chapter Test}
\providecommand{\prelimcoverline}{\prelimsubject\ Test}
\newcommand{\pp@htwo}{\begin{tikzpicture}[baseline=(b.base)]
  \node[draw, line width=0.9pt, minimum width=19mm, minimum height=15mm, inner sep=0pt] (b)
    {\fontsize{38}{38}\selectfont\special{pdf:literal direct 1 Tr 0.7 w}H2\special{pdf:literal direct 0 Tr}};
  \end{tikzpicture}}
\newcommand{\prelimcover}[3]{% #1 total marks, #2 duration, #3 date line
\thispagestyle{prelimcover}%
{\coverfont\fontsize{11}{13.5}\selectfont\linespread{1}\selectfont\setlength{\parskip}{0pt}%
\vspace*{4mm}
\noindent\begin{tabular}{@{}l@{\hspace{5mm}}l@{}}\raisebox{-0.5\height}{\pp@htwo}&
  \raisebox{-0.5\height}{\begin{minipage}{120mm}\fontsize{12}{14.5}\selectfont\prelimcoverhead\end{minipage}}\end{tabular}\par
\vspace{9mm}
\noindent\rule{\linewidth}{0.8pt}\par\vspace{1.2mm}
\noindent{\fontsize{12}{14.5}\selectfont\bfseries MATHEMATICS\hfill 9758}\par\vspace{1.2mm}
\noindent\prelimcoverline\hfill\textbf{#3}\par\vspace{1.2mm}
\setbox0=\hbox{#2}\ifdim\wd0>0pt \noindent\hfill\textbf{#2}\par\fi
\vspace{3mm}
\noindent You must answer on writing paper.\par\vspace{3mm}
\noindent\begin{tabular}{@{}l@{\hspace{2mm}}l@{}}You will need:&Writing paper\\ &List of formulae and results (MF27)\end{tabular}\par
\vspace{2mm}
\noindent\rule{\linewidth}{0.8pt}\par\vspace{1.5mm}
\noindent{\bfseries INSTRUCTIONS}\par\vspace{0.5mm}
\begin{list}{\textbullet}{\setlength\leftmargin{7mm}\setlength\labelwidth{4mm}\setlength\labelsep{3mm}\setlength\itemsep{0.6mm}\setlength\topsep{0pt}\setlength\parsep{0pt}}
\item Answer \textbf{all} questions.
\item Give non-exact numerical answers correct to 3 significant figures, or 1 decimal place for angles in degrees, unless a different level of accuracy is specified in the question.
\item You should use an approved graphing calculator.
\item Unsupported answers from a graphing calculator are allowed unless a question specifically states otherwise.
\item Where unsupported answers from a graphing calculator are \textbf{not} allowed in a question, you must present the mathematical steps using mathematical notations and not calculator commands.
\item You must show all working clearly.
\end{list}
\vspace{6mm}
\noindent{\bfseries INFORMATION}\par\vspace{0.5mm}
\begin{list}{\textbullet}{\setlength\leftmargin{7mm}\setlength\labelwidth{4mm}\setlength\labelsep{3mm}\setlength\itemsep{0.6mm}\setlength\topsep{0pt}\setlength\parsep{0pt}}
\item The total mark for this paper is #1.
\item The number of marks for each question or part question is shown in brackets [\ ].
\end{list}
\vfill
\noindent\rule{\linewidth}{0.8pt}\par\vspace{1mm}
\centerline{This document has \textbf{\pageref{LastPage}} pages.}\par\vspace{2mm}
}\clearpage}
\makeatother

% ---- solutions header block (page 1 of the worked solutions; same cover style) ----
\newcommand{\prelimsolhead}[1]{% #1 total marks
{\coverfont\fontsize{10}{12}\selectfont\linespread{1}\selectfont
\noindent{\bfseries \prelimtitle}\\[0.2ex]{\bfseries Compiled by Maximus Pung Jun}\\[1.6ex]Higher 2\par\vspace{2ex}
\noindent\rule{\linewidth}{1.2pt}\par\vspace{0.4ex}
\noindent{\Large\bfseries MATHEMATICS}\hfill{\Large\bfseries 9758}\par\vspace{1.2ex}
\noindent \prelimsubject\hfill\textbf{2026}\\
WORKED SOLUTIONS\hfill\textbf{Total marks: #1}\par\vspace{1ex}
\noindent\rule{\linewidth}{0.6pt}\par}\vspace{4ex}}
% house_sol.tex: HOUSE COPY of definite_integrals/common/ri_solstyle.tex (unchanged)
% =====================================================================
% ri_solstyle.tex: worked solutions in the layout of the RI 2026 Year 6 Prelim
% "Solutions with comments" (grail.moe 26612/26613), loaded AFTER prelim_style.
%   * page 1: bold school-style block + bold underlined title; later pages: running
%     title top-left with a rule under it; footer "Page N of M" bottom right (Calibri).
%   * each question: prelim-format question text (bold number, marks right, source tag
%     bottom right), then a bordered table:  Solution | Comments  (bold header row);
%     body rows  label + [marks] (bold, narrow left column) | working | comments,
%     a full-width rule between parts, working left-aligned (fleqn).
% Rows come from the existing sources: \pt in qparts (Ch10 HW), \slabel (Ch10 Test),
% spart (Ch11). Marks under each label are taken from the question's \mk.
% =====================================================================
\makeatletter
\usepackage{lastpage}
\usepackage{adjustbox}
\usepackage{zref-savepos,zref-abspage}
\zref@addprop{savepos}{abspage}
\geometry{a4paper, top=22mm, bottom=24mm, left=22mm, right=22mm, headsep=8mm, footskip=12mm, headheight=24pt}
\providecommand{\risoltitle}{\prelimtitle: Solutions with comments}
\providecommand{\risolschool}{Compiled by Maximus Pung Jun}
\fancypagestyle{risol}{\fancyhf{}%
  \lhead{\fontsize{11}{12}\selectfont 2026 \risoltitle\\[-1.2ex]\rule{\textwidth}{0.6pt}}%
  \rfoot{\calibrifont\fontsize{11}{13}\selectfont Page \textbf{\thepage} of \textbf{\pageref*{LastPage}}}%
  \renewcommand{\headrulewidth}{0pt}\renewcommand{\footrulewidth}{0pt}}
\fancypagestyle{risolfirst}{\fancyhf{}%
  \rfoot{\calibrifont\fontsize{11}{13}\selectfont Page \textbf{\thepage} of \textbf{\pageref*{LastPage}}}%
  \renewcommand{\headrulewidth}{0pt}\renewcommand{\footrulewidth}{0pt}}
\pagestyle{risol}
\newif\ifri@fresh
\newcommand{\risolfront}{\thispagestyle{risolfirst}\global\ri@freshtrue%
  {\noindent\bfseries\linespread{1}\selectfont \risolschool\\ H2 Mathematics (9758)\\ 2026\par}%
  \vspace{4ex}{\noindent\bfseries\underline{2026 \risoltitle}\par}\vspace{3ex}}
% ---- table geometry (RI: label 42pt | working 284pt | comments 142pt of 468pt) ----
\newlength{\ri@L}\setlength{\ri@L}{15mm}
\newlength{\ri@C}\setlength{\ri@C}{40mm}
\newlength{\ri@pad}\setlength{\ri@pad}{1.8mm}
\newif\ifri@sol \newif\ifri@first
\newcount\ri@depth
\def\ri@outer{}
\newtcolorbox{risol}{enhanced, breakable, sharp corners, frame hidden, colback=white,
  boxrule=0pt, boxsep=0pt, left=\dimexpr\ri@L+\ri@pad\relax, right=\dimexpr\ri@C+\ri@pad\relax,
  top=1.2mm, bottom=2.5mm, toptitle=0.6mm, bottomtitle=0.6mm, lefttitle=\ri@pad, righttitle=0pt,
  colbacktitle=white, coltitle=black, fonttitle=\bfseries,
  title={Solution\hfill\makebox[\ri@C][l]{\hspace{\ri@pad}Comments}},
  before skip=2ex, after skip=2ex,
  overlay unbroken and first={\draw[line width=0.5pt] (frame.south west) rectangle (frame.north east);
     \draw[line width=0.5pt] (title.south west) -- (title.south east);
     \draw[line width=0.5pt] ([xshift=\ri@L]interior.north west) -- ([xshift=\ri@L]frame.south west);
     \draw[line width=0.5pt] ([xshift=-\ri@C]frame.north east) -- ([xshift=-\ri@C]frame.south east);},
  overlay middle and last={\draw[line width=0.5pt] (frame.south west) rectangle (frame.north east);
     \draw[line width=0.5pt] ([xshift=\ri@L]frame.north west) -- ([xshift=\ri@L]frame.south west);
     \draw[line width=0.5pt] ([xshift=-\ri@C]frame.north east) -- ([xshift=-\ri@C]frame.south east);},
  after upper={\ri@endrow\protected@write\@auxout{}{\string\ri@settotal{\ri@qid}{\the\ri@sumH}}},
  before upper={\global\ri@sumH=14mm\ri@soltrue\ri@firsttrue\ri@depth=0 \parindent=0pt \parskip=0.6ex
     \predisplaypenalty=100 
     \linespread{1.12}\selectfont\raggedright\@fleqntrue\@mathmargin=0pt
     \abovedisplayskip=5pt plus 2pt\belowdisplayskip=5pt plus 2pt
     \abovedisplayshortskip=3pt plus 1pt\belowdisplayshortskip=4pt plus 2pt
     \ifcsname ri@hp@\ri@qid\endcsname\else\ri@row{}\fi}}
\def\ri@qid{}
% one row: full-width rule (glue leaders, so it vanishes at a page break where the frame
% edge takes its place; not before the first row), then bold label + bold [marks]
% in the left column, on the baseline of the first line of the working
\ExplSyntaxOn
\tl_new:N \g__ri_tmp_tl
\cs_new_protected:Npn \ri_setkey:Nn #1#2
  { \group_begin: \cs_set_eq:NN \chg \use_ii:nn \cs_set_eq:NN \chgm \use_ii:nn
    \cs_set_eq:NN \textbf \use:n \cs_set_eq:NN \mathrm \use:n
    \tl_gset:Nx \g__ri_tmp_tl {#2} \group_end:
    \tl_set_eq:NN #1 \g__ri_tmp_tl \tl_remove_all:Nn #1 {~} }
\cs_set_eq:NN \ri@setkey \ri_setkey:Nn
\ExplSyntaxOff
% ---- page-break rules (positions from the previous run, zref labels in the .aux) ----
% * a row (one part) never breaks across pages: if it does not fit on the rest of the
%   page it starts the next page (forced break between rows); only a row taller than a
%   whole page may break inside;
% * a whole solution that fits on one page is kept together (needspace before the table);
%   a longer one starts where at least its first row fits;
% * a row is at least as tall as its label + [marks].
\newcount\ri@n \newbox\ri@labbox \def\ri@open{0}
\newdimen\ri@ybot \newdimen\ri@ytop \newdimen\ri@h \newdimen\ri@sumH
\AtBeginDocument{\ri@ybot=\dimexpr\paperheight-1in-\voffset-\topmargin-\headheight-\headsep-\textheight\relax
  \ri@ytop=\dimexpr\ri@ybot+\textheight\relax}
\def\ri@pg#1{\zref@extractdefault{#1}{abspage}{-1}}
% \ri@rowh{n}: height of row n in the previous run -> \ri@h (0pt if unknown)
\newcommand{\ri@rowh}[1]{\ri@h=0pt
  \ifnum\ri@pg{ri-s#1}>0 \ifnum\ri@pg{ri-e#1}>0
    \ifnum\ri@pg{ri-s#1}=\ri@pg{ri-e#1}\relax \ri@h=\dimexpr\zposy{ri-s#1}sp-\zposy{ri-e#1}sp\relax
    \else \ri@h=\dimexpr\zposy{ri-s#1}sp-\ri@ybot+\ri@ytop-\zposy{ri-e#1}sp-6mm\relax\fi\fi\fi}
\newcommand{\ri@setfirst}[2]{\expandafter\gdef\csname ri@first@#1\endcsname{#2}}
\newcommand{\ri@setkept}[1]{\expandafter\gdef\csname ri@kept@#1\endcsname{}}
\newcommand{\ri@settotal}[2]{\expandafter\gdef\csname ri@total@#1\endcsname{#2}}
\newcommand{\ri@endrow}{\ifnum\ri@open=1 \par\zsavepos{ri-e\the\ri@n}%
  \ifnum\ri@pg{ri-s\the\ri@n}=\ri@pg{ri-e\the\ri@n}\relax
    \dimen@=\dimexpr\ri@curH sp+1.6mm-\zposy{ri-s\the\ri@n}sp+\zposy{ri-e\the\ri@n}sp\relax
    \ifdim\dimen@>0pt \nobreak\kern\dimen@\penalty\z@\fi\fi
  \ri@rowh{\the\ri@n}\global\advance\ri@sumH\dimexpr\ri@h+3.5mm\relax
  \gdef\ri@open{0}\fi}
\newcommand{\ri@row}[1]{\par\ri@endrow\global\advance\ri@n\@ne\gdef\ri@open{1}\@beginparpenalty\@M
  \ifri@first\global\ri@firstfalse\ri@firstfalse
    \protected@write\@auxout{}{\string\ri@setfirst{\ri@qid}{\the\ri@n}}%
  \else
    \zsavepos{ri-p\the\ri@n}%
    % push the row to the next page if it would not fit on the rest of this one
    \ri@rowh{\the\ri@n}%
    \ifnum\ri@pg{ri-p\the\ri@n}>0 \ifdim\ri@h>0pt
      \ifdim\dimexpr\ri@h+4mm\relax<\dimexpr\textheight-10mm\relax
        \ifdim\dimexpr\ri@h+7mm\relax>\dimexpr\zposy{ri-p\the\ri@n}sp-\ri@ybot\relax
          \penalty-\@M\fi\fi\fi\fi
    \vskip0.6ex\nointerlineskip
    \leaders\hbox{\hskip-\dimexpr\ri@L+\ri@pad\relax\vrule width\dimexpr\linewidth+\ri@L+\ri@C+2\ri@pad\relax height0.5pt depth0pt}\vskip0.5pt
    \nobreak\vskip1.4mm\fi
  \ri@setkey\ri@key{\ri@qid#1}\def\ri@lab{#1}%
  \setbox\ri@labbox=\vtop{\hsize=\dimexpr\ri@L-1.6mm\relax\bfseries\raggedright\parskip=0pt
      \strut\ri@lab\ifcsname ri@mk@\ri@key\endcsname\ifx\ri@lab\@empty\else\\\fi
      [\csname ri@mk@\ri@key\endcsname]\fi\strut\par}%
  \xdef\ri@curH{\number\dimexpr\ht\ri@labbox+\dp\ri@labbox\relax}%
  \nointerlineskip\zsavepos{ri-s\the\ri@n}\nointerlineskip\nobreak
  \vbox to 0pt{\hbox{\hskip-\dimexpr\ri@L+\ri@pad-1.2mm\relax\box\ri@labbox}\vss}\nobreak\prevdepth=\dp\strutbox
  \vskip-\parskip\ignorespaces}
% ---- qparts / \pt: question mode records labels (for the marks); solution mode makes rows
\NewCommandCopy\ri@qparts@begin\qparts \let\ri@qparts@end\endqparts
\renewenvironment{qparts}[1][2.6em]{\advance\ri@depth\@ne
  \ifri@sol\else\ri@qparts@begin[#1]\fi}{\ifri@sol\par\else\ri@qparts@end\fi}
\def\ri@qpartsname{qparts}
\renewcommand{\pt}[1]{\ifnum\ri@depth<2 \gdef\ri@outer{#1}\gdef\ri@cur{#1}\else\expandafter\gdef\expandafter\ri@cur\expandafter{\ri@outer#1}\fi
  \ifri@sol\expandafter\ri@solpt\else\expandafter\ri@qpt\fi{#1}}
\def\ri@qpt#1{\expandafter\gdef\csname ri@hp@\ri@qid\endcsname{}\item[#1]}
\def\ri@solpt#1{\@ifnextchar\begin\ri@ptpeek{\ri@row{\ri@cur}}}
\def\ri@ptpeek\begin#1{\def\ri@tmp{#1}\ifx\ri@tmp\ri@qpartsname
   \expandafter\@firstoftwo\else\expandafter\@secondoftwo\fi
  {\begin{qparts}}{\ri@row{\ri@cur}\begin{#1}}}
% marks: record [n] against the current question/part
\let\ri@mk\mk
\renewcommand{\mk}[1]{\ri@setkey\ri@k{\ri@qid\ri@cur}\ri@setkey\ri@mv{#1}\expandafter\xdef\csname ri@mk@\ri@k\endcsname{\ri@mv}\ri@mk{#1}}
% \slabel (Ch10 Test answers): letter labels set the outer part, roman ones are appended;
% consecutive \slabel calls merge into one row label
\def\ri@romans{(i)(ii)(iii)(iv)(v)(vi)(vii)}
\def\ri@firstpart(#1)#2\@nil{(#1)}
\def\slabel#1{\ifri@sol\expandafter\ri@solslabel\else\expandafter\ri@qslabel\fi{#1}}
\def\ri@qslabel#1{\par\smallskip\noindent\textbf{#1}\par\nobreak}
\def\ri@solslabel#1{\@expandtwoargs\in@{#1}{\ri@romans}%
  \ifin@\expandafter\gdef\expandafter\ri@cur\expandafter{\ri@outer#1}\else\gdef\ri@cur{#1}\xdef\ri@outer{\ri@firstpart#1\@nil}\fi
  \@ifnextchar\slabel{}{\ri@row{\ri@cur}}}
% Ch11 spart{label}: a row; soln: the table
\def\spart#1{\gdef\ri@cur{#1}\ri@row{#1}}\def\endspart{\par}
\def\soln{\@ifnextchar[{\ri@solnopt}{\risol}}\def\ri@solnopt[#1]{\risol}\def\endsoln{\endrisol}
% ---- working that started on the label line is left-aligned in the working column
\AtBeginDocument{\@ifundefined{topalign}{}{%
  \renewenvironment{topalign}{\par\noindent$\displaystyle\begin{aligned}[t]}{\end{aligned}$\par\smallskip}%
  \renewcommand{\topdisplay}[1]{\par\noindent$\displaystyle #1$\par\smallskip}}}
% ---- Ch11 side-by-side (text | sketch): the sketch shrinks to fit its half of the working column
\AtBeginDocument{\@ifundefined{sbsplain}{}{%
  \renewenvironment{sbsplain}[1]{%
    \sbs@start\begin{minipage}[t]{#1\linewidth}\vspace{0pt}\sbs@shift\raggedright%
    \def\sbsnext{\end{minipage}\hfill\begin{minipage}[t]{\dimexpr0.97\linewidth-#1\linewidth\relax}\vspace{0pt}\sbs@shift\centering
      \begin{adjustbox}{max width=\linewidth}}%
  }{\end{adjustbox}\end{minipage}\par}}}
% ---- figures in the working column shrink to fit it (only if wider), colours kept
\newcommand{\rifitfig}[1]{\expandafter\let\csname ri@orig@\string#1\endcsname#1%
  \def#1{\@ifnextchar\bgroup{\ri@fitarg#1}{\ri@fitnoarg#1}}}
\def\ri@fitarg#1#2{\ifri@sol\adjustbox{max width=\linewidth}{\csname ri@orig@\string#1\endcsname{#2}}\else\csname ri@orig@\string#1\endcsname{#2}\fi}
\def\ri@fitnoarg#1{\ifri@sol\adjustbox{max width=\linewidth}{\csname ri@orig@\string#1\endcsname}\else\csname ri@orig@\string#1\endcsname\fi}
\@ifundefined{solfig}{}{\renewcommand{\solfig}[2][]{\par\begin{center}\adjustbox{max width=\linewidth}{\begin{tikzpicture}\begin{axis}[qfig, width=6.6cm, height=5cm, #1]#2\end{axis}\end{tikzpicture}}\end{center}}}
% \ricomment{text}: text in the Comments column of the current row (call right after \pt,
% \slabel or \spart). The row is made at least as tall as the comment.
\newbox\ri@cmtbox
\newcommand{\ricomment}[1]{\par\setbox\ri@cmtbox=\vtop{\hsize=\dimexpr\ri@C-2\ri@pad\relax
    \parskip=0pt\raggedright\small{\color{black!60}\strut#1\strut\par}}%
  \ifdim\dimexpr\ht\ri@cmtbox+\dp\ri@cmtbox\relax>\ri@curH sp\xdef\ri@curH{\number\dimexpr\ht\ri@cmtbox+\dp\ri@cmtbox\relax}\fi
  \nointerlineskip\vbox to 0pt{\hbox to\linewidth{\hskip\dimexpr\linewidth+2\ri@pad\relax\box\ri@cmtbox\hss}\vss}\nobreak\prevdepth=\dp\strutbox
  \vskip-\parskip\ignorespaces}
% \putTFlong{n}: the True/False question n with its answer table allowed to break across
% pages (the body of \defQ{n} must be: intro text, then \begin{tftable}...\end{tftable});
% the solutions copy prints it with tflongtable.
\renewcommand{\putTFlong}[1]{\ri@newq\begingroup\let\tftable\tflongtable\let\endtftable\endtflongtable
  \stepcounter{diq}\expandafter\xdef\csname diLab@#1\endcsname{\thediq}%
  \par\noindent\hbox to 2.6em{\textbf{\thediq}\hfil}\setbox0=\hbox{\di@tag{#1}}\ifdim\wd0=0pt\else{\small\textcolor{black!60}{[\di@tag{#1}]}}\par\noindent\hskip2.6em\fi\ignorespaces\csname diQbody@#1\endcsname\par
  \endgroup\vspace{3ex}}
% ---- question + solution; each question starts on a new page (as RI)
\def\ri@qid{}
\let\ri@putQ\putQ
\renewcommand{\putQ}[1]{\gdef\ri@qid{#1}\global\ri@depth=0 \gdef\ri@cur{}\ri@putQ{#1}}
% \risolQ{n}{body macro name}: question n + its solution body in the table;
% \risolQraw: body already contains its own soln environment (Ch11 Test)
\newcommand{\ri@newq}{\ifri@fresh\global\ri@freshfalse\else\clearpage\fi}
% keep the whole table together if it fits on a page; otherwise make sure at least its
% header and first row fit before starting it
\newcommand{\ri@keep}[1]{\par
  \ifcsname ri@total@#1\endcsname
    % hysteresis: a table kept whole in the previous run stays whole up to textheight-4.5mm,
    % so a total near the limit cannot flip between runs
    \ifdim\csname ri@total@#1\endcsname<\dimexpr\textheight-\ifcsname ri@kept@#1\endcsname4.5\else8\fi mm\relax
      \protected@write\@auxout{}{\string\ri@setkept{#1}}%
      \needspace{\dimexpr\csname ri@total@#1\endcsname+4mm\relax}%
    \else\ifcsname ri@first@#1\endcsname
      \ri@rowh{\csname ri@first@#1\endcsname}%
      \ifdim\ri@h<\dimexpr\textheight-20mm\relax\needspace{\dimexpr\ri@h+18mm\relax}\fi\fi\fi\fi}
\newcommand{\risolQ}[2]{\ri@newq\putQ{#1}\vspace{-3ex}\ri@keep{#1}\begin{risol}\csname #2@#1\endcsname\par\end{risol}}
\newcommand{\risolQraw}[2]{\ri@newq\putQ{#1}\vspace{-3ex}\ri@keep{#1}\csname #2@#1\endcsname\par}
\renewcommand{\putQS}[1]{\risolQ{#1}{diS}}
% end marker only if it fits on the last page (no page holding just the marker)
\renewcommand{\diend}[1][END OF SOLUTIONS]{\par\ifdim\dimexpr\pagegoal-\pagetotal\relax<4\baselineskip\else\vspace{1.5ex}\begin{center}\textbf{#1}\end{center}\fi}
\makeatother
\tfanswerstrue\fi
  \ifx\housexm\housextest % house_prelim.tex: HOUSE COPY of definite_integrals/common/prelim_style.tex; [HOUSE] titles from \housechapter
% =====================================================================
% prelim_style.tex: makes the Test look like a Singapore JC H2 Mathematics
% 9758 prelim paper (studied from ACJC/RI 2026 P1/P2 on the box:
% body Times New Roman + Times-style maths, cover in Arial).
% Fonts: prelim_fonts.tex (real Times New Roman/Arial + Times-style maths).
% Shared by Ch10 (di_style) and Ch11 (de_style; sets \prelimtitle/\prelimsubject). Loaded AFTER the chapter style (overrides its layout for the test only).
% =====================================================================
% [HOUSE] fonts loaded by house.tex
\providecommand{\prelimtitle}{\MakeUppercase{\housechap} TEST}
\providecommand{\prelimsubject}{\housechap}
\geometry{a4paper, top=22mm, bottom=24mm, left=25mm, right=25mm, headsep=6mm, footskip=12mm}
\setlength{\parindent}{0pt}
\setlength{\parskip}{0.9ex plus 0.3ex minus 0.2ex}
\linespread{1.3}
\setlength{\emergencystretch}{3em}
\clubpenalty=10000 \widowpenalty=10000 \interlinepenalty=10000
\AtBeginDocument{\fontsize{12}{14.4}\selectfont}
\renewcommand{\normalsize}{\fontsize{12}{14.4}\selectfont}

% ---- figures: plain black/grey like a printed paper ----------------------
\ifdefined\prelimkeepcolour\else % question paper: black/grey figures; solutions keep colour
\pgfplotsset{curve/.style={black, semithick, smooth, samples=150},
  line2/.style={black, semithick, samples=2},
  sketch/.append style={axis line style={thin, black}, every axis label/.append style={black}}}
\tikzset{region/.style={fill=black!15, draw=none}, regionB/.style={fill=black!30, draw=none},
  rect/.style={fill=black!8, draw=black, thin}, ax3/.style={-{Stealth[length=1.5mm]}, thin, black}}
\colorlet{blue}{black}\colorlet{red}{black}\colorlet{orange}{black}
\fi

% ---- header / footer: page number centred top (bold), [Turn over at bottom right
\makeatletter
\providecommand{\prelimfootlabel}{\prelimtitle\ 2026}
\newcommand{\pp@footer}{{\coverfont\scriptsize \prelimfootlabel\hspace{3em}H2 MATHEMATICS 9758}}
\fancypagestyle{prelim}{\fancyhf{}%
  \chead{\bfseries\thepage}%
  \lfoot{\pp@footer}%
  \rfoot{\ifnum\value{page}<\getpagerefnumber{LastPage}\relax\textbf{[Turn over}\fi}%
  \renewcommand{\headrulewidth}{0pt}\renewcommand{\footrulewidth}{0pt}}
\fancypagestyle{prelimcover}{\fancyhf{}%
  \rfoot{\coverfont\textbf{[Turn over}}%
  \renewcommand{\headrulewidth}{0pt}\renewcommand{\footrulewidth}{0pt}}
\pagestyle{prelim}

% ---- marks: [n] regular weight, right-aligned at the end of the part ----------
\renewcommand{\mk}[1]{{\unskip\nobreak\hfil\penalty50\hskip1.5em\hbox{}\nobreak\hfill[#1]%
  \parfillskip=0pt\finalhyphendemerits=0\par}}

% ---- questions: bold number at the left margin, text indented; grey [tag] straight after the number (3B)
\newsavebox{\pp@qmeas}
\newcommand{\pp@qlist}[1]{%
  \begin{list}{}{\setlength\leftmargin{2.6em}\setlength\labelwidth{2.1em}\setlength\labelsep{0.5em}%
     \setlength\topsep{0pt}\setlength\itemsep{0pt}\setlength\parsep{0.9ex}%
     \renewcommand\makelabel[1]{\textbf{##1}\hfil}}%
  \item[\thediq]\@endparpenalty\@M\setbox0=\hbox{\di@tag{#1}}\ifdim\wd0=0pt\else{\small\textcolor{black!60}{[\di@tag{#1}]}}\par\nobreak\fi\csname diQbody@#1\endcsname\par\penalty\@M
  \end{list}}
% a question starts on a new page unless it fits whole on the current one
% (questions taller than a page may break, but only between parts)
\renewcommand{\putQ}[1]{\stepcounter{diq}%
  \expandafter\xdef\csname diLab@#1\endcsname{\thediq}%
  \par\setbox\pp@qmeas=\vbox{\pp@qlist{#1}}%
  \ifdim\dimexpr\ht\pp@qmeas+\dp\pp@qmeas\relax<0.95\textheight
    \needspace{\dimexpr\ht\pp@qmeas+\dp\pp@qmeas\relax}\else\needspace{14\baselineskip}\fi
  \zsavepos{q-s\thediq}\pp@qlist{#1}\par\zsavepos{q-e\thediq}\vspace{5ex plus 2ex minus 1ex}}
% part labels (a)/(i) bold, indented under the question text
\renewenvironment{qparts}[1][2.6em]{\begin{list}{}{\setlength\leftmargin{#1}%
  \setlength\labelwidth{\dimexpr#1-0.5em\relax}\setlength\labelsep{0.5em}%
  \setlength\itemsep{1.4ex}\setlength\topsep{1.2ex}\setlength\parsep{0.8ex}%
  \renewcommand\makelabel[1]{\textbf{##1}\hfil}}}{\end{list}}
\makeatother

% ---- cover page --------------------------------------------------------------
% A-level H2 style cover (9 Nov 2026; reference house_format/ref/alevel_cover_ref.png), no logos/barcodes.
% \prelimcover{total marks}{duration}{date}; heading block from \housechap; override with
% \prelimcoverhead (4 lines) and \prelimcoverline (the line under MATHEMATICS).
\makeatletter
\providecommand{\prelimcoverhead}{MAXIMUS PUNG JUN\\ H2 Mathematics (9758)\\ \housechap\\ Chapter Test}
\providecommand{\prelimcoverline}{\prelimsubject\ Test}
\newcommand{\pp@htwo}{\begin{tikzpicture}[baseline=(b.base)]
  \node[draw, line width=0.9pt, minimum width=19mm, minimum height=15mm, inner sep=0pt] (b)
    {\fontsize{38}{38}\selectfont\special{pdf:literal direct 1 Tr 0.7 w}H2\special{pdf:literal direct 0 Tr}};
  \end{tikzpicture}}
\newcommand{\prelimcover}[3]{% #1 total marks, #2 duration, #3 date line
\thispagestyle{prelimcover}%
{\coverfont\fontsize{11}{13.5}\selectfont\linespread{1}\selectfont\setlength{\parskip}{0pt}%
\vspace*{4mm}
\noindent\begin{tabular}{@{}l@{\hspace{5mm}}l@{}}\raisebox{-0.5\height}{\pp@htwo}&
  \raisebox{-0.5\height}{\begin{minipage}{120mm}\fontsize{12}{14.5}\selectfont\prelimcoverhead\end{minipage}}\end{tabular}\par
\vspace{9mm}
\noindent\rule{\linewidth}{0.8pt}\par\vspace{1.2mm}
\noindent{\fontsize{12}{14.5}\selectfont\bfseries MATHEMATICS\hfill 9758}\par\vspace{1.2mm}
\noindent\prelimcoverline\hfill\textbf{#3}\par\vspace{1.2mm}
\setbox0=\hbox{#2}\ifdim\wd0>0pt \noindent\hfill\textbf{#2}\par\fi
\vspace{3mm}
\noindent You must answer on writing paper.\par\vspace{3mm}
\noindent\begin{tabular}{@{}l@{\hspace{2mm}}l@{}}You will need:&Writing paper\\ &List of formulae and results (MF27)\end{tabular}\par
\vspace{2mm}
\noindent\rule{\linewidth}{0.8pt}\par\vspace{1.5mm}
\noindent{\bfseries INSTRUCTIONS}\par\vspace{0.5mm}
\begin{list}{\textbullet}{\setlength\leftmargin{7mm}\setlength\labelwidth{4mm}\setlength\labelsep{3mm}\setlength\itemsep{0.6mm}\setlength\topsep{0pt}\setlength\parsep{0pt}}
\item Answer \textbf{all} questions.
\item Give non-exact numerical answers correct to 3 significant figures, or 1 decimal place for angles in degrees, unless a different level of accuracy is specified in the question.
\item You should use an approved graphing calculator.
\item Unsupported answers from a graphing calculator are allowed unless a question specifically states otherwise.
\item Where unsupported answers from a graphing calculator are \textbf{not} allowed in a question, you must present the mathematical steps using mathematical notations and not calculator commands.
\item You must show all working clearly.
\end{list}
\vspace{6mm}
\noindent{\bfseries INFORMATION}\par\vspace{0.5mm}
\begin{list}{\textbullet}{\setlength\leftmargin{7mm}\setlength\labelwidth{4mm}\setlength\labelsep{3mm}\setlength\itemsep{0.6mm}\setlength\topsep{0pt}\setlength\parsep{0pt}}
\item The total mark for this paper is #1.
\item The number of marks for each question or part question is shown in brackets [\ ].
\end{list}
\vfill
\noindent\rule{\linewidth}{0.8pt}\par\vspace{1mm}
\centerline{This document has \textbf{\pageref{LastPage}} pages.}\par\vspace{2mm}
}\clearpage}
\makeatother

% ---- solutions header block (page 1 of the worked solutions; same cover style) ----
\newcommand{\prelimsolhead}[1]{% #1 total marks
{\coverfont\fontsize{10}{12}\selectfont\linespread{1}\selectfont
\noindent{\bfseries \prelimtitle}\\[0.2ex]{\bfseries Compiled by Maximus Pung Jun}\\[1.6ex]Higher 2\par\vspace{2ex}
\noindent\rule{\linewidth}{1.2pt}\par\vspace{0.4ex}
\noindent{\Large\bfseries MATHEMATICS}\hfill{\Large\bfseries 9758}\par\vspace{1.2ex}
\noindent \prelimsubject\hfill\textbf{2026}\\
WORKED SOLUTIONS\hfill\textbf{Total marks: #1}\par\vspace{1ex}
\noindent\rule{\linewidth}{0.6pt}\par}\vspace{4ex}}
\fi
  \ifx\housexm\housextestsol \def\prelimkeepcolour{1}% house_prelim.tex: HOUSE COPY of definite_integrals/common/prelim_style.tex; [HOUSE] titles from \housechapter
% =====================================================================
% prelim_style.tex: makes the Test look like a Singapore JC H2 Mathematics
% 9758 prelim paper (studied from ACJC/RI 2026 P1/P2 on the box:
% body Times New Roman + Times-style maths, cover in Arial).
% Fonts: prelim_fonts.tex (real Times New Roman/Arial + Times-style maths).
% Shared by Ch10 (di_style) and Ch11 (de_style; sets \prelimtitle/\prelimsubject). Loaded AFTER the chapter style (overrides its layout for the test only).
% =====================================================================
% [HOUSE] fonts loaded by house.tex
\providecommand{\prelimtitle}{\MakeUppercase{\housechap} TEST}
\providecommand{\prelimsubject}{\housechap}
\geometry{a4paper, top=22mm, bottom=24mm, left=25mm, right=25mm, headsep=6mm, footskip=12mm}
\setlength{\parindent}{0pt}
\setlength{\parskip}{0.9ex plus 0.3ex minus 0.2ex}
\linespread{1.3}
\setlength{\emergencystretch}{3em}
\clubpenalty=10000 \widowpenalty=10000 \interlinepenalty=10000
\AtBeginDocument{\fontsize{12}{14.4}\selectfont}
\renewcommand{\normalsize}{\fontsize{12}{14.4}\selectfont}

% ---- figures: plain black/grey like a printed paper ----------------------
\ifdefined\prelimkeepcolour\else % question paper: black/grey figures; solutions keep colour
\pgfplotsset{curve/.style={black, semithick, smooth, samples=150},
  line2/.style={black, semithick, samples=2},
  sketch/.append style={axis line style={thin, black}, every axis label/.append style={black}}}
\tikzset{region/.style={fill=black!15, draw=none}, regionB/.style={fill=black!30, draw=none},
  rect/.style={fill=black!8, draw=black, thin}, ax3/.style={-{Stealth[length=1.5mm]}, thin, black}}
\colorlet{blue}{black}\colorlet{red}{black}\colorlet{orange}{black}
\fi

% ---- header / footer: page number centred top (bold), [Turn over at bottom right
\makeatletter
\providecommand{\prelimfootlabel}{\prelimtitle\ 2026}
\newcommand{\pp@footer}{{\coverfont\scriptsize \prelimfootlabel\hspace{3em}H2 MATHEMATICS 9758}}
\fancypagestyle{prelim}{\fancyhf{}%
  \chead{\bfseries\thepage}%
  \lfoot{\pp@footer}%
  \rfoot{\ifnum\value{page}<\getpagerefnumber{LastPage}\relax\textbf{[Turn over}\fi}%
  \renewcommand{\headrulewidth}{0pt}\renewcommand{\footrulewidth}{0pt}}
\fancypagestyle{prelimcover}{\fancyhf{}%
  \rfoot{\coverfont\textbf{[Turn over}}%
  \renewcommand{\headrulewidth}{0pt}\renewcommand{\footrulewidth}{0pt}}
\pagestyle{prelim}

% ---- marks: [n] regular weight, right-aligned at the end of the part ----------
\renewcommand{\mk}[1]{{\unskip\nobreak\hfil\penalty50\hskip1.5em\hbox{}\nobreak\hfill[#1]%
  \parfillskip=0pt\finalhyphendemerits=0\par}}

% ---- questions: bold number at the left margin, text indented; grey [tag] straight after the number (3B)
\newsavebox{\pp@qmeas}
\newcommand{\pp@qlist}[1]{%
  \begin{list}{}{\setlength\leftmargin{2.6em}\setlength\labelwidth{2.1em}\setlength\labelsep{0.5em}%
     \setlength\topsep{0pt}\setlength\itemsep{0pt}\setlength\parsep{0.9ex}%
     \renewcommand\makelabel[1]{\textbf{##1}\hfil}}%
  \item[\thediq]\@endparpenalty\@M\setbox0=\hbox{\di@tag{#1}}\ifdim\wd0=0pt\else{\small\textcolor{black!60}{[\di@tag{#1}]}}\par\nobreak\fi\csname diQbody@#1\endcsname\par\penalty\@M
  \end{list}}
% a question starts on a new page unless it fits whole on the current one
% (questions taller than a page may break, but only between parts)
\renewcommand{\putQ}[1]{\stepcounter{diq}%
  \expandafter\xdef\csname diLab@#1\endcsname{\thediq}%
  \par\setbox\pp@qmeas=\vbox{\pp@qlist{#1}}%
  \ifdim\dimexpr\ht\pp@qmeas+\dp\pp@qmeas\relax<0.95\textheight
    \needspace{\dimexpr\ht\pp@qmeas+\dp\pp@qmeas\relax}\else\needspace{14\baselineskip}\fi
  \zsavepos{q-s\thediq}\pp@qlist{#1}\par\zsavepos{q-e\thediq}\vspace{5ex plus 2ex minus 1ex}}
% part labels (a)/(i) bold, indented under the question text
\renewenvironment{qparts}[1][2.6em]{\begin{list}{}{\setlength\leftmargin{#1}%
  \setlength\labelwidth{\dimexpr#1-0.5em\relax}\setlength\labelsep{0.5em}%
  \setlength\itemsep{1.4ex}\setlength\topsep{1.2ex}\setlength\parsep{0.8ex}%
  \renewcommand\makelabel[1]{\textbf{##1}\hfil}}}{\end{list}}
\makeatother

% ---- cover page --------------------------------------------------------------
% A-level H2 style cover (9 Nov 2026; reference house_format/ref/alevel_cover_ref.png), no logos/barcodes.
% \prelimcover{total marks}{duration}{date}; heading block from \housechap; override with
% \prelimcoverhead (4 lines) and \prelimcoverline (the line under MATHEMATICS).
\makeatletter
\providecommand{\prelimcoverhead}{MAXIMUS PUNG JUN\\ H2 Mathematics (9758)\\ \housechap\\ Chapter Test}
\providecommand{\prelimcoverline}{\prelimsubject\ Test}
\newcommand{\pp@htwo}{\begin{tikzpicture}[baseline=(b.base)]
  \node[draw, line width=0.9pt, minimum width=19mm, minimum height=15mm, inner sep=0pt] (b)
    {\fontsize{38}{38}\selectfont\special{pdf:literal direct 1 Tr 0.7 w}H2\special{pdf:literal direct 0 Tr}};
  \end{tikzpicture}}
\newcommand{\prelimcover}[3]{% #1 total marks, #2 duration, #3 date line
\thispagestyle{prelimcover}%
{\coverfont\fontsize{11}{13.5}\selectfont\linespread{1}\selectfont\setlength{\parskip}{0pt}%
\vspace*{4mm}
\noindent\begin{tabular}{@{}l@{\hspace{5mm}}l@{}}\raisebox{-0.5\height}{\pp@htwo}&
  \raisebox{-0.5\height}{\begin{minipage}{120mm}\fontsize{12}{14.5}\selectfont\prelimcoverhead\end{minipage}}\end{tabular}\par
\vspace{9mm}
\noindent\rule{\linewidth}{0.8pt}\par\vspace{1.2mm}
\noindent{\fontsize{12}{14.5}\selectfont\bfseries MATHEMATICS\hfill 9758}\par\vspace{1.2mm}
\noindent\prelimcoverline\hfill\textbf{#3}\par\vspace{1.2mm}
\setbox0=\hbox{#2}\ifdim\wd0>0pt \noindent\hfill\textbf{#2}\par\fi
\vspace{3mm}
\noindent You must answer on writing paper.\par\vspace{3mm}
\noindent\begin{tabular}{@{}l@{\hspace{2mm}}l@{}}You will need:&Writing paper\\ &List of formulae and results (MF27)\end{tabular}\par
\vspace{2mm}
\noindent\rule{\linewidth}{0.8pt}\par\vspace{1.5mm}
\noindent{\bfseries INSTRUCTIONS}\par\vspace{0.5mm}
\begin{list}{\textbullet}{\setlength\leftmargin{7mm}\setlength\labelwidth{4mm}\setlength\labelsep{3mm}\setlength\itemsep{0.6mm}\setlength\topsep{0pt}\setlength\parsep{0pt}}
\item Answer \textbf{all} questions.
\item Give non-exact numerical answers correct to 3 significant figures, or 1 decimal place for angles in degrees, unless a different level of accuracy is specified in the question.
\item You should use an approved graphing calculator.
\item Unsupported answers from a graphing calculator are allowed unless a question specifically states otherwise.
\item Where unsupported answers from a graphing calculator are \textbf{not} allowed in a question, you must present the mathematical steps using mathematical notations and not calculator commands.
\item You must show all working clearly.
\end{list}
\vspace{6mm}
\noindent{\bfseries INFORMATION}\par\vspace{0.5mm}
\begin{list}{\textbullet}{\setlength\leftmargin{7mm}\setlength\labelwidth{4mm}\setlength\labelsep{3mm}\setlength\itemsep{0.6mm}\setlength\topsep{0pt}\setlength\parsep{0pt}}
\item The total mark for this paper is #1.
\item The number of marks for each question or part question is shown in brackets [\ ].
\end{list}
\vfill
\noindent\rule{\linewidth}{0.8pt}\par\vspace{1mm}
\centerline{This document has \textbf{\pageref{LastPage}} pages.}\par\vspace{2mm}
}\clearpage}
\makeatother

% ---- solutions header block (page 1 of the worked solutions; same cover style) ----
\newcommand{\prelimsolhead}[1]{% #1 total marks
{\coverfont\fontsize{10}{12}\selectfont\linespread{1}\selectfont
\noindent{\bfseries \prelimtitle}\\[0.2ex]{\bfseries Compiled by Maximus Pung Jun}\\[1.6ex]Higher 2\par\vspace{2ex}
\noindent\rule{\linewidth}{1.2pt}\par\vspace{0.4ex}
\noindent{\Large\bfseries MATHEMATICS}\hfill{\Large\bfseries 9758}\par\vspace{1.2ex}
\noindent \prelimsubject\hfill\textbf{2026}\\
WORKED SOLUTIONS\hfill\textbf{Total marks: #1}\par\vspace{1ex}
\noindent\rule{\linewidth}{0.6pt}\par}\vspace{4ex}}
% house_sol.tex: HOUSE COPY of definite_integrals/common/ri_solstyle.tex (unchanged)
% =====================================================================
% ri_solstyle.tex: worked solutions in the layout of the RI 2026 Year 6 Prelim
% "Solutions with comments" (grail.moe 26612/26613), loaded AFTER prelim_style.
%   * page 1: bold school-style block + bold underlined title; later pages: running
%     title top-left with a rule under it; footer "Page N of M" bottom right (Calibri).
%   * each question: prelim-format question text (bold number, marks right, source tag
%     bottom right), then a bordered table:  Solution | Comments  (bold header row);
%     body rows  label + [marks] (bold, narrow left column) | working | comments,
%     a full-width rule between parts, working left-aligned (fleqn).
% Rows come from the existing sources: \pt in qparts (Ch10 HW), \slabel (Ch10 Test),
% spart (Ch11). Marks under each label are taken from the question's \mk.
% =====================================================================
\makeatletter
\usepackage{lastpage}
\usepackage{adjustbox}
\usepackage{zref-savepos,zref-abspage}
\zref@addprop{savepos}{abspage}
\geometry{a4paper, top=22mm, bottom=24mm, left=22mm, right=22mm, headsep=8mm, footskip=12mm, headheight=24pt}
\providecommand{\risoltitle}{\prelimtitle: Solutions with comments}
\providecommand{\risolschool}{Compiled by Maximus Pung Jun}
\fancypagestyle{risol}{\fancyhf{}%
  \lhead{\fontsize{11}{12}\selectfont 2026 \risoltitle\\[-1.2ex]\rule{\textwidth}{0.6pt}}%
  \rfoot{\calibrifont\fontsize{11}{13}\selectfont Page \textbf{\thepage} of \textbf{\pageref*{LastPage}}}%
  \renewcommand{\headrulewidth}{0pt}\renewcommand{\footrulewidth}{0pt}}
\fancypagestyle{risolfirst}{\fancyhf{}%
  \rfoot{\calibrifont\fontsize{11}{13}\selectfont Page \textbf{\thepage} of \textbf{\pageref*{LastPage}}}%
  \renewcommand{\headrulewidth}{0pt}\renewcommand{\footrulewidth}{0pt}}
\pagestyle{risol}
\newif\ifri@fresh
\newcommand{\risolfront}{\thispagestyle{risolfirst}\global\ri@freshtrue%
  {\noindent\bfseries\linespread{1}\selectfont \risolschool\\ H2 Mathematics (9758)\\ 2026\par}%
  \vspace{4ex}{\noindent\bfseries\underline{2026 \risoltitle}\par}\vspace{3ex}}
% ---- table geometry (RI: label 42pt | working 284pt | comments 142pt of 468pt) ----
\newlength{\ri@L}\setlength{\ri@L}{15mm}
\newlength{\ri@C}\setlength{\ri@C}{40mm}
\newlength{\ri@pad}\setlength{\ri@pad}{1.8mm}
\newif\ifri@sol \newif\ifri@first
\newcount\ri@depth
\def\ri@outer{}
\newtcolorbox{risol}{enhanced, breakable, sharp corners, frame hidden, colback=white,
  boxrule=0pt, boxsep=0pt, left=\dimexpr\ri@L+\ri@pad\relax, right=\dimexpr\ri@C+\ri@pad\relax,
  top=1.2mm, bottom=2.5mm, toptitle=0.6mm, bottomtitle=0.6mm, lefttitle=\ri@pad, righttitle=0pt,
  colbacktitle=white, coltitle=black, fonttitle=\bfseries,
  title={Solution\hfill\makebox[\ri@C][l]{\hspace{\ri@pad}Comments}},
  before skip=2ex, after skip=2ex,
  overlay unbroken and first={\draw[line width=0.5pt] (frame.south west) rectangle (frame.north east);
     \draw[line width=0.5pt] (title.south west) -- (title.south east);
     \draw[line width=0.5pt] ([xshift=\ri@L]interior.north west) -- ([xshift=\ri@L]frame.south west);
     \draw[line width=0.5pt] ([xshift=-\ri@C]frame.north east) -- ([xshift=-\ri@C]frame.south east);},
  overlay middle and last={\draw[line width=0.5pt] (frame.south west) rectangle (frame.north east);
     \draw[line width=0.5pt] ([xshift=\ri@L]frame.north west) -- ([xshift=\ri@L]frame.south west);
     \draw[line width=0.5pt] ([xshift=-\ri@C]frame.north east) -- ([xshift=-\ri@C]frame.south east);},
  after upper={\ri@endrow\protected@write\@auxout{}{\string\ri@settotal{\ri@qid}{\the\ri@sumH}}},
  before upper={\global\ri@sumH=14mm\ri@soltrue\ri@firsttrue\ri@depth=0 \parindent=0pt \parskip=0.6ex
     \predisplaypenalty=100 
     \linespread{1.12}\selectfont\raggedright\@fleqntrue\@mathmargin=0pt
     \abovedisplayskip=5pt plus 2pt\belowdisplayskip=5pt plus 2pt
     \abovedisplayshortskip=3pt plus 1pt\belowdisplayshortskip=4pt plus 2pt
     \ifcsname ri@hp@\ri@qid\endcsname\else\ri@row{}\fi}}
\def\ri@qid{}
% one row: full-width rule (glue leaders, so it vanishes at a page break where the frame
% edge takes its place; not before the first row), then bold label + bold [marks]
% in the left column, on the baseline of the first line of the working
\ExplSyntaxOn
\tl_new:N \g__ri_tmp_tl
\cs_new_protected:Npn \ri_setkey:Nn #1#2
  { \group_begin: \cs_set_eq:NN \chg \use_ii:nn \cs_set_eq:NN \chgm \use_ii:nn
    \cs_set_eq:NN \textbf \use:n \cs_set_eq:NN \mathrm \use:n
    \tl_gset:Nx \g__ri_tmp_tl {#2} \group_end:
    \tl_set_eq:NN #1 \g__ri_tmp_tl \tl_remove_all:Nn #1 {~} }
\cs_set_eq:NN \ri@setkey \ri_setkey:Nn
\ExplSyntaxOff
% ---- page-break rules (positions from the previous run, zref labels in the .aux) ----
% * a row (one part) never breaks across pages: if it does not fit on the rest of the
%   page it starts the next page (forced break between rows); only a row taller than a
%   whole page may break inside;
% * a whole solution that fits on one page is kept together (needspace before the table);
%   a longer one starts where at least its first row fits;
% * a row is at least as tall as its label + [marks].
\newcount\ri@n \newbox\ri@labbox \def\ri@open{0}
\newdimen\ri@ybot \newdimen\ri@ytop \newdimen\ri@h \newdimen\ri@sumH
\AtBeginDocument{\ri@ybot=\dimexpr\paperheight-1in-\voffset-\topmargin-\headheight-\headsep-\textheight\relax
  \ri@ytop=\dimexpr\ri@ybot+\textheight\relax}
\def\ri@pg#1{\zref@extractdefault{#1}{abspage}{-1}}
% \ri@rowh{n}: height of row n in the previous run -> \ri@h (0pt if unknown)
\newcommand{\ri@rowh}[1]{\ri@h=0pt
  \ifnum\ri@pg{ri-s#1}>0 \ifnum\ri@pg{ri-e#1}>0
    \ifnum\ri@pg{ri-s#1}=\ri@pg{ri-e#1}\relax \ri@h=\dimexpr\zposy{ri-s#1}sp-\zposy{ri-e#1}sp\relax
    \else \ri@h=\dimexpr\zposy{ri-s#1}sp-\ri@ybot+\ri@ytop-\zposy{ri-e#1}sp-6mm\relax\fi\fi\fi}
\newcommand{\ri@setfirst}[2]{\expandafter\gdef\csname ri@first@#1\endcsname{#2}}
\newcommand{\ri@setkept}[1]{\expandafter\gdef\csname ri@kept@#1\endcsname{}}
\newcommand{\ri@settotal}[2]{\expandafter\gdef\csname ri@total@#1\endcsname{#2}}
\newcommand{\ri@endrow}{\ifnum\ri@open=1 \par\zsavepos{ri-e\the\ri@n}%
  \ifnum\ri@pg{ri-s\the\ri@n}=\ri@pg{ri-e\the\ri@n}\relax
    \dimen@=\dimexpr\ri@curH sp+1.6mm-\zposy{ri-s\the\ri@n}sp+\zposy{ri-e\the\ri@n}sp\relax
    \ifdim\dimen@>0pt \nobreak\kern\dimen@\penalty\z@\fi\fi
  \ri@rowh{\the\ri@n}\global\advance\ri@sumH\dimexpr\ri@h+3.5mm\relax
  \gdef\ri@open{0}\fi}
\newcommand{\ri@row}[1]{\par\ri@endrow\global\advance\ri@n\@ne\gdef\ri@open{1}\@beginparpenalty\@M
  \ifri@first\global\ri@firstfalse\ri@firstfalse
    \protected@write\@auxout{}{\string\ri@setfirst{\ri@qid}{\the\ri@n}}%
  \else
    \zsavepos{ri-p\the\ri@n}%
    % push the row to the next page if it would not fit on the rest of this one
    \ri@rowh{\the\ri@n}%
    \ifnum\ri@pg{ri-p\the\ri@n}>0 \ifdim\ri@h>0pt
      \ifdim\dimexpr\ri@h+4mm\relax<\dimexpr\textheight-10mm\relax
        \ifdim\dimexpr\ri@h+7mm\relax>\dimexpr\zposy{ri-p\the\ri@n}sp-\ri@ybot\relax
          \penalty-\@M\fi\fi\fi\fi
    \vskip0.6ex\nointerlineskip
    \leaders\hbox{\hskip-\dimexpr\ri@L+\ri@pad\relax\vrule width\dimexpr\linewidth+\ri@L+\ri@C+2\ri@pad\relax height0.5pt depth0pt}\vskip0.5pt
    \nobreak\vskip1.4mm\fi
  \ri@setkey\ri@key{\ri@qid#1}\def\ri@lab{#1}%
  \setbox\ri@labbox=\vtop{\hsize=\dimexpr\ri@L-1.6mm\relax\bfseries\raggedright\parskip=0pt
      \strut\ri@lab\ifcsname ri@mk@\ri@key\endcsname\ifx\ri@lab\@empty\else\\\fi
      [\csname ri@mk@\ri@key\endcsname]\fi\strut\par}%
  \xdef\ri@curH{\number\dimexpr\ht\ri@labbox+\dp\ri@labbox\relax}%
  \nointerlineskip\zsavepos{ri-s\the\ri@n}\nointerlineskip\nobreak
  \vbox to 0pt{\hbox{\hskip-\dimexpr\ri@L+\ri@pad-1.2mm\relax\box\ri@labbox}\vss}\nobreak\prevdepth=\dp\strutbox
  \vskip-\parskip\ignorespaces}
% ---- qparts / \pt: question mode records labels (for the marks); solution mode makes rows
\NewCommandCopy\ri@qparts@begin\qparts \let\ri@qparts@end\endqparts
\renewenvironment{qparts}[1][2.6em]{\advance\ri@depth\@ne
  \ifri@sol\else\ri@qparts@begin[#1]\fi}{\ifri@sol\par\else\ri@qparts@end\fi}
\def\ri@qpartsname{qparts}
\renewcommand{\pt}[1]{\ifnum\ri@depth<2 \gdef\ri@outer{#1}\gdef\ri@cur{#1}\else\expandafter\gdef\expandafter\ri@cur\expandafter{\ri@outer#1}\fi
  \ifri@sol\expandafter\ri@solpt\else\expandafter\ri@qpt\fi{#1}}
\def\ri@qpt#1{\expandafter\gdef\csname ri@hp@\ri@qid\endcsname{}\item[#1]}
\def\ri@solpt#1{\@ifnextchar\begin\ri@ptpeek{\ri@row{\ri@cur}}}
\def\ri@ptpeek\begin#1{\def\ri@tmp{#1}\ifx\ri@tmp\ri@qpartsname
   \expandafter\@firstoftwo\else\expandafter\@secondoftwo\fi
  {\begin{qparts}}{\ri@row{\ri@cur}\begin{#1}}}
% marks: record [n] against the current question/part
\let\ri@mk\mk
\renewcommand{\mk}[1]{\ri@setkey\ri@k{\ri@qid\ri@cur}\ri@setkey\ri@mv{#1}\expandafter\xdef\csname ri@mk@\ri@k\endcsname{\ri@mv}\ri@mk{#1}}
% \slabel (Ch10 Test answers): letter labels set the outer part, roman ones are appended;
% consecutive \slabel calls merge into one row label
\def\ri@romans{(i)(ii)(iii)(iv)(v)(vi)(vii)}
\def\ri@firstpart(#1)#2\@nil{(#1)}
\def\slabel#1{\ifri@sol\expandafter\ri@solslabel\else\expandafter\ri@qslabel\fi{#1}}
\def\ri@qslabel#1{\par\smallskip\noindent\textbf{#1}\par\nobreak}
\def\ri@solslabel#1{\@expandtwoargs\in@{#1}{\ri@romans}%
  \ifin@\expandafter\gdef\expandafter\ri@cur\expandafter{\ri@outer#1}\else\gdef\ri@cur{#1}\xdef\ri@outer{\ri@firstpart#1\@nil}\fi
  \@ifnextchar\slabel{}{\ri@row{\ri@cur}}}
% Ch11 spart{label}: a row; soln: the table
\def\spart#1{\gdef\ri@cur{#1}\ri@row{#1}}\def\endspart{\par}
\def\soln{\@ifnextchar[{\ri@solnopt}{\risol}}\def\ri@solnopt[#1]{\risol}\def\endsoln{\endrisol}
% ---- working that started on the label line is left-aligned in the working column
\AtBeginDocument{\@ifundefined{topalign}{}{%
  \renewenvironment{topalign}{\par\noindent$\displaystyle\begin{aligned}[t]}{\end{aligned}$\par\smallskip}%
  \renewcommand{\topdisplay}[1]{\par\noindent$\displaystyle #1$\par\smallskip}}}
% ---- Ch11 side-by-side (text | sketch): the sketch shrinks to fit its half of the working column
\AtBeginDocument{\@ifundefined{sbsplain}{}{%
  \renewenvironment{sbsplain}[1]{%
    \sbs@start\begin{minipage}[t]{#1\linewidth}\vspace{0pt}\sbs@shift\raggedright%
    \def\sbsnext{\end{minipage}\hfill\begin{minipage}[t]{\dimexpr0.97\linewidth-#1\linewidth\relax}\vspace{0pt}\sbs@shift\centering
      \begin{adjustbox}{max width=\linewidth}}%
  }{\end{adjustbox}\end{minipage}\par}}}
% ---- figures in the working column shrink to fit it (only if wider), colours kept
\newcommand{\rifitfig}[1]{\expandafter\let\csname ri@orig@\string#1\endcsname#1%
  \def#1{\@ifnextchar\bgroup{\ri@fitarg#1}{\ri@fitnoarg#1}}}
\def\ri@fitarg#1#2{\ifri@sol\adjustbox{max width=\linewidth}{\csname ri@orig@\string#1\endcsname{#2}}\else\csname ri@orig@\string#1\endcsname{#2}\fi}
\def\ri@fitnoarg#1{\ifri@sol\adjustbox{max width=\linewidth}{\csname ri@orig@\string#1\endcsname}\else\csname ri@orig@\string#1\endcsname\fi}
\@ifundefined{solfig}{}{\renewcommand{\solfig}[2][]{\par\begin{center}\adjustbox{max width=\linewidth}{\begin{tikzpicture}\begin{axis}[qfig, width=6.6cm, height=5cm, #1]#2\end{axis}\end{tikzpicture}}\end{center}}}
% \ricomment{text}: text in the Comments column of the current row (call right after \pt,
% \slabel or \spart). The row is made at least as tall as the comment.
\newbox\ri@cmtbox
\newcommand{\ricomment}[1]{\par\setbox\ri@cmtbox=\vtop{\hsize=\dimexpr\ri@C-2\ri@pad\relax
    \parskip=0pt\raggedright\small{\color{black!60}\strut#1\strut\par}}%
  \ifdim\dimexpr\ht\ri@cmtbox+\dp\ri@cmtbox\relax>\ri@curH sp\xdef\ri@curH{\number\dimexpr\ht\ri@cmtbox+\dp\ri@cmtbox\relax}\fi
  \nointerlineskip\vbox to 0pt{\hbox to\linewidth{\hskip\dimexpr\linewidth+2\ri@pad\relax\box\ri@cmtbox\hss}\vss}\nobreak\prevdepth=\dp\strutbox
  \vskip-\parskip\ignorespaces}
% \putTFlong{n}: the True/False question n with its answer table allowed to break across
% pages (the body of \defQ{n} must be: intro text, then \begin{tftable}...\end{tftable});
% the solutions copy prints it with tflongtable.
\renewcommand{\putTFlong}[1]{\ri@newq\begingroup\let\tftable\tflongtable\let\endtftable\endtflongtable
  \stepcounter{diq}\expandafter\xdef\csname diLab@#1\endcsname{\thediq}%
  \par\noindent\hbox to 2.6em{\textbf{\thediq}\hfil}\setbox0=\hbox{\di@tag{#1}}\ifdim\wd0=0pt\else{\small\textcolor{black!60}{[\di@tag{#1}]}}\par\noindent\hskip2.6em\fi\ignorespaces\csname diQbody@#1\endcsname\par
  \endgroup\vspace{3ex}}
% ---- question + solution; each question starts on a new page (as RI)
\def\ri@qid{}
\let\ri@putQ\putQ
\renewcommand{\putQ}[1]{\gdef\ri@qid{#1}\global\ri@depth=0 \gdef\ri@cur{}\ri@putQ{#1}}
% \risolQ{n}{body macro name}: question n + its solution body in the table;
% \risolQraw: body already contains its own soln environment (Ch11 Test)
\newcommand{\ri@newq}{\ifri@fresh\global\ri@freshfalse\else\clearpage\fi}
% keep the whole table together if it fits on a page; otherwise make sure at least its
% header and first row fit before starting it
\newcommand{\ri@keep}[1]{\par
  \ifcsname ri@total@#1\endcsname
    % hysteresis: a table kept whole in the previous run stays whole up to textheight-4.5mm,
    % so a total near the limit cannot flip between runs
    \ifdim\csname ri@total@#1\endcsname<\dimexpr\textheight-\ifcsname ri@kept@#1\endcsname4.5\else8\fi mm\relax
      \protected@write\@auxout{}{\string\ri@setkept{#1}}%
      \needspace{\dimexpr\csname ri@total@#1\endcsname+4mm\relax}%
    \else\ifcsname ri@first@#1\endcsname
      \ri@rowh{\csname ri@first@#1\endcsname}%
      \ifdim\ri@h<\dimexpr\textheight-20mm\relax\needspace{\dimexpr\ri@h+18mm\relax}\fi\fi\fi\fi}
\newcommand{\risolQ}[2]{\ri@newq\putQ{#1}\vspace{-3ex}\ri@keep{#1}\begin{risol}\csname #2@#1\endcsname\par\end{risol}}
\newcommand{\risolQraw}[2]{\ri@newq\putQ{#1}\vspace{-3ex}\ri@keep{#1}\csname #2@#1\endcsname\par}
\renewcommand{\putQS}[1]{\risolQ{#1}{diS}}
% end marker only if it fits on the last page (no page holding just the marker)
\renewcommand{\diend}[1][END OF SOLUTIONS]{\par\ifdim\dimexpr\pagegoal-\pagetotal\relax<4\baselineskip\else\vspace{1.5ex}\begin{center}\textbf{#1}\end{center}\fi}
\makeatother
\tfanswerstrue\fi
\fi
):
%   notes: \documentclass[a4paper]{book}
%   others: \documentclass[a4paper,11pt,oneside]{book}
% A HW file uses \providecommand\housemode{hw}; its solutions file has NO \documentclass:
%   \def\housemode{sol}\def\diquestionsonly{1}\input{HW1}  (as DI/DE HWn_solutions.tex)
% Compile: TEXINPUTS=.:/workspace/house_format: xelatex (2 runs; solutions until .aux stable)
% =====================================================================
\providecommand\housemode{hw}
\providecommand\housedir{/workspace/house_format}
\newcommand\housechap{}
\newcommand\housechapter[1]{\renewcommand\housechap{#1}}
\def\housexnotes{notes}\def\housexhw{hw}\def\housexsol{sol}\def\housextest{test}\def\housextestsol{testsol}
\edef\housexm{\housemode}
% house_base.tex: HOUSE COPY of definite_integrals/common/ws_preamble_build.tex
% (= Maximus's own preamble.tex + the DE v3 / AoI additions). Only change: GC font path. [HOUSE]
% WORKSHEET COPY of ../notes/preamble_build.tex (2026-10-09 version), unchanged.
% preamble_build.tex: copy of Maximus's preamble.tex for the DE v3 build.
% The blocks marked [BUILD COPY CHANGE] are the earlier build fixes.
% The blocks marked [V3 ADDITION] are new in this version (example box,
% spacing). Everything else is unchanged from preamble.tex.

% For margins
\usepackage[a4paper, margin = 20mm]{geometry}

\usepackage[Lenny]{fncychap}
% Rejne, Glenn, Lenny

\usepackage{etoc}

% [V3 ADDITION] cancel: used to show the dx's "cancelling" in Section 11.3
\usepackage{cancel}
\renewcommand{\CancelColor}{\color{red}}

% For enumeration & math symbols
\usepackage{enumerate}
\usepackage{amsmath,amsfonts,amssymb,amsthm,graphicx,lipsum}

% For diagrams, plots and boxes
\usepackage{pgfplots}
\pgfplotsset{compat=1.18}
\usepackage{tikz}
\usepackage[most]{tcolorbox}
\tcbuselibrary{skins}
\usepackage{xcolor}
\pgfplotsset{
    standard/.style={clip = false,
    axis line style = thick,
    trig format=rad,
    axis lines = middle,
    every axis x label/.style={at={(current axis.right of origin)},anchor= west},
    every axis y label/.style={at={(current axis.above origin)},anchor=south}
    }
}

% For GC font
\usepackage{fontspec}
% [BUILD COPY CHANGE 1] The TI-84 key font is loaded only if the file exists
% (checked in ./Fonts/... and ../Fonts/...). Otherwise \ti{} falls back to a
% small boxed sans-serif key. Known letters of the TI font are mapped to their
% key names in the fallback: o = [y=], z = [mode], f = [L3].
\newcommand{\tikey}[1]{\tcbox[on line, enhanced, boxsep=0pt, left=2pt, right=2pt,
  top=1pt, bottom=1pt, boxrule=0.5pt, arc=1.5pt, colback=white,
  colframe=black!70]{\sffamily\footnotesize\strut #1}}
\ExplSyntaxOn
\NewDocumentCommand{\tifallback}{m}{
  \str_case:nnF {#1}
    { {o}{\tikey{y=}} {z}{\tikey{mode}} {f}{\tikey{L3}} }
    { \tikey{#1} }
}
\ExplSyntaxOff
\def\tifontcase{0}
\IfFileExists{Fonts/TI-84 Plus CE Key Fonts/TI84PlusCEKeys.ttf}
  {\def\tifontcase{1}}{}
\IfFileExists{../Fonts/TI-84 Plus CE Key Fonts/TI84PlusCEKeys.ttf}
  {\def\tifontcase{2}}{}
\ifcase\tifontcase\relax
  % 0: font not found, use boxed fallback keys
  \newcommand{\ti}[1]{\tifallback{#1}}
\or
  % 1: original location (as in preamble.tex)
  \newfontfamily\tifont[
    Path = {Fonts/TI-84 Plus CE Key Fonts/},
    UprightFont = {TI84PlusCEKeys.ttf}
  ]{}
  \newcommand{\ti}[1]{{\tifont #1}}
\or
  % 2: one level up (when compiling inside draft-ch1/)
  \newfontfamily\tifont[
    Path = {../Fonts/TI-84 Plus CE Key Fonts/},
    UprightFont = {TI84PlusCEKeys.ttf}
  ]{}
  \newcommand{\ti}[1]{{\tifont #1}}
\fi

% [AoI items 17/18/37] GC DISPLAY FONT. Every GC key, command, menu item and
% screen line goes through \gcfont, so the font can be swapped HERE in one place.
%   \gc{fnInt(}      command / menu item / stored variable in the display font
%   \gckey{2nd}      a calculator key (display font in a small key outline)
%   \gcscreen{...}   a calculator screen (display font on a grey-green LCD panel)
% The font file lives in ../Fonts (shared by notes, slides and worksheet).
\def\gcfontfile{VT323-Regular.ttf}
% [HOUSE] GC font: always the absolute copy in \housedir/Fonts (relative Path= breaks under TEXINPUTS)
\providecommand\housedir{/workspace/house_format}
\IfFileExists{\housedir/Fonts/\gcfontfile}{\newfontfamily\gcfont[Path=\housedir/Fonts/, Scale=1.32]{\gcfontfile}}{%
  \newcommand\gcfont{\ttfamily}}
\newcommand{\gc}[1]{\mbox{\gcfont #1}}
\newcommand{\gckey}[1]{\tcbox[on line, enhanced, boxsep=0pt, left=2pt, right=2pt,
  top=0.6pt, bottom=0.2pt, boxrule=0.5pt, arc=1.5pt, colback=black!4,
  colframe=black!65]{\gcfont\strut #1}}
\definecolor{lcd}{RGB}{214,222,204}
\newcommand{\gcscreen}[1]{\tcbox[on line, enhanced, boxsep=1pt, left=4pt, right=4pt,
  top=2pt, bottom=2pt, boxrule=1.2pt, arc=1pt, colback=lcd, colframe=black!70]{%
  \gcfont\begin{tabular}{@{}l@{}}#1\end{tabular}}}
% integral template as shown on the TI-84 Plus CE (MathPrint): \gcint{0}{2}{e^(-X²)}{X}
\newcommand{\gcint}[4]{{\gcfont\raisebox{-0.35ex}{\Large ∫}\hspace{-0.15em}%
  \raisebox{-0.2ex}{\begin{tabular}{@{}l@{}}\footnotesize #2\\[-1.3ex]\footnotesize #1\end{tabular}}%
  \hspace{0.1em}(#3)\mathrm{d}#4}}
% old key macro kept as an alias so nothing breaks
\renewcommand{\ti}[1]{\gckey{#1}}

% SI units
\usepackage{siunitx}

% setspace
%\usepackage{setspace}
%\onehalfspacing

% For hyperlinks within and out of document
\usepackage{hyperref}
\hypersetup{
    colorlinks = true,
    linkcolor = black,
    urlcolor = blue,
}   

% For captions for pics
\usepackage{capt-of}


% tcolorbox custom environments
% Definition ========================================================
\newtcolorbox[auto counter,number within=section]{definition}[2][drop shadow]{%
enhanced, before skip=5mm, after skip=6mm, % [V3 ADDITION] spacing
colback=blue!5!white,colframe=blue!50!black,fonttitle=\bfseries,
title= Definition~\thetcbcounter. #2,#1}

% Theorem ==========================================================
\newtcolorbox[auto counter,number within=section]{theorem}[2][drop shadow]{%
enhanced, before skip=5mm, after skip=6mm, % [V3 ADDITION] spacing
colback=black!5!white,colframe=black,fonttitle=\bfseries,
title= Theorem~\thetcbcounter. #2,#1}

% Caution box ======================================================
\newtcolorbox{caution}[1][]{enhanced,
before skip=5mm,after skip=6mm, % [V3 ADDITION] was 2mm/3mm
boxrule=1pt,left=11mm,right=2mm,top=3mm,bottom=3mm,
colback=yellow!50,
colframe=yellow!20!black,
sharp corners,rounded corners=southeast,arc is angular,arc=3mm,
underlay={%
\path[fill=tcbcolback!80!black] ([yshift=3mm]interior.south east)--++(-0.4,-0.1)--++(0.1,-0.2);
\path[draw=tcbcolframe,shorten <=-0.05mm,shorten >=-0.05mm] ([yshift=3mm]interior.south east)--++(-0.4,-0.1)--++(0.1,-0.2);
\path[fill=yellow!50!black,draw=none] (interior.south west) rectangle node[white]{\Huge\bfseries !} ([xshift=9mm]interior.north west);
},
drop fuzzy shadow,#1}

% Question box ======================================================
\newtcolorbox{question}[1][]{enhanced,
before skip=5mm,after skip=6mm, % [V3 ADDITION] was 2mm/3mm
boxrule=1pt,left=11mm,right=2mm,top=3mm,bottom=3mm,
colback=blue!20!white,
colframe=blue!20!black,
sharp corners,rounded corners=southeast,arc is angular,arc=3mm,
underlay={%
\path[fill=tcbcolback!80!black] ([yshift=3mm]interior.south east)--++(-0.4,-0.1)--++(0.1,-0.2);
\path[draw=tcbcolframe,shorten <=-0.05mm,shorten >=-0.05mm] ([yshift=3mm]interior.south east)--++(-0.4,-0.1)--++(0.1,-0.2);
\path[fill=blue!50!black,draw=none] (interior.south west) rectangle node[white]{\Huge\bfseries ?} ([xshift=9mm]interior.north west);
},
drop fuzzy shadow,#1}

% Information box ======================================================
\newtcolorbox{info}[1][]{enhanced,
before skip=5mm,after skip=6mm, % [V3 ADDITION] was 2mm/3mm
boxrule=1pt,left=11mm,right=2mm,top=3mm,bottom=3mm,
colback=cyan!20!white,
colframe=cyan!20!black,
sharp corners,rounded corners=southeast,arc is angular,arc=3mm,
underlay={%
\path[fill=tcbcolback!80!black] ([yshift=3mm]interior.south east)--++(-0.4,-0.1)--++(0.1,-0.2);
\path[draw=tcbcolframe,shorten <=-0.05mm,shorten >=-0.05mm] ([yshift=3mm]interior.south east)--++(-0.4,-0.1)--++(0.1,-0.2);
\path[fill=cyan!50!black,draw=none] (interior.south west) rectangle node[white]{\Huge\bfseries i} ([xshift=9mm]interior.north west);
},
drop fuzzy shadow,#1}

% Mistake box ======================================================
\newtcolorbox{mistake}[1][]{enhanced,
before skip=5mm,after skip=6mm, % [V3 ADDITION] was 2mm/3mm
boxrule=1pt,left=11mm,right=2mm,top=3mm,bottom=3mm,
colback=red!20!white,
colframe=red!20!black,
sharp corners,rounded corners=southeast,arc is angular,arc=3mm,
underlay={%
\path[fill=tcbcolback!80!black] ([yshift=3mm]interior.south east)--++(-0.4,-0.1)--++(0.1,-0.2);
\path[draw=tcbcolframe,shorten <=-0.05mm,shorten >=-0.05mm] ([yshift=3mm]interior.south east)--++(-0.4,-0.1)--++(0.1,-0.2);
\path[fill=red!50!black,draw=none] (interior.south west) rectangle node[white]{\Huge\bfseries $\times$} ([xshift=9mm]interior.north west);
},
drop fuzzy shadow,#1}



% [V3 ADDITION] Worked example box ==================================
% Own counter, numbered within the section (Example 11.2.1, 11.2.2, ...).
% For per-subsection numbering (11.2.1.1) change `section' to `subsection'. Blue `question' boxes are kept for thought-provoking questions
% only; worked examples always use this box. Usage:
%   \begin{example} ...question text... \end{example}
%   \solution{...}
\newtcolorbox[auto counter,number within=section]{example}[1][]{%
enhanced, breakable=false, before skip=5mm, after skip=4mm,
colback=teal!5!white, colframe=teal!55!black, boxrule=0.8pt,
left=3mm, right=3mm, top=2mm, bottom=2mm,
sharp corners, rounded corners=northwest, arc=3mm,
fonttitle=\bfseries, coltitle=white,
title=Example~\thetcbcounter, attach boxed title to top left={yshift=-2mm, xshift=3mm},
boxed title style={colback=teal!55!black, sharp corners, rounded corners=northwest, arc=2mm},
drop fuzzy shadow, #1}

% [V3 ADDITION] Spacing between sections, subsections, parts, paragraphs =
% (fncychap's chapter heading is untouched; only \section and below.)
\makeatletter
\renewcommand\section{\@startsection{section}{1}{\z@}%
  {-2ex \@plus -.5ex \@minus -.2ex}{3.2ex \@plus .3ex}%
  {\normalfont\LARGE\bfseries\color{blue!50!black}}}
\renewcommand\subsection{\@startsection{subsection}{2}{\z@}%
  {-1.5ex \@plus -.5ex \@minus -.2ex}{2.6ex \@plus .3ex}%
  {\normalfont\Large\bfseries}}
\renewcommand\subsubsection{\@startsection{subsubsection}{3}{\z@}%
  {-4ex \@plus -1ex \@minus -.5ex}{1.6ex \@plus .2ex}%
  {\normalfont\large\bfseries\color{teal!50!black}}}
\makeatother
\setlength{\parindent}{0pt}
\setlength{\parskip}{0.9ex plus 0.3ex minus 0.2ex}

% Headers and Footers
\usepackage{fancyhdr,lastpage}
% [V3 R2] grey left-margin line on the out-of-syllabus pages (item 13)
\usepackage{eso-pic}
% [AoI item 63] thicker line (2.6pt, was 1.6pt) that starts just BELOW the header rule
% (top 14mm of the page left clear). Flag: copy this line into every chapter's preamble.
\newcommand{\oosmargin}{\AddToShipoutPictureBG{\put(12mm,0){\color{gray!75}\rule{2.6pt}{\dimexpr\paperheight-14mm\relax}}}}
\pagestyle{fancy}
\lhead{Maximus Pung Jun}
% [BUILD COPY CHANGE 2] Header bug fix: \chaptername is the word "Chapter",
% so the original header printed "Chapter 1: Chapter". \chaptermark now stores
% the bare chapter title in \leftmark, and the header uses \leftmark.
\renewcommand{\chaptermark}[1]{\markboth{#1}{}}
\rhead{Chapter \thechapter: \leftmark}
% \cfoot{Page \thepage\  of \pageref{LastPage}} % buggy, does not appear on chapter pages.

% [V3 R2] tick and cross markers for mistake boxes (G3)
\newcommand{\wrong}{\textcolor{red!70!black}{\ensuremath{\boldsymbol{\times}}}}
\newcommand{\cmark}{\textcolor{green!50!black}{\ensuremath{\checkmark}}}

\newcommand{\xmark}{\wrong}   % alias: red cross
% [V3 R2] out-of-syllabus helpers (DE notes item 13)
%   \oosmargin    : grey line in the left margin on every following page
%   \oosmarginoff : stop the margin line
%   \oosnote      : grey banner "Out of syllabus. For your interest only."
%   outofsyllabus : environment = new page + margin line + banner; ends with
%                   \clearpage and switches the margin line off again
\newcommand{\oosmarginoff}{\ClearShipoutPictureBG}
\newcommand{\oosnote}[1][Nothing in this section is examined in H2 Mathematics.]{\par\noindent\colorbox{gray!15}{\parbox{\dimexpr\textwidth-2\fboxsep}{\textbf{Out of syllabus. For your interest only.} #1}}\par}
\newenvironment{outofsyllabus}{\clearpage\oosmargin}{\clearpage\oosmarginoff}

% Notation as in Maximus's notes: italic d, display style everywhere (G1).
% Fractions inside exponents stay small, so never use these macros inside ^{}.
\newcommand{\dydx}{\dfrac{\mathrm{d}y}{\mathrm{d}x}}
\newcommand{\dvdx}{\dfrac{\mathrm{d}v}{\mathrm{d}x}}
\newcommand{\dudx}{\dfrac{\mathrm{d}u}{\mathrm{d}x}}
\newcommand{\ddydx}{\dfrac{\mathrm{d}^{2}y}{\mathrm{d}x^{2}}}
\newcommand{\ddt}[1]{\dfrac{\mathrm{d}#1}{\mathrm{d}t}}
% Small syllabus flag (no longer used in the DE notes, kept for other chapters)
\newcommand{\syllabusflag}[1]{\par\noindent{\footnotesize\color{red!70!black}\textbf{Syllabus note.} #1}\par}


% Teacher's copy toggle ================================================
% [BUILD COPY CHANGE 3] \ifx compares the whole word (the original
% \if\conform test was always false). DE_notes_student.tex predefines
% \conform as HideSolutions.
% A student copy defines \conform as HideSolutions BEFORE loading this preamble.
\providecommand\conform{ShowSolutions}
\newcommand\ShowSolutionsKey{ShowSolutions}
% [V3 ADDITION] In the student copy the hidden solution leaves stretchable
% blank space (\vfill), so the free space on each page is shared out as
% working space below the examples instead of collecting at the bottom.
% [AoI item 10] "Solution." on its own line, so that (a), (b), (c) line up below it.
\newcommand\solution[1]{\ifx\conform\ShowSolutionsKey
  \par\noindent\textbf{Solution.}\par\nobreak\vspace{0.5ex}\noindent #1\par\medskip\else\par\vspace{\workspace}\fi}
% [V3 R2] minimum working space left under each example in the student copy
\newlength\workspace\setlength\workspace{3cm plus 1fill}


\newcommand{\NN}{\mathbb{N}}
\newcommand{\ZZ}{\mathbb{Z}}
\newcommand{\QQ}{\mathbb{Q}}
\newcommand{\RR}{\mathbb{R}}
\newcommand{\CC}{\mathbb{C}}


% house_local.tex: HOUSE COPY of definite_integrals/common/ws_preamble_local.tex ([HOUSE] figstyles renamed)
% WORKSHEET COPY of ../notes/preamble_local.tex (version the worksheet was last built with,
% from ../_backup_pre_limitsum, 2026-10-09 17:31). Local so the worksheet no longer depends
% on ../notes; changes: % house_figstyles.tex: HOUSE COPY of definite_integrals/common/ws_figstyles.tex (unchanged)
% WORKSHEET COPY of ../notes/figstyles.tex (2026-10-09 17:31 version). Only change (follow-up):
% thin light axes in `sketch' and `ax3' (marked [2026-10-09 follow-up]).
% =====================================================================
% figstyles.tex: plot styles and figure macros shared by the notes, the
% slides (../slides) and the worksheets (../worksheet). Needs pgfplots,
% the fillbetween library and tikz libraries arrows.meta and calc.
% House rules built in:
%   * `axis on top' in every sketch: shading never covers an axis, and the
%     fill starts at the axis line (items 31, 32, 34);
%   * regions are filled from the curve itself (fill between / closed plot);
%   * 3D solids in the page-13 style (translucent blue, light end caps).
% =====================================================================
\pgfplotsset{
  sketch/.style={clip=false, axis line style={thin, black!45}, every axis label/.append style={black!60}, trig format=rad, % [2026-10-09 follow-up] was axis line style=thick
    axis lines=middle, axis on top,
    every axis x label/.style={at={(current axis.right of origin)},anchor=west},
    every axis y label/.style={at={(current axis.above origin)},anchor=south},
    ticks=none, xlabel=$x$, ylabel=$y$, width=7.6cm, height=6cm,
    every node/.append style={font=\footnotesize}},
  curve/.style={blue, semithick, smooth, samples=150},
  line2/.style={red!75!black, semithick, samples=2},
  % 3D solids of revolution (angles in DEGREES: no `trig format=rad' here)
  solid3d/.style={hide axis, trig format=deg, width=8.4cm, height=6.4cm, clip=false,
    z buffer=sort, colormap={bl}{rgb255=(150,185,235) rgb255=(40,90,180)},
    every node/.append style={font=\footnotesize}},
  solidsurf/.style={surf, shader=interp, opacity=0.75, point meta=y},
  red3d/.style={surf, shader=flat, draw=none, point meta=1,
    colormap={r}{rgb255=(215,50,50) rgb255=(215,50,50)}},
  orange3d/.style={surf, shader=interp, point meta=z,
    colormap={o}{rgb255=(250,210,160) rgb255=(220,120,30)}},
  cap3d/.style={surf, shader=flat, draw=none, point meta=1,
    colormap={c}{rgb255=(185,205,238) rgb255=(185,205,238)}},
  mini3d/.style={solid3d, width=2.6cm, height=2.4cm},
}
\tikzset{
  region/.style={fill=blue!22, draw=none},
  regionB/.style={fill=orange!35, draw=none},
  strip/.style={fill=red!55, draw=red!70!black, thin},
  asym/.style={black, dashed, thin},
  pt/.style={circle, fill=black, inner sep=1.1pt},
  dimarrow/.style={<->, >=Stealth, thin},
  ax3/.style={-{Stealth[length=1.5mm]}, thin, black!50}, % [2026-10-09 follow-up] was -{Stealth[length=2mm]}, thick
}
\definecolor{hl}{RGB}{30,110,200}

% ---------------------------------------------------------------------
% Mini 3D solids for SOLUTIONS (items 54, 62, 66, 74): solid only, no
% washer, every picture 1.5 cm tall; a "big minus small" pair is 4.4 cm wide.
%  \solX[x unit cm]{R(x)}{a}{b}{R(a)}{R(b)}{x-axis end}{y-axis top}{curve colour}
%     solid swept out by y = R(x), a <= x <= b, about the x-axis
%  \solY[y unit cm]{R(y)}{c}{d}{R(c)}{R(d)}{x-axis end}{y-axis top}{curve colour}
%     solid swept out by x = R(y), c <= y <= d, about the y-axis
%     (R is written in terms of the letter x, which here stands for y)
% ---------------------------------------------------------------------
\newcommand{\solX}[9][0.5]{%
\begin{tikzpicture}[baseline=(current bounding box.center)]
\begin{axis}[mini3d, x={(#1cm,-0.06cm)}, y={(0.12cm,0.09cm)}, z={(0cm,0.4cm)}]
  \addplot3[cap3d, domain=0:1, y domain=0:360, samples=2, samples y=40] ({#3}, {x*(#5)*cos(y)}, {x*(#5)*sin(y)});
  \addplot3[solidsurf, opacity=0.5, point meta=z, domain=#3:#4, y domain=0:360, samples=24, samples y=40] ({x}, {(#2)*cos(y)}, {(#2)*sin(y)});
  \addplot3[cap3d, domain=0:1, y domain=0:360, samples=2, samples y=40] ({#4}, {x*(#6)*cos(y)}, {x*(#6)*sin(y)});
  \addplot3[#9, thick, domain=#3:#4, samples=40, samples y=1] ({x}, {0}, {#2});
  \draw[ax3] (axis cs:{min(0,#3)-0.3},0,0) -- (axis cs:#7,0,0);
  \draw[ax3] (axis cs:0,0,{-0.15*(#8)}) -- (axis cs:0,0,#8);
\end{axis}
\end{tikzpicture}}
\newcommand{\solY}[9][0.6]{%
\begin{tikzpicture}[baseline=(current bounding box.center)]
\begin{axis}[mini3d, x={(0.3cm,-0.03cm)}, y={(0.13cm,0.11cm)}, z={(0cm,#1cm)}]
  \addplot3[cap3d, domain=0:1, y domain=0:360, samples=2, samples y=40] ({x*(#5)*cos(y)}, {x*(#5)*sin(y)}, {#3});
  \addplot3[solidsurf, opacity=0.5, point meta=z, domain=#3:#4, y domain=0:360, samples=24, samples y=40] ({(#2)*cos(y)}, {(#2)*sin(y)}, {x});
  \addplot3[cap3d, domain=0:1, y domain=0:360, samples=2, samples y=40] ({x*(#6)*cos(y)}, {x*(#6)*sin(y)}, {#4});
  \addplot3[#9, thick, domain=#3:#4, samples=40, samples y=1] ({#2}, {0}, {x});
  \draw[ax3] (axis cs:0,0,{min(0,#3)-0.15*(#8)}) -- (axis cs:0,0,#8);
  \draw[ax3] (axis cs:-0.3,0,0) -- (axis cs:#7,0,0);
\end{axis}
\end{tikzpicture}}
\newcommand{\minisolo}[1]{\adjustbox{totalheight=1.5cm, max width=4.6cm}{#1}}
\newcommand{\minipair}[2]{\adjustbox{totalheight=1.5cm, max width=4.6cm}{#1\;\raisebox{0.15ex}{\Large$-$}\;#2}}
\newcommand{\minilabel}[1]{\raisebox{0.65cm}{\small #1}\;}
; [2026-10-09 follow-up] \mk marks now bold, as in the notes.
% =====================================================================
% preamble_local.tex: extra packages, styles and macros for the
% Definite Integrals chapter ONLY. Loaded after preamble_build.tex.
% Plot styles and the mini 3D macros are in figstyles.tex (shared with the
% slides and worksheets).
% =====================================================================
\usepgfplotslibrary{fillbetween}
\usepackage{adjustbox} % mini 3D pictures (figstyles.tex)
\usetikzlibrary{arrows.meta, decorations.pathreplacing, calc}
% house_figstyles.tex: HOUSE COPY of definite_integrals/common/ws_figstyles.tex (unchanged)
% WORKSHEET COPY of ../notes/figstyles.tex (2026-10-09 17:31 version). Only change (follow-up):
% thin light axes in `sketch' and `ax3' (marked [2026-10-09 follow-up]).
% =====================================================================
% figstyles.tex: plot styles and figure macros shared by the notes, the
% slides (../slides) and the worksheets (../worksheet). Needs pgfplots,
% the fillbetween library and tikz libraries arrows.meta and calc.
% House rules built in:
%   * `axis on top' in every sketch: shading never covers an axis, and the
%     fill starts at the axis line (items 31, 32, 34);
%   * regions are filled from the curve itself (fill between / closed plot);
%   * 3D solids in the page-13 style (translucent blue, light end caps).
% =====================================================================
\pgfplotsset{
  sketch/.style={clip=false, axis line style={thin, black!45}, every axis label/.append style={black!60}, trig format=rad, % [2026-10-09 follow-up] was axis line style=thick
    axis lines=middle, axis on top,
    every axis x label/.style={at={(current axis.right of origin)},anchor=west},
    every axis y label/.style={at={(current axis.above origin)},anchor=south},
    ticks=none, xlabel=$x$, ylabel=$y$, width=7.6cm, height=6cm,
    every node/.append style={font=\footnotesize}},
  curve/.style={blue, semithick, smooth, samples=150},
  line2/.style={red!75!black, semithick, samples=2},
  % 3D solids of revolution (angles in DEGREES: no `trig format=rad' here)
  solid3d/.style={hide axis, trig format=deg, width=8.4cm, height=6.4cm, clip=false,
    z buffer=sort, colormap={bl}{rgb255=(150,185,235) rgb255=(40,90,180)},
    every node/.append style={font=\footnotesize}},
  solidsurf/.style={surf, shader=interp, opacity=0.75, point meta=y},
  red3d/.style={surf, shader=flat, draw=none, point meta=1,
    colormap={r}{rgb255=(215,50,50) rgb255=(215,50,50)}},
  orange3d/.style={surf, shader=interp, point meta=z,
    colormap={o}{rgb255=(250,210,160) rgb255=(220,120,30)}},
  cap3d/.style={surf, shader=flat, draw=none, point meta=1,
    colormap={c}{rgb255=(185,205,238) rgb255=(185,205,238)}},
  mini3d/.style={solid3d, width=2.6cm, height=2.4cm},
}
\tikzset{
  region/.style={fill=blue!22, draw=none},
  regionB/.style={fill=orange!35, draw=none},
  strip/.style={fill=red!55, draw=red!70!black, thin},
  asym/.style={black, dashed, thin},
  pt/.style={circle, fill=black, inner sep=1.1pt},
  dimarrow/.style={<->, >=Stealth, thin},
  ax3/.style={-{Stealth[length=1.5mm]}, thin, black!50}, % [2026-10-09 follow-up] was -{Stealth[length=2mm]}, thick
}
\definecolor{hl}{RGB}{30,110,200}

% ---------------------------------------------------------------------
% Mini 3D solids for SOLUTIONS (items 54, 62, 66, 74): solid only, no
% washer, every picture 1.5 cm tall; a "big minus small" pair is 4.4 cm wide.
%  \solX[x unit cm]{R(x)}{a}{b}{R(a)}{R(b)}{x-axis end}{y-axis top}{curve colour}
%     solid swept out by y = R(x), a <= x <= b, about the x-axis
%  \solY[y unit cm]{R(y)}{c}{d}{R(c)}{R(d)}{x-axis end}{y-axis top}{curve colour}
%     solid swept out by x = R(y), c <= y <= d, about the y-axis
%     (R is written in terms of the letter x, which here stands for y)
% ---------------------------------------------------------------------
\newcommand{\solX}[9][0.5]{%
\begin{tikzpicture}[baseline=(current bounding box.center)]
\begin{axis}[mini3d, x={(#1cm,-0.06cm)}, y={(0.12cm,0.09cm)}, z={(0cm,0.4cm)}]
  \addplot3[cap3d, domain=0:1, y domain=0:360, samples=2, samples y=40] ({#3}, {x*(#5)*cos(y)}, {x*(#5)*sin(y)});
  \addplot3[solidsurf, opacity=0.5, point meta=z, domain=#3:#4, y domain=0:360, samples=24, samples y=40] ({x}, {(#2)*cos(y)}, {(#2)*sin(y)});
  \addplot3[cap3d, domain=0:1, y domain=0:360, samples=2, samples y=40] ({#4}, {x*(#6)*cos(y)}, {x*(#6)*sin(y)});
  \addplot3[#9, thick, domain=#3:#4, samples=40, samples y=1] ({x}, {0}, {#2});
  \draw[ax3] (axis cs:{min(0,#3)-0.3},0,0) -- (axis cs:#7,0,0);
  \draw[ax3] (axis cs:0,0,{-0.15*(#8)}) -- (axis cs:0,0,#8);
\end{axis}
\end{tikzpicture}}
\newcommand{\solY}[9][0.6]{%
\begin{tikzpicture}[baseline=(current bounding box.center)]
\begin{axis}[mini3d, x={(0.3cm,-0.03cm)}, y={(0.13cm,0.11cm)}, z={(0cm,#1cm)}]
  \addplot3[cap3d, domain=0:1, y domain=0:360, samples=2, samples y=40] ({x*(#5)*cos(y)}, {x*(#5)*sin(y)}, {#3});
  \addplot3[solidsurf, opacity=0.5, point meta=z, domain=#3:#4, y domain=0:360, samples=24, samples y=40] ({(#2)*cos(y)}, {(#2)*sin(y)}, {x});
  \addplot3[cap3d, domain=0:1, y domain=0:360, samples=2, samples y=40] ({x*(#6)*cos(y)}, {x*(#6)*sin(y)}, {#4});
  \addplot3[#9, thick, domain=#3:#4, samples=40, samples y=1] ({#2}, {0}, {x});
  \draw[ax3] (axis cs:0,0,{min(0,#3)-0.15*(#8)}) -- (axis cs:0,0,#8);
  \draw[ax3] (axis cs:-0.3,0,0) -- (axis cs:#7,0,0);
\end{axis}
\end{tikzpicture}}
\newcommand{\minisolo}[1]{\adjustbox{totalheight=1.5cm, max width=4.6cm}{#1}}
\newcommand{\minipair}[2]{\adjustbox{totalheight=1.5cm, max width=4.6cm}{#1\;\raisebox{0.15ex}{\Large$-$}\;#2}}
\newcommand{\minilabel}[1]{\raisebox{0.65cm}{\small #1}\;}


% Small fractions are allowed only in exponents and in the limits of an
% integral (which are superscripts/subscripts); everywhere else use \dfrac.
\newcommand{\dx}{\,\mathrm{d}x}
\newcommand{\dy}{\,\mathrm{d}y}
\newcommand{\evalat}[3]{\left[\,#1\,\right]_{#2}^{#3}}

% [item 78] right-aligned A-level style marks at the end of a question part
\newcommand{\mk}[1]{{\unskip\nobreak\hfil\penalty50\hskip1.5em\hbox{}\nobreak\hfill\textbf{[#1]}%
  \parfillskip=0pt\finalhyphendemerits=0\par}}
% [item 11] visually emphasised step labels in worked solutions
\newcommand{\step}[1]{\textcolor{teal!55!black}{\textbf{Step #1}}\quad}
% (a), (b), (c) part lists inside questions and solutions
\newenvironment{parts}{\begin{enumerate}[(a)]\setlength{\itemsep}{0.4ex}}{\end{enumerate}}
% SEAB-style solution layout (approved 10.5.8 template): working on the
% left, region sketch on the right (mini 3D pictures go under the sketch)
\newcommand{\solpair}[3][0.56]{\noindent\begin{minipage}[t]{#1\textwidth}\vspace{0pt}#2\end{minipage}\hfill
  \begin{minipage}[t]{\dimexpr0.98\textwidth-#1\textwidth\relax}\vspace{0pt}\centering #3\end{minipage}\par}
% minus-diagram style (drafts/minus_style_template.tex, items 29/31/62)
\pgfplotsset{minus/.style={axis lines=middle, width=5.0cm, height=4.2cm, ticks=none,
  xlabel={$x$}, ylabel={$y$}, every axis x label/.style={at={(ticklabel* cs:1)},anchor=west},
  every axis y label/.style={at={(ticklabel* cs:1)},anchor=south}, clip=false, axis on top,
  axis line style={gray!70}, every axis label/.append style={gray},
  every node/.append style={text=gray!80!black, font=\footnotesize}}}
\tikzset{hlreg/.style={fill=hl!35, draw=hl, thick}}

\ifx\housexm\housexnotes
  \ifdefined\housenotesLM\else% house_fonts.tex: HOUSE COPY of definite_integrals/common/prelim_fonts.tex (unchanged)
% =====================================================================
% prelim_fonts.tex: fonts of the real 9758 prelim papers (Word + MathType; pdffonts on
% ACJC 2026 P1/P2, ACJC 2025 P1, JPJC 2026, RVHS 2026, CJC/DHS/TJC/NYJC/HCI/EJC 2025):
%   text Times New Roman; cover/footer Arial; RI page footer Calibri (here Carlito);
%   maths = MathType: Times New Roman letters/digits (italic variables; upright digits,
%   function names, d, e, f), SymbolMT for the integral, sigma, pi and = + < >,
%   scripts 7 pt on 12 pt (about 58 %), tight operator spacing.
% Candidate C (teacher-approved 13 Oct 2026): the real Times New Roman and Arial
% (ttf-mscorefonts-installer) + URW Standard Symbols PS (fonts-urw-base35, a SymbolMT
% clone with the same Symbol encoding) for the integral, sigma, pi, = + < >;
% TeX Gyre Termes Math only for what MathType draws itself (radicals, stretchy brackets,
% fraction bars, arrows). House style: upright d (\mathrm{d}x), upright e and f.
% Loaded once by di_style / de_core (HW sheets) and prelim_style (Tests, solutions).
% =====================================================================
\ifdefined\prelimfontsloaded\endinput\fi
\def\prelimfontsloaded{1}
\edef\pfrestoreat{\catcode`\noexpand\@=\the\catcode`\@\relax}\makeatletter
% MathType script sizes: 7 pt superscripts on 12 pt text (58 %), 5 pt second level
\DeclareMathSizes{12}{12}{7}{5}
\DeclareMathSizes{11}{11}{6.5}{5}
\DeclareMathSizes{10.95}{10.95}{6.5}{5}
\DeclareMathSizes{10}{10}{6}{5}
\DeclareMathSizes{9}{9}{5.5}{5}
\DeclareMathSizes{14.4}{14.4}{8.5}{6}
\setmainfont{Times New Roman}
\setsansfont{Arial}
\newfontfamily\coverfont{Arial}
\newfontfamily\calibrifont{Carlito}
\usepackage{unicode-math}
\setmathfont{TeX Gyre Termes Math}
\setmathfont[range={up/{latin,Latin,num}}]{Times New Roman}
\setmathfont[range={it/{latin,Latin}}]{Times New Roman Italic}
\setmathfont[range={bfup/{latin,Latin,num}}]{Times New Roman Bold}
\setmathfont[range={"002B,"003D,"003C,"003E}]{Standard Symbols PS}
\newfontfamily\pf@sym{Standard Symbols PS}
\newfontfamily\pf@symi[FakeSlant=0.21]{Standard Symbols PS}
% MathType-like spacing (tighter than TeX around = + -)
\AtBeginDocument{\thickmuskip=2.5mu plus 1mu\medmuskip=1.5mu plus 1mu\thinmuskip=3mu
% Symbol integral and sigma (18 pt at 12 pt text, as MathType), slanted Symbol pi
 \renewcommand{\int}{\mathop{\vcenter{\hbox{\pf@sym\fontsize{\fpeval{1.5*\f@size}}{0pt}\selectfont\char"F2}}}\nolimits}%
 \renewcommand{\sum}{\mathop{\vcenter{\hbox{\pf@sym\fontsize{\fpeval{1.5*\f@size}}{0pt}\selectfont\char"E5}}}\limits}%
 \renewcommand{\pi}{\mathord{\text{\pf@symi\char"70}}}}
\pfrestoreat
\fi
  % =====================================================================
% house_tf.tex: True/False table (standing rules 2026-11-05/06/07):
%   * a T/F set is a NUMBERED QUESTION (no "Check Your Understanding" heading);
%   * table: row label (a),(b),... | Statement (left) | T/F column (right);
%   * row spacing slightly larger than default (extrarowheight 6pt, arraystretch 1.18,
%     as the Vectors Basics Worksheet after the 2026-11-07 change), font size unchanged;
%   * question doc: T/F column blank; solutions doc (\tfanswerstrue): bold T or F,
%     optional short reason under the statement in small grey.
% Usage:
%   \begin{tftable}  \tfrow{statement}{T}{optional reason}  ...  \end{tftable}
% Plain tabular: a question never splits across pages (house rule), so the table is kept whole.
% =====================================================================
\RequirePackage{array}
\makeatletter
\newif\iftfanswers \tfanswersfalse
\newcounter{tfrow}
\newlength{\tf@lab}\setlength{\tf@lab}{0.9cm}
\newlength{\tf@ans}\setlength{\tf@ans}{1.6cm}
\newenvironment{tftable}{\setcounter{tfrow}{0}\par\smallskip
  \setlength{\extrarowheight}{6pt}\renewcommand{\arraystretch}{1.18}%
  \noindent\begin{tabular}{|>{\centering\arraybackslash}p{\tf@lab}|p{\dimexpr\linewidth-\tf@lab-\tf@ans-6\tabcolsep-4\arrayrulewidth\relax}|>{\centering\arraybackslash}p{\tf@ans}|}
  \hline &\textbf{Statement}&\textbf{T/F}\\\hline}
 {\end{tabular}\par\smallskip}
\newcommand{\tfrow}[3]{\stepcounter{tfrow}\textbf{(\alph{tfrow})} & #1%
  \iftfanswers\if\relax\detokenize{#3}\relax\else\newline{\small\color{black!60}#3}\fi\fi
  & \iftfanswers\textbf{#2}\fi\\\hline}
% explicit row label (a set that keeps the Basics Worksheet letters, e.g. (e) omitted)
\newcommand{\tfrowl}[4]{\textbf{#1} & #2%
  \iftfanswers\if\relax\detokenize{#4}\relax\else\newline{\small\color{black!60}#4}\fi\fi
  & \iftfanswers\textbf{#3}\fi\\\hline}
% \tfcompact: allowed exception (Vectors I HW2/HW3) - compressed rows so the whole set
% stays on page 1 of the question sheet (extrarowheight 1.5pt, tabcolsep 4pt, 0.85/1.45 cm)
\newif\iftf@compact
\newcommand{\tfcompact}{\tf@compacttrue\setlength{\tf@lab}{0.85cm}\setlength{\tf@ans}{1.45cm}}
\renewenvironment{tftable}{\setcounter{tfrow}{0}\par\smallskip
  \iftf@compact\setlength{\extrarowheight}{1.5pt}\setlength{\tabcolsep}{4pt}\else\setlength{\extrarowheight}{6pt}\fi
  \renewcommand{\arraystretch}{1.18}%
  \noindent\begin{tabular}{|>{\centering\arraybackslash}p{\tf@lab}|p{\dimexpr\linewidth-\tf@lab-\tf@ans-6\tabcolsep-4\arrayrulewidth\relax}|>{\centering\arraybackslash}p{\tf@ans}|}
  \hline &\textbf{Statement}&\textbf{T/F}\\\hline}
 {\end{tabular}\par\smallskip}
% tflongtable: same table, may break across pages with the header repeated (solutions
% copies whose reasons make the set longer than a page). Needs longtable.
\RequirePackage{longtable}
\newenvironment{tflongtable}{\setcounter{tfrow}{0}\par\smallskip
  \setlength{\extrarowheight}{6pt}\renewcommand{\arraystretch}{1.18}%
  \setlength{\LTleft}{0pt}\setlength{\LTright}{0pt}%
  \begin{longtable}{|>{\centering\arraybackslash}p{\tf@lab}|p{\dimexpr\linewidth-\tf@lab-\tf@ans-6\tabcolsep-4\arrayrulewidth\relax}|>{\centering\arraybackslash}p{\tf@ans}|}
  \hline &\textbf{Statement}&\textbf{T/F}\\\hline\endhead}
 {\end{longtable}\par}
% \putTFlong{n}: True/False question n (body of \defQ{n}: intro, then tftable) printed with
% tflongtable, so a long set may run over a page (house_sol redefines it for solutions).
\newcommand{\putTFlong}[1]{\begingroup\let\tftable\tflongtable\let\endtftable\endtflongtable
  \stepcounter{diq}\expandafter\xdef\csname diLab@#1\endcsname{\thediq}%
  \Qhead{\thediq}{\csname diQtag@#1\endcsname}\csname diQbody@#1\endcsname\par\endgroup}
\makeatother

  \providecommand{\units}[1]{\ \text{units}^{#1}}
  \raggedbottom
\else
  % house_fonts.tex: HOUSE COPY of definite_integrals/common/prelim_fonts.tex (unchanged)
% =====================================================================
% prelim_fonts.tex: fonts of the real 9758 prelim papers (Word + MathType; pdffonts on
% ACJC 2026 P1/P2, ACJC 2025 P1, JPJC 2026, RVHS 2026, CJC/DHS/TJC/NYJC/HCI/EJC 2025):
%   text Times New Roman; cover/footer Arial; RI page footer Calibri (here Carlito);
%   maths = MathType: Times New Roman letters/digits (italic variables; upright digits,
%   function names, d, e, f), SymbolMT for the integral, sigma, pi and = + < >,
%   scripts 7 pt on 12 pt (about 58 %), tight operator spacing.
% Candidate C (teacher-approved 13 Oct 2026): the real Times New Roman and Arial
% (ttf-mscorefonts-installer) + URW Standard Symbols PS (fonts-urw-base35, a SymbolMT
% clone with the same Symbol encoding) for the integral, sigma, pi, = + < >;
% TeX Gyre Termes Math only for what MathType draws itself (radicals, stretchy brackets,
% fraction bars, arrows). House style: upright d (\mathrm{d}x), upright e and f.
% Loaded once by di_style / de_core (HW sheets) and prelim_style (Tests, solutions).
% =====================================================================
\ifdefined\prelimfontsloaded\endinput\fi
\def\prelimfontsloaded{1}
\edef\pfrestoreat{\catcode`\noexpand\@=\the\catcode`\@\relax}\makeatletter
% MathType script sizes: 7 pt superscripts on 12 pt text (58 %), 5 pt second level
\DeclareMathSizes{12}{12}{7}{5}
\DeclareMathSizes{11}{11}{6.5}{5}
\DeclareMathSizes{10.95}{10.95}{6.5}{5}
\DeclareMathSizes{10}{10}{6}{5}
\DeclareMathSizes{9}{9}{5.5}{5}
\DeclareMathSizes{14.4}{14.4}{8.5}{6}
\setmainfont{Times New Roman}
\setsansfont{Arial}
\newfontfamily\coverfont{Arial}
\newfontfamily\calibrifont{Carlito}
\usepackage{unicode-math}
\setmathfont{TeX Gyre Termes Math}
\setmathfont[range={up/{latin,Latin,num}}]{Times New Roman}
\setmathfont[range={it/{latin,Latin}}]{Times New Roman Italic}
\setmathfont[range={bfup/{latin,Latin,num}}]{Times New Roman Bold}
\setmathfont[range={"002B,"003D,"003C,"003E}]{Standard Symbols PS}
\newfontfamily\pf@sym{Standard Symbols PS}
\newfontfamily\pf@symi[FakeSlant=0.21]{Standard Symbols PS}
% MathType-like spacing (tighter than TeX around = + -)
\AtBeginDocument{\thickmuskip=2.5mu plus 1mu\medmuskip=1.5mu plus 1mu\thinmuskip=3mu
% Symbol integral and sigma (18 pt at 12 pt text, as MathType), slanted Symbol pi
 \renewcommand{\int}{\mathop{\vcenter{\hbox{\pf@sym\fontsize{\fpeval{1.5*\f@size}}{0pt}\selectfont\char"F2}}}\nolimits}%
 \renewcommand{\sum}{\mathop{\vcenter{\hbox{\pf@sym\fontsize{\fpeval{1.5*\f@size}}{0pt}\selectfont\char"E5}}}\limits}%
 \renewcommand{\pi}{\mathord{\text{\pf@symi\char"70}}}}
\pfrestoreat

  % house_hw.tex: HOUSE COPY of definite_integrals/common/di_style.tex; [HOUSE] chapter name from \housechapter{...}
% =====================================================================
% di_style.tex: Definite Integrals homework / test layout.
% Loads the house preamble (local copies of the AoI worksheet files:
% ws_preamble_build.tex, ws_preamble_local.tex -> ws_figstyles.tex),
% then adds the homework layout used by HW1-HW3 and the Test.
% Usage (before \begin{document}):  \ditle{Homework 1}
% =====================================================================
% [HOUSE] base, local and fonts are loaded by house.tex
\usepackage{needspace}
\usepackage{zref-savepos,zref-abspage}\makeatletter\zref@addprop{savepos}{abspage}\makeatother
\usepackage{array}

% ---- shortlist toggle ------------------------------------------------
% \ifshortlist true  -> only the questions in SHORTLIST.md (+ challenges)
% A wrapper file sets it with \def\setshortlist{1} before \input-ing the
% main file; the main file can also be edited to say \shortlisttrue.
\newif\ifshortlist
% Task 5: the FINAL homework (= former shortlist) is the default build. HWn.tex -> HWn.pdf.
% A wrapper that sets \def\setfull{1} builds the archive with all original questions.
\ifdefined\setfull\else\shortlisttrue\fi

\makeatletter
% ---- REVIEW build: highlight amended content ---------------------------
% A wrapper sets \def\setreview{1}. In review mode amended text is shaded
% yellow and carries a hover note (PDF tooltip / sticky note); otherwise the
% macros print the content plainly.
%   \chg{note}{inline content}        short inline change (no line breaks), text mode
%   \chgm{note}{math}                 the same inside math mode
%   \begin{chgblock}{note} ... \end{chgblock}   paragraph/display-level change
\newif\ifreview
\ifdefined\setreview \reviewtrue \fi
% hover notes are raw PDF annotations (xdvipdfmx \special; no extra packages).
% Keep notes plain ASCII; balanced parentheses are fine.
\newsavebox{\chg@box}
\newcommand{\chg@ann}[4]{% width, height, depth, note -> invisible markup annotation with a hover note
  \special{pdf:ann width #1 height #2 depth #3 << /Type /Annot /Subtype /Square /Border [0 0 0] /BS << /W 0 >> /C [1 0.85 0] /CA 0.01
    /T (Review) /Subj (Amendment) /Contents (#4) >>}}
\newcommand{\chg@sticky}[1]{\special{pdf:ann width 12pt height 12pt depth 0pt << /Type /Annot /Subtype /Text /Name /Comment
    /C [1 0.85 0] /T (Review) /Subj (Amendment) /Contents (#1) >>}}
\newcommand{\chg}[2]{\ifreview\sbox\chg@box{\colorbox{yellow!60}{#2}}%
  \leavevmode\rlap{\chg@ann{\the\wd\chg@box}{\the\ht\chg@box}{\the\dp\chg@box}{#1}}\usebox\chg@box\else#2\fi}
% \chgm{note}{math}: the same, for use INSIDE math mode (display style kept)
\newcommand{\chgm}[2]{\ifreview\mathchoice{\chg{#1}{$\displaystyle#2$}}{\chg{#1}{$\textstyle#2$}}{\chg{#1}{$\scriptstyle#2$}}{\chg{#1}{$\scriptscriptstyle#2$}}\else#2\fi}
\newtcolorbox{chgbox}[1]{enhanced, breakable, colback=yellow!35, colframe=yellow!35, boxrule=0pt, sharp corners,
  left=2pt, right=14pt, top=1pt, bottom=1pt, before skip=2pt, after skip=2pt,
  overlay first and middle={}, overlay last={}, overlay unbroken and first={\node[anchor=north east, inner sep=1pt] at (frame.north east) {\chg@sticky{#1}\hbox to 12pt{\vrule width 0pt height 12pt}};}}
% \chgnote{note}: just the sticky-note icon (e.g. inside a figure); nothing outside review mode
\newcommand{\chgnote}[1]{\ifreview\chg@sticky{#1}\hbox to 12pt{\vrule width 0pt height 12pt}\fi}
\newenvironment{chgblock}[1]{\ifreview\begin{chgbox}{#1}\fi}{\ifreview\end{chgbox}\fi}

% ---- header ---------------------------------------------------------
\newcommand\di@title{}
\makeatletter
\newcommand{\ditle}[1]{\gdef\di@title{#1}}
\fancypagestyle{di}{\fancyhf{}%
  \lhead{\small \housechap: \di@title}%
  \rhead{\small Compiled by Maximus Pung Jun}%
  \cfoot{\small\thepage}%
  \renewcommand{\headrulewidth}{0.4pt}}
\pagestyle{di}
% title block on page 1, as in his Typst sheets
\newcommand{\dititleblock}[1][]{\thispagestyle{di}%
  {\LARGE\scshape \housechap: \di@title}\par\smallskip
  {Compiled by Maximus Pung Jun}\par\ifx&#1&\else\smallskip{\small\textcolor{black!60}{#1}}\par\fi\medskip}
\makeatother
\newcommand{\disection}[1]{\par\bigskip\noindent{\Large\scshape #1}\par\medskip}

% ---- questions --------------------------------------------------------
% \Qhead{label}{source tag}: question heading; grey [tag] straight after the number, on the left (rule 3B, 9 Nov 2026)
\newcommand{\Qhead}[2]{\par\needspace{8\baselineskip}\bigskip\noindent
  {\large\bfseries #1}\setbox0=\hbox{#2}\ifdim\wd0=0pt\else\enspace{\small\textcolor{black!60}{[#2]}}\fi\par\smallskip}
% part lists with bold labels as he writes them: \pt{(a)}, \pt{(i)}, \pt{(b) (i)}
\newenvironment{qparts}[1][2.4em]{\begin{list}{}{\setlength\leftmargin{#1}%
  \setlength\labelwidth{\dimexpr#1-0.5em\relax}\setlength\labelsep{0.5em}%
  \setlength\itemsep{0.8ex}\setlength\topsep{0.6ex}\setlength\parsep{0.6ex}%
  \renewcommand\makelabel[1]{\textbf{##1}\hfil}}}{\end{list}}
\newcommand{\pt}[1]{\item[#1]}
% thin horizontal label for a challenge question
\newcommand{\challengebanner}{\par\bigskip\noindent\textcolor{black!50}{\rule[0.5ex]{\linewidth}{0.4pt}}\par
  \noindent{\large\scshape Challenge (optional)}\par}
% upright "units"
\newcommand{\units}[1]{\ \text{units}^{#1}}
% END line
\newcommand{\diend}[1][END OF WORKSHEET]{\par\vspace{2\baselineskip}\begin{center}\textbf{#1}\end{center}}

% ---- question registry (full / shortlist from one source) -------------
% \defQ{n}{tag}{body}   question n of the original sheet
% \defA{n}{answer body} its final answer, as printed in his answer section
% \putQ{n} / \putC{n}   print question n as the next (challenge) question
% \putA{n}              print its answer under the same new label
\newcounter{diq}
\makeatletter
\newcommand{\defQ}[3]{\expandafter\gdef\csname diQtag@#1\endcsname{#2}%
  \expandafter\long\expandafter\gdef\csname diQbody@#1\endcsname{#3}}
\newcommand{\defA}[2]{\expandafter\long\expandafter\gdef\csname diA@#1\endcsname{#2}}
\newcommand{\di@tag}[1]{\csname diQtag@#1\endcsname}% source tag only (no original Q number)
% A question never breaks across pages: it is measured and starts on a new page unless it
% fits whole on the rest of this one (only a question taller than a page may break, and then
% only between parts). Page positions of its start/end (zref labels q-s<n>/q-e<n> in the .aux)
% let the checker verify this.
\newsavebox{\di@qmeas}
\newcommand{\di@placeQ}[1]{\par\setbox\di@qmeas=\vbox{#1}%
  \ifdim\dimexpr\ht\di@qmeas+\dp\di@qmeas\relax<0.95\textheight
    \needspace{\dimexpr\ht\di@qmeas+\dp\di@qmeas+\bigskipamount\relax}\else\needspace{14\baselineskip}\fi
  \zsavepos{q-s\thediq}#1\par\zsavepos{q-e\thediq}}
\newcommand{\putQ}[1]{\stepcounter{diq}%
  \expandafter\xdef\csname diLab@#1\endcsname{\thediq}%
  \di@placeQ{\Qhead{\thediq}{\di@tag{#1}}\csname diQbody@#1\endcsname\par}}
\newcommand{\putC}[1]{\stepcounter{diq}%
  \expandafter\xdef\csname diLab@#1\endcsname{\thediq}%
  \expandafter\gdef\csname diChal@#1\endcsname{1}%
  \di@placeQ{\Qhead{\thediq\quad\normalsize\mdseries\scshape Challenge}{\di@tag{#1}}\csname diQbody@#1\endcsname\par}}
\newcommand{\putA}[1]{\par\needspace{4\baselineskip}%
  \begin{qparts}[2.4em]\item[{\csname diLab@#1\endcsname}]\csname diA@#1\endcsname\end{qparts}\medskip}
% ---- worked solutions: \defS{n}{body}, \putQS{n} prints question n + its solution ----
\newcommand{\defS}[2]{\expandafter\long\expandafter\gdef\csname diS@#1\endcsname{#2}}
\newcommand{\disolhead}{\par\medskip\noindent{\bfseries Solution}\par\nobreak\vspace{-0.6ex}\noindent\textcolor{black!35}{\rule{\linewidth}{0.4pt}}\par\nobreak}
\newcommand{\putQS}[1]{\putQ{#1}\disolhead\csname diS@#1\endcsname\par}
% source of a worked solution; visible only in REVIEW builds. #1: his / pub / pip
\newcommand{\solsrc}[2]{\ifreview\par\noindent\def\di@a{#1}\def\di@b{pip}%
  \ifx\di@a\di@b\colorbox{yellow!60}{\parbox{\dimexpr\linewidth-2\fboxsep-14pt}{\small\textbf{Pip-worked:} #2}}\chgnote{Pip-worked solution. No worked solution by the teacher and no published solution was found. Please check it.}%
  \else{\small\textcolor{black!55}{Source of this solution: #2}}\fi\par\smallskip\fi}
\makeatother
% compact sketch for worked solutions: \solfig[axis options]{plot commands}
\newcommand{\solfig}[2][]{\par\begin{center}\begin{tikzpicture}\begin{axis}[qfig, width=6.6cm, height=5cm, #1]#2\end{axis}\end{tikzpicture}\end{center}}
% a sketch added to an existing worked solution (teacher rule), shaded in REVIEW
\newenvironment{addedsketch}{\begin{chgblock}{Sketch added: teacher rule - every area or volume solution has a sketch.}}{\end{chgblock}}
% answer part list (compact)
\newenvironment{aparts}[1][2.4em]{\begin{qparts}[#1]\setlength\itemsep{0.3ex}}{\end{qparts}}

% ---- figure helpers ---------------------------------------------------------
\tikzset{rect/.style={fill=blue!10, draw=blue!55!black, thin}}
\pgfplotsset{qfig/.style={sketch, every node/.append style={font=\scriptsize}}}
% aligned working that starts on the same line as a part label: \pt{(iii)} \begin{topalign} ... \end{topalign}
\newenvironment{topalign}{\hfill$\displaystyle\begin{aligned}[t]}{\end{aligned}$\hfill\null\par\smallskip}
% one display that starts on the same line as a part label
\newcommand{\topdisplay}[1]{\hfill$\displaystyle #1$\hfill\null\par\smallskip}

  % house_solmacros.tex: HOUSE COPY of de-docs/common/de_style.tex minus \input{de_core} and the DE test titles.
% Adds the Ch11 solution helpers: soln/spart tables, sbs/sbsplain side-by-side, \note, \award, \tagm, \dstart.
% =====================================================================
% de_style.tex: Differential Equations (Ch11) homework / test layout.
% = Ch10 di_style (copied as de_core.tex; header text changed) + the DE
% notation and the framed solution boxes of the v1 DE documents
% (soln / spart, used by the Test worked solutions).
% =====================================================================
\providecommand{\e}{\mathrm{e}}
\providecommand{\cosec}{\operatorname{cosec}}
\providecommand{\dt}{\,\mathrm{d}t}
\providecommand{\abs}[1]{\left|#1\right|}
\providecommand{\dd}[2]{\dfrac{\mathrm{d}#1}{\mathrm{d}#2}}
\providecommand{\dvdx}{\dfrac{\mathrm{d}v}{\mathrm{d}x}}
\providecommand{\dudx}{\dfrac{\mathrm{d}u}{\mathrm{d}x}}
\colorlet{notegrey}{black!60}
\makeatletter
\protected\def\note#1{\ifmmode\text{\small\color{notegrey}\itshape #1}\else{\small\color{notegrey}\itshape #1}\fi}
\newcommand{\award}[1]{\unskip\nobreak\hfill\textbf{\small #1}}
\newcommand{\tagm}[1]{\tag*{\textbf{\small #1}}}
\newcommand{\mfAminus}{\[\begin{aligned}&\int\dfrac{1}{a^{2}-x^{2}}\dx\\&\quad=\dfrac{1}{2a}\ln\abs{\dfrac{a+x}{a-x}}+c\end{aligned}\]}
\newcommand{\mfXminus}{\[\begin{aligned}&\int\dfrac{1}{x^{2}-a^{2}}\dx\\&\quad=\dfrac{1}{2a}\ln\abs{\dfrac{x-a}{x+a}}+c\end{aligned}\]}
\newcommand{\sbs@start}{\if@inlabel\def\sbs@shift{\vspace{-\ht\strutbox}}\else
  \ifvmode\else\par\fi\addvspace{0.4ex}\def\sbs@shift{}\fi\noindent}
\newsavebox{\sbs@box}
\newenvironment{sbs}[1]{%
  \sbs@start\begin{minipage}[t]{#1\linewidth}\vspace{0pt}\sbs@shift\raggedright%
  \def\sbsnext{\end{minipage}\hfill\begin{minipage}[t]{\dimexpr0.97\linewidth-#1\linewidth\relax}\vspace{0pt}\sbs@shift\setlength{\fboxsep}{1.5mm}\begin{lrbox}{\sbs@box}\begin{minipage}[t]{\dimexpr\linewidth-2\fboxsep-1pt\relax}\small\raggedright}%
}{\end{minipage}\end{lrbox}\setlength{\fboxrule}{0.5pt}\fbox{\usebox{\sbs@box}}\end{minipage}\par}
\newenvironment{sbsplain}[1]{%
  \sbs@start\begin{minipage}[t]{#1\linewidth}\vspace{0pt}\sbs@shift\raggedright%
  \def\sbsnext{\end{minipage}\hfill\begin{minipage}[t]{\dimexpr0.97\linewidth-#1\linewidth\relax}\vspace{0pt}\sbs@shift\centering}%
}{\end{minipage}\par}
% ---- framed worked-solution box (v1 DE mark-scheme layout, kept for the Test) ----
\newlength{\solcol}\setlength{\solcol}{15mm}
\newcounter{spartno}
\newcommand{\solskips}{\abovedisplayskip=4pt plus 2pt\belowdisplayskip=4pt plus 2pt\abovedisplayshortskip=2pt\belowdisplayshortskip=3pt}
\newtcolorbox{soln}[1][]{enhanced, breakable, sharp corners,
  colback=white, colframe=black, boxrule=0.6pt,
  boxsep=0pt, left=0pt, right=3mm, top=1.5mm, bottom=2mm,
  before skip=1.2ex, after skip=1.5ex,
  before upper={\setcounter{spartno}{0}\parskip=0.5ex\solskips\@fleqntrue\@mathmargin=0pt\raggedright},
  overlay={\draw[line width=0.4pt] ([xshift=\solcol]frame.north west) -- ([xshift=\solcol]frame.south west);},
  #1}
\newenvironment{spart}[1]{%
  \stepcounter{spartno}%
  \ifnum\value{spartno}>1 \par\vspace{0.2ex}\noindent
    \makebox[\linewidth][l]{\rule{\dimexpr\linewidth+3mm\relax}{0.4pt}}\par\nobreak\vspace{0.3ex}\fi
  \begin{list}{}{\leftmargin=\dimexpr\solcol+2.5mm\relax
    \labelwidth=\solcol \labelsep=2.5mm \itemindent=0pt
    \topsep=0.3ex \parsep=0.5ex \rightmargin=0pt \@beginparpenalty=\@M}%
  \item[{\makebox[\solcol][l]{\strut\hspace{1.5mm}\textbf{#1}}}]%
  \ignorespaces}{\end{list}}
\makeatother
% ---- DE sketches: thin light axes (teacher rule), as the Ch10 'sketch' style ----
\pgfplotsset{dsketch/.style={sketch, axis lines=middle, axis line style={thin, black!45, -{Stealth[length=1.6mm]}},
    ticks=none, clip=false, width=6.5cm, height=4.6cm,
    every axis x label/.style={at={(current axis.right of origin)},anchor=west},
    every axis y label/.style={at={(current axis.above origin)},anchor=south},
    every node/.append style={font=\small}},
  dcurve/.style={black, semithick, smooth, samples=120}}
\tikzset{asym/.style={black, dashed, thin}}
\makeatletter
\newcommand{\dstart}{\leavevmode\vadjust{\nobreak\vskip\dimexpr-\baselineskip-\abovedisplayskip+1pt\relax}}
\newcommand{\sdstart}{\leavevmode\vadjust{\nobreak\vskip\dimexpr-\baselineskip-\abovedisplayskip+4pt\relax}}
\makeatother

  % =====================================================================
% house_tf.tex: True/False table (standing rules 2026-11-05/06/07):
%   * a T/F set is a NUMBERED QUESTION (no "Check Your Understanding" heading);
%   * table: row label (a),(b),... | Statement (left) | T/F column (right);
%   * row spacing slightly larger than default (extrarowheight 6pt, arraystretch 1.18,
%     as the Vectors Basics Worksheet after the 2026-11-07 change), font size unchanged;
%   * question doc: T/F column blank; solutions doc (\tfanswerstrue): bold T or F,
%     optional short reason under the statement in small grey.
% Usage:
%   \begin{tftable}  \tfrow{statement}{T}{optional reason}  ...  \end{tftable}
% Plain tabular: a question never splits across pages (house rule), so the table is kept whole.
% =====================================================================
\RequirePackage{array}
\makeatletter
\newif\iftfanswers \tfanswersfalse
\newcounter{tfrow}
\newlength{\tf@lab}\setlength{\tf@lab}{0.9cm}
\newlength{\tf@ans}\setlength{\tf@ans}{1.6cm}
\newenvironment{tftable}{\setcounter{tfrow}{0}\par\smallskip
  \setlength{\extrarowheight}{6pt}\renewcommand{\arraystretch}{1.18}%
  \noindent\begin{tabular}{|>{\centering\arraybackslash}p{\tf@lab}|p{\dimexpr\linewidth-\tf@lab-\tf@ans-6\tabcolsep-4\arrayrulewidth\relax}|>{\centering\arraybackslash}p{\tf@ans}|}
  \hline &\textbf{Statement}&\textbf{T/F}\\\hline}
 {\end{tabular}\par\smallskip}
\newcommand{\tfrow}[3]{\stepcounter{tfrow}\textbf{(\alph{tfrow})} & #1%
  \iftfanswers\if\relax\detokenize{#3}\relax\else\newline{\small\color{black!60}#3}\fi\fi
  & \iftfanswers\textbf{#2}\fi\\\hline}
% explicit row label (a set that keeps the Basics Worksheet letters, e.g. (e) omitted)
\newcommand{\tfrowl}[4]{\textbf{#1} & #2%
  \iftfanswers\if\relax\detokenize{#4}\relax\else\newline{\small\color{black!60}#4}\fi\fi
  & \iftfanswers\textbf{#3}\fi\\\hline}
% \tfcompact: allowed exception (Vectors I HW2/HW3) - compressed rows so the whole set
% stays on page 1 of the question sheet (extrarowheight 1.5pt, tabcolsep 4pt, 0.85/1.45 cm)
\newif\iftf@compact
\newcommand{\tfcompact}{\tf@compacttrue\setlength{\tf@lab}{0.85cm}\setlength{\tf@ans}{1.45cm}}
\renewenvironment{tftable}{\setcounter{tfrow}{0}\par\smallskip
  \iftf@compact\setlength{\extrarowheight}{1.5pt}\setlength{\tabcolsep}{4pt}\else\setlength{\extrarowheight}{6pt}\fi
  \renewcommand{\arraystretch}{1.18}%
  \noindent\begin{tabular}{|>{\centering\arraybackslash}p{\tf@lab}|p{\dimexpr\linewidth-\tf@lab-\tf@ans-6\tabcolsep-4\arrayrulewidth\relax}|>{\centering\arraybackslash}p{\tf@ans}|}
  \hline &\textbf{Statement}&\textbf{T/F}\\\hline}
 {\end{tabular}\par\smallskip}
% tflongtable: same table, may break across pages with the header repeated (solutions
% copies whose reasons make the set longer than a page). Needs longtable.
\RequirePackage{longtable}
\newenvironment{tflongtable}{\setcounter{tfrow}{0}\par\smallskip
  \setlength{\extrarowheight}{6pt}\renewcommand{\arraystretch}{1.18}%
  \setlength{\LTleft}{0pt}\setlength{\LTright}{0pt}%
  \begin{longtable}{|>{\centering\arraybackslash}p{\tf@lab}|p{\dimexpr\linewidth-\tf@lab-\tf@ans-6\tabcolsep-4\arrayrulewidth\relax}|>{\centering\arraybackslash}p{\tf@ans}|}
  \hline &\textbf{Statement}&\textbf{T/F}\\\hline\endhead}
 {\end{longtable}\par}
% \putTFlong{n}: True/False question n (body of \defQ{n}: intro, then tftable) printed with
% tflongtable, so a long set may run over a page (house_sol redefines it for solutions).
\newcommand{\putTFlong}[1]{\begingroup\let\tftable\tflongtable\let\endtftable\endtflongtable
  \stepcounter{diq}\expandafter\xdef\csname diLab@#1\endcsname{\thediq}%
  \Qhead{\thediq}{\csname diQtag@#1\endcsname}\csname diQbody@#1\endcsname\par\endgroup}
\makeatother

  \ifx\housexm\housexsol \def\prelimkeepcolour{1}% house_prelim.tex: HOUSE COPY of definite_integrals/common/prelim_style.tex; [HOUSE] titles from \housechapter
% =====================================================================
% prelim_style.tex: makes the Test look like a Singapore JC H2 Mathematics
% 9758 prelim paper (studied from ACJC/RI 2026 P1/P2 on the box:
% body Times New Roman + Times-style maths, cover in Arial).
% Fonts: prelim_fonts.tex (real Times New Roman/Arial + Times-style maths).
% Shared by Ch10 (di_style) and Ch11 (de_style; sets \prelimtitle/\prelimsubject). Loaded AFTER the chapter style (overrides its layout for the test only).
% =====================================================================
% [HOUSE] fonts loaded by house.tex
\providecommand{\prelimtitle}{\MakeUppercase{\housechap} TEST}
\providecommand{\prelimsubject}{\housechap}
\geometry{a4paper, top=22mm, bottom=24mm, left=25mm, right=25mm, headsep=6mm, footskip=12mm}
\setlength{\parindent}{0pt}
\setlength{\parskip}{0.9ex plus 0.3ex minus 0.2ex}
\linespread{1.3}
\setlength{\emergencystretch}{3em}
\clubpenalty=10000 \widowpenalty=10000 \interlinepenalty=10000
\AtBeginDocument{\fontsize{12}{14.4}\selectfont}
\renewcommand{\normalsize}{\fontsize{12}{14.4}\selectfont}

% ---- figures: plain black/grey like a printed paper ----------------------
\ifdefined\prelimkeepcolour\else % question paper: black/grey figures; solutions keep colour
\pgfplotsset{curve/.style={black, semithick, smooth, samples=150},
  line2/.style={black, semithick, samples=2},
  sketch/.append style={axis line style={thin, black}, every axis label/.append style={black}}}
\tikzset{region/.style={fill=black!15, draw=none}, regionB/.style={fill=black!30, draw=none},
  rect/.style={fill=black!8, draw=black, thin}, ax3/.style={-{Stealth[length=1.5mm]}, thin, black}}
\colorlet{blue}{black}\colorlet{red}{black}\colorlet{orange}{black}
\fi

% ---- header / footer: page number centred top (bold), [Turn over at bottom right
\makeatletter
\providecommand{\prelimfootlabel}{\prelimtitle\ 2026}
\newcommand{\pp@footer}{{\coverfont\scriptsize \prelimfootlabel\hspace{3em}H2 MATHEMATICS 9758}}
\fancypagestyle{prelim}{\fancyhf{}%
  \chead{\bfseries\thepage}%
  \lfoot{\pp@footer}%
  \rfoot{\ifnum\value{page}<\getpagerefnumber{LastPage}\relax\textbf{[Turn over}\fi}%
  \renewcommand{\headrulewidth}{0pt}\renewcommand{\footrulewidth}{0pt}}
\fancypagestyle{prelimcover}{\fancyhf{}%
  \rfoot{\coverfont\textbf{[Turn over}}%
  \renewcommand{\headrulewidth}{0pt}\renewcommand{\footrulewidth}{0pt}}
\pagestyle{prelim}

% ---- marks: [n] regular weight, right-aligned at the end of the part ----------
\renewcommand{\mk}[1]{{\unskip\nobreak\hfil\penalty50\hskip1.5em\hbox{}\nobreak\hfill[#1]%
  \parfillskip=0pt\finalhyphendemerits=0\par}}

% ---- questions: bold number at the left margin, text indented; grey [tag] straight after the number (3B)
\newsavebox{\pp@qmeas}
\newcommand{\pp@qlist}[1]{%
  \begin{list}{}{\setlength\leftmargin{2.6em}\setlength\labelwidth{2.1em}\setlength\labelsep{0.5em}%
     \setlength\topsep{0pt}\setlength\itemsep{0pt}\setlength\parsep{0.9ex}%
     \renewcommand\makelabel[1]{\textbf{##1}\hfil}}%
  \item[\thediq]\@endparpenalty\@M\setbox0=\hbox{\di@tag{#1}}\ifdim\wd0=0pt\else{\small\textcolor{black!60}{[\di@tag{#1}]}}\par\nobreak\fi\csname diQbody@#1\endcsname\par\penalty\@M
  \end{list}}
% a question starts on a new page unless it fits whole on the current one
% (questions taller than a page may break, but only between parts)
\renewcommand{\putQ}[1]{\stepcounter{diq}%
  \expandafter\xdef\csname diLab@#1\endcsname{\thediq}%
  \par\setbox\pp@qmeas=\vbox{\pp@qlist{#1}}%
  \ifdim\dimexpr\ht\pp@qmeas+\dp\pp@qmeas\relax<0.95\textheight
    \needspace{\dimexpr\ht\pp@qmeas+\dp\pp@qmeas\relax}\else\needspace{14\baselineskip}\fi
  \zsavepos{q-s\thediq}\pp@qlist{#1}\par\zsavepos{q-e\thediq}\vspace{5ex plus 2ex minus 1ex}}
% part labels (a)/(i) bold, indented under the question text
\renewenvironment{qparts}[1][2.6em]{\begin{list}{}{\setlength\leftmargin{#1}%
  \setlength\labelwidth{\dimexpr#1-0.5em\relax}\setlength\labelsep{0.5em}%
  \setlength\itemsep{1.4ex}\setlength\topsep{1.2ex}\setlength\parsep{0.8ex}%
  \renewcommand\makelabel[1]{\textbf{##1}\hfil}}}{\end{list}}
\makeatother

% ---- cover page --------------------------------------------------------------
% A-level H2 style cover (9 Nov 2026; reference house_format/ref/alevel_cover_ref.png), no logos/barcodes.
% \prelimcover{total marks}{duration}{date}; heading block from \housechap; override with
% \prelimcoverhead (4 lines) and \prelimcoverline (the line under MATHEMATICS).
\makeatletter
\providecommand{\prelimcoverhead}{MAXIMUS PUNG JUN\\ H2 Mathematics (9758)\\ \housechap\\ Chapter Test}
\providecommand{\prelimcoverline}{\prelimsubject\ Test}
\newcommand{\pp@htwo}{\begin{tikzpicture}[baseline=(b.base)]
  \node[draw, line width=0.9pt, minimum width=19mm, minimum height=15mm, inner sep=0pt] (b)
    {\fontsize{38}{38}\selectfont\special{pdf:literal direct 1 Tr 0.7 w}H2\special{pdf:literal direct 0 Tr}};
  \end{tikzpicture}}
\newcommand{\prelimcover}[3]{% #1 total marks, #2 duration, #3 date line
\thispagestyle{prelimcover}%
{\coverfont\fontsize{11}{13.5}\selectfont\linespread{1}\selectfont\setlength{\parskip}{0pt}%
\vspace*{4mm}
\noindent\begin{tabular}{@{}l@{\hspace{5mm}}l@{}}\raisebox{-0.5\height}{\pp@htwo}&
  \raisebox{-0.5\height}{\begin{minipage}{120mm}\fontsize{12}{14.5}\selectfont\prelimcoverhead\end{minipage}}\end{tabular}\par
\vspace{9mm}
\noindent\rule{\linewidth}{0.8pt}\par\vspace{1.2mm}
\noindent{\fontsize{12}{14.5}\selectfont\bfseries MATHEMATICS\hfill 9758}\par\vspace{1.2mm}
\noindent\prelimcoverline\hfill\textbf{#3}\par\vspace{1.2mm}
\setbox0=\hbox{#2}\ifdim\wd0>0pt \noindent\hfill\textbf{#2}\par\fi
\vspace{3mm}
\noindent You must answer on writing paper.\par\vspace{3mm}
\noindent\begin{tabular}{@{}l@{\hspace{2mm}}l@{}}You will need:&Writing paper\\ &List of formulae and results (MF27)\end{tabular}\par
\vspace{2mm}
\noindent\rule{\linewidth}{0.8pt}\par\vspace{1.5mm}
\noindent{\bfseries INSTRUCTIONS}\par\vspace{0.5mm}
\begin{list}{\textbullet}{\setlength\leftmargin{7mm}\setlength\labelwidth{4mm}\setlength\labelsep{3mm}\setlength\itemsep{0.6mm}\setlength\topsep{0pt}\setlength\parsep{0pt}}
\item Answer \textbf{all} questions.
\item Give non-exact numerical answers correct to 3 significant figures, or 1 decimal place for angles in degrees, unless a different level of accuracy is specified in the question.
\item You should use an approved graphing calculator.
\item Unsupported answers from a graphing calculator are allowed unless a question specifically states otherwise.
\item Where unsupported answers from a graphing calculator are \textbf{not} allowed in a question, you must present the mathematical steps using mathematical notations and not calculator commands.
\item You must show all working clearly.
\end{list}
\vspace{6mm}
\noindent{\bfseries INFORMATION}\par\vspace{0.5mm}
\begin{list}{\textbullet}{\setlength\leftmargin{7mm}\setlength\labelwidth{4mm}\setlength\labelsep{3mm}\setlength\itemsep{0.6mm}\setlength\topsep{0pt}\setlength\parsep{0pt}}
\item The total mark for this paper is #1.
\item The number of marks for each question or part question is shown in brackets [\ ].
\end{list}
\vfill
\noindent\rule{\linewidth}{0.8pt}\par\vspace{1mm}
\centerline{This document has \textbf{\pageref{LastPage}} pages.}\par\vspace{2mm}
}\clearpage}
\makeatother

% ---- solutions header block (page 1 of the worked solutions; same cover style) ----
\newcommand{\prelimsolhead}[1]{% #1 total marks
{\coverfont\fontsize{10}{12}\selectfont\linespread{1}\selectfont
\noindent{\bfseries \prelimtitle}\\[0.2ex]{\bfseries Compiled by Maximus Pung Jun}\\[1.6ex]Higher 2\par\vspace{2ex}
\noindent\rule{\linewidth}{1.2pt}\par\vspace{0.4ex}
\noindent{\Large\bfseries MATHEMATICS}\hfill{\Large\bfseries 9758}\par\vspace{1.2ex}
\noindent \prelimsubject\hfill\textbf{2026}\\
WORKED SOLUTIONS\hfill\textbf{Total marks: #1}\par\vspace{1ex}
\noindent\rule{\linewidth}{0.6pt}\par}\vspace{4ex}}
% house_sol.tex: HOUSE COPY of definite_integrals/common/ri_solstyle.tex (unchanged)
% =====================================================================
% ri_solstyle.tex: worked solutions in the layout of the RI 2026 Year 6 Prelim
% "Solutions with comments" (grail.moe 26612/26613), loaded AFTER prelim_style.
%   * page 1: bold school-style block + bold underlined title; later pages: running
%     title top-left with a rule under it; footer "Page N of M" bottom right (Calibri).
%   * each question: prelim-format question text (bold number, marks right, source tag
%     bottom right), then a bordered table:  Solution | Comments  (bold header row);
%     body rows  label + [marks] (bold, narrow left column) | working | comments,
%     a full-width rule between parts, working left-aligned (fleqn).
% Rows come from the existing sources: \pt in qparts (Ch10 HW), \slabel (Ch10 Test),
% spart (Ch11). Marks under each label are taken from the question's \mk.
% =====================================================================
\makeatletter
\usepackage{lastpage}
\usepackage{adjustbox}
\usepackage{zref-savepos,zref-abspage}
\zref@addprop{savepos}{abspage}
\geometry{a4paper, top=22mm, bottom=24mm, left=22mm, right=22mm, headsep=8mm, footskip=12mm, headheight=24pt}
\providecommand{\risoltitle}{\prelimtitle: Solutions with comments}
\providecommand{\risolschool}{Compiled by Maximus Pung Jun}
\fancypagestyle{risol}{\fancyhf{}%
  \lhead{\fontsize{11}{12}\selectfont 2026 \risoltitle\\[-1.2ex]\rule{\textwidth}{0.6pt}}%
  \rfoot{\calibrifont\fontsize{11}{13}\selectfont Page \textbf{\thepage} of \textbf{\pageref*{LastPage}}}%
  \renewcommand{\headrulewidth}{0pt}\renewcommand{\footrulewidth}{0pt}}
\fancypagestyle{risolfirst}{\fancyhf{}%
  \rfoot{\calibrifont\fontsize{11}{13}\selectfont Page \textbf{\thepage} of \textbf{\pageref*{LastPage}}}%
  \renewcommand{\headrulewidth}{0pt}\renewcommand{\footrulewidth}{0pt}}
\pagestyle{risol}
\newif\ifri@fresh
\newcommand{\risolfront}{\thispagestyle{risolfirst}\global\ri@freshtrue%
  {\noindent\bfseries\linespread{1}\selectfont \risolschool\\ H2 Mathematics (9758)\\ 2026\par}%
  \vspace{4ex}{\noindent\bfseries\underline{2026 \risoltitle}\par}\vspace{3ex}}
% ---- table geometry (RI: label 42pt | working 284pt | comments 142pt of 468pt) ----
\newlength{\ri@L}\setlength{\ri@L}{15mm}
\newlength{\ri@C}\setlength{\ri@C}{40mm}
\newlength{\ri@pad}\setlength{\ri@pad}{1.8mm}
\newif\ifri@sol \newif\ifri@first
\newcount\ri@depth
\def\ri@outer{}
\newtcolorbox{risol}{enhanced, breakable, sharp corners, frame hidden, colback=white,
  boxrule=0pt, boxsep=0pt, left=\dimexpr\ri@L+\ri@pad\relax, right=\dimexpr\ri@C+\ri@pad\relax,
  top=1.2mm, bottom=2.5mm, toptitle=0.6mm, bottomtitle=0.6mm, lefttitle=\ri@pad, righttitle=0pt,
  colbacktitle=white, coltitle=black, fonttitle=\bfseries,
  title={Solution\hfill\makebox[\ri@C][l]{\hspace{\ri@pad}Comments}},
  before skip=2ex, after skip=2ex,
  overlay unbroken and first={\draw[line width=0.5pt] (frame.south west) rectangle (frame.north east);
     \draw[line width=0.5pt] (title.south west) -- (title.south east);
     \draw[line width=0.5pt] ([xshift=\ri@L]interior.north west) -- ([xshift=\ri@L]frame.south west);
     \draw[line width=0.5pt] ([xshift=-\ri@C]frame.north east) -- ([xshift=-\ri@C]frame.south east);},
  overlay middle and last={\draw[line width=0.5pt] (frame.south west) rectangle (frame.north east);
     \draw[line width=0.5pt] ([xshift=\ri@L]frame.north west) -- ([xshift=\ri@L]frame.south west);
     \draw[line width=0.5pt] ([xshift=-\ri@C]frame.north east) -- ([xshift=-\ri@C]frame.south east);},
  after upper={\ri@endrow\protected@write\@auxout{}{\string\ri@settotal{\ri@qid}{\the\ri@sumH}}},
  before upper={\global\ri@sumH=14mm\ri@soltrue\ri@firsttrue\ri@depth=0 \parindent=0pt \parskip=0.6ex
     \predisplaypenalty=100 
     \linespread{1.12}\selectfont\raggedright\@fleqntrue\@mathmargin=0pt
     \abovedisplayskip=5pt plus 2pt\belowdisplayskip=5pt plus 2pt
     \abovedisplayshortskip=3pt plus 1pt\belowdisplayshortskip=4pt plus 2pt
     \ifcsname ri@hp@\ri@qid\endcsname\else\ri@row{}\fi}}
\def\ri@qid{}
% one row: full-width rule (glue leaders, so it vanishes at a page break where the frame
% edge takes its place; not before the first row), then bold label + bold [marks]
% in the left column, on the baseline of the first line of the working
\ExplSyntaxOn
\tl_new:N \g__ri_tmp_tl
\cs_new_protected:Npn \ri_setkey:Nn #1#2
  { \group_begin: \cs_set_eq:NN \chg \use_ii:nn \cs_set_eq:NN \chgm \use_ii:nn
    \cs_set_eq:NN \textbf \use:n \cs_set_eq:NN \mathrm \use:n
    \tl_gset:Nx \g__ri_tmp_tl {#2} \group_end:
    \tl_set_eq:NN #1 \g__ri_tmp_tl \tl_remove_all:Nn #1 {~} }
\cs_set_eq:NN \ri@setkey \ri_setkey:Nn
\ExplSyntaxOff
% ---- page-break rules (positions from the previous run, zref labels in the .aux) ----
% * a row (one part) never breaks across pages: if it does not fit on the rest of the
%   page it starts the next page (forced break between rows); only a row taller than a
%   whole page may break inside;
% * a whole solution that fits on one page is kept together (needspace before the table);
%   a longer one starts where at least its first row fits;
% * a row is at least as tall as its label + [marks].
\newcount\ri@n \newbox\ri@labbox \def\ri@open{0}
\newdimen\ri@ybot \newdimen\ri@ytop \newdimen\ri@h \newdimen\ri@sumH
\AtBeginDocument{\ri@ybot=\dimexpr\paperheight-1in-\voffset-\topmargin-\headheight-\headsep-\textheight\relax
  \ri@ytop=\dimexpr\ri@ybot+\textheight\relax}
\def\ri@pg#1{\zref@extractdefault{#1}{abspage}{-1}}
% \ri@rowh{n}: height of row n in the previous run -> \ri@h (0pt if unknown)
\newcommand{\ri@rowh}[1]{\ri@h=0pt
  \ifnum\ri@pg{ri-s#1}>0 \ifnum\ri@pg{ri-e#1}>0
    \ifnum\ri@pg{ri-s#1}=\ri@pg{ri-e#1}\relax \ri@h=\dimexpr\zposy{ri-s#1}sp-\zposy{ri-e#1}sp\relax
    \else \ri@h=\dimexpr\zposy{ri-s#1}sp-\ri@ybot+\ri@ytop-\zposy{ri-e#1}sp-6mm\relax\fi\fi\fi}
\newcommand{\ri@setfirst}[2]{\expandafter\gdef\csname ri@first@#1\endcsname{#2}}
\newcommand{\ri@setkept}[1]{\expandafter\gdef\csname ri@kept@#1\endcsname{}}
\newcommand{\ri@settotal}[2]{\expandafter\gdef\csname ri@total@#1\endcsname{#2}}
\newcommand{\ri@endrow}{\ifnum\ri@open=1 \par\zsavepos{ri-e\the\ri@n}%
  \ifnum\ri@pg{ri-s\the\ri@n}=\ri@pg{ri-e\the\ri@n}\relax
    \dimen@=\dimexpr\ri@curH sp+1.6mm-\zposy{ri-s\the\ri@n}sp+\zposy{ri-e\the\ri@n}sp\relax
    \ifdim\dimen@>0pt \nobreak\kern\dimen@\penalty\z@\fi\fi
  \ri@rowh{\the\ri@n}\global\advance\ri@sumH\dimexpr\ri@h+3.5mm\relax
  \gdef\ri@open{0}\fi}
\newcommand{\ri@row}[1]{\par\ri@endrow\global\advance\ri@n\@ne\gdef\ri@open{1}\@beginparpenalty\@M
  \ifri@first\global\ri@firstfalse\ri@firstfalse
    \protected@write\@auxout{}{\string\ri@setfirst{\ri@qid}{\the\ri@n}}%
  \else
    \zsavepos{ri-p\the\ri@n}%
    % push the row to the next page if it would not fit on the rest of this one
    \ri@rowh{\the\ri@n}%
    \ifnum\ri@pg{ri-p\the\ri@n}>0 \ifdim\ri@h>0pt
      \ifdim\dimexpr\ri@h+4mm\relax<\dimexpr\textheight-10mm\relax
        \ifdim\dimexpr\ri@h+7mm\relax>\dimexpr\zposy{ri-p\the\ri@n}sp-\ri@ybot\relax
          \penalty-\@M\fi\fi\fi\fi
    \vskip0.6ex\nointerlineskip
    \leaders\hbox{\hskip-\dimexpr\ri@L+\ri@pad\relax\vrule width\dimexpr\linewidth+\ri@L+\ri@C+2\ri@pad\relax height0.5pt depth0pt}\vskip0.5pt
    \nobreak\vskip1.4mm\fi
  \ri@setkey\ri@key{\ri@qid#1}\def\ri@lab{#1}%
  \setbox\ri@labbox=\vtop{\hsize=\dimexpr\ri@L-1.6mm\relax\bfseries\raggedright\parskip=0pt
      \strut\ri@lab\ifcsname ri@mk@\ri@key\endcsname\ifx\ri@lab\@empty\else\\\fi
      [\csname ri@mk@\ri@key\endcsname]\fi\strut\par}%
  \xdef\ri@curH{\number\dimexpr\ht\ri@labbox+\dp\ri@labbox\relax}%
  \nointerlineskip\zsavepos{ri-s\the\ri@n}\nointerlineskip\nobreak
  \vbox to 0pt{\hbox{\hskip-\dimexpr\ri@L+\ri@pad-1.2mm\relax\box\ri@labbox}\vss}\nobreak\prevdepth=\dp\strutbox
  \vskip-\parskip\ignorespaces}
% ---- qparts / \pt: question mode records labels (for the marks); solution mode makes rows
\NewCommandCopy\ri@qparts@begin\qparts \let\ri@qparts@end\endqparts
\renewenvironment{qparts}[1][2.6em]{\advance\ri@depth\@ne
  \ifri@sol\else\ri@qparts@begin[#1]\fi}{\ifri@sol\par\else\ri@qparts@end\fi}
\def\ri@qpartsname{qparts}
\renewcommand{\pt}[1]{\ifnum\ri@depth<2 \gdef\ri@outer{#1}\gdef\ri@cur{#1}\else\expandafter\gdef\expandafter\ri@cur\expandafter{\ri@outer#1}\fi
  \ifri@sol\expandafter\ri@solpt\else\expandafter\ri@qpt\fi{#1}}
\def\ri@qpt#1{\expandafter\gdef\csname ri@hp@\ri@qid\endcsname{}\item[#1]}
\def\ri@solpt#1{\@ifnextchar\begin\ri@ptpeek{\ri@row{\ri@cur}}}
\def\ri@ptpeek\begin#1{\def\ri@tmp{#1}\ifx\ri@tmp\ri@qpartsname
   \expandafter\@firstoftwo\else\expandafter\@secondoftwo\fi
  {\begin{qparts}}{\ri@row{\ri@cur}\begin{#1}}}
% marks: record [n] against the current question/part
\let\ri@mk\mk
\renewcommand{\mk}[1]{\ri@setkey\ri@k{\ri@qid\ri@cur}\ri@setkey\ri@mv{#1}\expandafter\xdef\csname ri@mk@\ri@k\endcsname{\ri@mv}\ri@mk{#1}}
% \slabel (Ch10 Test answers): letter labels set the outer part, roman ones are appended;
% consecutive \slabel calls merge into one row label
\def\ri@romans{(i)(ii)(iii)(iv)(v)(vi)(vii)}
\def\ri@firstpart(#1)#2\@nil{(#1)}
\def\slabel#1{\ifri@sol\expandafter\ri@solslabel\else\expandafter\ri@qslabel\fi{#1}}
\def\ri@qslabel#1{\par\smallskip\noindent\textbf{#1}\par\nobreak}
\def\ri@solslabel#1{\@expandtwoargs\in@{#1}{\ri@romans}%
  \ifin@\expandafter\gdef\expandafter\ri@cur\expandafter{\ri@outer#1}\else\gdef\ri@cur{#1}\xdef\ri@outer{\ri@firstpart#1\@nil}\fi
  \@ifnextchar\slabel{}{\ri@row{\ri@cur}}}
% Ch11 spart{label}: a row; soln: the table
\def\spart#1{\gdef\ri@cur{#1}\ri@row{#1}}\def\endspart{\par}
\def\soln{\@ifnextchar[{\ri@solnopt}{\risol}}\def\ri@solnopt[#1]{\risol}\def\endsoln{\endrisol}
% ---- working that started on the label line is left-aligned in the working column
\AtBeginDocument{\@ifundefined{topalign}{}{%
  \renewenvironment{topalign}{\par\noindent$\displaystyle\begin{aligned}[t]}{\end{aligned}$\par\smallskip}%
  \renewcommand{\topdisplay}[1]{\par\noindent$\displaystyle #1$\par\smallskip}}}
% ---- Ch11 side-by-side (text | sketch): the sketch shrinks to fit its half of the working column
\AtBeginDocument{\@ifundefined{sbsplain}{}{%
  \renewenvironment{sbsplain}[1]{%
    \sbs@start\begin{minipage}[t]{#1\linewidth}\vspace{0pt}\sbs@shift\raggedright%
    \def\sbsnext{\end{minipage}\hfill\begin{minipage}[t]{\dimexpr0.97\linewidth-#1\linewidth\relax}\vspace{0pt}\sbs@shift\centering
      \begin{adjustbox}{max width=\linewidth}}%
  }{\end{adjustbox}\end{minipage}\par}}}
% ---- figures in the working column shrink to fit it (only if wider), colours kept
\newcommand{\rifitfig}[1]{\expandafter\let\csname ri@orig@\string#1\endcsname#1%
  \def#1{\@ifnextchar\bgroup{\ri@fitarg#1}{\ri@fitnoarg#1}}}
\def\ri@fitarg#1#2{\ifri@sol\adjustbox{max width=\linewidth}{\csname ri@orig@\string#1\endcsname{#2}}\else\csname ri@orig@\string#1\endcsname{#2}\fi}
\def\ri@fitnoarg#1{\ifri@sol\adjustbox{max width=\linewidth}{\csname ri@orig@\string#1\endcsname}\else\csname ri@orig@\string#1\endcsname\fi}
\@ifundefined{solfig}{}{\renewcommand{\solfig}[2][]{\par\begin{center}\adjustbox{max width=\linewidth}{\begin{tikzpicture}\begin{axis}[qfig, width=6.6cm, height=5cm, #1]#2\end{axis}\end{tikzpicture}}\end{center}}}
% \ricomment{text}: text in the Comments column of the current row (call right after \pt,
% \slabel or \spart). The row is made at least as tall as the comment.
\newbox\ri@cmtbox
\newcommand{\ricomment}[1]{\par\setbox\ri@cmtbox=\vtop{\hsize=\dimexpr\ri@C-2\ri@pad\relax
    \parskip=0pt\raggedright\small{\color{black!60}\strut#1\strut\par}}%
  \ifdim\dimexpr\ht\ri@cmtbox+\dp\ri@cmtbox\relax>\ri@curH sp\xdef\ri@curH{\number\dimexpr\ht\ri@cmtbox+\dp\ri@cmtbox\relax}\fi
  \nointerlineskip\vbox to 0pt{\hbox to\linewidth{\hskip\dimexpr\linewidth+2\ri@pad\relax\box\ri@cmtbox\hss}\vss}\nobreak\prevdepth=\dp\strutbox
  \vskip-\parskip\ignorespaces}
% \putTFlong{n}: the True/False question n with its answer table allowed to break across
% pages (the body of \defQ{n} must be: intro text, then \begin{tftable}...\end{tftable});
% the solutions copy prints it with tflongtable.
\renewcommand{\putTFlong}[1]{\ri@newq\begingroup\let\tftable\tflongtable\let\endtftable\endtflongtable
  \stepcounter{diq}\expandafter\xdef\csname diLab@#1\endcsname{\thediq}%
  \par\noindent\hbox to 2.6em{\textbf{\thediq}\hfil}\setbox0=\hbox{\di@tag{#1}}\ifdim\wd0=0pt\else{\small\textcolor{black!60}{[\di@tag{#1}]}}\par\noindent\hskip2.6em\fi\ignorespaces\csname diQbody@#1\endcsname\par
  \endgroup\vspace{3ex}}
% ---- question + solution; each question starts on a new page (as RI)
\def\ri@qid{}
\let\ri@putQ\putQ
\renewcommand{\putQ}[1]{\gdef\ri@qid{#1}\global\ri@depth=0 \gdef\ri@cur{}\ri@putQ{#1}}
% \risolQ{n}{body macro name}: question n + its solution body in the table;
% \risolQraw: body already contains its own soln environment (Ch11 Test)
\newcommand{\ri@newq}{\ifri@fresh\global\ri@freshfalse\else\clearpage\fi}
% keep the whole table together if it fits on a page; otherwise make sure at least its
% header and first row fit before starting it
\newcommand{\ri@keep}[1]{\par
  \ifcsname ri@total@#1\endcsname
    % hysteresis: a table kept whole in the previous run stays whole up to textheight-4.5mm,
    % so a total near the limit cannot flip between runs
    \ifdim\csname ri@total@#1\endcsname<\dimexpr\textheight-\ifcsname ri@kept@#1\endcsname4.5\else8\fi mm\relax
      \protected@write\@auxout{}{\string\ri@setkept{#1}}%
      \needspace{\dimexpr\csname ri@total@#1\endcsname+4mm\relax}%
    \else\ifcsname ri@first@#1\endcsname
      \ri@rowh{\csname ri@first@#1\endcsname}%
      \ifdim\ri@h<\dimexpr\textheight-20mm\relax\needspace{\dimexpr\ri@h+18mm\relax}\fi\fi\fi\fi}
\newcommand{\risolQ}[2]{\ri@newq\putQ{#1}\vspace{-3ex}\ri@keep{#1}\begin{risol}\csname #2@#1\endcsname\par\end{risol}}
\newcommand{\risolQraw}[2]{\ri@newq\putQ{#1}\vspace{-3ex}\ri@keep{#1}\csname #2@#1\endcsname\par}
\renewcommand{\putQS}[1]{\risolQ{#1}{diS}}
% end marker only if it fits on the last page (no page holding just the marker)
\renewcommand{\diend}[1][END OF SOLUTIONS]{\par\ifdim\dimexpr\pagegoal-\pagetotal\relax<4\baselineskip\else\vspace{1.5ex}\begin{center}\textbf{#1}\end{center}\fi}
\makeatother
\tfanswerstrue\fi
  \ifx\housexm\housextest % house_prelim.tex: HOUSE COPY of definite_integrals/common/prelim_style.tex; [HOUSE] titles from \housechapter
% =====================================================================
% prelim_style.tex: makes the Test look like a Singapore JC H2 Mathematics
% 9758 prelim paper (studied from ACJC/RI 2026 P1/P2 on the box:
% body Times New Roman + Times-style maths, cover in Arial).
% Fonts: prelim_fonts.tex (real Times New Roman/Arial + Times-style maths).
% Shared by Ch10 (di_style) and Ch11 (de_style; sets \prelimtitle/\prelimsubject). Loaded AFTER the chapter style (overrides its layout for the test only).
% =====================================================================
% [HOUSE] fonts loaded by house.tex
\providecommand{\prelimtitle}{\MakeUppercase{\housechap} TEST}
\providecommand{\prelimsubject}{\housechap}
\geometry{a4paper, top=22mm, bottom=24mm, left=25mm, right=25mm, headsep=6mm, footskip=12mm}
\setlength{\parindent}{0pt}
\setlength{\parskip}{0.9ex plus 0.3ex minus 0.2ex}
\linespread{1.3}
\setlength{\emergencystretch}{3em}
\clubpenalty=10000 \widowpenalty=10000 \interlinepenalty=10000
\AtBeginDocument{\fontsize{12}{14.4}\selectfont}
\renewcommand{\normalsize}{\fontsize{12}{14.4}\selectfont}

% ---- figures: plain black/grey like a printed paper ----------------------
\ifdefined\prelimkeepcolour\else % question paper: black/grey figures; solutions keep colour
\pgfplotsset{curve/.style={black, semithick, smooth, samples=150},
  line2/.style={black, semithick, samples=2},
  sketch/.append style={axis line style={thin, black}, every axis label/.append style={black}}}
\tikzset{region/.style={fill=black!15, draw=none}, regionB/.style={fill=black!30, draw=none},
  rect/.style={fill=black!8, draw=black, thin}, ax3/.style={-{Stealth[length=1.5mm]}, thin, black}}
\colorlet{blue}{black}\colorlet{red}{black}\colorlet{orange}{black}
\fi

% ---- header / footer: page number centred top (bold), [Turn over at bottom right
\makeatletter
\providecommand{\prelimfootlabel}{\prelimtitle\ 2026}
\newcommand{\pp@footer}{{\coverfont\scriptsize \prelimfootlabel\hspace{3em}H2 MATHEMATICS 9758}}
\fancypagestyle{prelim}{\fancyhf{}%
  \chead{\bfseries\thepage}%
  \lfoot{\pp@footer}%
  \rfoot{\ifnum\value{page}<\getpagerefnumber{LastPage}\relax\textbf{[Turn over}\fi}%
  \renewcommand{\headrulewidth}{0pt}\renewcommand{\footrulewidth}{0pt}}
\fancypagestyle{prelimcover}{\fancyhf{}%
  \rfoot{\coverfont\textbf{[Turn over}}%
  \renewcommand{\headrulewidth}{0pt}\renewcommand{\footrulewidth}{0pt}}
\pagestyle{prelim}

% ---- marks: [n] regular weight, right-aligned at the end of the part ----------
\renewcommand{\mk}[1]{{\unskip\nobreak\hfil\penalty50\hskip1.5em\hbox{}\nobreak\hfill[#1]%
  \parfillskip=0pt\finalhyphendemerits=0\par}}

% ---- questions: bold number at the left margin, text indented; grey [tag] straight after the number (3B)
\newsavebox{\pp@qmeas}
\newcommand{\pp@qlist}[1]{%
  \begin{list}{}{\setlength\leftmargin{2.6em}\setlength\labelwidth{2.1em}\setlength\labelsep{0.5em}%
     \setlength\topsep{0pt}\setlength\itemsep{0pt}\setlength\parsep{0.9ex}%
     \renewcommand\makelabel[1]{\textbf{##1}\hfil}}%
  \item[\thediq]\@endparpenalty\@M\setbox0=\hbox{\di@tag{#1}}\ifdim\wd0=0pt\else{\small\textcolor{black!60}{[\di@tag{#1}]}}\par\nobreak\fi\csname diQbody@#1\endcsname\par\penalty\@M
  \end{list}}
% a question starts on a new page unless it fits whole on the current one
% (questions taller than a page may break, but only between parts)
\renewcommand{\putQ}[1]{\stepcounter{diq}%
  \expandafter\xdef\csname diLab@#1\endcsname{\thediq}%
  \par\setbox\pp@qmeas=\vbox{\pp@qlist{#1}}%
  \ifdim\dimexpr\ht\pp@qmeas+\dp\pp@qmeas\relax<0.95\textheight
    \needspace{\dimexpr\ht\pp@qmeas+\dp\pp@qmeas\relax}\else\needspace{14\baselineskip}\fi
  \zsavepos{q-s\thediq}\pp@qlist{#1}\par\zsavepos{q-e\thediq}\vspace{5ex plus 2ex minus 1ex}}
% part labels (a)/(i) bold, indented under the question text
\renewenvironment{qparts}[1][2.6em]{\begin{list}{}{\setlength\leftmargin{#1}%
  \setlength\labelwidth{\dimexpr#1-0.5em\relax}\setlength\labelsep{0.5em}%
  \setlength\itemsep{1.4ex}\setlength\topsep{1.2ex}\setlength\parsep{0.8ex}%
  \renewcommand\makelabel[1]{\textbf{##1}\hfil}}}{\end{list}}
\makeatother

% ---- cover page --------------------------------------------------------------
% A-level H2 style cover (9 Nov 2026; reference house_format/ref/alevel_cover_ref.png), no logos/barcodes.
% \prelimcover{total marks}{duration}{date}; heading block from \housechap; override with
% \prelimcoverhead (4 lines) and \prelimcoverline (the line under MATHEMATICS).
\makeatletter
\providecommand{\prelimcoverhead}{MAXIMUS PUNG JUN\\ H2 Mathematics (9758)\\ \housechap\\ Chapter Test}
\providecommand{\prelimcoverline}{\prelimsubject\ Test}
\newcommand{\pp@htwo}{\begin{tikzpicture}[baseline=(b.base)]
  \node[draw, line width=0.9pt, minimum width=19mm, minimum height=15mm, inner sep=0pt] (b)
    {\fontsize{38}{38}\selectfont\special{pdf:literal direct 1 Tr 0.7 w}H2\special{pdf:literal direct 0 Tr}};
  \end{tikzpicture}}
\newcommand{\prelimcover}[3]{% #1 total marks, #2 duration, #3 date line
\thispagestyle{prelimcover}%
{\coverfont\fontsize{11}{13.5}\selectfont\linespread{1}\selectfont\setlength{\parskip}{0pt}%
\vspace*{4mm}
\noindent\begin{tabular}{@{}l@{\hspace{5mm}}l@{}}\raisebox{-0.5\height}{\pp@htwo}&
  \raisebox{-0.5\height}{\begin{minipage}{120mm}\fontsize{12}{14.5}\selectfont\prelimcoverhead\end{minipage}}\end{tabular}\par
\vspace{9mm}
\noindent\rule{\linewidth}{0.8pt}\par\vspace{1.2mm}
\noindent{\fontsize{12}{14.5}\selectfont\bfseries MATHEMATICS\hfill 9758}\par\vspace{1.2mm}
\noindent\prelimcoverline\hfill\textbf{#3}\par\vspace{1.2mm}
\setbox0=\hbox{#2}\ifdim\wd0>0pt \noindent\hfill\textbf{#2}\par\fi
\vspace{3mm}
\noindent You must answer on writing paper.\par\vspace{3mm}
\noindent\begin{tabular}{@{}l@{\hspace{2mm}}l@{}}You will need:&Writing paper\\ &List of formulae and results (MF27)\end{tabular}\par
\vspace{2mm}
\noindent\rule{\linewidth}{0.8pt}\par\vspace{1.5mm}
\noindent{\bfseries INSTRUCTIONS}\par\vspace{0.5mm}
\begin{list}{\textbullet}{\setlength\leftmargin{7mm}\setlength\labelwidth{4mm}\setlength\labelsep{3mm}\setlength\itemsep{0.6mm}\setlength\topsep{0pt}\setlength\parsep{0pt}}
\item Answer \textbf{all} questions.
\item Give non-exact numerical answers correct to 3 significant figures, or 1 decimal place for angles in degrees, unless a different level of accuracy is specified in the question.
\item You should use an approved graphing calculator.
\item Unsupported answers from a graphing calculator are allowed unless a question specifically states otherwise.
\item Where unsupported answers from a graphing calculator are \textbf{not} allowed in a question, you must present the mathematical steps using mathematical notations and not calculator commands.
\item You must show all working clearly.
\end{list}
\vspace{6mm}
\noindent{\bfseries INFORMATION}\par\vspace{0.5mm}
\begin{list}{\textbullet}{\setlength\leftmargin{7mm}\setlength\labelwidth{4mm}\setlength\labelsep{3mm}\setlength\itemsep{0.6mm}\setlength\topsep{0pt}\setlength\parsep{0pt}}
\item The total mark for this paper is #1.
\item The number of marks for each question or part question is shown in brackets [\ ].
\end{list}
\vfill
\noindent\rule{\linewidth}{0.8pt}\par\vspace{1mm}
\centerline{This document has \textbf{\pageref{LastPage}} pages.}\par\vspace{2mm}
}\clearpage}
\makeatother

% ---- solutions header block (page 1 of the worked solutions; same cover style) ----
\newcommand{\prelimsolhead}[1]{% #1 total marks
{\coverfont\fontsize{10}{12}\selectfont\linespread{1}\selectfont
\noindent{\bfseries \prelimtitle}\\[0.2ex]{\bfseries Compiled by Maximus Pung Jun}\\[1.6ex]Higher 2\par\vspace{2ex}
\noindent\rule{\linewidth}{1.2pt}\par\vspace{0.4ex}
\noindent{\Large\bfseries MATHEMATICS}\hfill{\Large\bfseries 9758}\par\vspace{1.2ex}
\noindent \prelimsubject\hfill\textbf{2026}\\
WORKED SOLUTIONS\hfill\textbf{Total marks: #1}\par\vspace{1ex}
\noindent\rule{\linewidth}{0.6pt}\par}\vspace{4ex}}
\fi
  \ifx\housexm\housextestsol \def\prelimkeepcolour{1}% house_prelim.tex: HOUSE COPY of definite_integrals/common/prelim_style.tex; [HOUSE] titles from \housechapter
% =====================================================================
% prelim_style.tex: makes the Test look like a Singapore JC H2 Mathematics
% 9758 prelim paper (studied from ACJC/RI 2026 P1/P2 on the box:
% body Times New Roman + Times-style maths, cover in Arial).
% Fonts: prelim_fonts.tex (real Times New Roman/Arial + Times-style maths).
% Shared by Ch10 (di_style) and Ch11 (de_style; sets \prelimtitle/\prelimsubject). Loaded AFTER the chapter style (overrides its layout for the test only).
% =====================================================================
% [HOUSE] fonts loaded by house.tex
\providecommand{\prelimtitle}{\MakeUppercase{\housechap} TEST}
\providecommand{\prelimsubject}{\housechap}
\geometry{a4paper, top=22mm, bottom=24mm, left=25mm, right=25mm, headsep=6mm, footskip=12mm}
\setlength{\parindent}{0pt}
\setlength{\parskip}{0.9ex plus 0.3ex minus 0.2ex}
\linespread{1.3}
\setlength{\emergencystretch}{3em}
\clubpenalty=10000 \widowpenalty=10000 \interlinepenalty=10000
\AtBeginDocument{\fontsize{12}{14.4}\selectfont}
\renewcommand{\normalsize}{\fontsize{12}{14.4}\selectfont}

% ---- figures: plain black/grey like a printed paper ----------------------
\ifdefined\prelimkeepcolour\else % question paper: black/grey figures; solutions keep colour
\pgfplotsset{curve/.style={black, semithick, smooth, samples=150},
  line2/.style={black, semithick, samples=2},
  sketch/.append style={axis line style={thin, black}, every axis label/.append style={black}}}
\tikzset{region/.style={fill=black!15, draw=none}, regionB/.style={fill=black!30, draw=none},
  rect/.style={fill=black!8, draw=black, thin}, ax3/.style={-{Stealth[length=1.5mm]}, thin, black}}
\colorlet{blue}{black}\colorlet{red}{black}\colorlet{orange}{black}
\fi

% ---- header / footer: page number centred top (bold), [Turn over at bottom right
\makeatletter
\providecommand{\prelimfootlabel}{\prelimtitle\ 2026}
\newcommand{\pp@footer}{{\coverfont\scriptsize \prelimfootlabel\hspace{3em}H2 MATHEMATICS 9758}}
\fancypagestyle{prelim}{\fancyhf{}%
  \chead{\bfseries\thepage}%
  \lfoot{\pp@footer}%
  \rfoot{\ifnum\value{page}<\getpagerefnumber{LastPage}\relax\textbf{[Turn over}\fi}%
  \renewcommand{\headrulewidth}{0pt}\renewcommand{\footrulewidth}{0pt}}
\fancypagestyle{prelimcover}{\fancyhf{}%
  \rfoot{\coverfont\textbf{[Turn over}}%
  \renewcommand{\headrulewidth}{0pt}\renewcommand{\footrulewidth}{0pt}}
\pagestyle{prelim}

% ---- marks: [n] regular weight, right-aligned at the end of the part ----------
\renewcommand{\mk}[1]{{\unskip\nobreak\hfil\penalty50\hskip1.5em\hbox{}\nobreak\hfill[#1]%
  \parfillskip=0pt\finalhyphendemerits=0\par}}

% ---- questions: bold number at the left margin, text indented; grey [tag] straight after the number (3B)
\newsavebox{\pp@qmeas}
\newcommand{\pp@qlist}[1]{%
  \begin{list}{}{\setlength\leftmargin{2.6em}\setlength\labelwidth{2.1em}\setlength\labelsep{0.5em}%
     \setlength\topsep{0pt}\setlength\itemsep{0pt}\setlength\parsep{0.9ex}%
     \renewcommand\makelabel[1]{\textbf{##1}\hfil}}%
  \item[\thediq]\@endparpenalty\@M\setbox0=\hbox{\di@tag{#1}}\ifdim\wd0=0pt\else{\small\textcolor{black!60}{[\di@tag{#1}]}}\par\nobreak\fi\csname diQbody@#1\endcsname\par\penalty\@M
  \end{list}}
% a question starts on a new page unless it fits whole on the current one
% (questions taller than a page may break, but only between parts)
\renewcommand{\putQ}[1]{\stepcounter{diq}%
  \expandafter\xdef\csname diLab@#1\endcsname{\thediq}%
  \par\setbox\pp@qmeas=\vbox{\pp@qlist{#1}}%
  \ifdim\dimexpr\ht\pp@qmeas+\dp\pp@qmeas\relax<0.95\textheight
    \needspace{\dimexpr\ht\pp@qmeas+\dp\pp@qmeas\relax}\else\needspace{14\baselineskip}\fi
  \zsavepos{q-s\thediq}\pp@qlist{#1}\par\zsavepos{q-e\thediq}\vspace{5ex plus 2ex minus 1ex}}
% part labels (a)/(i) bold, indented under the question text
\renewenvironment{qparts}[1][2.6em]{\begin{list}{}{\setlength\leftmargin{#1}%
  \setlength\labelwidth{\dimexpr#1-0.5em\relax}\setlength\labelsep{0.5em}%
  \setlength\itemsep{1.4ex}\setlength\topsep{1.2ex}\setlength\parsep{0.8ex}%
  \renewcommand\makelabel[1]{\textbf{##1}\hfil}}}{\end{list}}
\makeatother

% ---- cover page --------------------------------------------------------------
% A-level H2 style cover (9 Nov 2026; reference house_format/ref/alevel_cover_ref.png), no logos/barcodes.
% \prelimcover{total marks}{duration}{date}; heading block from \housechap; override with
% \prelimcoverhead (4 lines) and \prelimcoverline (the line under MATHEMATICS).
\makeatletter
\providecommand{\prelimcoverhead}{MAXIMUS PUNG JUN\\ H2 Mathematics (9758)\\ \housechap\\ Chapter Test}
\providecommand{\prelimcoverline}{\prelimsubject\ Test}
\newcommand{\pp@htwo}{\begin{tikzpicture}[baseline=(b.base)]
  \node[draw, line width=0.9pt, minimum width=19mm, minimum height=15mm, inner sep=0pt] (b)
    {\fontsize{38}{38}\selectfont\special{pdf:literal direct 1 Tr 0.7 w}H2\special{pdf:literal direct 0 Tr}};
  \end{tikzpicture}}
\newcommand{\prelimcover}[3]{% #1 total marks, #2 duration, #3 date line
\thispagestyle{prelimcover}%
{\coverfont\fontsize{11}{13.5}\selectfont\linespread{1}\selectfont\setlength{\parskip}{0pt}%
\vspace*{4mm}
\noindent\begin{tabular}{@{}l@{\hspace{5mm}}l@{}}\raisebox{-0.5\height}{\pp@htwo}&
  \raisebox{-0.5\height}{\begin{minipage}{120mm}\fontsize{12}{14.5}\selectfont\prelimcoverhead\end{minipage}}\end{tabular}\par
\vspace{9mm}
\noindent\rule{\linewidth}{0.8pt}\par\vspace{1.2mm}
\noindent{\fontsize{12}{14.5}\selectfont\bfseries MATHEMATICS\hfill 9758}\par\vspace{1.2mm}
\noindent\prelimcoverline\hfill\textbf{#3}\par\vspace{1.2mm}
\setbox0=\hbox{#2}\ifdim\wd0>0pt \noindent\hfill\textbf{#2}\par\fi
\vspace{3mm}
\noindent You must answer on writing paper.\par\vspace{3mm}
\noindent\begin{tabular}{@{}l@{\hspace{2mm}}l@{}}You will need:&Writing paper\\ &List of formulae and results (MF27)\end{tabular}\par
\vspace{2mm}
\noindent\rule{\linewidth}{0.8pt}\par\vspace{1.5mm}
\noindent{\bfseries INSTRUCTIONS}\par\vspace{0.5mm}
\begin{list}{\textbullet}{\setlength\leftmargin{7mm}\setlength\labelwidth{4mm}\setlength\labelsep{3mm}\setlength\itemsep{0.6mm}\setlength\topsep{0pt}\setlength\parsep{0pt}}
\item Answer \textbf{all} questions.
\item Give non-exact numerical answers correct to 3 significant figures, or 1 decimal place for angles in degrees, unless a different level of accuracy is specified in the question.
\item You should use an approved graphing calculator.
\item Unsupported answers from a graphing calculator are allowed unless a question specifically states otherwise.
\item Where unsupported answers from a graphing calculator are \textbf{not} allowed in a question, you must present the mathematical steps using mathematical notations and not calculator commands.
\item You must show all working clearly.
\end{list}
\vspace{6mm}
\noindent{\bfseries INFORMATION}\par\vspace{0.5mm}
\begin{list}{\textbullet}{\setlength\leftmargin{7mm}\setlength\labelwidth{4mm}\setlength\labelsep{3mm}\setlength\itemsep{0.6mm}\setlength\topsep{0pt}\setlength\parsep{0pt}}
\item The total mark for this paper is #1.
\item The number of marks for each question or part question is shown in brackets [\ ].
\end{list}
\vfill
\noindent\rule{\linewidth}{0.8pt}\par\vspace{1mm}
\centerline{This document has \textbf{\pageref{LastPage}} pages.}\par\vspace{2mm}
}\clearpage}
\makeatother

% ---- solutions header block (page 1 of the worked solutions; same cover style) ----
\newcommand{\prelimsolhead}[1]{% #1 total marks
{\coverfont\fontsize{10}{12}\selectfont\linespread{1}\selectfont
\noindent{\bfseries \prelimtitle}\\[0.2ex]{\bfseries Compiled by Maximus Pung Jun}\\[1.6ex]Higher 2\par\vspace{2ex}
\noindent\rule{\linewidth}{1.2pt}\par\vspace{0.4ex}
\noindent{\Large\bfseries MATHEMATICS}\hfill{\Large\bfseries 9758}\par\vspace{1.2ex}
\noindent \prelimsubject\hfill\textbf{2026}\\
WORKED SOLUTIONS\hfill\textbf{Total marks: #1}\par\vspace{1ex}
\noindent\rule{\linewidth}{0.6pt}\par}\vspace{4ex}}
% house_sol.tex: HOUSE COPY of definite_integrals/common/ri_solstyle.tex (unchanged)
% =====================================================================
% ri_solstyle.tex: worked solutions in the layout of the RI 2026 Year 6 Prelim
% "Solutions with comments" (grail.moe 26612/26613), loaded AFTER prelim_style.
%   * page 1: bold school-style block + bold underlined title; later pages: running
%     title top-left with a rule under it; footer "Page N of M" bottom right (Calibri).
%   * each question: prelim-format question text (bold number, marks right, source tag
%     bottom right), then a bordered table:  Solution | Comments  (bold header row);
%     body rows  label + [marks] (bold, narrow left column) | working | comments,
%     a full-width rule between parts, working left-aligned (fleqn).
% Rows come from the existing sources: \pt in qparts (Ch10 HW), \slabel (Ch10 Test),
% spart (Ch11). Marks under each label are taken from the question's \mk.
% =====================================================================
\makeatletter
\usepackage{lastpage}
\usepackage{adjustbox}
\usepackage{zref-savepos,zref-abspage}
\zref@addprop{savepos}{abspage}
\geometry{a4paper, top=22mm, bottom=24mm, left=22mm, right=22mm, headsep=8mm, footskip=12mm, headheight=24pt}
\providecommand{\risoltitle}{\prelimtitle: Solutions with comments}
\providecommand{\risolschool}{Compiled by Maximus Pung Jun}
\fancypagestyle{risol}{\fancyhf{}%
  \lhead{\fontsize{11}{12}\selectfont 2026 \risoltitle\\[-1.2ex]\rule{\textwidth}{0.6pt}}%
  \rfoot{\calibrifont\fontsize{11}{13}\selectfont Page \textbf{\thepage} of \textbf{\pageref*{LastPage}}}%
  \renewcommand{\headrulewidth}{0pt}\renewcommand{\footrulewidth}{0pt}}
\fancypagestyle{risolfirst}{\fancyhf{}%
  \rfoot{\calibrifont\fontsize{11}{13}\selectfont Page \textbf{\thepage} of \textbf{\pageref*{LastPage}}}%
  \renewcommand{\headrulewidth}{0pt}\renewcommand{\footrulewidth}{0pt}}
\pagestyle{risol}
\newif\ifri@fresh
\newcommand{\risolfront}{\thispagestyle{risolfirst}\global\ri@freshtrue%
  {\noindent\bfseries\linespread{1}\selectfont \risolschool\\ H2 Mathematics (9758)\\ 2026\par}%
  \vspace{4ex}{\noindent\bfseries\underline{2026 \risoltitle}\par}\vspace{3ex}}
% ---- table geometry (RI: label 42pt | working 284pt | comments 142pt of 468pt) ----
\newlength{\ri@L}\setlength{\ri@L}{15mm}
\newlength{\ri@C}\setlength{\ri@C}{40mm}
\newlength{\ri@pad}\setlength{\ri@pad}{1.8mm}
\newif\ifri@sol \newif\ifri@first
\newcount\ri@depth
\def\ri@outer{}
\newtcolorbox{risol}{enhanced, breakable, sharp corners, frame hidden, colback=white,
  boxrule=0pt, boxsep=0pt, left=\dimexpr\ri@L+\ri@pad\relax, right=\dimexpr\ri@C+\ri@pad\relax,
  top=1.2mm, bottom=2.5mm, toptitle=0.6mm, bottomtitle=0.6mm, lefttitle=\ri@pad, righttitle=0pt,
  colbacktitle=white, coltitle=black, fonttitle=\bfseries,
  title={Solution\hfill\makebox[\ri@C][l]{\hspace{\ri@pad}Comments}},
  before skip=2ex, after skip=2ex,
  overlay unbroken and first={\draw[line width=0.5pt] (frame.south west) rectangle (frame.north east);
     \draw[line width=0.5pt] (title.south west) -- (title.south east);
     \draw[line width=0.5pt] ([xshift=\ri@L]interior.north west) -- ([xshift=\ri@L]frame.south west);
     \draw[line width=0.5pt] ([xshift=-\ri@C]frame.north east) -- ([xshift=-\ri@C]frame.south east);},
  overlay middle and last={\draw[line width=0.5pt] (frame.south west) rectangle (frame.north east);
     \draw[line width=0.5pt] ([xshift=\ri@L]frame.north west) -- ([xshift=\ri@L]frame.south west);
     \draw[line width=0.5pt] ([xshift=-\ri@C]frame.north east) -- ([xshift=-\ri@C]frame.south east);},
  after upper={\ri@endrow\protected@write\@auxout{}{\string\ri@settotal{\ri@qid}{\the\ri@sumH}}},
  before upper={\global\ri@sumH=14mm\ri@soltrue\ri@firsttrue\ri@depth=0 \parindent=0pt \parskip=0.6ex
     \predisplaypenalty=100 
     \linespread{1.12}\selectfont\raggedright\@fleqntrue\@mathmargin=0pt
     \abovedisplayskip=5pt plus 2pt\belowdisplayskip=5pt plus 2pt
     \abovedisplayshortskip=3pt plus 1pt\belowdisplayshortskip=4pt plus 2pt
     \ifcsname ri@hp@\ri@qid\endcsname\else\ri@row{}\fi}}
\def\ri@qid{}
% one row: full-width rule (glue leaders, so it vanishes at a page break where the frame
% edge takes its place; not before the first row), then bold label + bold [marks]
% in the left column, on the baseline of the first line of the working
\ExplSyntaxOn
\tl_new:N \g__ri_tmp_tl
\cs_new_protected:Npn \ri_setkey:Nn #1#2
  { \group_begin: \cs_set_eq:NN \chg \use_ii:nn \cs_set_eq:NN \chgm \use_ii:nn
    \cs_set_eq:NN \textbf \use:n \cs_set_eq:NN \mathrm \use:n
    \tl_gset:Nx \g__ri_tmp_tl {#2} \group_end:
    \tl_set_eq:NN #1 \g__ri_tmp_tl \tl_remove_all:Nn #1 {~} }
\cs_set_eq:NN \ri@setkey \ri_setkey:Nn
\ExplSyntaxOff
% ---- page-break rules (positions from the previous run, zref labels in the .aux) ----
% * a row (one part) never breaks across pages: if it does not fit on the rest of the
%   page it starts the next page (forced break between rows); only a row taller than a
%   whole page may break inside;
% * a whole solution that fits on one page is kept together (needspace before the table);
%   a longer one starts where at least its first row fits;
% * a row is at least as tall as its label + [marks].
\newcount\ri@n \newbox\ri@labbox \def\ri@open{0}
\newdimen\ri@ybot \newdimen\ri@ytop \newdimen\ri@h \newdimen\ri@sumH
\AtBeginDocument{\ri@ybot=\dimexpr\paperheight-1in-\voffset-\topmargin-\headheight-\headsep-\textheight\relax
  \ri@ytop=\dimexpr\ri@ybot+\textheight\relax}
\def\ri@pg#1{\zref@extractdefault{#1}{abspage}{-1}}
% \ri@rowh{n}: height of row n in the previous run -> \ri@h (0pt if unknown)
\newcommand{\ri@rowh}[1]{\ri@h=0pt
  \ifnum\ri@pg{ri-s#1}>0 \ifnum\ri@pg{ri-e#1}>0
    \ifnum\ri@pg{ri-s#1}=\ri@pg{ri-e#1}\relax \ri@h=\dimexpr\zposy{ri-s#1}sp-\zposy{ri-e#1}sp\relax
    \else \ri@h=\dimexpr\zposy{ri-s#1}sp-\ri@ybot+\ri@ytop-\zposy{ri-e#1}sp-6mm\relax\fi\fi\fi}
\newcommand{\ri@setfirst}[2]{\expandafter\gdef\csname ri@first@#1\endcsname{#2}}
\newcommand{\ri@setkept}[1]{\expandafter\gdef\csname ri@kept@#1\endcsname{}}
\newcommand{\ri@settotal}[2]{\expandafter\gdef\csname ri@total@#1\endcsname{#2}}
\newcommand{\ri@endrow}{\ifnum\ri@open=1 \par\zsavepos{ri-e\the\ri@n}%
  \ifnum\ri@pg{ri-s\the\ri@n}=\ri@pg{ri-e\the\ri@n}\relax
    \dimen@=\dimexpr\ri@curH sp+1.6mm-\zposy{ri-s\the\ri@n}sp+\zposy{ri-e\the\ri@n}sp\relax
    \ifdim\dimen@>0pt \nobreak\kern\dimen@\penalty\z@\fi\fi
  \ri@rowh{\the\ri@n}\global\advance\ri@sumH\dimexpr\ri@h+3.5mm\relax
  \gdef\ri@open{0}\fi}
\newcommand{\ri@row}[1]{\par\ri@endrow\global\advance\ri@n\@ne\gdef\ri@open{1}\@beginparpenalty\@M
  \ifri@first\global\ri@firstfalse\ri@firstfalse
    \protected@write\@auxout{}{\string\ri@setfirst{\ri@qid}{\the\ri@n}}%
  \else
    \zsavepos{ri-p\the\ri@n}%
    % push the row to the next page if it would not fit on the rest of this one
    \ri@rowh{\the\ri@n}%
    \ifnum\ri@pg{ri-p\the\ri@n}>0 \ifdim\ri@h>0pt
      \ifdim\dimexpr\ri@h+4mm\relax<\dimexpr\textheight-10mm\relax
        \ifdim\dimexpr\ri@h+7mm\relax>\dimexpr\zposy{ri-p\the\ri@n}sp-\ri@ybot\relax
          \penalty-\@M\fi\fi\fi\fi
    \vskip0.6ex\nointerlineskip
    \leaders\hbox{\hskip-\dimexpr\ri@L+\ri@pad\relax\vrule width\dimexpr\linewidth+\ri@L+\ri@C+2\ri@pad\relax height0.5pt depth0pt}\vskip0.5pt
    \nobreak\vskip1.4mm\fi
  \ri@setkey\ri@key{\ri@qid#1}\def\ri@lab{#1}%
  \setbox\ri@labbox=\vtop{\hsize=\dimexpr\ri@L-1.6mm\relax\bfseries\raggedright\parskip=0pt
      \strut\ri@lab\ifcsname ri@mk@\ri@key\endcsname\ifx\ri@lab\@empty\else\\\fi
      [\csname ri@mk@\ri@key\endcsname]\fi\strut\par}%
  \xdef\ri@curH{\number\dimexpr\ht\ri@labbox+\dp\ri@labbox\relax}%
  \nointerlineskip\zsavepos{ri-s\the\ri@n}\nointerlineskip\nobreak
  \vbox to 0pt{\hbox{\hskip-\dimexpr\ri@L+\ri@pad-1.2mm\relax\box\ri@labbox}\vss}\nobreak\prevdepth=\dp\strutbox
  \vskip-\parskip\ignorespaces}
% ---- qparts / \pt: question mode records labels (for the marks); solution mode makes rows
\NewCommandCopy\ri@qparts@begin\qparts \let\ri@qparts@end\endqparts
\renewenvironment{qparts}[1][2.6em]{\advance\ri@depth\@ne
  \ifri@sol\else\ri@qparts@begin[#1]\fi}{\ifri@sol\par\else\ri@qparts@end\fi}
\def\ri@qpartsname{qparts}
\renewcommand{\pt}[1]{\ifnum\ri@depth<2 \gdef\ri@outer{#1}\gdef\ri@cur{#1}\else\expandafter\gdef\expandafter\ri@cur\expandafter{\ri@outer#1}\fi
  \ifri@sol\expandafter\ri@solpt\else\expandafter\ri@qpt\fi{#1}}
\def\ri@qpt#1{\expandafter\gdef\csname ri@hp@\ri@qid\endcsname{}\item[#1]}
\def\ri@solpt#1{\@ifnextchar\begin\ri@ptpeek{\ri@row{\ri@cur}}}
\def\ri@ptpeek\begin#1{\def\ri@tmp{#1}\ifx\ri@tmp\ri@qpartsname
   \expandafter\@firstoftwo\else\expandafter\@secondoftwo\fi
  {\begin{qparts}}{\ri@row{\ri@cur}\begin{#1}}}
% marks: record [n] against the current question/part
\let\ri@mk\mk
\renewcommand{\mk}[1]{\ri@setkey\ri@k{\ri@qid\ri@cur}\ri@setkey\ri@mv{#1}\expandafter\xdef\csname ri@mk@\ri@k\endcsname{\ri@mv}\ri@mk{#1}}
% \slabel (Ch10 Test answers): letter labels set the outer part, roman ones are appended;
% consecutive \slabel calls merge into one row label
\def\ri@romans{(i)(ii)(iii)(iv)(v)(vi)(vii)}
\def\ri@firstpart(#1)#2\@nil{(#1)}
\def\slabel#1{\ifri@sol\expandafter\ri@solslabel\else\expandafter\ri@qslabel\fi{#1}}
\def\ri@qslabel#1{\par\smallskip\noindent\textbf{#1}\par\nobreak}
\def\ri@solslabel#1{\@expandtwoargs\in@{#1}{\ri@romans}%
  \ifin@\expandafter\gdef\expandafter\ri@cur\expandafter{\ri@outer#1}\else\gdef\ri@cur{#1}\xdef\ri@outer{\ri@firstpart#1\@nil}\fi
  \@ifnextchar\slabel{}{\ri@row{\ri@cur}}}
% Ch11 spart{label}: a row; soln: the table
\def\spart#1{\gdef\ri@cur{#1}\ri@row{#1}}\def\endspart{\par}
\def\soln{\@ifnextchar[{\ri@solnopt}{\risol}}\def\ri@solnopt[#1]{\risol}\def\endsoln{\endrisol}
% ---- working that started on the label line is left-aligned in the working column
\AtBeginDocument{\@ifundefined{topalign}{}{%
  \renewenvironment{topalign}{\par\noindent$\displaystyle\begin{aligned}[t]}{\end{aligned}$\par\smallskip}%
  \renewcommand{\topdisplay}[1]{\par\noindent$\displaystyle #1$\par\smallskip}}}
% ---- Ch11 side-by-side (text | sketch): the sketch shrinks to fit its half of the working column
\AtBeginDocument{\@ifundefined{sbsplain}{}{%
  \renewenvironment{sbsplain}[1]{%
    \sbs@start\begin{minipage}[t]{#1\linewidth}\vspace{0pt}\sbs@shift\raggedright%
    \def\sbsnext{\end{minipage}\hfill\begin{minipage}[t]{\dimexpr0.97\linewidth-#1\linewidth\relax}\vspace{0pt}\sbs@shift\centering
      \begin{adjustbox}{max width=\linewidth}}%
  }{\end{adjustbox}\end{minipage}\par}}}
% ---- figures in the working column shrink to fit it (only if wider), colours kept
\newcommand{\rifitfig}[1]{\expandafter\let\csname ri@orig@\string#1\endcsname#1%
  \def#1{\@ifnextchar\bgroup{\ri@fitarg#1}{\ri@fitnoarg#1}}}
\def\ri@fitarg#1#2{\ifri@sol\adjustbox{max width=\linewidth}{\csname ri@orig@\string#1\endcsname{#2}}\else\csname ri@orig@\string#1\endcsname{#2}\fi}
\def\ri@fitnoarg#1{\ifri@sol\adjustbox{max width=\linewidth}{\csname ri@orig@\string#1\endcsname}\else\csname ri@orig@\string#1\endcsname\fi}
\@ifundefined{solfig}{}{\renewcommand{\solfig}[2][]{\par\begin{center}\adjustbox{max width=\linewidth}{\begin{tikzpicture}\begin{axis}[qfig, width=6.6cm, height=5cm, #1]#2\end{axis}\end{tikzpicture}}\end{center}}}
% \ricomment{text}: text in the Comments column of the current row (call right after \pt,
% \slabel or \spart). The row is made at least as tall as the comment.
\newbox\ri@cmtbox
\newcommand{\ricomment}[1]{\par\setbox\ri@cmtbox=\vtop{\hsize=\dimexpr\ri@C-2\ri@pad\relax
    \parskip=0pt\raggedright\small{\color{black!60}\strut#1\strut\par}}%
  \ifdim\dimexpr\ht\ri@cmtbox+\dp\ri@cmtbox\relax>\ri@curH sp\xdef\ri@curH{\number\dimexpr\ht\ri@cmtbox+\dp\ri@cmtbox\relax}\fi
  \nointerlineskip\vbox to 0pt{\hbox to\linewidth{\hskip\dimexpr\linewidth+2\ri@pad\relax\box\ri@cmtbox\hss}\vss}\nobreak\prevdepth=\dp\strutbox
  \vskip-\parskip\ignorespaces}
% \putTFlong{n}: the True/False question n with its answer table allowed to break across
% pages (the body of \defQ{n} must be: intro text, then \begin{tftable}...\end{tftable});
% the solutions copy prints it with tflongtable.
\renewcommand{\putTFlong}[1]{\ri@newq\begingroup\let\tftable\tflongtable\let\endtftable\endtflongtable
  \stepcounter{diq}\expandafter\xdef\csname diLab@#1\endcsname{\thediq}%
  \par\noindent\hbox to 2.6em{\textbf{\thediq}\hfil}\setbox0=\hbox{\di@tag{#1}}\ifdim\wd0=0pt\else{\small\textcolor{black!60}{[\di@tag{#1}]}}\par\noindent\hskip2.6em\fi\ignorespaces\csname diQbody@#1\endcsname\par
  \endgroup\vspace{3ex}}
% ---- question + solution; each question starts on a new page (as RI)
\def\ri@qid{}
\let\ri@putQ\putQ
\renewcommand{\putQ}[1]{\gdef\ri@qid{#1}\global\ri@depth=0 \gdef\ri@cur{}\ri@putQ{#1}}
% \risolQ{n}{body macro name}: question n + its solution body in the table;
% \risolQraw: body already contains its own soln environment (Ch11 Test)
\newcommand{\ri@newq}{\ifri@fresh\global\ri@freshfalse\else\clearpage\fi}
% keep the whole table together if it fits on a page; otherwise make sure at least its
% header and first row fit before starting it
\newcommand{\ri@keep}[1]{\par
  \ifcsname ri@total@#1\endcsname
    % hysteresis: a table kept whole in the previous run stays whole up to textheight-4.5mm,
    % so a total near the limit cannot flip between runs
    \ifdim\csname ri@total@#1\endcsname<\dimexpr\textheight-\ifcsname ri@kept@#1\endcsname4.5\else8\fi mm\relax
      \protected@write\@auxout{}{\string\ri@setkept{#1}}%
      \needspace{\dimexpr\csname ri@total@#1\endcsname+4mm\relax}%
    \else\ifcsname ri@first@#1\endcsname
      \ri@rowh{\csname ri@first@#1\endcsname}%
      \ifdim\ri@h<\dimexpr\textheight-20mm\relax\needspace{\dimexpr\ri@h+18mm\relax}\fi\fi\fi\fi}
\newcommand{\risolQ}[2]{\ri@newq\putQ{#1}\vspace{-3ex}\ri@keep{#1}\begin{risol}\csname #2@#1\endcsname\par\end{risol}}
\newcommand{\risolQraw}[2]{\ri@newq\putQ{#1}\vspace{-3ex}\ri@keep{#1}\csname #2@#1\endcsname\par}
\renewcommand{\putQS}[1]{\risolQ{#1}{diS}}
% end marker only if it fits on the last page (no page holding just the marker)
\renewcommand{\diend}[1][END OF SOLUTIONS]{\par\ifdim\dimexpr\pagegoal-\pagetotal\relax<4\baselineskip\else\vspace{1.5ex}\begin{center}\textbf{#1}\end{center}\fi}
\makeatother
\tfanswerstrue\fi
\fi

\usepackage{pifont}
\usetikzlibrary{arrows.meta, calc, positioning, 3d, decorations.pathreplacing, angles, quotes}
% =====================================================================
% vec_common.tex: shared by the notes (via notes/preamble_local.tex) and
% the slides (via slides/slides_local.tex) for the Vectors chapter.
%
%   CHAPTER NUMBER: change the 13 below and rebuild. The notes set the
%   chapter counter from it, and every slide frame title, Example/Quiz
%   number and box number on the slides uses \chapnum.
% =====================================================================
\def\chapnum{13}

% ---- answers -------------------------------------------------------------
% \ans{key}{value} just prints the value. checks/verify.py reads every
% \ans in the notes and slides and checks it with sympy.
\newcommand{\ans}[2]{#2}

% Column vector, the form used in MF27 and in the examples.
% pmatrix cells are text style, so \dfrac alone still looks small and the
% rows collide. Force display style in every cell and open the rows up.
% Each cell gets a strut as tall as its own content plus \colpad above and
% below, so rows of \dfrac never touch, while integer rows stay compact.
\newlength{\colpad}\setlength{\colpad}{0.55ex}
\newcommand{\colcell}[1]{%
  \setbox0=\hbox{$\displaystyle #1$}%
  \rule[-\dimexpr\dp0+\colpad\relax]{0pt}{\dimexpr\ht0+\dp0+2\colpad\relax}%
  \box0}
\newcommand{\col}[3]{%
  \begingroup
  \renewcommand{\arraystretch}{1.35}%
  \begin{pmatrix}
    \colcell{#1}\\
    \colcell{#2}\\
    \colcell{#3}
  \end{pmatrix}%
  \endgroup
}

% \place{Name}{x}{y}{z} defines a TikZ coordinate. The parentheses are
% inside the macro, so TikZ sees them only after expansion, and
% checks/verify.py can read the drawn coordinates back out.
\newcommand{\place}[4]{\coordinate (#1) at (#2,#3,#4);}

% ---- 3D pictures ----------------------------------------------------------
% Oblique axes: x to the right and slightly down, y out of the page,
% z straight up. Use scale to fit the frame.
\tikzset{
  threeD/.style={
    x={(0.78cm,-0.20cm)},
    y={(0.36cm,0.26cm)},
    z={(0cm,0.78cm)},
    line cap=round, line join=round,
    >=Stealth
  },
  varr/.style={->, thick},
  pt/.style={circle, fill=black, inner sep=1.15pt},
  plane/.style={fill=blue!12, draw=blue!45!black, thick, fill opacity=0.85},
  hid/.style={dashed, black!55}
}
% vecfig.tex: shared 3D diagram style (tikz-3dplot) for Vectors I / II notes and slides.
\usepackage{tikz-3dplot}
\tdplotsetmaincoords{62}{55}
\definecolor{vA}{RGB}{31,90,180}   % first vector / line 1
\definecolor{vB}{RGB}{210,90,20}   % second vector / line 2
\definecolor{vC}{RGB}{0,128,110}   % result / third
\definecolor{vN}{RGB}{140,40,150}  % normals
\definecolor{vP}{RGB}{200,215,240} % plane 1
\definecolor{vQ}{RGB}{245,215,190} % plane 2
\definecolor{vR}{RGB}{200,235,215} % plane 3
\tikzset{
  v3/.style={tdplot_main_coords, line cap=round, line join=round, >=Stealth, scale=#1},
  v3/.default=1,
  ax/.style={->, black!50, thin},
  va/.style={->, very thick, vA},
  vb/.style={->, very thick, vB},
  vc/.style={->, very thick, vC},
  vn/.style={->, very thick, vN},
  ln/.style={very thick},
  dsh/.style={dashed, black!55},
  pp/.style={fill=vP, fill opacity=0.55, draw=vA!70!black, thin},
  pq/.style={fill=vQ, fill opacity=0.55, draw=vB!70!black, thin},
  pr/.style={fill=vR, fill opacity=0.55, draw=vC!70!black, thin},
  dot/.style={circle, fill=black, inner sep=1.2pt},
}
% \vaxes{xlen}{ylen}{zlen}
\newcommand{\vaxes}[3]{%
  \draw[ax] (0,0,0) -- (#1,0,0) node[anchor=west]{$x$};
  \draw[ax] (0,0,0) -- (0,#2,0) node[anchor=south west]{$y$};
  \draw[ax] (0,0,0) -- (0,0,#3) node[anchor=south]{$z$};
  \node[anchor=east] at (0,0,0) {\scriptsize $O$};}
% \rightmark{corner}{along1}{along2} small right-angle square
\newcommand{\rtang}[3]{\draw[thin] ($(#1)!0.25cm!(#2)$) -- ($($(#1)!0.25cm!(#2)$)+($(#1)!0.25cm!(#3)$)-(#1)$) -- ($(#1)!0.25cm!(#3)$);}
\newenvironment{vis}{\par\medskip\begin{center}}{\end{center}\medskip}
% \vang{vertex}{p1}{p2}{label}: small angle arc drawn on the page
\newcommand{\vang}[4]{%
  \coordinate (vg1) at ($(#1)!0.55cm!(#2)$); \coordinate (vg2) at ($(#1)!0.55cm!(#3)$);
  \coordinate (vgm) at ($(vg1)!0.5!(vg2)$);
  \draw[vC, thick] (vg1) .. controls ($(#1)!0.72cm!(vgm)$) .. (vg2);
  \node[vC] at ($(#1)!0.95cm!(vgm)$) {#4};}

\providecommand{\vv}[1]{\overrightarrow{#1}}\providecommand{\cvv}[3]{\begin{pmatrix}#1\\#2\\#3\end{pmatrix}}
\housechapter{Vectors I}
\ditle{Homework 3}% topic is in the grey note line (rule 1A, 9 Nov 2026)
\tfcompact% allowed exception: the True/False set stays on page 1
\defQ{1}{}{%
For the following statements, state whether it is \textbf{true} or \textbf{false}.
\begin{tftable}
\tfrowl{(a)}{$\mathbf a\times\mathbf a=\mathbf 0$.}{T}{$\displaystyle |\mathbf a\times\mathbf a|=|\mathbf a|^2\sin0^\circ=0$}
\tfrowl{(b)}{$\mathbf a\times\mathbf a=0$.}{F}{$\displaystyle \mathbf a\times\mathbf a=\mathbf 0\neq0\ \text{(left side is the zero vector)}$}
\tfrowl{(c)}{$\mathbf i\times\mathbf j=\mathbf k$.}{T}{$\displaystyle \cvv{1}{0}{0}\times\cvv{0}{1}{0}=\cvv{0}{0}{1}$}
\tfrowl{(d)}{$\mathbf a\times\mathbf b=\mathbf b\times\mathbf a$.}{F}{$\displaystyle \mathbf i\times\mathbf j=\mathbf k,\ \mathbf j\times\mathbf i=-\mathbf k$}
\tfrowl{(e)}{$\mathbf a\times\mathbf b=-(\mathbf b\times\mathbf a)$.}{T}{$\displaystyle \mathbf a\times\mathbf b=-\mathbf b\times\mathbf a$}
\tfrowl{(f)}{$\mathbf a\times\mathbf b$ is a scalar.}{F}{$\displaystyle \mathbf i\times\mathbf j=\mathbf k\ \text{(a vector)}$}
\tfrowl{(g)}{$\mathbf k\times\mathbf j=\mathbf i$.}{F}{$\displaystyle \mathbf j\times\mathbf k=\mathbf i\ \therefore\ \mathbf k\times\mathbf j=-\mathbf i$}
\tfrowl{(h)}{$\cvv{1}{2}{3}\times\cvv{2}{4}{6}=\mathbf 0$.}{T}{$\displaystyle \cvv{2}{4}{6}=2\cvv{1}{2}{3}$}
\tfrowl{(i)}{$\mathbf a\times(\mathbf b+\mathbf c)=\mathbf a\times\mathbf b+\mathbf a\times\mathbf c$.}{T}{$\displaystyle \text{Distributive law}$}
\tfrowl{(j)}{$(\lambda\mathbf a)\times\mathbf b=\lambda(\mathbf a\times\mathbf b)$ for all $\lambda\in\mathbb R$.}{T}{$\displaystyle (\lambda\mathbf a)\times\mathbf b=\mathbf a\times(\lambda\mathbf b)=\lambda(\mathbf a\times\mathbf b)$}
\tfrowl{(k)}{For non-zero and non-parallel $\mathbf a$ and $\mathbf b$, $\mathbf a\times\mathbf b$ is perpendicular to both $\mathbf a$ and $\mathbf b$.}{T}{$\displaystyle (\mathbf a\times\mathbf b)\cdot\mathbf a=0,\ (\mathbf a\times\mathbf b)\cdot\mathbf b=0$}
\tfrowl{(l)}{$(\mathbf a\times\mathbf b)\cdot\mathbf c$ is a scalar.}{T}{$\displaystyle (\mathbf a\times\mathbf b)\cdot\mathbf c=\text{vector}\cdot\text{vector}$}
\tfrowl{(m)}{$|\mathbf a\times\mathbf b|=|\mathbf a||\mathbf b|\sin\theta$, where $\theta$ is the angle between $\mathbf a$ and $\mathbf b$.}{T}{$\displaystyle \text{Definition of vector product}$}
\tfrowl{(n)}{$|\mathbf a\times\mathbf b|=|\mathbf a||\mathbf b|$.}{F}{$\displaystyle |\mathbf i\times\mathbf i|=0\neq1$}
\tfrowl{(o)}{If $\mathbf a\times\mathbf b=\mathbf 0$, then $\mathbf a=\mathbf 0$ or $\mathbf b=\mathbf 0$.}{F}{$\displaystyle \mathbf i\times2\mathbf i=\mathbf 0,\ \mathbf i\neq\mathbf 0,\ 2\mathbf i\neq\mathbf 0$}
\tfrowl{(p)}{For non-zero $\mathbf a$ and $\mathbf b$, $\mathbf a\times\mathbf b=\mathbf 0$ means that $\mathbf a$ is parallel to $\mathbf b$.}{T}{$\displaystyle |\mathbf a||\mathbf b|\sin\theta=0\ \therefore\ \theta=0^\circ\ \text{or}\ 180^\circ$}
\tfrowl{(q)}{The area of the parallelogram with adjacent sides $\mathbf a$ and $\mathbf b$ is $|\mathbf a\times\mathbf b|$.}{T}{$\displaystyle \text{Area}=|\mathbf a||\mathbf b|\sin\theta=|\mathbf a\times\mathbf b|$}
\tfrowl{(r)}{The area of triangle $OAB$ is $|\vv{OA}\times\vv{OB}|$.}{F}{$\displaystyle \text{Area}=\tfrac12|\vv{OA}\times\vv{OB}|$}
\tfrowl{(s)}{$\mathbf a\times(\mathbf b\cdot\mathbf c)$ is a vector.}{F}{$\displaystyle \mathbf b\cdot\mathbf c\ \text{is a scalar}\ \therefore\ \text{meaningless}$}
\tfrowl{(t)}{For non-zero $\mathbf a$ and $\mathbf b$, if $|\mathbf a\times\mathbf b|=|\mathbf a||\mathbf b|$, then $\mathbf a$ is perpendicular to $\mathbf b$.}{T}{$\displaystyle \sin\theta=1\ \therefore\ \theta=90^\circ$}
\tfrowl{(u)}{$(\mathbf a\times\mathbf b)\cdot\mathbf c=\mathbf a\cdot(\mathbf b\times\mathbf c)$.}{T}{$\displaystyle \text{both}=\det\left(\mathbf a\ \mathbf b\ \mathbf c\right)$}
\tfrowl{(v)}{$|\mathbf a\times\mathbf b|^2+(\mathbf a\cdot\mathbf b)^2=|\mathbf a|^2|\mathbf b|^2$.}{T}{$\displaystyle |\mathbf a|^2|\mathbf b|^2(\sin^2\theta+\cos^2\theta)=|\mathbf a|^2|\mathbf b|^2$}
\tfrowl{(w)}{$\mathbf a\times\mathbf b=\mathbf a\times\mathbf c$ means that $\mathbf b=\mathbf c$.}{F}{$\displaystyle \mathbf i\times\mathbf j=\mathbf k=\mathbf i\times(\mathbf i+\mathbf j),\ \mathbf j\neq\mathbf i+\mathbf j$}
\tfrowl{(x)}{$\mathbf a\times(\mathbf b\times\mathbf c)=(\mathbf a\times\mathbf b)\times\mathbf c$.}{F}{$\displaystyle \mathbf i\times(\mathbf i\times\mathbf j)=\mathbf i\times\mathbf k=-\mathbf j,\ (\mathbf i\times\mathbf i)\times\mathbf j=\mathbf 0$}
\tfrowl{(y)}{The shortest distance from $B$ to the line through $A$ parallel to $\mathbf d$ is $|\vv{AB}\times\mathbf d|$.}{F}{$\displaystyle \text{distance}=|\vv{AB}\times\hat{\mathbf d}|=\dfrac{|\vv{AB}\times\mathbf d|}{|\mathbf d|}$}
\tfrowl{(z)}{If $\mathbf a\times\mathbf b=\mathbf 0$ and $\mathbf a\cdot\mathbf b=0$, then $\mathbf a=\mathbf 0$ or $\mathbf b=\mathbf 0$.}{T}{$\displaystyle |\mathbf a|^2|\mathbf b|^2=|\mathbf a\times\mathbf b|^2+(\mathbf a\cdot\mathbf b)^2=0$}
\end{tftable}}

\defQ{2}{JPJC 2024 Promo Q11}{%
Referred to the origin $O$, points $A$ and $B$ have position vectors $\mathbf a$ and $\mathbf b$ respectively, where $\mathbf a$ and $\mathbf b$ are non-zero and non-parallel vectors. Point $C$ is the midpoint of $OA$ and point $D$ lies on $BC$ such that $BD:DC=2:3$.
\begin{qparts}
\pt{(i)} Find the position vectors $\vv{OC}$ and $\vv{OD}$, giving your answers in terms of $\mathbf a$ and $\mathbf b$.\mk{2}
\pt{(ii)} Find the area of triangle $OCD$ in the form of $k|\mathbf a\times\mathbf b|$, where $k$ is a constant to be found.\mk{3}
\pt{(iii)} Given that $|\mathbf a|=2$, $|\mathbf b|=\sqrt2$ and angle $AOB$ is $\dfrac\pi4$ radian, using scalar product, show that $CD$ is perpendicular to $OA$.\mk{3}
\pt{(iv)} Given instead that $\mathbf a$ is a unit vector, give a geometrical interpretation of $|\mathbf a\cdot\mathbf b|$.\mk{1}
\end{qparts}}

\defS{2}{%
\pt{(i)} Given $\vv{OA}=\mathbf a$, $\vv{OB}=\mathbf b$, $\vv{OC}=\dfrac12\mathbf a$.
\begin{align*}
\vv{OD}&=\frac{2\mathbf c+3\mathbf b}{5}
=\frac15\left[2\left(\frac12\mathbf a\right)+3\mathbf b\right]
=\frac15(\mathbf a+3\mathbf b).
\end{align*}
\pt{(ii)}
\begin{align*}
\text{Area of triangle }OCD
&=\frac12|\vv{OC}\times\vv{OD}|\\
&=\frac12\left|\frac12\mathbf a\times\frac15(\mathbf a+3\mathbf b)\right|\\
&=\frac12\left|\frac1{10}(\mathbf a\times\mathbf a)+\frac3{10}(\mathbf a\times\mathbf b)\right|\\
&=\frac12\left|\frac3{10}(\mathbf a\times\mathbf b)\right|\qquad(\mathbf a\times\mathbf a=\mathbf0)\\
&=\frac3{20}|\mathbf a\times\mathbf b|,\qquad k=\frac3{20}.
\end{align*}
\pt{(iii)}
\begin{align*}
\vv{CD}&=\vv{OD}-\vv{OC}=\frac15(\mathbf a+3\mathbf b)-\frac12\mathbf a
=\frac35\mathbf b-\frac3{10}\mathbf a\\
\vv{CD}\cdot\vv{OA}
&=\left(\frac35\mathbf b-\frac3{10}\mathbf a\right)\cdot\mathbf a\\
&=\frac35\mathbf b\cdot\mathbf a-\frac3{10}\mathbf a\cdot\mathbf a\\
&=\frac35|\mathbf b||\mathbf a|\cos\frac\pi4-\frac3{10}|\mathbf a|^2\\
&=\frac35(\sqrt2)(2)\frac{\sqrt2}{2}-\frac3{10}(2)^2
=\frac65-\frac{12}{10}=0.
\end{align*}
Since $\vv{CD}\cdot\vv{OA}=0$, $CD$ is perpendicular to $OA$. (Shown)\par
\pt{(iv)} $|\mathbf a\cdot\mathbf b|$ is the length of projection of $\vv{OB}$ onto $\vv{OA}$ (or a line or vector parallel to $\vv{OA}$).}

\defQ{3}{MI 2024 Promo Q4}{%
Referred to the origin $O$, the points $A$, $B$ and $C$ have position vectors $\mathbf a$, $\mathbf b$ and $\mathbf c$. $C$ lies on $AB$ such that $AC:AB=1:3$.
\begin{qparts}
\pt{(i)} Find $\mathbf c$ in terms of $\mathbf a$ and $\mathbf b$.\mk{2}
\pt{(ii)} Show that the area of triangle $OAC$ is given by $k|\mathbf a\times\mathbf b|$, where $k$ is a constant to be determined.\mk{2}
\end{qparts}
It is now given that $\mathbf a$ and $\mathbf b$ are unit vectors.
\begin{qparts}
\pt{(iii)} Using the result in part (i), find the value of $(2\mathbf a-\mathbf b)\cdot\mathbf c$.\mk{3}
\end{qparts}}

\defS{3}{%
\pt{(i)}
\begin{align*}
AC:AB&=1:3\quad\Rightarrow\quad AC:CB=1:2\\
\mathbf c&=\frac{(1)\mathbf b+(2)\mathbf a}{1+2}=\frac{2\mathbf a+\mathbf b}{3}.
\end{align*}
\pt{(ii)}
\begin{align*}
\text{Area of triangle }OAC
&=\frac13\text{ Area of triangle }OAB\\
&=\frac13\times\frac12|\vv{OA}\times\vv{OB}|\\
&=\frac16|\mathbf a\times\mathbf b|,\qquad k=\frac16.
\end{align*}
\pt{(iii)}
\begin{align*}
(2\mathbf a-\mathbf b)\cdot\mathbf c
&=(2\mathbf a-\mathbf b)\cdot\left(\frac{2\mathbf a+\mathbf b}{3}\right)\\
&=\frac13(2\mathbf a-\mathbf b)\cdot(2\mathbf a+\mathbf b)\\
&=\frac13(4\mathbf a\cdot\mathbf a+2\mathbf a\cdot\mathbf b-2\mathbf b\cdot\mathbf a-\mathbf b\cdot\mathbf b)\\
&=\frac13(4|\mathbf a|^2-|\mathbf b|^2)\\
&=\frac13\left(4(1)^2-(1)^2\right)\quad\text{since $\mathbf a$, $\mathbf b$ are unit vectors}\\
&=1.
\end{align*}}

\defQ{4}{ASRJC 2024 Promo Q5}{%
In triangle $PQR$, the vectors $\mathbf a$, $\mathbf b$ and $\mathbf c$ represent the vectors $\vv{PQ}$, $\vv{QR}$ and $\vv{PR}$ respectively.
\begin{center}\begin{tikzpicture}[scale=0.85,font=\small,>=Stealth]
\coordinate(P)at(0,0);\coordinate(Q)at(5,0);\coordinate(R)at(4.15,2.2);
\draw[semithick](P)--(Q)--(R)--cycle;
\draw[->,semithick](P)--(2.7,0);\draw[->,semithick](Q)--(4.57,1.12);\draw[->,semithick](P)--(2.3,1.219);
\node[below left]at(P){$P$};\node[below right]at(Q){$Q$};\node[above]at(R){$R$};
\node[below]at(2.6,0){$\mathbf a$};\node[right]at(4.7,1.15){$\mathbf b$};\node[above left]at(2.1,1.1){$\mathbf c$};
\end{tikzpicture}\end{center}
\begin{qparts}
\pt{(a)} State a geometrical interpretation of $|\mathbf c\times\hat{\mathbf a}|$ and hence show that the area of the triangle $PQR$ can be written as $\dfrac12|\mathbf c\times\mathbf a|$ units$^2$.\mk{3}
\pt{(b)} With clear explanations, show that $|\mathbf c\times\mathbf a|=|\mathbf b\times\mathbf a|$ and deduce that
\[\frac{|\mathbf c|}{\sin\angle PQR}=\frac{|\mathbf b|}{\sin\angle QPR}\quad\text{(Sine Rule).}\]\mk{2}
\end{qparts}}

\defS{4}{%
\pt{(a)} $|\mathbf c\times\hat{\mathbf a}|$ represents the perpendicular distance from the point $R$ to the line $PQ$.
\begin{align*}
\text{Area of }\triangle PQR&=\frac12\times\text{base}\times\text{height}\\
&=\frac12|\mathbf a|\,|\mathbf c\times\hat{\mathbf a}|\\
&=\frac12|\mathbf a|\left|\mathbf c\times\frac{\mathbf a}{|\mathbf a|}\right|\\
&=\frac12|\mathbf c\times\mathbf a|\text{ units}^2.\quad\text{(Shown)}
\end{align*}
\pt{(b)} By replacing $\mathbf a$ with $-\mathbf a$ and $\mathbf c$ with $\mathbf b$ in part (a),
\begin{align*}
\text{Area of }\triangle PQR&=\frac12|\mathbf b\times(-\mathbf a)|=\frac12|\mathbf b\times\mathbf a|\\
\frac12|\mathbf c\times\mathbf a|&=\frac12|\mathbf b\times\mathbf a|\\
|\mathbf c\times\mathbf a|&=|\mathbf b\times\mathbf a|\quad\text{(Shown)}\\
|\mathbf c||\mathbf a|\sin\angle QPR&=|\mathbf b||\mathbf a|\sin\angle PQR\\
\frac{|\mathbf c|}{\sin\angle PQR}&=\frac{|\mathbf b|}{\sin\angle QPR}.\quad\text{(Shown)}
\end{align*}}

\defQ{5}{CJC 2024 Promo Q7}{%
\begin{qparts}
\pt{(a)} Three non-zero vectors $\mathbf p$, $\mathbf q$ and $\mathbf r$ are such that $3\mathbf p\times\mathbf q=\mathbf r\times\mathbf p$.\par
Find a linear relationship between $\mathbf p$, $\mathbf q$ and $\mathbf r$.\mk{3}
\pt{(b)} Referred to the origin $O$, the points $A$ and $B$ have position vectors $\mathbf a$ and $\mathbf b$.\par
It is given that $\mathbf a$ and $\mathbf b$ are non-parallel unit vectors.\par
The point $B$ divides $AC$ in the ratio $1:3$.\par
The point $D$ lies on $OB$ produced such that $OB:OD=1:m$, where $m\in\mathbb Z$, $m\ge2$.\par
Given that $\vv{AB}$ is perpendicular to $\vv{CD}$, find the numerical value of $m$.\mk{4}
\end{qparts}}

\defS{5}{%
\pt{(a)}
\begin{align*}
3\mathbf p\times\mathbf q&=\mathbf r\times\mathbf p\\
(3\mathbf p\times\mathbf q)-(\mathbf r\times\mathbf p)&=\mathbf0\\
(3\mathbf p\times\mathbf q)+(\mathbf p\times\mathbf r)&=\mathbf0\\
\mathbf p\times(3\mathbf q+\mathbf r)&=\mathbf0\\
\mathbf p&\parallel3\mathbf q+\mathbf r\\
\mathbf p&=k(3\mathbf q+\mathbf r),\quad k\in\mathbb R\setminus\{0\}.
\end{align*}
{\small\textit{Verification note (not part of the official scheme):} the printed last line assumes $3\mathbf q+\mathbf r\ne\mathbf0$. The question also permits $3\mathbf q+\mathbf r=\mathbf0$. The unrestricted relationship is $3\mathbf q+\mathbf r=\lambda\mathbf p$, $\lambda\in\mathbb R$, including $\lambda=0$.}\par\smallskip
\pt{(b)}
\begin{align*}
\vv{OB}&=\frac{3\vv{OA}+\vv{OC}}4\\
\vv{OC}&=4\mathbf b-3\mathbf a\\
\vv{OD}&=m\vv{OB}=m\mathbf b\\
\vv{AB}&=\mathbf b-\mathbf a\\
\vv{CD}&=m\mathbf b-(4\mathbf b-3\mathbf a)=3\mathbf a+(m-4)\mathbf b.
\end{align*}
Given that $\vv{AB}$ is perpendicular to $\vv{CD}$,
\begin{align*}
\vv{AB}\cdot\vv{CD}&=0\\
(\mathbf b-\mathbf a)\cdot[3\mathbf a+(m-4)\mathbf b]&=0\\
3\mathbf a\cdot\mathbf b+(m-4)\mathbf b\cdot\mathbf b-3\mathbf a\cdot\mathbf a-(m-4)\mathbf a\cdot\mathbf b&=0\\
3\mathbf a\cdot\mathbf b+(m-4)|\mathbf b|^2-3|\mathbf a|^2-(m-4)\mathbf a\cdot\mathbf b&=0\\
(m-4)-3-(m-7)\mathbf a\cdot\mathbf b&=0\\
m(1-\mathbf a\cdot\mathbf b)&=3-7\mathbf a\cdot\mathbf b+4\\
m&=\frac{3-7\mathbf a\cdot\mathbf b+4}{1-\mathbf a\cdot\mathbf b}\\
&=\frac{7(1-\mathbf a\cdot\mathbf b)}{1-\mathbf a\cdot\mathbf b}=7.
\end{align*}}

\defQ{6}{SAJC 2024 Promo Q9}{%
\begin{qparts}
\pt{(a)} Relative to the origin $O$, the points $A$ and $B$ have position vectors $\mathbf a$ and $\mathbf b$ respectively, where $\mathbf a$ and $\mathbf b$ are non-parallel. It is given that $|\mathbf a|=9$, $|\mathbf b|=1$ and $|2\mathbf a+\mathbf b|=2\sqrt{74}$.
\begin{qparts}[2.2em]
\pt{(i)} By considering $(2\mathbf a+\mathbf b)\cdot(2\mathbf a+\mathbf b)$, show that $\mathbf a\cdot\mathbf b=-\dfrac{29}{4}$.\mk{3}
\pt{(ii)} Give a geometrical meaning of $|\mathbf a\cdot\mathbf b|$.\mk{1}
\pt{(iii)} The points $P$, $Q$ and $R$ have position vectors $7\mathbf a-5\mathbf b$, $6\mathbf a+5\mathbf b$ and $9\mathbf a+\lambda\mathbf b$ respectively. Given that these three points are collinear, find the value of $\lambda$.\mk{3}
\end{qparts}
\pt{(b)} In the triangle $OUV$ where $O$ is the origin, the position vectors of the points $U$ and $V$ are $\mathbf u$ and $\mathbf v$ respectively. The point $W$ is the midpoint of $OU$, and the point $X$ has position vector given by $\dfrac1{16}\mathbf u+\dfrac34\mathbf v$. Show that the area of triangle $UWX$ can be written as $m|\mathbf u\times\mathbf v|$ where $m$ is a constant to be found.\mk{3}
\end{qparts}}

\defS{6}{%
\pt{(a)(i)}
\begin{align*}
(2\mathbf a+\mathbf b)\cdot(2\mathbf a+\mathbf b)
&=4(\mathbf a\cdot\mathbf a)+2(\mathbf b\cdot\mathbf a)+2(\mathbf a\cdot\mathbf b)+(\mathbf b\cdot\mathbf b)\\
&=4|\mathbf a|^2+4(\mathbf a\cdot\mathbf b)+|\mathbf b|^2\\
|2\mathbf a+\mathbf b|^2&=4|\mathbf a|^2+4(\mathbf a\cdot\mathbf b)+|\mathbf b|^2\\
(2\sqrt{74})^2&=4(9^2)+4(\mathbf a\cdot\mathbf b)+(1^2)\\
296&=4(9^2)+4(\mathbf a\cdot\mathbf b)+(1^2)\\
4(\mathbf a\cdot\mathbf b)&=296-4(9^2)-(1^2)\\
\mathbf a\cdot\mathbf b&=-\frac{29}4.
\end{align*}
\pt{(a)(ii)} $|\mathbf a\cdot\mathbf b|$ is the length of projection of $\vv{OA}$ onto $\vv{OB}$.\par
\pt{(a)(iii)} If $P$, $Q$ and $R$ are collinear, then $k\vv{PQ}=\vv{PR}$.
\begin{align*}
k[(6\mathbf a+5\mathbf b)-(7\mathbf a-5\mathbf b)]&=[(9\mathbf a+\lambda\mathbf b)-(7\mathbf a-5\mathbf b)]\\
k[-\mathbf a+10\mathbf b]&=[2\mathbf a+(\lambda+5)\mathbf b].
\end{align*}
Since $\mathbf a$ and $\mathbf b$ are non-parallel, non-zero vectors,
\begin{align*}
k&=-2\\
\lambda+5&=10k=-20\\
\lambda&=-25.
\end{align*}
\pt{(b)}
\begin{align*}
\vv{UX}&=\vv{OX}-\vv{OU}=\frac1{16}\mathbf u+\frac34\mathbf v-\mathbf u=\frac34\mathbf v-\frac{15}{16}\mathbf u\\
\vv{UW}&=\vv{OW}-\vv{OU}=\frac12\mathbf u-\mathbf u=-\frac12\mathbf u\\
\text{Area of triangle }UWX
&=\frac12|\vv{UW}\times\vv{UX}|\\
&=\frac12\left|\frac12\mathbf u\times\left(\frac34\mathbf v-\frac{15}{16}\mathbf u\right)\right|\\
&=\frac3{16}|\mathbf u\times\mathbf v|\qquad(\mathbf u\times\mathbf u=\mathbf0)\\
m&=\frac3{16}.
\end{align*}}

\defQ{7}{CJC 2025 Promo Q6}{%
The diagram shows a unit circle with centre $O$ and diameter $AB$. The point $C$ lies on the circumference of the circle. The position vectors of $A$, $B$ and $C$ are $\mathbf a$, $\mathbf b$ and $\mathbf c$ respectively.
\begin{center}\begin{tikzpicture}[scale=.72,font=\small,>=Stealth]
\coordinate(O)at(0,0);\coordinate(A)at(2,0);\coordinate(B)at(-2,0);\coordinate(C)at(-.9,1.786);
\draw[semithick](O)circle(2);\draw[semithick](B)--(A)--(C)--cycle;
\draw[->,semithick](O)--(1.2,0);\draw[->,semithick](O)--(-1.2,0);\draw[->,semithick](O)--(C);
\node[below]at(O){$O$};\node[right]at(A){$A$};\node[left]at(B){$B$};\node[above]at(C){$C$};
\node[below]at(1,0){$\mathbf a$};\node[below]at(-1,0){$\mathbf b$};\node[left]at(-.45,.9){$\mathbf c$};
\end{tikzpicture}\end{center}
\begin{qparts}
\pt{(a)} By considering an appropriate scalar product, show that $AC$ and $BC$ are perpendicular.\mk{4}
\pt{(b)} Give a geometrical interpretation of $|\mathbf a\cdot(\mathbf c-\mathbf a)|$.\mk{1}
\end{qparts}
The point $D$ is such that $ABCD$ is a parallelogram. The point $M$ is the mid-point of $AD$ and the point $R$ lies on $OA$ produced such that $OA:AR=1:2$.
\begin{center}\begin{tikzpicture}[scale=.64,font=\small]
\coordinate(O)at(0,0);\coordinate(A)at(2,0);\coordinate(B)at(-2,0);\coordinate(C)at(-.9,1.786);\coordinate(D)at(3.1,1.786);\coordinate(M)at(2.55,.893);\coordinate(R)at(6,0);
\draw[semithick](O)circle(2);\draw[semithick](B)--(R);\draw[semithick](B)--(C)--(D)--(A)--(C);
\fill(O)circle(1.6pt);\fill(M)circle(1.6pt);
\node[below]at(O){$O$};\node[below]at(A){$A$};\node[left]at(B){$B$};\node[above]at(C){$C$};\node[above]at(D){$D$};\node[right]at(M){$M$};\node[below]at(R){$R$};
\end{tikzpicture}\end{center}
\begin{qparts}
\pt{(c)} By considering $\vv{CM}\times\vv{CR}$ or otherwise, show that the points $C$, $M$ and $R$ are collinear.\mk{3}
\end{qparts}}

\defS{7}{%
\pt{(a)}
\begin{align*}
\vv{AC}&=\mathbf c-\mathbf a,\qquad\vv{BC}=\mathbf c-\mathbf b\\
\vv{AC}\cdot\vv{BC}
&=(\mathbf c-\mathbf a)\cdot(\mathbf c-\mathbf b)\\
&=(\mathbf c-\mathbf a)\cdot(\mathbf c-(-\mathbf a))\\
&=(\mathbf c-\mathbf a)\cdot(\mathbf c+\mathbf a)\\
&=\mathbf c\cdot\mathbf c-\mathbf a\cdot\mathbf a\\
&=|\mathbf c|^2-|\mathbf a|^2\\
&=0\qquad(|\mathbf c|=|\mathbf a|=\text{radius of circle}).
\end{align*}
Hence $AC$ and $BC$ are perpendicular.\par
\pt{(b)} $|\mathbf a\cdot(\mathbf c-\mathbf a)|$ is the length of projection of $\vv{AC}$ onto $\vv{OA}$.\par
\pt{(c)} By Ratio Theorem,
\begin{align*}
\vv{OM}&=\frac{\vv{OD}+\vv{OA}}2\\
&=\frac12(\mathbf c+2\mathbf a+\mathbf a)=\frac12\mathbf c+\frac32\mathbf a\\
\vv{CM}\times\vv{CR}
&=(\vv{OM}-\vv{OC})\times(\vv{OR}-\vv{OC})\\
&=\left[\frac12\mathbf c+\frac32\mathbf a-\mathbf c\right]\times(3\mathbf a-\mathbf c)\\
&=\left(\frac32\mathbf a-\frac12\mathbf c\right)\times(3\mathbf a-\mathbf c)\\
&=\frac92\mathbf a\times\mathbf a-\frac32\mathbf a\times\mathbf c-\frac32\mathbf c\times\mathbf a+\frac12\mathbf c\times\mathbf c\\
&=\mathbf0-\frac32\mathbf a\times\mathbf c+\frac32\mathbf a\times\mathbf c+\mathbf0\\
&=\mathbf0.
\end{align*}
$\vv{CM}$ and $\vv{CR}$ are parallel. Since $CR$ and $CM$ are parallel and $C$ is a common point, hence $C$, $M$ and $R$ are collinear.}

% ---- Answers (rule 8A, 9 Nov 2026): copied from the final lines of the worked solutions
\defA{1}{\textbf{(a)} T\quad \textbf{(b)} F\quad \textbf{(c)} T\quad \textbf{(d)} F\quad \textbf{(e)} T\quad \textbf{(f)} F\quad \textbf{(g)} F\quad \textbf{(h)} T\quad \textbf{(i)} T\quad \textbf{(j)} T\quad \textbf{(k)} T\quad \textbf{(l)} T\quad \textbf{(m)} T\quad \textbf{(n)} F\quad \textbf{(o)} F\quad \textbf{(p)} T\quad \textbf{(q)} T\quad \textbf{(r)} F\quad \textbf{(s)} F\quad \textbf{(t)} T\quad \textbf{(u)} T\quad \textbf{(v)} T\quad \textbf{(w)} F\quad \textbf{(x)} F\quad \textbf{(y)} F\quad \textbf{(z)} T}
\defA{2}{\begin{aparts}
\pt{(i)} $\vv{OC}=\dfrac12\mathbf a$, $\vv{OD}=\dfrac15(\mathbf a+3\mathbf b)$
\pt{(ii)} $\dfrac3{20}|\mathbf a\times\mathbf b|$, $k=\dfrac3{20}$
\pt{(iii)} Shown.
\pt{(iv)} $|\mathbf a\cdot\mathbf b|$ is the length of projection of $\vv{OB}$ onto $\vv{OA}$ (or a line or vector parallel to $\vv{OA}$).
\end{aparts}}
\defA{3}{\begin{aparts}
\pt{(i)} $\mathbf c=\dfrac{2\mathbf a+\mathbf b}{3}$
\pt{(ii)} Shown. $k=\dfrac16$
\pt{(iii)} $1$
\end{aparts}}
\defA{4}{\begin{aparts}
\pt{(a)} $|\mathbf c\times\hat{\mathbf a}|$ represents the perpendicular distance from the point $R$ to the line $PQ$. Shown.
\pt{(b)} Shown.
\end{aparts}}
\defA{5}{\begin{aparts}
\pt{(a)} $\mathbf p=k(3\mathbf q+\mathbf r)$, $k\in\mathbb R\setminus\{0\}$
\pt{(b)} $m=7$
\end{aparts}}
\defA{6}{\begin{aparts}
\pt{(a)(i)} Shown.
\pt{(a)(ii)} $|\mathbf a\cdot\mathbf b|$ is the length of projection of $\vv{OA}$ onto $\vv{OB}$.
\pt{(a)(iii)} $\lambda=-25$
\pt{(b)} $m=\dfrac3{16}$
\end{aparts}}
\defA{7}{\begin{aparts}
\pt{(a)} Shown.
\pt{(b)} $|\mathbf a\cdot(\mathbf c-\mathbf a)|$ is the length of projection of $\vv{AC}$ onto $\vv{OA}$.
\pt{(c)} Shown.
\end{aparts}}

\ifdefined\diquestionsonly\expandafter\endinput\fi
\begin{document}
\dititleblock[Vector product: properties, parallel vectors, areas, collinearity. Q1 True/False from the Basics Worksheet. Q2--Q7 are complete past-paper questions on the vector product. Total: 46 marks (Q1 unmarked).]
\disection{Questions}
\putQ{1}
\clearpage
\putQ{2}\putQ{3}\putQ{4}\putQ{5}\putQ{6}\putQ{7}
\clearpage
\disection{Answers}
\putA{1}\putA{2}\putA{3}\putA{4}\putA{5}\putA{6}\putA{7}
\diend
\end{document}
